%% ===========================================================================
%%  Mathilda -- A User's Manual and a Course in Computational Mathematics
%%
%%  Master file.  Build with the book Makefile (which regenerates the verified
%%  transcripts and the builtin-link table first, then runs latexmk):
%%
%%      make -C book pdf
%%
%%  Every In[]/Out[] transcript in this book is produced by running the inputs
%%  through the real ./Mathilda binary at build time (see book/CONTEXT.md).
%% ===========================================================================
\documentclass[11pt,oneside]{memoir}

% --- encoding & fonts ------------------------------------------------------
\usepackage[T1]{fontenc}
\usepackage[utf8]{inputenc}
\usepackage{lmodern}
\usepackage{inconsolata}          % a readable monospace for the transcripts
\usepackage{microtype}

% --- mathematics -----------------------------------------------------------
\usepackage{amsmath,amssymb,amsthm,mathtools}

% --- graphics, tables, colour ---------------------------------------------
\usepackage{graphicx}
\usepackage{booktabs}
\usepackage{xcolor}

% --- page layout (memoir) --------------------------------------------------
\setlrmarginsandblock{1.1in}{1.1in}{*}
\setulmarginsandblock{1.1in}{1.2in}{*}
\checkandfixthelayout
% Custom chapter style: a large, LEFT-anchored number (so two-digit chapter
% numbers never overflow the page as veelo's fixed-width black bar did) with a
% thin rule, echoing the running-head rule below.
\makeatletter
\makechapterstyle{mathildachap}{%
  \setlength{\beforechapskip}{40pt}%
  \setlength{\afterchapskip}{26pt}%
  \renewcommand*{\chapnamefont}{\normalfont\large\scshape}%
  \renewcommand*{\chapnumfont}{\normalfont\fontsize{48}{48}\selectfont\bfseries}%
  \renewcommand*{\chaptitlefont}{\normalfont\Huge\bfseries}%
  \renewcommand*{\printchaptername}{\chapnamefont\@chapapp\space}%
  \renewcommand*{\chapternamenum}{}%
  \renewcommand*{\printchapternum}{\chapnumfont\thechapter}%
  \renewcommand*{\afterchapternum}{\par\nobreak\vskip 10pt\hrule height 0.4pt\vskip 16pt}%
  \renewcommand*{\printchaptertitle}[1]{\raggedright\chaptitlefont ##1}%
}
\makeatother
\chapterstyle{mathildachap}
\setsecnumdepth{subsection}
\settocdepth{section}

% --- running heads: mixed case (not ALL CAPS), with a thin rule under them ---
\nouppercaseheads
\makepagestyle{mathilda}
\makepsmarks{mathilda}{%
  \createmark{chapter}{both}{shownumber}{}{. }%
  \createplainmark{toc}{both}{\contentsname}%
  \createplainmark{bib}{both}{\bibname}%
  \createplainmark{index}{both}{\indexname}%
}
\makeevenhead{mathilda}{\small\thepage}{}{\small\itshape\leftmark}
\makeoddhead{mathilda}{\small\itshape\leftmark}{}{\small\thepage}
\makeevenfoot{mathilda}{}{}{}
\makeoddfoot{mathilda}{}{}{}
\makeheadrule{mathilda}{\textwidth}{0.4pt}
\pagestyle{mathilda}

% --- index -----------------------------------------------------------------
\makeindex

% --- bibliography ----------------------------------------------------------
\usepackage[backend=biber,style=numeric,sorting=nyt]{biblatex}
\addbibresource{references.bib}

% --- hyperref (late), then the book's own support package ------------------
\usepackage{hyperref}
\usepackage{mathilda}
\hypersetup{%
  colorlinks=true,
  linkcolor=mLink, citecolor=mLink, urlcolor=mLink, filecolor=mLink,
  pdftitle={A User's Manual for the Mathilda Computer Algebra System},
  pdfauthor={Sam Blake}}

% --- theorem environments --------------------------------------------------
\theoremstyle{plain}
\newtheorem{theorem}{Theorem}[chapter]
\newtheorem{proposition}[theorem]{Proposition}
\newtheorem{lemma}[theorem]{Lemma}
\newtheorem{corollary}[theorem]{Corollary}
\theoremstyle{definition}
\newtheorem{definition}[theorem]{Definition}
\newtheorem{example}[theorem]{Example}
\theoremstyle{remark}
\newtheorem*{remark}{Remark}

% ---------------------------------------------------------------------------
\begin{document}

\frontmatter
%% Title page.
\begin{titlingpage}
\centering
\vspace*{\fill}

{\Huge\bfseries The Mathilda Book\par}
\vspace{0.6cm}
{\Large A User's Manual for the\\[0.2em] Mathilda Computer Algebra System\par}

\vspace{1.4cm}
{\large\itshape The freely available computer algebra system\par}

\vspace{2.2cm}
{\large Sam Blake, PhD\par}

\vfill

{\small
Built and verified against Mathilda \texttt{\mathildaversion}.\\
Every transcript in this book was produced by the Mathilda binary itself.\par}

\vspace{0.6cm}
{\small GNU General Public License, version 3\par}
\vspace*{\fill}
\end{titlingpage}

%% Dedication page (the second page, immediately after the title page).
%% Dedicates both the Mathilda computer algebra system and this book to the
%% author's daughter.
\thispagestyle{empty}
\vspace*{0.3\textheight}
\begin{center}
{\Large\itshape
For my daughter, Matilda,\\[0.45em]
to whom the Mathilda computer algebra system\\[0.45em]
and this book are both dedicated.\par}
\end{center}
\vspace*{\fill}
\clearpage

%% Copyright / colophon-lite page (follows the title and dedication pages).
\thispagestyle{empty}
\noindent
\textit{A User's Manual for the Mathilda Computer Algebra System.}

\vspace{1em}
\noindent
Copyright \copyright\ 2026 Sam Blake.

\vspace{1em}
\noindent
Mathilda is free software: you can redistribute it and/or modify it under the
terms of the \textbf{GNU General Public License} as published by the Free
Software Foundation, either version~3 of the License, or (at your option) any
later version. Mathilda is distributed in the hope that it will be useful, but
\textsc{without any warranty}; without even the implied warranty of
\textsc{merchantability} or \textsc{fitness for a particular purpose}. See the
GNU General Public License for more details:
\url{https://www.gnu.org/licenses/gpl-3.0.html}.

\vspace{1em}
\noindent
This book documents Mathilda and shares its free-software spirit: you are
encouraged to read it, run every example yourself, and improve both the software
and this text.

\vspace{1em}
\noindent
Typeset in \LaTeX\ with the \textsf{memoir} class. The source for this book, its
example programs, and the tools that verify every transcript against the running
system live alongside Mathilda in the \texttt{book/} directory of the source
tree.

\vfill
\noindent
\textbf{Reproducibility.} The transcripts in this edition were generated by
Mathilda version \texttt{0.076}. Because the examples are re-run from source at
build time, a later build against a newer Mathilda will faithfully show whatever
that version prints---see \texttt{book/CONTEXT.md}.

\tableofcontents*
% \chapter{Preface}

This is a book about doing mathematics with a computer.

It has two ambitions at once. The first is to be a \emph{user's manual}: a
practical, example-driven guide to Mathilda, a computer algebra system you can
download, read, and run for free. The second is to be a \emph{textbook on
computational mathematics}---an account of how genuine, sometimes advanced,
mathematics is actually carried out inside a computer algebra system: how a
machine differentiates and integrates, factors an integer or a polynomial,
solves an equation exactly, sums a series in closed form, or races through a
million-element array without ever leaving machine precision.

Those two goals reinforce each other. You cannot really understand a computer
algebra system without seeing the mathematics it embodies, and you cannot
appreciate the mathematics without watching a real system do it. So this book is
built around transcripts of Mathilda at work, and it treats every one of them as
a small experiment worth explaining.

\section*{Every example is real}

There is a promise behind every transcript in this book:

\begin{quote}
Nothing you see printed as \texttt{Out[\,]} was typed by a human. Each example is
run through the Mathilda binary at the moment the book is built, and the output
is inserted automatically.
\end{quote}

This is not a stylistic flourish; it is a design decision baked into how the book
is produced. The inputs live in small program files; a build step feeds them to
Mathilda and captures exactly what it prints; and the typesetting includes that
captured text verbatim. There is, quite literally, nowhere for a hand-written
answer to hide. If Mathilda's behaviour changes, the next build tells the truth
about it. The mechanics are described in \texttt{book/CONTEXT.md} for the curious
and for anyone who wants to contribute.

One consequence is worth stating plainly: Mathilda's output form is not always
the form a textbook would choose. Terms may be reordered, an exponential may be
written \texttt{E\^{}x} rather than $e^{x}$, and an indefinite integral may differ
from a table by a constant. Everything printed is nonetheless mathematically
correct, and part of learning a computer algebra system is learning to read what
it says.

\section*{Links to the reference}

Mathilda has hundreds of built-in functions, each with its own reference page in
the online documentation. Whenever this book mentions a built-in for the first
time in a discussion---\B{Integrate}, \B{Solve}, \B{FactorInteger}, and so
on---the name is a live hyperlink to that reference page. The book teaches; the
reference pages enumerate every option and corner case. Read this for the story
and follow the links for the details.

\section*{How to read this book}

You do not need Mathematica, a licence, or a network connection to follow along;
you need only a copy of Mathilda, which you can build from source. Type the
\texttt{In[\,]} lines yourself and you will see the same \texttt{Out[\,]}. Later
chapters assume the vocabulary of earlier ones, but the mathematical chapters in
Part~II are largely independent---dip into calculus or number theory as your
interest takes you.

Welcome. Let us compute.
   % preface deferred until the book is further along

\mainmatter

\chapter{About Mathilda}
\label{ch:about}

\index{Mathilda!goals}%
Mathilda is a computer algebra system---a program\index{computer algebra system} that does mathematics
\emph{symbolically}\index{symbolic computation}, computing with $x$ and $\pi$ and $\sqrt{2}$ as exact objects
rather than as approximate decimals, while remaining fast enough to crunch a
million floating-point numbers when that is what you need. It is modelled closely
on the core architecture and evaluation semantics of the Mathematica\index{Mathematica}, and it
is written in C, and in its own language.

This chapter is about \emph{why} Mathilda exists, what it is built on, and why it
is written the way it is. The next chapter is practical---how to obtain, build, and
run the system on your own machine---and the mathematics begins in earnest in
Chapter~\ref{ch:intro}; here we set the stage.

\section{The mission}
\label{sec:mission}

The goal of Mathilda is simple to state and hard to achieve: \emph{to make the
most advanced, most efficient computer algebra system freely available to the
world.}

Each word in that sentence is load-bearing.

\emph{Advanced}, because a computer algebra system earns its keep at the frontier
---factoring polynomials over algebraic number fields, integrating in closed form
by the Risch algorithm\index{Risch algorithm}, solving Diophantine equations, computing with Gr\"obner\index{Grobner basis@Gr\"obner basis}
bases~\cite{gcl}. Anyone can add two fractions; the interesting question is how far
into real mathematics a system will carry you.

\emph{Efficient}, because a symbolic answer you cannot get before lunch is of
little use. Mathilda is built to be quick: exact arithmetic on arbitrary-precision\index{arbitrary-precision arithmetic}
integers, a bytecode compiler for numeric inner loops\index{bytecode compiler}, and transparently packed\index{array!packed}
arrays that let element-wise work run at the speed of machine code.

\emph{Freely available}, because mathematics belongs to everyone. The dominant
computer algebra systems are proprietary and expensive; their algorithms are
sealed. A student, a researcher in a country without site licences, or a curious
amateur should be able to compute at the highest level without asking anyone's
permission and without paying a toll.

\section{Free software, and why it matters}
\label{sec:license}

\index{license!GPLv3}%
Mathilda is released under the \textbf{GNU General Public License, version~3}
(GPLv3). You may run it for any purpose, read every line of its source, modify it,
and redistribute your modifications---provided you extend the same freedoms to
whomever you pass it on to. This is not an incidental legal choice; it follows
directly from the mission.

For a computer algebra system in particular, openness is a matter of
\emph{mathematical integrity}\index{mathematical integrity}. When a system tells you that
\[
  \int \frac{dx}{1+x^{2}} = \arctan x,
\]
you can check it by differentiating. But when it returns a page-long closed form
for a hard integral, or announces that a two-hundred-digit number is prime, you
are trusting the algorithm. With proprietary software that trust is an act of
faith; with free software it is, in principle, verifiable. The algorithm is right
there to be read, tested, and---when it is wrong---fixed. Scientific results
computed with a black box are not fully reproducible; results computed with an
open system are.

Openness also buys \emph{longevity}\index{longevity}. Proprietary formats and products are
abandoned; a licence server goes dark and a decade of work becomes unreadable.
Free software cannot be taken away. As long as one copy of the source survives, so
does the ability to build it, study it, and carry it forward. A system meant to
serve mathematics---which measures its lifetimes in centuries---should be built to
last, and only free software can make that promise.

\section{Standing on the shoulders of giants}
\label{sec:libraries}

Mathilda does not reinvent the deep numerical kernels that decades of careful
engineering have already perfected. Instead it builds on a small set of
best-in-class libraries, each the state of the art in its niche. You can see them
listed the moment you start the program, in the version banner:

\mtranscript{01-about/version}

Behind that one line stands an enormous amount of work by other people:

\begin{itemize}
  \item \textbf{GMP}~\cite{gmp}\index{GMP}, the GNU Multiple Precision Arithmetic Library, provides the
        arbitrary-precision integers on which all of Mathilda's exact arithmetic
        rests. When a machine integer would overflow, the value is promoted to a
        GMP ``bignum'' automatically and silently.
  \item \textbf{MPFR}~\cite{mpfr}\index{MPFR} extends that idea to real numbers, giving correctly rounded
        arbitrary-precision floating point---so \B{N}\texttt{[expr, 100]} can be
        trusted to a hundred digits.
  \item \textbf{FLINT}~\cite{flint}\index{FLINT} accelerates polynomial arithmetic, factoring, and the
        rigorous numerics behind many special functions.
  \item \textbf{LAPACK}\index{LAPACK} and \textbf{BLAS}\index{BLAS}
        power dense machine-precision linear algebra---the decompositions behind
        \B{Det}, \B{Inverse}, \B{Eigenvalues}, and
        \B{SingularValueDecomposition}.
  \item \textbf{GMP-ECM}\index{GMP-ECM} supplies the Elliptic Curve Method for the harder cases
        of integer factorization.
  \item \textbf{Raylib}\index{Raylib} renders the interactive graphics produced by \B{Plot} and
        \B{Show}, and \textbf{GNU~Readline}\index{GNU Readline} gives the interactive session its
        line editing and history.
\end{itemize}

The art is not in rewriting these libraries but in \emph{orchestrating} them: in
knowing when an exact computation should hand off to GMP, when a numeric one
should drop into LAPACK, and in making those transitions invisible so that you can
think about the mathematics instead of the plumbing. Much of this book is about
that orchestration.

\section{Why Mathilda is implemented in C}
\label{sec:why-c}

A computer algebra system could be written in almost any language. Mathilda is
written in C99\index{C99}, and the choice is deliberate.

\emph{Performance.} C imposes essentially no overhead between your program and the
machine. The libraries Mathilda depends on---GMP, MPFR, FLINT, LAPACK---are
themselves written in C and Fortran, and calling them from C means no marshalling,
no foreign-function boundary, no garbage-collector pauses in the middle of a tight
loop. When Mathilda packs an array of doubles and maps a function over it, the
inner loop is exactly the loop a C programmer would write by hand.

\emph{Language design and predictability.} C is small. Its semantics fit in one
head; its cost model is transparent---you can look at a line and know roughly what
the machine will do. For a system whose entire job is to be trustworthy about
computation, that transparency is a feature, not a limitation.

\emph{Longevity.} C is perhaps the most stable programming language in wide use.
Code written to the C standard decades ago still compiles and runs today, on
hardware that did not exist when it was written. A project that intends to be
useful for a very long time---see \S\ref{sec:license}---should be written in a
language that will still be understood, and still be compilable, a very long time
from now. C is that language.

\emph{An unfair advantage.} There is also a modern, practical reason. Large
language models are exceptionally strong\index{large language model} at writing and reasoning about C: it is
one of the most abundant languages in their training, its patterns are regular,
and its errors are mechanically checkable. Much of Mathilda is developed with such
models in the loop, and C's combination of ubiquity and simplicity makes that
collaboration unusually productive. A language that is easy for both humans and
machines to reason about is exactly the right substrate for building an ambitious
system quickly and carefully.

\section{The long-term vision of Mathilda}
\label{sec:vision}
\index{vision}%

The mission of \S\ref{sec:mission} sets a deliberately immodest bar, and it is worth
being explicit about how high. Mathilda's long-term goal is not merely to be \emph{a}
capable computer algebra system, but to be the \emph{best} system in the world at
every domain of computational mathematics---to gather, in a single free program, the
capabilities that today are scattered across a dozen specialised and often
proprietary tools, and to match or surpass the finest of each on its own ground.

That standard is measured against the best there is, domain by domain. To make it
concrete, here is some of what ``best in each domain'' is meant to mean:

\begin{itemize}
  \item \textbf{Symbolic integration.} A complete implementation of the
        Risch--Trager--Bronstein algorithm\index{Risch--Trager--Bronstein algorithm}
        for integration in closed form, rivalling the Scratchpad / AXIOM / FriCAS\index{FriCAS}
        lineage that pioneered it---so that when an elementary antiderivative exists
        Mathilda finds it, and when none exists Mathilda proves that too.
  \item \textbf{Algebraic geometry and commutative algebra.} Gr\"obner bases,
        ideals, varieties, and the wider machinery of computational algebra at a level
        that rivals Magma\index{Magma}, Singular\index{Singular}, and Macaulay2\index{Macaulay2}.
  \item \textbf{Number theory.} Factorization, primality proving, elliptic curves,
        and class-group computation to rival PARI/GP\index{PARI/GP}.
  \item \textbf{Finite fields and discrete computation.} Arithmetic and linear
        algebra over finite fields at the scale that coding theory and cryptography
        demand, rivalling Magma.
  \item \textbf{Tensor computation.} Manipulation of tensors and the objects of
        differential geometry that is as expressive and as efficient as Cadabra\index{Cadabra}.
  \item \textbf{Arbitrary-precision numerics.} Correctly rounded and rigorously
        bounded computation---ball arithmetic in the spirit of Arb\index{Arb} and mpmath\index{mpmath}---so a
        numerical answer comes with a guarantee, not merely a value.
  \item \textbf{Numerical computing.} Dense linear algebra, quadrature, and
        differential-equation solving at the level of LAPACK, SciPy\index{SciPy}, and Julia\index{Julia}.
  \item \textbf{Optimization.} Local and global optimization---constrained,
        unconstrained, and mixed-integer---to rival SciPy and NLopt\index{NLopt}.
  \item \textbf{Array numerics.} Element-wise and reduction operations on dense,
        packed numerical arrays that are at least as fast as NumPy\index{NumPy} on a single
        core---and, because Mathilda owns its memory layout and kernels end to end,
        that parallelise cleanly across many cores.
  \item \textbf{Compilation.} A bytecode compiler whose numeric throughput rivals
        Pipi, so that an inner loop\index{Pipi} written in Mathilda runs at the speed of compiled
        code rather than interpreted rewriting.
\end{itemize}

\noindent
---and so on, across special functions, polynomial factorization, graph theory, and
every other branch. In each the bar is the same: not ``adequate'' but ``as good as
the best there is,'' and the whole brought together under one roof, freely available.

This is an ambition, stated honestly, not a finished boast. Mathilda does not yet
match every one of those systems on every problem, and this book is candid about where
it falls short: throughout, we hold Mathilda's answers and its speed up against the
best available tools and say plainly where there is still ground to make up. The
vision is the direction of travel. What follows is an honest account of how far along
that road Mathilda has come---and, just as importantly, of \emph{how} it does what it
does, because the point of this book is not only to compute, but to understand the
computation.

\chapter{Compiling and Running Mathilda}
\label{ch:building}

Before you can compute anything, you have to get Mathilda onto your machine and
build it. For most newcomers this is the single largest hurdle---not the
mathematics, but the toolchain\index{building!prerequisites}---so this chapter treats it carefully and
completely. By the end you will have obtained the source, installed the handful of
libraries it needs, compiled the program, started an interactive session, and
learned your way around the read--eval--print loop. If you have built open-source C
programs before, much of this will be familiar and you can skim; if you have not,
follow it line by line and you will be fine.

Mathilda is distributed as \emph{source code}, not as a pre-built binary. That is a
deliberate consequence of the mission of Chapter~\ref{ch:about}: a system meant to
be read, checked, and carried forward is one you compile yourself, on your own
hardware, from source you can inspect. The good news is that the build is ordinary
and undramatic---a C project with a plain \texttt{makefile}, a few well-known
libraries, and no exotic build system to learn.

\section{Obtaining the source}
\label{sec:obtaining}
\index{building!obtaining the source}

Mathilda lives in a public Git repository. To get a copy, clone it and step
inside:

\begin{shell}
$ git clone https://github.com/stblake/Mathilda.git
$ cd Mathilda
\end{shell}

\noindent
Everything the rest of this chapter does happens from inside that top-level
\texttt{Mathilda/} directory. Two files there are worth knowing about from the
start: \texttt{README.md} is the quick-start and the authoritative list of
dependencies, and \texttt{SPEC.md} describes the system's architecture. You do not
need to read either to build the program, but they are where to turn when something
is unusual about your platform.

\section{The toolchain and its dependencies}
\label{sec:prerequisites}

Mathilda is written in C99 and leans on a small set of best-in-class numerical
libraries (introduced in \S\ref{sec:libraries}). Building it therefore needs a C
compiler and those libraries' development headers. Only three of the libraries are
required; the rest are optional and the build \emph{degrades gracefully}\index{graceful degradation} without
them---it detects what is present, switches on the corresponding features, and
switches off the ones whose libraries are missing, always producing a working
binary.

\subsection{A real GCC, not clang}
\label{sec:gcc-not-clang}
\index{GCC!required, not clang}

The one platform rule that trips people up: Mathilda must be built with a genuine
\textbf{GCC}, never Apple's or LLVM's \texttt{clang}. The build enforces this for
you. Its makefile auto-detects a real GCC---trying \texttt{gcc-16}, \texttt{gcc-15},
\ldots\ down to a plain \texttt{gcc}---and if the compiler it is about to use reports
itself as clang, it stops immediately with a message telling you how to fix it,
rather than silently producing an unsupported binary.%
\calloutnote{Under the hood}{The rule is not arbitrary. On macOS the plain
\texttt{gcc} command is a symlink to Apple's clang, which both miscompiles a few of
Mathilda's constructs and---more importantly---hides the strict-C99 portability
warnings that the project's Linux continuous-integration build treats as errors.
Building with real GCC on every platform keeps one standard of correctness. The fix
is always the same: install GCC from Homebrew.}
On Linux the plain \texttt{gcc} is already real GCC, so there is nothing to do. On
macOS, install one from Homebrew (see below); the makefile will then find and prefer
it automatically. If you ever need to point it at a specific compiler, pass it
explicitly, for example \texttt{make CC=gcc-16}.

\subsection{The libraries}
\label{sec:libraries-list}

\paragraph{Required.} These three must be present or the build cannot proceed:
\begin{itemize}
  \item \textbf{GMP}\index{GMP}, the GNU Multiple Precision Arithmetic Library---arbitrary-precision
        integers, the foundation of all of Mathilda's exact arithmetic.
  \item \textbf{MPFR}\index{MPFR}---correctly rounded arbitrary-precision reals, so that \B{N}
        can give you a hundred trustworthy digits. It is on by default; you can build
        without it, but almost no one should.
  \item \textbf{GNU~Readline}\index{GNU Readline}---the line editing, history, and command recall of the
        interactive session.
\end{itemize}

\paragraph{Optional.} Each of these unlocks a family of features when present and is
simply omitted when absent:
\begin{itemize}
  \item \textbf{FLINT}\index{FLINT} (version~3.0 or newer)---fast, rigorous polynomial arithmetic
        over algebraic number fields and rigorous numerics for many special
        functions. Distributions that ship only FLINT~2.x are detected as too old
        and the accelerated paths are switched off automatically.
  \item \textbf{GMP-ECM}\index{GMP-ECM}---the Elliptic Curve Method for the harder cases of integer
        factorization.
  \item \textbf{LAPACK}\index{LAPACK} and \textbf{BLAS}---fast machine-precision linear algebra. On
        macOS the build finds Apple's \textbf{Accelerate} framework and you need
        install nothing.
  \item \textbf{FFTW}---fast Fourier transforms; without it \B{Fourier} falls back to
        a slower $O(n^2)$ transform.
  \item \textbf{Raylib}\index{Raylib}---the interactive window that \B{Plot}, \B{Plot3D}, and \B{Show}
        draw into. Without it those functions print a text placeholder and return
        normally (see Chapter~\ref{ch:graphics}).
\end{itemize}

\noindent
POSIX threads, used to spread the packed-array kernels across your cores, are on by
default on macOS and Linux and need no extra package. \textbf{CMake} is needed only
to build the test suite (\S\ref{sec:tests}).

Every optional backend is governed by a build-time flag you can force on or off. The
defaults auto-detect, so you rarely need to touch them, but they are worth knowing:

\begin{table}[htbp]
  \centering
  \small
  \begin{tabular}{@{}llp{0.52\linewidth}@{}}
    \toprule
    \textbf{Flag} & \textbf{Default} & \textbf{What it enables} \\
    \midrule
    \mcode{USE\_MPFR}     & \mcode{1} & Arbitrary-precision reals: \mcode{N[expr, prec]}, precision tracking. \\
    \mcode{USE\_FLINT}    & \mcode{1} & Fast FLINT ($\geq 3.0$) kernels for number fields and rigorous numerics. \\
    \mcode{USE\_LAPACK}   & \mcode{1} & Fast machine-precision linear algebra (Accelerate on macOS). \\
    \mcode{USE\_ECM}      & \mcode{1} & Elliptic Curve Method factorization via GMP-ECM. \\
    \mcode{USE\_FFTW}     & \mcode{1} & FFTW-backed \mcode{Fourier}; naive $O(n^2)$ transform when off. \\
    \mcode{USE\_THREADS}  & \mcode{1} & POSIX-thread parallelism for the packed-array kernels. \\
    \mcode{USE\_GRAPHICS} & \mcode{1} & Interactive 2D/3D plot windows via Raylib. \\
    \bottomrule
  \end{tabular}
  \caption{Optional-backend flags. Each auto-detects; append e.g.\ \mcode{USE\_LAPACK=0}
  to a \mcode{make} command to force one off.}
  \label{tab:use-flags}
\end{table}

\subsection{Installing the dependencies}
\label{sec:installing-deps}
\index{building!installing dependencies}

On \textbf{Debian or Ubuntu}, install the three required libraries and as many of the
optional ones as you want:

\begin{shell}
# Required
$ sudo apt install libgmp-dev libmpfr-dev libreadline-dev

# Optional, each auto-detected by the build
$ sudo apt install libflint-dev          # FLINT >= 3.0 (Ubuntu 24.04+/Debian 12+)
$ sudo apt install libecm-dev            # GMP-ECM factorization
$ sudo apt install liblapacke-dev libopenblas-dev   # LAPACK/BLAS
$ sudo apt install libfftw3-dev          # FFTW
$ sudo apt install libraylib-dev         # interactive graphics
$ sudo apt install cmake                 # only for the test suite
\end{shell}

\noindent
On \textbf{Fedora or RHEL} the package names are \texttt{gmp-devel},
\texttt{mpfr-devel}, \texttt{readline-devel}, and then \texttt{flint-devel},
\texttt{gmp-ecm-devel}, \texttt{lapack-devel}/\texttt{openblas-devel},
\texttt{fftw-devel}, and \texttt{cmake}.

On \textbf{macOS} with Homebrew, install a real GCC (\S\ref{sec:gcc-not-clang}) and
the libraries; Accelerate is already part of the system, so you do not install
LAPACK:

\begin{shell}
$ brew install gcc gmp mpfr readline cmake
# Optional
$ brew install flint fftw raylib gmp-ecm
\end{shell}

\section{Building}
\label{sec:building}
\index{building!compiling}

With the dependencies in place, build the program with a single command. The
makefile auto-discovers every source file, configures the internal dependencies, and
links the executable with \texttt{-std=c99 -O3}:

\begin{shell}
$ make -j
\end{shell}

\noindent
The \texttt{-j} flag builds in parallel across your cores; on Linux you can be
explicit with \texttt{make -j\$(nproc)}. The first build compiles several hundred
source files and takes a little while; afterwards, \texttt{make} rebuilds only what
changed. When it finishes you will have an executable called \texttt{Mathilda} in the
current directory.

Want a leaner binary, or missing a library you do not care about? Turn any optional
backend off explicitly (Table~\ref{tab:use-flags}). For instance, to build without
MPFR and without the LAPACK acceleration:

\begin{shell}
$ make -j USE_MPFR=0 USE_LAPACK=0
\end{shell}

\paragraph{If the build stops.} The three failures new users hit, and their fixes:
\begin{itemize}
  \item \emph{``\,'gcc' is clang, not GCC\,''}---you are on macOS with only Apple's
        compiler. Run \texttt{brew install gcc} and build again; the makefile will
        find the real GCC on its own (\S\ref{sec:gcc-not-clang}).
  \item \emph{A missing header such as \texttt{gmp.h} or \texttt{readline/readline.h}}---%
        a required library's development package is not installed. Install it from
        \S\ref{sec:installing-deps} and rebuild.
  \item \emph{FLINT features quietly absent}---your distribution ships FLINT~2.x.
        The build detects this, falls back to \texttt{USE\_FLINT=0}, and still
        succeeds; install FLINT~$\geq 3.0$ from source or a backport if you want the
        accelerated paths.
\end{itemize}

\section{Your first session}
\label{sec:first-session}
\index{REPL!starting}

Start the interactive session by running the binary with no arguments:

\begin{shell}
$ ./Mathilda

Mathilda <version> - A small, open source computer algebra system.

This program is free, open source software and comes with ABSOLUTELY NO WARRANTY.

Press Return to evaluate. An open bracket, string or comment continues
on the next line; press Esc then Return to force evaluation.
Exit by evaluating Quit[] or CONTROL-C.

In[1]:=
\end{shell}

\noindent
The banner's first line names the version you just built and the libraries it was
compiled against; the same information is available at any time from
\B{\$Version}\index{version!\$Version}, and the version number alone from \B{\$VersionNumber}:

\begin{codepairs}
\pair{02-building/version/1}{The full banner: Mathilda's version and every library
it was built against---exactly the build you are running.}
\pair{02-building/version/2}{Just the number, as a machine-readable value you can
compare against in a script.}
\end{codepairs}

The \texttt{In[1]:=} is the \emph{prompt}\index{REPL!In[] and Out[]}: Mathilda is waiting for you to type an
expression. Type one, press Return, and it prints the result on an
\texttt{Out[1]=} line and prompts again. Here is the whole of a first session---exact
arithmetic, an exact symbolic value, and a machine-precision number on demand:

\begin{codepairs}
\pair{02-building/first-session/1}{Ordinary arithmetic. The answer is the exact
integer~$4$.}
\pair{02-building/first-session/2}{\B{Sin} of $\pi/6$ is returned \emph{exactly}, as
the rational $1/2$---not a decimal.}
\pair{02-building/first-session/3}{Ask for a number with \B{N} and you get one: here
a quick machine-precision value of $\pi$.}
\end{codepairs}

\noindent
That is the entire rhythm of the system: you type an expression, Mathilda evaluates
it to a fixed point, and it shows you the result. Everything else in this book is
variations on that loop.

\section{How the REPL works}
\label{sec:repl}
\index{REPL}

The interactive loop---read an expression, evaluate it, print it, repeat---is where
you will spend most of your time, so it is worth knowing its conveniences.

\subsection{Results persist, and you can name them}
\label{sec:session-state}

A session has \emph{memory}. Everything you define stays in force until you leave, so
you can name a result on one line and use it on the next:

\begin{codepairs}
\pair{02-building/session-state/1}{Assign with \mcode{=}. The assignment returns the
value, so you see it echoed.}
\pair{02-building/session-state/2}{\mcode{x} now stands for $1024$ everywhere.}
\pair{02-building/session-state/3}{And keeps standing for it until you reassign or
clear it.}
\end{codepairs}

Beyond your own names, the session numbers everything for you.
Each result is stored as \mcode{Out[n]}, matching the \texttt{Out[n]=} label beside
it, and each input as \mcode{In[n]}. The special symbol \mcode{\%} is shorthand for
the most recent result and \mcode{\%\%} for the one before it, so in an interactive
session \mcode{\% + 1} adds one to whatever was just printed, and \mcode{Out[3]}
recalls the third result exactly.%
\calloutnote{Under the hood}{\mcode{In[n]}, \mcode{Out[n]}, and \mcode{\%} are a
feature of the \emph{interactive} loop---the REPL records each result as it prints
it. A script run with \mcode{-file} (\S\ref{sec:scripts}) and the machine pipe used
by front-ends do not keep this history, so reach for a named variable, as above, for
anything you want to reuse in a script.}

A trailing semicolon suppresses the printed output. The expression is still
evaluated and its value still stored---you just do not see it, which is exactly what
you want for setup lines and large intermediate results:

\begin{codepairs}
\pair{02-building/suppression/1}{The semicolon suppresses the \texttt{Out[]} line;
the assignment still happens.}
\pair{02-building/suppression/2}{Proof that it did: \mcode{a} is $3$, so \mcode{a\^{}2}
is $9$.}
\end{codepairs}

\subsection{Editing, history, and multiline input}
\label{sec:editing}
\index{REPL!line editing and history}

Because the session is built on GNU~Readline, it behaves like a modern shell prompt.
You can move along the line with the arrow keys, edit in place, and---most
usefully---press the up arrow to recall previous inputs and re-run or tweak them. The
Emacs-style key bindings work too: \texttt{Ctrl-A} to the start of the line,
\texttt{Ctrl-E} to the end, \texttt{Ctrl-R} to search backwards through history.

For an expression too long or too structured to sit comfortably on one line, you need
do nothing special: just keep typing.\index{REPL!multiline input} Pressing Return
evaluates the buffer \emph{only} once it forms a complete
expression---every \mcode{(}, \mcode{[} and \mcode{\{} closed, and every string and
comment finished. Until then, Return simply opens a fresh line and waits, so a
bracket you have opened but not yet closed carries you onto the next line
automatically:

\begin{shell}
In[1]:= Table[i^2,
{i, 1, 4}]
Out[1]= {1, 4, 9, 16}
\end{shell}

\noindent
You close the final brace, press Return, and the whole thing runs as one
expression---the line breaks inside it are just whitespace, exactly as they are in a
script. Should you ever want to submit an \emph{incomplete} line as it stands---to
see the parser's error, say, rather than being carried onto a new line
forever---press \texttt{Esc} then \texttt{Return} (or \texttt{Alt-Return}), which
forces evaluation regardless.%
\calloutnote{Under the hood}{This in-place multiline editing needs a genuine
GNU~Readline. macOS ships a \texttt{libedit} stand-in under the same header name that
cannot do it, so the build prefers a Homebrew \texttt{readline} when one is present
and otherwise falls back to plain one-line-at-a-time entry.}

\subsection{Getting help}
\label{sec:help}
\index{REPL!getting help}

To see what a function does, type a question mark before its name---\mcode{?Sin},
say---and Mathilda prints its usage message: a one-line signature and a short
description, drawn from the function's own built-in documentation. Every builtin
carries one. This is the same text that titles the function's reference card; for
\B{Sin} it reads:

\usagebox{Sin}

To \emph{discover} functions rather than look one up, \B{Names} lists every symbol
matching a string pattern---a quick way to see what a fresh session already knows:

\begin{codepairs}
\pair{02-building/help/1}{Every built-in name beginning with \mcode{Bessel}. The
\mcode{*} is a wildcard.}
\end{codepairs}

\noindent
For the fuller machine-readable story about a symbol---its attributes and
definitions as well as its usage---there is also \B{Information}, of which
\mcode{?name} is the everyday shorthand.

\subsection{Leaving the session}
\label{sec:leaving}

To quit, evaluate \mcode{Quit[]} or press \texttt{Ctrl-D} (end-of-input) at the
prompt. \texttt{Ctrl-C} interrupts as well. Any of them returns you to your shell.

\section{Running scripts and other modes}
\label{sec:scripts}
\index{building!running scripts}

The REPL is one of three ways Mathilda can run. It is chosen automatically: if you
name a file, Mathilda runs it as a script; otherwise, if its input is an interactive
terminal, it starts the REPL you have just met; and if its input is a pipe, it speaks
a machine protocol instead (below).

\paragraph{Scripts.} Point Mathilda at a file of expressions with \texttt{-file} (a
bare filename means the same thing) and it evaluates each one in order and exits.
Crucially, \emph{nothing is echoed}: there are no \texttt{In[]}/\texttt{Out[]} lines,
so the output is exactly and only what the script chooses to \B{Print}. Given a file
\texttt{area.m}:

\begin{shell}
$ cat area.m
(* area.m -- a tiny script. Nothing is echoed; only what we Print appears. *)
r = 3;
Print["Area of a disk of radius ", r, " is ", N[Pi r^2]]

$ ./Mathilda -file area.m
Area of a disk of radius 3 is 28.2743
\end{shell}

\noindent
This is the mode to use for batch work, for reproducible computations kept under
version control, and for anything you want to pipe into another program. (It is also
how this very book is built: every transcript you see was produced by feeding its
inputs to the real binary.)

\paragraph{The command line.} A few options round out the picture; \texttt{--help}
lists them all:

\begin{shell}
$ ./Mathilda --help
Mathilda <version> - a small, open source computer algebra system.

Usage: ./Mathilda [options] [file]

  -file <path>    evaluate every expression in <path>, then exit
  -h, --help      show this message and exit
  -v, --version   print version information and exit

A bare <file> argument is equivalent to -file <file>. With no file,
Mathilda starts the interactive REPL when stdin is a terminal, and
otherwise speaks the NDJSON pipe protocol on stdio.
\end{shell}

\paragraph{The pipe protocol.} That last sentence describes the third mode. When
Mathilda's input is not a terminal and no script was named, it speaks a line-oriented
JSON protocol on its standard input and output, designed for a front-end---a
notebook, an editor plugin, a web interface---to drive it programmatically.%
\calloutnote{Under the hood}{The protocol is \emph{NDJSON}: one JSON request per
line in, one or more JSON responses per line out, each result carrying both a text
form and a \LaTeX{} form so a graphical front-end can render it. You will not use
this mode by hand, but it is what lets Mathilda sit behind a richer interface than a
terminal.}

\section{Building and running the tests}
\label{sec:tests}
\index{building!test suite}

Mathilda ships a large suite of C-based unit tests. You do not need them to use the
system, but running them is the surest way to confirm your build is sound---and if
you go on to modify Mathilda (Chapter~\ref{ch:contributing}), they are how you check
you have not broken anything. The suite is built with CMake, which auto-detects the
same optional backends as the main build:

\begin{shell}
$ cd tests
$ mkdir -p build && cd build
$ cmake ..
$ make -j

# Run every test binary
$ for t in *_tests; do ./$t; done
\end{shell}

\bigskip
\noindent
That is everything you need to obtain, build, and run Mathilda, and to find your way
around a session. With a working binary in front of you and the prompt waiting, we
can turn to the mathematics. The next chapter is a whirlwind tour of what the system
can do---arithmetic, number theory, algebra, calculus, and more---so that the
detailed chapters afterwards feel like a return to places you have already seen.

\chapter{A Brief Introduction to Mathilda}
\label{ch:intro}

This chapter is a whirlwind tour. Rather than explain any one thing thoroughly---%
that is the job of Chapter~\ref{ch:mathematics} and beyond---it moves quickly across
the landscape, showing a small, concrete example of each of the things Mathilda can
do: arithmetic and numerics, number theory, algebra, calculus, numerical calculus,
linear algebra, and programming. The aim is to give you a feel for the system and a sense of how much
ground it covers, so that the detailed chapters later feel like a return to places
you have already visited.

Everything below is real: each \texttt{In[\,]} was fed to the Mathilda binary\index{In[] and Out[]} as this
book was built and each \texttt{Out[\,]} is exactly what came back. Type them yourself
and you will see the same results. Throughout, each input and its result stand on the
left, with a few words about that line alongside on the right. This chapter is a
starting point and will grow over time; think of it as the first, shortest pass over
material the rest of the book returns to in depth.

\section{Numbers, exact and approximate}
\label{sec:intro-numbers}

Mathilda distinguishes sharply between \emph{exact} numbers\index{number!exact vs approximate} and \emph{approximate}
ones, and it knows how they combine. Watch what happens as we add numbers of different
kinds:

\begin{codepairs}
\pair{03-introduction/numbers/1}{Two integers add to an integer.}
\pair{03-introduction/numbers/2}{An integer plus a rational\index{number!rational} stays \emph{exact}: the
fraction $7/3$, not a decimal.}
\pair{03-introduction/numbers/3}{A rational plus a \emph{real} gives a real\index{number!real}. One
approximate number makes the whole result approximate---a rule called
\emph{contagion}\index{numeric contagion}.}
\pair{03-introduction/numbers/4}{Integers grow without bound; $2^{64}$ is kept
exactly, far beyond what a single machine word can hold.}
\pair{03-introduction/numbers/5}{To get a decimal you ask, with \B{N}: here a quick
machine-precision value of $\sqrt2$, six figures by default.}
\pair{03-introduction/numbers/6}{Ask \B{N} for more digits and it obliges: here is
$\sqrt2$ to thirty places.}
\pair{03-introduction/numbers/7}{And here is $\pi$ to fifty digits, computed to
arbitrary precision.}
\end{codepairs}

Section~\ref{sec:arithmetic} takes up how these three worlds---machine numbers, exact
integers and rationals, and arbitrary-precision reals---fit together.

\section{Elementary functions}
\label{sec:intro-elementary}

Mathilda knows the elementary functions\index{elementary functions}---powers and roots, the exponential and
logarithm, and the trigonometric functions and their inverses---and treats them the
same way it treats numbers: exactly when it can, symbolically when it must, and
numerically when you ask. At arguments where a simple closed form exists, it returns
that exact value:

\begin{codepairs}
\pair{03-introduction/elementary-exact/1}{A perfect square gives an integer.}
\pair{03-introduction/elementary-exact/2}{\B{Sin} of $\pi/6$ is the rational $1/2$.}
\pair{03-introduction/elementary-exact/3}{\B{Cos} of $\pi/4$ comes back as an exact
surd.}
\pair{03-introduction/elementary-exact/4}{So does \B{Tan} of $\pi/3$.}
\pair{03-introduction/elementary-exact/5}{\B{Log} and \B{Exp} undo each other, so
$\log e^{5}=5$.}
\pair{03-introduction/elementary-exact/6}{\B{ArcTan} of $1$ is the exact angle
$\pi/4$.}
\pair{03-introduction/elementary-exact/7}{At an exact argument with \emph{no} closed\index{closed form}
form, \B{Sin}$[1]$ is held as \texttt{Sin[1]}---not silently turned into a decimal.}
\end{codepairs}

When you do want a number, you get one either by supplying an approximate argument or
by applying \B{N}. By default the result is \emph{machine precision}\index{machine precision}---the fast,
fixed-size floating point of the hardware---but \B{N} will also work to
\emph{arbitrary} precision\index{precision!arbitrary}:

\begin{codepairs}
\pair{03-introduction/elementary-numeric/1}{An approximate argument \texttt{1.0}
forces a machine-precision result.}
\pair{03-introduction/elementary-numeric/2}{Applying \B{N} to the exact
\texttt{Sin[1]} gives the same value.}
\pair{03-introduction/elementary-numeric/3}{Ask for forty digits and you get $\sin 1$
to that precision.}
\pair{03-introduction/elementary-numeric/4}{Here is the constant $e$ to forty digits.}
\pair{03-introduction/elementary-numeric/5}{And here is $\ln 2$ to forty digits. The
elementary functions return in Section~\ref{sec:calculus}.}
\end{codepairs}

\section{Special functions}
\label{sec:intro-special}

Beyond the elementary functions lies a second tier---the \emph{special functions}\index{special functions} of
analysis and mathematical physics: the gamma and zeta functions, the error function,
the Bessel and Airy functions, and many more. Mathilda handles them with exactly the
same discipline, returning exact values wherever they exist:

\begin{codepairs}
\pair{03-introduction/special-exact/1}{\B{Gamma} extends the factorial: $\Gamma(5)=4!
=24$.}
\pair{03-introduction/special-exact/2}{At the half-integers it takes the celebrated
value $\Gamma(1/2)=\sqrt\pi$.}
\pair{03-introduction/special-exact/3}{\B{Zeta} gives Euler's $\zeta(2)=\pi^{2}/6$.}
\pair{03-introduction/special-exact/4}{By analytic continuation\index{analytic continuation}, the notorious
$\zeta(-1)=-\tfrac{1}{12}$.}
\pair{03-introduction/special-exact/5}{\B{Beta} reduces to a rational.}
\pair{03-introduction/special-exact/6}{\B{PolyGamma}, the digamma at $1$, is exactly
$-\gamma$, the negative Euler--Mascheroni constant\index{Euler--Mascheroni constant} \B{EulerGamma}.}
\pair{03-introduction/special-exact/7}{Even \B{AiryAi} at the origin has a closed
form, in terms of \B{Gamma}.}
\end{codepairs}

Where no closed form exists, these functions are computed numerically, to machine or
arbitrary precision just like the elementary ones:

\begin{codepairs}
\pair{03-introduction/special-numeric/1}{$\zeta(3)$, Ap\'ery's constant, to machine
precision.}
\pair{03-introduction/special-numeric/2}{The same constant to forty digits.}
\pair{03-introduction/special-numeric/3}{\B{Erf} at an approximate argument.}
\pair{03-introduction/special-numeric/4}{\B{BesselJ} likewise.}
\pair{03-introduction/special-numeric/5}{$\Gamma(1/2)$ to thirty digits, agreeing with
$\sqrt\pi$. Section~\ref{sec:specialfns} is devoted to these functions.}
\end{codepairs}

\section{Number theory}
\label{sec:intro-numbertheory}

Whole numbers have a rich structure, and Mathilda knows a great deal of it. Primes,
factorizations, divisors, greatest common divisors\index{greatest common divisor}, and modular arithmetic\index{modular arithmetic} are all a
single call away.

\begin{codepairs}
\pair{03-introduction/number-theory/1}{\B{Prime} gives the hundredth prime.}
\pair{03-introduction/number-theory/2}{\B{FactorInteger} breaks $360$ into prime
powers $2^{3}\cdot 3^{2}\cdot 5$.}
\pair{03-introduction/number-theory/3}{\B{Divisors} lists every divisor of $36$.}
\pair{03-introduction/number-theory/4}{\B{GCD} is the greatest common divisor.}
\pair{03-introduction/number-theory/5}{\B{ExtendedGCD} adds B\'ezout's coefficients\index{Bezout's coefficients@B\'ezout's coefficients}:
$12 = 1\cdot 48 + (-1)\cdot 36$.}
\pair{03-introduction/number-theory/6}{\B{PowerMod} finds $7^{100}\bmod 13$ without
ever forming $7^{100}$.}
\pair{03-introduction/number-theory/7}{\B{LCM} is the least common multiple;
$\mathrm{lcm}(12,18,20)=180$.}
\pair{03-introduction/number-theory/8}{\B{EulerPhi} counts the integers from $1$ to $36$
that are coprime to it: $\varphi(36)=12$.}
\pair{03-introduction/number-theory/9}{\B{MultiplicativeOrder} gives the least $k$ with
$2^{k}\equiv 1\pmod{11}$: here that order is $10$.}
\pair{03-introduction/number-theory/10}{\B{PrimitiveRoot} returns a generator of the
units modulo~$11$; the least is $2$, which is exactly why its order above is the full
$\varphi(11)=10$.}
\pair{03-introduction/number-theory/11}{\B{Fibonacci} returns the fiftieth Fibonacci
number exactly. Section~\ref{sec:numbertheory} takes up this material.}
\end{codepairs}

\section{Algebra}
\label{sec:intro-algebra}

Algebra is where a computer algebra system\index{computer algebra system} earns its name: it manipulates
expressions containing symbols\index{expression!symbolic}, not just numbers. Two everyday operations are
multiplying out and its partial inverse, collecting like terms\index{collecting like terms}.

\begin{codepairs}
\pair{03-introduction/algebra-expand/1}{\B{Expand} multiplies out a power.}
\pair{03-introduction/algebra-expand/2}{It works in several variables too.}
\pair{03-introduction/algebra-expand/3}{\B{Collect} gathers terms by the powers of a
chosen variable---here $x$, pulling the common factor out of the linear term.}
\end{codepairs}

The deeper operation is \B{Factor}, which writes a polynomial as a product of\index{polynomial!factorization}
irreducible pieces. What counts as irreducible\index{polynomial!irreducible} depends on the number system you allow,
and Mathilda lets you choose:

\begin{codepairs}
\pair{03-introduction/algebra-factor/1}{Over the rationals, $x^{6}-1$ splits into four
factors.}
\pair{03-introduction/algebra-factor/2}{But $x^{2}+1$ will not factor at all---it is
irreducible over~$\mathbb{Q}$.}
\pair{03-introduction/algebra-factor/3}{Allow the imaginary unit, with
\texttt{Extension}, and it becomes\index{field extension} $(x-i)(x+i)$.}
\pair{03-introduction/algebra-factor/4}{Allow $\sqrt2$ and $x^{4}+1$ splits into two
real quadratics. Extensions return in Section~\ref{sec:algebra}.}
\end{codepairs}

Closely related to factoring is the polynomial greatest common divisor\index{polynomial!greatest common divisor}. Just as \B{GCD}
finds the largest integer dividing two whole numbers, \B{PolynomialGCD} finds the
highest-degree polynomial dividing two polynomials---their shared factor:

\begin{codepairs}
\pair{03-introduction/algebra-gcd/1}{These two polynomials share the factor $x-1$.}
\pair{03-introduction/algebra-gcd/2}{These share $x^{2}-1$, found directly, without
factoring either polynomial.}
\end{codepairs}

That common factor is exactly what simplifies a \emph{rational function}\index{rational function}, a ratio of
polynomials. Mathilda provides the three everyday operations:

\begin{codepairs}
\pair{03-introduction/algebra-rational/1}{\B{Cancel} divides out the common factor
(the $x-1$ above), reducing the fraction to lowest terms.}
\pair{03-introduction/algebra-rational/2}{\B{Together} combines a sum of fractions
over one common denominator.}
\pair{03-introduction/algebra-rational/3}{\B{Apart} does the reverse---the partial-%
fraction decomposition you reach for\index{partial fraction decomposition} before integrating. More in
Section~\ref{sec:algebra}.}
\end{codepairs}

Beyond a single polynomial lies a whole \emph{system} of them, and the tool for those is
the \B{GroebnerBasis}: a canonical set of generators for the ideal\index{ideal} a system defines. It
is the multivariate, nonlinear analogue of Gaussian elimination\index{Gaussian elimination}---it rewrites a system
into an equivalent triangular one that is far easier to solve, and it can eliminate
chosen variables outright.

\begin{codepairs}
\pair{03-introduction/algebra-groebner/1}{The unit circle meets the parabola $y=x^{2}$.
\B{GroebnerBasis} triangulates the system: the first polynomial involves only $y$---a
quadratic to solve---and the second then recovers $x$ from $y$.}
\pair{03-introduction/algebra-groebner/2}{Naming a variable to eliminate---here the
parameter $t$ of the curve $(t^{2},t^{3})$---collapses the system to the curve's implicit
equation $x^{3}=y^{2}$, the cuspidal cubic. Section~\ref{sec:algebra} returns to
Gr\"obner bases.}
\end{codepairs}

\section{Solving equations}
\label{sec:intro-solve}

Solving equations is among the oldest tasks in mathematics, and Mathilda does it both
exactly and numerically.

\subsection{Exact solutions of equations}

\B{Solve} finds \emph{exact} solutions---as integers, radicals, or general closed
forms---for a single equation or a whole system:

\begin{codepairs}
\pair{03-introduction/solve-exact/1}{A quadratic with two integer roots. Note
\texttt{==}, which builds an equation\index{equation vs assignment}, versus \texttt{=}, which assigns.}
\pair{03-introduction/solve-exact/2}{$x^{2}=2$: the two exact surds $\pm\sqrt2$.}
\pair{03-introduction/solve-exact/3}{A linear system in two unknowns\index{system of equations!linear}, uniquely solved.}
\pair{03-introduction/solve-exact/4}{A \emph{nonlinear} system---a circle meeting the
line $y=x$---with both intersection points, exactly.}
\pair{03-introduction/solve-exact/5}{A cubic with no rational roots: the exact solutions
come back as \B{Root} objects, each an implicit algebraic number\index{algebraic number}---the $k$th root of the
given polynomial---rather than an unwieldy nest of radicals.}
\pair{03-introduction/solve-exact/6}{A \B{Root} object is still an exact number, so \B{N}
evaluates it to as many digits as you ask---here to machine precision.}
\end{codepairs}

\subsection{Numerical solutions of equations}

When an exact solution is unavailable---the general quintic has no formula in radicals,
by Abel's theorem\index{Abel's theorem}---or simply not wanted, the numerical solvers take over. \B{NRoots}
and \B{NSolve} return every root of a polynomial equation or system as machine numbers,
while \B{FindRoot} follows Newton's method\index{Newton's method} from a starting point and handles
transcendental problems too:

\begin{codepairs}
\pair{03-introduction/solve-numerical/1}{\B{NRoots} finds all five roots of $x^{5}-x-1$:
one real and two complex-conjugate pairs, none expressible in radicals.}
\pair{03-introduction/solve-numerical/2}{\B{NSolve} handles systems. Eliminating $y$ here
leaves that very same quintic, so the five solutions carry its unsolvable roots in their
$x$-coordinate.}
\pair{03-introduction/solve-numerical/3}{\B{FindRoot}, started near $x=1$, converges to
the root of $\cos x = x$, the ``Dottie number''\index{Dottie number}---a transcendental equation with no
closed form.}
\pair{03-introduction/solve-numerical/4}{\B{FindRoot} takes systems as well: a
transcendental pair in $x$ and $y$, solved from the starting point $(1,1)$.}
\end{codepairs}

\subsection{Equations with a continuum of solutions}

\B{Solve} expects a finite list of solutions to name. But an equation built from
absolute values or nested radicals can hold across a whole interval, and then what
is wanted is again a \emph{description} of the solution set\index{solution
set}\index{equation!continuum of solutions} rather than a list of points.
\B{Reduce}---the general real-equation engine taken up in the next
section---provides it, solving over the reals and returning the exact region on
which the equation holds:

\begin{codepairs}
\pair{03-introduction/solve-region/1}{Each nested radical\index{radical!nested} is a
perfect square in disguise---$x+1-2\sqrt x=(\sqrt x-1)^{2}$ and
$x+16-8\sqrt x=(\sqrt x-4)^{2}$---so the left side is $|\sqrt x-1|+|\sqrt x-4|$,
which equals~$3$ across the entire interval $1\le x\le16$.}
\pair{03-introduction/solve-region/2}{A sum of absolute values\index{absolute value}
over~$x$: the equation is satisfied not at isolated points but along the whole ray
$x\ge3$.}
\pair{03-introduction/solve-region/3}{Nested absolute values, even in~$x$: the
solution set is a \emph{union} of two intervals, one mirrored across the origin from
the other.}
\end{codepairs}

Exact solving belongs to the algebra of Section~\ref{sec:algebra}; the numerical
solvers to Section~\ref{sec:numcalc}.

\section{Solving inequalities}
\label{sec:intro-inequalities}

An equation has a handful of solutions to name; an inequality\index{inequality} usually
has infinitely many, so what you want back is not a list but a \emph{description} of the
solution set\index{solution set}. \B{Reduce} provides it: given one or more polynomial,
rational, or radical inequalities it returns an equivalent, exhaustive statement---an
interval, a union of intervals, or a region of the plane---computed by \emph{cylindrical
algebraic
decomposition}\index{cylindrical algebraic decomposition}, a complete decision
procedure\index{decision procedure} for systems of polynomial equations and inequalities
over the reals.

\begin{codepairs}
\pair{03-introduction/inequalities/1}{A quadratic inequality: the solution set is the
closed interval where the parabola lies on or below the axis.}
\pair{03-introduction/inequalities/2}{A quartic can cut the line into several
pieces---here a union of two open intervals.}
\pair{03-introduction/inequalities/3}{A rational inequality. The pole at $x=-2$ is
excluded automatically---\B{Reduce} keeps track of where the expression is undefined.}
\pair{03-introduction/inequalities/4}{The endpoints need not be rational: here they are
the exact surds $\pm\sqrt2$.}
\pair{03-introduction/inequalities/5}{A rational inequality whose numerator carries the
denominator's factors---$x^{3}-x=x(x^{2}-1)$---so after cancellation it collapses to the
plain linear condition $x\le2$.}
\pair{03-introduction/inequalities/6}{A \emph{radical}
inequality\index{inequality!radical}, not polynomial at all. The square root's domain
$x\ge-4$ meets the single crossing at $x=5$, where $\sqrt{x+4}=x-2$, to give
$-4\le x<5$.}
\end{codepairs}

Because it \emph{decides} the theory rather than merely searching for solutions,
\B{Reduce} settles the universal questions too\calloutnote{Under the hood}{That the real
numbers admit \emph{any} such decision procedure is Tarski's theorem; cylindrical
algebraic decomposition is the algorithm that makes it practical, and it is what lets
\B{Reduce} answer a question about \emph{all} real $x$ with a definite \texttt{True} or
\texttt{False}.}, returning \texttt{True} or \texttt{False} outright when the solution
set turns out to be everything, or nothing:

\begin{codepairs}
\pair{03-introduction/inequalities/7}{$x^{2}+1>0$ holds for every real $x$---a one-line
proof that the expression is positive-definite.}
\pair{03-introduction/inequalities/8}{No real number makes $x^{2}+x+1$ negative (its
discriminant is negative), so the answer is \texttt{False}.}
\pair{03-introduction/inequalities/9}{The arithmetic--geometric-mean
inequality\index{AM--GM inequality} $x+1/x\ge2$: \B{Reduce} confirms it across the
\emph{whole} positive ray, the constraint collapsing to just $x>0$.}
\pair{03-introduction/inequalities/10}{With two variables the solution set becomes a
region of the plane---the open unit disk, described as $x$ ranging over $(-1,1)$ with $y$
between the two branches of its bounding circle.}
\pair{03-introduction/inequalities/11}{Absolute values\index{absolute value} in two
variables: the open \emph{diamond} $|x|+|y|<1$, the unit ball of the $\ell^{1}$
norm\index{norm!$\ell^1$}. \B{Reduce} case-splits each \B{Abs} on the sign of its
argument, so the region falls out as its four quadrant sectors.}
\end{codepairs}

\section{Simplification}
\label{sec:intro-simplify}

Expanding, factoring, and cancelling each drive an expression toward a \emph{particular}
form. Often, though, you simply want the \emph{simplest} form, whatever that turns out
to be. \B{Simplify} searches for it\index{simplification}, trying a battery of algebraic and trigonometric
transformations and keeping the smallest result:

\begin{codepairs}
\pair{03-introduction/simplification/1}{The Pythagorean identity\index{Pythagorean identity} collapses to~$1$.}
\pair{03-introduction/simplification/2}{A rational function reduces.}
\pair{03-introduction/simplification/3}{A double-angle identity is recognised:
$\cos^{2}x-\sin^{2}x=\cos 2x$.}
\pair{03-introduction/simplification/4}{And a nested radical
\emph{denests}\index{radical!denesting}: $\sqrt{3+2\sqrt2}=1+\sqrt2$.}
\end{codepairs}

Some simplifications hold only under conditions, and Mathilda will not assume what you
have not told it. Because $\sqrt{x^{2}}$ equals $|x|$, not $x$:

\begin{codepairs}
\pair{03-introduction/simplification-assumptions/1}{With nothing known about $x$,
$\sqrt{x^{2}}$ is left exactly as written.}
\pair{03-introduction/simplification-assumptions/2}{Told that $x>0$, it simplifies to
$x$.}
\end{codepairs}

An assumption can be more interesting\index{assumptions} than a sign condition---sometimes what matters is
the \emph{kind} of number involved:

\begin{codepairs}
\pair{03-introduction/simplification-integer/1}{For a general~$n$, $\sin(n\pi)$ is left
untouched.}
\pair{03-introduction/simplification-integer/2}{Told that $n$ is an integer, with
\B{Element}, Mathilda proves it must vanish. More in Section~\ref{sec:algebra}.}
\end{codepairs}

\section{Calculus}
\label{sec:intro-calculus}

Mathilda does symbolic calculus\index{symbolic calculus}: differentiation, limits, and integration on
expressions, returning exact results. Differentiation with \B{D} applies the rules of
calculus recursively, and by naming a second variable it takes partial derivatives\index{partial derivative}:

\begin{codepairs}
\pair{03-introduction/calculus-diff/1}{The power rule, applied term by term.}
\pair{03-introduction/calculus-diff/2}{The product and chain rules, automatically.}
\pair{03-introduction/calculus-diff/3}{Symbolic exponents are fine too.}
\pair{03-introduction/calculus-diff/4}{A \emph{partial} derivative in $x$, treating
$y$ as a constant.}
\pair{03-introduction/calculus-diff/5}{The same expression, now differentiated in
$y$.}
\end{codepairs}

Limits are computed with\index{limit} \B{Limit}, including the classic ones that define the
derivative and the number~$e$:

\begin{codepairs}
\pair{03-introduction/calculus-limits/1}{The limit at the heart of the derivative of
sine.}
\pair{03-introduction/calculus-limits/2}{One of the definitions of~$e$.}
\pair{03-introduction/calculus-limits/3}{The leading behaviour of $1-\cos x$ near the
origin.}
\end{codepairs}

Near a point a function is captured by its power series\index{power series}, and \B{Series} computes it to any
order, returning the polynomial part together with an $O[x]^{n}$ term that records where
the expansion was truncated:

\begin{codepairs}
\pair{03-introduction/calculus-series/1}{The exponential series, to sixth order.}
\pair{03-introduction/calculus-series/2}{$\arctan$: the alternating odd reciprocals whose
value at $x=1$ gives Leibniz's series for $\pi/4$.}
\pair{03-introduction/calculus-series/3}{At a pole the expansion is a \emph{Laurent}\index{Laurent series}
series---here $\cot x$ begins with a $1/x$ term.}
\pair{03-introduction/calculus-series/4}{With a \emph{symbolic} exponent the coefficients
come out as the generalized binomial coefficients.}
\pair{03-introduction/calculus-series/5}{\B{Normal} drops the $O[\,]$ term, turning the
\texttt{SeriesData} object\index{SeriesData} a series is into an ordinary polynomial.}
\end{codepairs}

Integration, the harder inverse of differentiation\index{integration}, is done with \B{Integrate}---%
indefinite integrals as symbolic antiderivatives, and definite integrals\index{definite integral}, including
improper ones:

\begin{codepairs}
\pair{03-introduction/calculus-integrate/1}{An indefinite integral by parts.}
\pair{03-introduction/calculus-integrate/2}{Another, giving an inverse tangent.}
\pair{03-introduction/calculus-integrate/3}{A definite integral over $[0,1]$.}
\pair{03-introduction/calculus-integrate/4}{The Gaussian integral over the whole real\index{Gaussian integral}
line: $\int_{-\infty}^{\infty} e^{-x^{2}}\,dx=\sqrt\pi$.}
\end{codepairs}

\B{Integrate} is not confined to one variable. Give it several iterators and it
computes a \emph{multiple} integral\index{multiple integral}, from the inside out, so the inner limits may
depend on an outer variable:

\begin{codepairs}
\pair{03-introduction/calculus-integrate-multi/1}{$x y$ integrated over the unit
square.}
\pair{03-introduction/calculus-integrate-multi/2}{Over the triangle $0\le y\le x\le 1$,
whose inner bound is $x$ itself.}
\end{codepairs}

And when the limits are \emph{complex}, \B{Integrate} evaluates a \emph{contour\index{contour integral}
integral} along the straight segments joining the points you list:

\begin{codepairs}
\pair{03-introduction/calculus-integrate-contour/1}{A line integral of $z^{2}$ along
the segment from $0$ to $1+i$.}
\pair{03-introduction/calculus-integrate-contour/2}{A closed loop around the origin:
$\oint \frac{dz}{z}=2\pi i$, the cornerstone of the residue calculus\index{residue calculus}.
Section~\ref{sec:calculus} has the full story.}
\end{codepairs}

\section{Numerical calculus}
\label{sec:intro-numcalc}

Not every problem has a closed form, and even when it does a number is sometimes all
you want. The \texttt{N}-prefixed family answers\index{numerical calculus} these questions numerically and
adaptively---to machine precision by default, or to any arbitrary precision you request
with the \texttt{WorkingPrecision} option\index{WorkingPrecision option}:

\begin{codepairs}
\pair{03-introduction/numerical-calculus/1}{\B{ND} takes a numerical derivative---here
of $\sin$ at $1$, giving $\cos 1$.}
\pair{03-introduction/numerical-calculus/2}{\B{NIntegrate} evaluates the Gaussian
integral numerically ($\approx\sqrt\pi$).}
\pair{03-introduction/numerical-calculus/3}{\B{NSum} sums a series numerically.}
\end{codepairs}

Differential equations are solved numerically\index{differential equation} by \B{NDSolve}. Take a damped harmonic
oscillator, $y'' + \tfrac{3}{10}\,y' + 4y = 0$ with $y(0)=1$ and $y'(0)=0$: the solution
comes back as an \B{InterpolatingFunction}, and \B{Plot} takes that
directly,\footnote{That plot is genuinely the \texttt{Out[2]} of the \B{Plot} command,
drawn by the running system. To keep it, \B{Export} writes any such graphic to PDF, PNG,
or JPEG (\texttt{Export["decay.pdf", plot]})---indeed every figure in this book is
produced exactly this way.} so the motion---a decaying oscillation settling to rest---is
one further command away:

\plotcell{03-introduction/ndsolve-solution}

Finally, \B{NLimit} evaluates a limit numerically, invaluable when a symbolic limit is
hard to obtain:

\begin{codepairs}
\pair{03-introduction/numerical-calculus-nlimit/1}{The familiar limit that
defines~$e$.}
\pair{03-introduction/numerical-calculus-nlimit/2}{More striking: $\zeta(s)$ (the
\B{Zeta} function) and $1/(s-1)$ both diverge at $s=1$, yet their difference tends to
the Euler--Mascheroni constant $\gamma$.}
\pair{03-introduction/numerical-calculus-nlimit/3}{A slower, non-obvious limit,
recovered numerically.}
\end{codepairs}

Section~\ref{sec:numcalc} covers numerical differentiation, integration, summation,
differential equations, and limits.

\section{Linear algebra}
\label{sec:intro-linalg}

Vectors and matrices are not a new kind of object in Mathilda---a
vector\index{vector} is just a list, and a matrix\index{matrix} a list of its
rows---so everything you already know about lists carries over unchanged. What
linear algebra adds is a vocabulary of operations on them, beginning with the
\emph{dot product}\index{dot product}, written with an ordinary period. The single
operator \texttt{.}\ (\B{Dot}) serves throughout: it contracts two vectors to a
scalar, applies a matrix to a vector, and multiplies two matrices.

\begin{codepairs}
\pair{03-introduction/linear-algebra-vectors/1}{The dot product of two vectors is a
single number.}
\pair{03-introduction/linear-algebra-vectors/2}{\B{Norm} is the Euclidean
length\index{norm!Euclidean} $\sqrt{v\cdot v}$, kept exact---here the surd
$\sqrt{14}$.}
\pair{03-introduction/linear-algebra-vectors/3}{\B{Cross} is the three-dimensional
cross product\index{cross product}: $\hat\imath\times\hat\jmath=\hat k$.}
\pair{03-introduction/linear-algebra-vectors/4}{The same \texttt{.}\ applies a matrix
to a vector\index{matrix!times vector}.}
\pair{03-introduction/linear-algebra-vectors/5}{And multiplies two
matrices\index{matrix!product}---here by the swap matrix, which exchanges the
columns.}
\end{codepairs}

The staple operations on a square matrix are each a single call away, and---this is
the point---they return \emph{exact} answers. Working over the rationals rather than
in floating point, \B{Inverse} produces an exact rational matrix, and multiplying it
back recovers the identity with no rounding error at all:

\begin{codepairs}
\pair{03-introduction/linear-algebra-matrices/1}{A matrix is a list of its rows; \texttt{A}
is defined once and reused below.}
\pair{03-introduction/linear-algebra-matrices/2}{\B{Det} is the
determinant\index{determinant}.}
\pair{03-introduction/linear-algebra-matrices/3}{\B{Tr} is the trace\index{trace}, the
sum of the diagonal entries.}
\pair{03-introduction/linear-algebra-matrices/4}{\B{Inverse} gives the exact inverse,
with rational entries.}
\pair{03-introduction/linear-algebra-matrices/5}{Multiplying a matrix by its inverse
returns the identity\index{identity matrix}, exactly.}
\pair{03-introduction/linear-algebra-matrices/6}{\B{Transpose} reflects a matrix across
its diagonal, turning rows into columns.}
\pair{03-introduction/linear-algebra-matrices/7}{\B{IdentityMatrix} builds the
$n\times n$ identity.}
\pair{03-introduction/linear-algebra-matrices/8}{\B{DiagonalMatrix} places a list along
the diagonal, zeros elsewhere.}
\end{codepairs}

Nothing here requires the entries to be numbers. Given symbols, the same functions
return the general formula\index{matrix!symbolic}---the textbook determinant and inverse
of a $2\times2$ matrix:

\begin{codepairs}
\pair{03-introduction/linear-algebra-symbolic/1}{The determinant $ad-bc$.}
\pair{03-introduction/linear-algebra-symbolic/2}{And the inverse, each entry divided by
that determinant.}
\end{codepairs}

\subsection{Linear systems}

Solving $A\mathbf{x}=\mathbf{b}$ is the oldest task in the subject. \B{LinearSolve}
does it directly; \B{RowReduce} performs the Gaussian elimination\index{Gaussian
elimination} behind it, returning the reduced row echelon form\index{row echelon form};
and \B{MatrixRank} and \B{NullSpace} report the structure of the solution set---the
rank\index{rank (of a matrix)} of the coefficient matrix and a basis for the vectors it
sends to zero.

\begin{codepairs}
\pair{03-introduction/linear-algebra-solve/1}{\B{LinearSolve} solves a $2\times2$ system
exactly.}
\pair{03-introduction/linear-algebra-solve/2}{\B{RowReduce} drives a matrix to reduced row
echelon form; this one is invertible, so the result is the identity.}
\pair{03-introduction/linear-algebra-solve/3}{\B{MatrixRank}: these two rows are
proportional, so the rank is~$1$.}
\pair{03-introduction/linear-algebra-solve/4}{\B{NullSpace} returns a basis for the
vectors the matrix annihilates.}
\end{codepairs}

\subsection{Eigenvalues and eigenvectors}

The deepest invariants of a matrix are its \emph{eigenvalues}\index{eigenvalue} and
\emph{eigenvectors}\index{eigenvector}---the scalars $\lambda$ and directions
$\mathbf{v}$ for which $A\mathbf{v}=\lambda\mathbf{v}$. \B{Eigenvalues} and
\B{Eigenvectors} compute them exactly:

\begin{codepairs}
\pair{03-introduction/linear-algebra-eigen/1}{A symmetric matrix with integer
eigenvalues.}
\pair{03-introduction/linear-algebra-eigen/2}{Its eigenvectors, orthogonal as symmetry
demands.}
\pair{03-introduction/linear-algebra-eigen/3}{Here the eigenvalues are the golden
ratio\index{golden ratio} $\varphi=(1+\sqrt5)/2$ and its conjugate---the growth rates
concealed in the Fibonacci recurrence.}
\pair{03-introduction/linear-algebra-eigen/4}{Indeed the powers of that same matrix
\emph{are} the Fibonacci numbers\index{Fibonacci numbers}: the tenth power displays
$F_{11}=89$ and $F_{10}=55$.}
\pair{03-introduction/linear-algebra-eigen/5}{When the eigenvalues have no radical form
they come back as \B{Root} objects\index{Root object}---the roots of $x^{3}-x-1$, exact
algebraic numbers, exactly as \B{Solve} returned a cubic in \S\ref{sec:intro-solve}.}
\pair{03-introduction/linear-algebra-eigen/6}{Numericalising those exact eigenvalues
with \texttt{N} resolves the real root to $\approx1.32472$ and the complex-conjugate
pair to $-0.662359\pm0.56228\,i$.}
\pair{03-introduction/linear-algebra-eigen/7}{Asking instead for the eigenvalues of the
floating-point matrix lands on the identical list---by a numerical algorithm that never
forms the characteristic polynomial at all.}
\end{codepairs}

Those last two lines reach the \emph{same} three numbers by opposite roads. The first
takes the exact \Bnoidx{Root} eigenvalues and rounds them to machine precision; the second
hands the floating-point matrix straight to a numerical eigenvalue algorithm, which
never builds a characteristic polynomial\index{characteristic polynomial} to
factor---though \B{CharacteristicPolynomial} forms that polynomial, $\det(m-xI)$
(or $\det(m-xa)$ for a pencil $\{m,a\}$), on demand. That the two agree is a reassuring
cross-check---and a first sight of a choice that runs through all of computational
linear algebra\index{linear algebra!exact versus numerical}: exact or numerical, and
what each costs.

Section~\ref{sec:linalg} takes up linear algebra---exact and numerical alike---in
depth.

\section{Programming}
\label{sec:intro-programming}

Mathilda is also a programming language\index{programming language}, and an unusually expressive one. It supports
three styles that blend freely: pattern-based rewriting, procedural code, and
functional programming.

\subsection*{Pattern matching}

The most characteristic style is definition by \emph{patterns}\index{pattern matching}. A name followed by an
underscore, like \texttt{x\_}, matches anything and binds it\index{pattern matching!Blank (underscore)}, and a pattern can be
narrowed: \texttt{x\_Integer} matches only an expression with head \texttt{Integer}\index{head (of expression)},
and \texttt{?Positive} adds a test\index{pattern matching!test} the match must pass. A definition with \texttt{:=}
applies whenever its pattern fits:

\begin{codepairs}
\pair{03-introduction/pattern-matching/1}{Define \texttt{f} only on positive
integers.}
\pair{03-introduction/pattern-matching/2}{A positive integer matches: $7^{2}=49$.}
\pair{03-introduction/pattern-matching/3}{A negative integer does not match.}
\pair{03-introduction/pattern-matching/4}{A real does not match either.}
\pair{03-introduction/pattern-matching/5}{Nor does a complex number; each call is
returned unevaluated.}
\pair{03-introduction/pattern-matching/6}{\B{Cases} keeps the elements matching a
pattern---the even ones.}
\pair{03-introduction/pattern-matching/7}{The rule \texttt{/.}\ (\B{ReplaceAll})
substitutes into an expression.}
\pair{03-introduction/pattern-matching/8}{A pattern can destructure a two-element
list.}
\pair{03-introduction/pattern-matching/9}{It binds both parts at once. More in
Chapter~\ref{ch:programming}.}
\end{codepairs}

\subsection*{Procedural programming}

When an ordinary loop is the clearest thing, Mathilda has loops. \B{Module} introduces
local variables\index{local variable (Module)}, and \B{Do}, \B{For}, and \B{While} iterate:

\begin{codepairs}
\pair{03-introduction/procedural/1}{A \B{Do} loop sums $1$ through $100$.}
\pair{03-introduction/procedural/2}{A \B{For} loop computes $5! = 120$.}
\pair{03-introduction/procedural/3}{A \B{While} loop counts the steps of the Collatz
iteration from $27$.}
\end{codepairs}

\subsection*{Functional programming}

Often the shortest code is functional\index{functional programming}---building a result by applying and composing
functions rather than by mutating variables. The \texttt{\&} notation makes an
anonymous function, with \texttt{\#} as its argument\index{anonymous function}:

\begin{codepairs}
\pair{03-introduction/functional/1}{\B{Map} applies a function to every element.}
\pair{03-introduction/functional/2}{\B{Select} keeps those passing a test---the
primes.}
\pair{03-introduction/functional/3}{\B{Fold} chains a function across a list.}
\pair{03-introduction/functional/4}{\B{Nest} applies one repeatedly: $\sqrt{\cdot}$ of
$2$, three times.}
\pair{03-introduction/functional/5}{\B{Table} builds a list from a formula.}
\pair{03-introduction/functional/6}{\B{Apply}, written \texttt{@@}, replaces the head:
\texttt{Times @@ list} is the product.}
\pair{03-introduction/functional/7}{\texttt{@@@} applies at level one, to each
sublist.}
\pair{03-introduction/functional/8}{\B{Cases} picks out the elements matching a
pattern---the integers.}
\end{codepairs}

\section{Graphics}
\label{sec:intro-graphics}

Mathilda draws its own graphics. Every plotting function returns a \texttt{Graphics}\index{Graphics object} (in two
dimensions) or \texttt{Graphics3D} (in three) object built from primitives---lines, points,
polygons, colours---which the built-in engine renders on screen and \B{Export} writes to
PDF, PNG, or JPEG. This closing tour gives one example of each plotting function; as in
\S\ref{sec:intro-numcalc}, every plot below is the genuine \texttt{Out} of the command
shown, produced by the running system.

\subsection{Functions and data}

\B{Plot} draws one or more functions of a real variable---here the first odd harmonics of
a square wave:

\plotcell{03-introduction/graphics/plot}

\noindent \B{ListPlot} draws discrete data as points:

\plotcell{03-introduction/graphics/listplot}

\noindent \B{ParametricPlot} traces a curve $(x(t), y(t))$---a Lissajous figure---and
\B{PolarPlot} draws $r(\theta)$, here a five-petal rose:

\plotcell{03-introduction/graphics/parametricplot}

\plotcell{03-introduction/graphics/polarplot}

\subsection{Vector fields}

\B{VectorPlot} shows a field as a grid of arrows\index{vector field}, \B{StreamPlot} as flow lines. Here is the
phase field of a pendulum, then a pure rotation:

\plotcell{03-introduction/graphics/vectorplot}

\plotcell{03-introduction/graphics/streamplot}

\subsection{Scalar fields and arrays}

\B{ContourPlot} draws level curves of $f(x,y)$, \B{DensityPlot} shades the plane by value,
and \B{ArrayPlot} renders a matrix as a grid of cells---here the greatest-common-divisor
table, whose diagonals of bright cells expose its number-theoretic structure:

\plotcell{03-introduction/graphics/contourplot}

\plotcell{03-introduction/graphics/densityplot}

\plotcell{03-introduction/graphics/arrayplot}

\subsection{Charts and graphs}

\B{BarChart} draws labelled bars; \B{GraphPlot} lays out a graph automatically:

\plotcell{03-introduction/graphics/barchart}

\plotcell{03-introduction/graphics/graphplot}

\subsection{Complex functions}

\B{ComplexPlot} colours the plane by the value of a complex function\index{domain coloring}---hue for the phase,
brightness for the magnitude---so its zeros and poles stand out at a glance:

\plotcell{03-introduction/graphics/complexplot}

\subsection{Three dimensions}

\B{Plot3D} draws a surface $z=f(x,y)$, \B{ParametricPlot3D} a curve or surface in space,
and \B{ComplexPlot3D} the modulus of a complex function as a surface coloured by phase.
These return \texttt{Graphics3D}, which the engine renders with lighting and a movable camera,
and which \B{Export} writes as PNG or JPEG:

\plotcell{03-introduction/graphics/plot3d}

\plotcell{03-introduction/graphics/parametricplot3d}

\plotcell{03-introduction/graphics/complexplot3d}

\bigskip
\noindent
That is the tour. Every area glimpsed here has a chapter of its own ahead, and this
introduction will itself grow as Mathilda does. With the shape of the system in view,
we can now begin in earnest---starting, in the next part, with numbers.

\chapter{Mathematics in Mathilda}
\label{ch:mathematics}

This is the heart of the book. The previous chapter was a quick tour; the sections
that follow treat, one domain at a time and in depth, the mathematics Mathilda
performs---arithmetic and its precision models, algebra, calculus, linear algebra,
numerical calculus, number theory, and the special functions---showing not just how
to use each but how the system does what it does.

\section{Arithmetic}
\label{sec:arithmetic}

Arithmetic is the floor everything else stands on, and Mathilda takes it seriously.
Rather than one kind of number, it keeps a small \emph{tower} of numeric
representations and moves between them automatically:

\begin{itemize}
  \item \textbf{Machine integers} — signed 64-bit words, the hardware's native
        integer, used while a value is small enough to fit.
  \item \textbf{Arbitrary-precision integers} — GMP~\cite{gmp} ``bignums'', into
        which a machine integer is promoted the instant it would overflow, and out
        of which a result is demoted the instant it fits again.
  \item \textbf{Exact rationals} — pairs of exact integers, kept in lowest terms.
  \item \textbf{Machine reals} — IEEE-754 double-precision floating point, the fast
        approximate numbers of the hardware.
  \item \textbf{Arbitrary-precision reals} — MPFR~\cite{mpfr} values, correctly
        rounded to as many digits as you ask for.
\end{itemize}

The whole of this section is about how these five representations behave and how
they hand off to one another. We begin with the way numbers are typed and tested, and
the operators that combine them; then we take the representations in three groups: the
machine numbers (fast, fixed-size), the arbitrary-precision numbers (exact integers and
rationals, and MPFR reals), and finally interval arithmetic, which computes with
\emph{guaranteed enclosures} rather than points.

\subsection{The types of number}
\label{ssec:number-types}

In Mathilda every value carries a \emph{head} that names its type, and \B{Head} reads
it off. The four numeric heads are \texttt{Integer}, \texttt{Real}, \texttt{Rational},
and \texttt{Complex}:

\begin{codepairs}
\pair{arithmetic/heads/1}{A whole number has head \texttt{Integer}.}
\pair{arithmetic/heads/2}{So does a bignum---an arbitrary-precision integer is still an
\texttt{Integer}, not a separate type.}
\pair{arithmetic/heads/3}{A machine real has head \texttt{Real}.}
\pair{arithmetic/heads/4}{A fraction has head \texttt{Rational}: internally it is
\texttt{Rational[p, q]}.}
\pair{arithmetic/heads/5}{And a complex number has head \texttt{Complex}, stored as
\texttt{Complex[re, im]}.}
\end{codepairs}

Rather than test the head by hand, you use the \emph{predicates}---the functions ending
in \texttt{Q} (for ``question'') that return \texttt{True} or \texttt{False}. Mathilda
provides a full family for numbers:

\begin{center}\small
\begin{tabular}{@{}ll@{}}
\toprule
Predicate & True when the argument is\ldots \\
\midrule
\B{IntegerQ}        & an integer (\texttt{Integer} or a bignum) \\
\B{MachineIntegerQ} & a machine-word (64-bit) integer, but not a bignum \\
\B{RationalQ}       & an exact rational (an integer or a \texttt{Rational}) \\
\B{ComplexQ}        & a \texttt{Complex} number \\
\B{ExactNumberQ}    & an exact number (integer, rational, or exact complex) \\
\B{InexactNumberQ}  & an inexact number (a machine or MPFR real) \\
\B{NumberQ}         & any explicit number \\
\B{NumericQ}        & a number \emph{or} a numeric constant such as $\pi$ \\
\bottomrule
\end{tabular}
\end{center}

\begin{codepairs}
\pair{arithmetic/predicates/1}{\B{IntegerQ} accepts any integer.}
\pair{arithmetic/predicates/2}{\B{MachineIntegerQ} is stricter: a bignum, having
overflowed the machine word, is not a machine integer.}
\pair{arithmetic/predicates/3}{\B{RationalQ} recognises a fraction.}
\pair{arithmetic/predicates/4}{A Gaussian integer is an \emph{exact} number, so
\B{ExactNumberQ} is true.}
\pair{arithmetic/predicates/5}{A fifty-digit MPFR value is \emph{inexact}, so
\B{InexactNumberQ} is true.}
\pair{arithmetic/predicates/6}{\B{ComplexQ} tests for the \texttt{Complex} head.}
\end{codepairs}

Each predicate is a small structural test: it looks at the head and type of its one
argument and returns \texttt{True} or \texttt{False}, never leaving the call
unevaluated for a concrete number. Every one of these functions carries its own
documentation, which you can read at the prompt with \texttt{?Name}; here, for
instance, is what \B{MachineIntegerQ} reports about itself:

\usagebox{MachineIntegerQ}

Throughout the rest of this section, boxes like the one above reproduce a function's
own docstring---exactly the text \texttt{?Name} prints---together with its attributes.

\subsection{The arithmetic operators}
\label{ssec:operators}

The operators you type---$+$, $-$, $\times$, $/$, $\wedge$---are surface syntax for a
few built-in heads, and \emph{only a few}. \B{Plus} and \B{Times} and \B{Power} are the
primitives; subtraction, division, and negation are not separate operations but
rewrites into them, as \B{FullForm} reveals:

\begin{codepairs}
\pair{arithmetic/operator-forms/1}{Subtraction is \B{Plus} against a negated term:
$a-b$ is \texttt{Plus[a, Times[-1, b]]}.}
\pair{arithmetic/operator-forms/2}{Division is \B{Times} against a reciprocal power:
$a/b$ is \texttt{Times[a, Power[b, -1]]}.}
\pair{arithmetic/operator-forms/3}{Unary minus is multiplication by $-1$.}
\pair{arithmetic/operator-forms/4}{Even $\sqrt{x}$ is a \B{Power}, with exponent the
exact rational $1/2$.}
\end{codepairs}

Reducing everything to \B{Plus}, \B{Times}, and \B{Power} is what keeps the evaluator
small: a single canonicalizer for each, driven by attributes\calloutnote{Under the
hood}{The attributes do the organising. \texttt{Flat} means nested calls are spliced
(\texttt{Plus[a, Plus[b, c]]} becomes \texttt{Plus[a, b, c]}); \texttt{Orderless} sorts
the arguments into a canonical order so that \texttt{a + b} and \texttt{b + a} are the
same expression and like terms sit together to be combined; \texttt{OneIdentity} lets a
one-argument call collapse; \texttt{Listable} threads the operation over lists; and
\texttt{NumericFunction} marks it for numeric evaluation. Underneath, the exact-integer
paths add and multiply in a $128$-bit intermediate and, on overflow, fall through to
GMP~\cite{gmp}, promoting the result to a bignum---so the answer is always exact, and
the fast machine-word path is taken whenever it fits.}, handles addition,
subtraction, multiplication, division, powers, and roots at once. Here are the three,
with their own docstrings and attributes:

\usagebox{Plus}
\usagebox{Times}
\usagebox{Power}

\subsection{Machine-precision computation}
\label{ssec:machine-precision}

At the bottom of the tower are the machine numbers: a signed 64-bit integer and an
IEEE-754 double. These are the types the processor implements directly, and Mathilda
uses them whenever it can, because they are fast.

An integer stays a machine word only while it fits in 64 bits. The largest such value
is $2^{63}-1$; one step beyond and Mathilda promotes it to an arbitrary-precision
integer, silently and without loss:

\begin{codepairs}
\pair{arithmetic/machine-integers/1}{$2^{62}$ fits in a signed 64-bit word and stays
a machine integer.}
\pair{arithmetic/machine-integers/2}{$2^{63}$ is one past the largest machine integer,
so it is promoted to a GMP bignum---still exact.}
\pair{arithmetic/machine-integers/3}{Addition promotes too: two machine integers whose
sum overflows produce a bignum. The promotion is detected with a wider (128-bit)
intermediate.}
\end{codepairs}

Machine reals are IEEE-754 binary64 (``double'') values. Their characteristics are
reported by a handful of read-only system constants, derived at start-up from the
platform's \texttt{<float.h>}:

\begin{codepairs}
\pair{arithmetic/machine-limits/1}{A double carries about $15.95$ decimal digits (a
53-bit mantissa).}
\pair{arithmetic/machine-limits/2}{\B{\$MachineEpsilon}: the gap between $1$ and the
next representable double.}
\pair{arithmetic/machine-limits/3}{The largest finite double.}
\pair{arithmetic/machine-limits/4}{And the smallest positive normal double.}
\end{codepairs}

Here is the single most important thing to understand about Mathilda's machine reals,
and a real difference from Mathematica: \textbf{Mathilda does not use significance
arithmetic}.\calloutnote{Theory}{\emph{Significance arithmetic} (which Mathematica uses
for arbitrary-precision numbers) attaches to each number an estimate of its error and
propagates that estimate through every operation, so that catastrophic cancellation
visibly reduces the reported precision. It is elegant but only a first-order model of
error, and it makes a number's behaviour depend on its history. Mathilda takes the other
well-trodden road: a machine real is a fixed-precision IEEE double, exactly as specified
by the floating-point standard~\cite{ieee754}. If you need to know how error propagates,
you either reason about it yourself---the classic reference is Goldberg~\cite{goldberg1991}---or
reach for the rigorous tool in \S\ref{ssec:interval}, interval arithmetic.} A machine
real is just an IEEE double, with a fixed 53-bit mantissa; it
carries no error term, and operations do not track how many digits remain trustworthy.
What you get is exactly what the hardware, C, or NumPy would give you---no more and no
less:

\begin{codepairs}
\pair{arithmetic/machine-doubles/1}{Dividing two machine reals gives a machine real,
printed to six figures by default.}
\pair{arithmetic/machine-doubles/2}{$0.1+0.2$ \emph{prints} as \texttt{0.3}, because
six figures round to that.}
\pair{arithmetic/machine-doubles/3}{But it is not exactly $0.3$: the residual is the
true IEEE rounding error, shown plainly rather than hidden.}
\pair{arithmetic/machine-doubles/4}{The precision of any machine real is the symbol
\texttt{MachinePrecision}.}
\end{codepairs}

This fixed-precision model is not a limitation so much as a deliberate choice, and for
the overwhelming majority of numerical work---linear algebra, differential equations,
optimization, plotting---it is exactly right: it is fast, it is predictable, and it
matches every other serious numerical system bit for bit. When you genuinely need more
than a double can hold, you move up the tower.

\subsection{Arbitrary-precision computation}
\label{ssec:arbitrary-precision}

Exactness in Mathilda is never an option you switch on; it is the default, and it never
silently degrades. Integers grow without bound, factorials and powers stay exact, and
the arithmetic behind them is GMP~\cite{gmp}, the same library that underpins much of
computational number theory:

\begin{codepairs}
\pair{arithmetic/bignums/1}{A thirty-one-digit power, kept exactly.}
\pair{arithmetic/bignums/2}{$100!$ is a $158$-digit integer, computed exactly by GMP.}
\pair{arithmetic/bignums/3}{Factoring the Mersenne number $2^{67}-1$---the number Cole
famously factored by hand in 1903.}
\end{codepairs}

Rational numbers are exact too. A ratio of integers is kept as a reduced fraction, its
numerator and denominator sharing no common factor, and it collapses to an integer the
moment its denominator becomes~$1$:

\begin{codepairs}
\pair{arithmetic/rationals/1}{A sum of fractions is another fraction, reduced to lowest
terms.}
\pair{arithmetic/rationals/2}{Reduction happens on construction, so $2/4$ is already
$1/2$.}
\pair{arithmetic/rationals/3}{When the denominator reduces to $1$, the rational
collapses to a plain integer.}
\end{codepairs}

For \emph{approximate} numbers beyond machine precision, Mathilda uses MPFR, which
computes correctly-rounded floating point to any requested precision. You reach it with
the two-argument form \B{N}\texttt{[expr, d]}, asking for \texttt{d} digits:

\begin{codepairs}
\pair{arithmetic/mpfr-N/1}{$\pi$ to fifty digits, via MPFR's \texttt{mpfr\_const\_pi}.}
\pair{arithmetic/mpfr-N/2}{$e$ to a hundred digits.}
\pair{arithmetic/mpfr-N/3}{$\sqrt2$ to forty digits---an MPFR power, not a table
lookup.}
\pair{arithmetic/mpfr-N/4}{Ap\'ery's constant $\zeta(3)$ to fifty digits. The zeta
function is evaluated by FLINT's rigorous ball arithmetic~\cite{flint,arb}, not by
MPFR.}
\end{codepairs}

So the numeric substrate is really three libraries working together: GMP for exact
integers, MPFR for correctly-rounded arbitrary-precision reals, and FLINT/Arb for the
harder special functions, where rigorous complex-ball arithmetic guarantees every
printed digit. \B{N} orchestrates them, replacing each exact leaf of an expression with
a numeric value of the requested precision and letting the ordinary evaluator re-run.

A number's precision is itself a value you can query, and here Mathilda is honest about
a small subtlety:

\begin{codepairs}
\pair{arithmetic/precision/1}{Exact numbers have \emph{infinite} precision.}
\pair{arithmetic/precision/2}{But \B{N}\texttt{[Pi, 50]} reports about $50.27$ digits,
not $50$: fifty decimal digits are stored as $\lceil 50\log_2 10\rceil = 167$ bits,
which is a hair more than fifty digits' worth.}
\pair{arithmetic/precision/3}{You can also state a number's precision on input: the
literal \texttt{1.5`20} is $1.5$ carried at twenty digits.}
\end{codepairs}

Precision is \emph{contagious}, and when representations of different precision meet, the
lowest wins---machine precision being the floor. This, too, follows Mathematica:

\begin{codepairs}
\pair{arithmetic/contagion/1}{Add a machine real to a hundred-digit number and the
result drops to machine precision: one approximate double poisons the lot.}
\pair{arithmetic/contagion/2}{Applying \B{N} to an already-precise number preserves its
precision rather than collapsing it to a double.}
\pair{arithmetic/contagion/3}{$1001!$ has $2{,}571$ digits, far past a double's range,
so \B{N} keeps it as a machine-precision value with an unbounded exponent---a real with
a mantissa but no overflow.}
\end{codepairs}

\subsection{Interval arithmetic}
\label{ssec:interval}

The tools so far give you either an exact answer or an approximate one whose error you
must reason about yourself. \emph{Interval arithmetic} gives you a third thing: an
approximate answer that comes with a \emph{proof}. Instead of a single number, you
compute with an interval $[a,b]$ that is guaranteed to contain the true value, and every
operation returns an interval guaranteed to contain the true result. The idea goes back
to Moore~\cite{moore2009}; Mathilda implements it as a first-class numeric object,
\B{Interval}\texttt{[\{lo, hi\}]}.\calloutnote{Pitfall}{Mathilda's intervals differ from
Mathematica's in two deliberate ways, both documented. A machine number becomes a
\emph{point} interval---\texttt{Interval[1.]} is \texttt{Interval[\{1., 1.\}]}---rather
than being preloaded with a half-ULP uncertainty band; if you want an uncertainty
propagated, you supply it. And \B{Interval} is a set, not a scalar, so it is deliberately
neither \texttt{Listable} nor a \texttt{NumericFunction}. The enclosures are rigorous in
the spirit of the interval standard, though Mathilda implements no formal decoration
system.}

The arithmetic operations and the elementary functions all propagate intervals:

\begin{codepairs}
\pair{arithmetic/interval-basic/1}{Interval addition: the sum of $[1,2]$ and $[3,4]$ is
$[4,6]$.}
\pair{arithmetic/interval-basic/2}{Multiplication takes the extremes over the four
corner products.}
\pair{arithmetic/interval-basic/3}{Squaring is tighter than multiplying an interval by
itself: $[-2,5]^2=[0,25]$, because the operation knows the two factors are equal.}
\pair{arithmetic/interval-basic/4}{A reciprocal that straddles zero splits into a
disjoint union of two rays---an interval may be several ranges at once.}
\pair{arithmetic/interval-basic/5}{Functions thread too: $\sqrt{[4,9]}=[2,3]$.}
\end{codepairs}

What makes these \emph{rigorous} rather than merely suggestive is \emph{outward
rounding}. Exact endpoints (integers, rationals, $\pi$) pass through untouched, but the
moment an endpoint becomes an inexact machine or MPFR real, Mathilda nudges it one unit
in the last place \emph{outward}---lower bounds down, upper bounds up---so the printed
interval is a mathematically certain enclosure, never an underestimate. There is no
dependence on an external interval library; the widening is done directly with
\texttt{nextafter} and MPFR's directed neighbours.

The classic application is certifying a floating-point computation. Recall that
$0.1+0.2-0.3$ was not exactly zero; interval arithmetic pins down exactly where the true
answer lies:

\begin{codepairs}
\pair{arithmetic/interval-rounding/1}{A certified enclosure of the rounding error: the
true value of $0.1+0.2-0.3$ lies somewhere in this tiny interval around zero.}
\pair{arithmetic/interval-rounding/2}{And zero is provably inside it---so the exact
answer really is $0$, rounding error notwithstanding.}
\end{codepairs}

A more striking use is proving that a solution \emph{exists}. If the image of a
continuous function over an interval straddles zero, the function must have a root there.
Evaluate $\sin$ over $[3, 7/2]$:

\begin{codepairs}
\pair{arithmetic/interval-root/1}{The enclosure of $\sin$ over $[3, 7/2]$, with exact
symbolic endpoints, runs from a negative value to a positive one.}
\pair{arithmetic/interval-root/2}{Zero lies inside it. This is a computer-assisted
proof that $\sin$ has a root in $[3, 7/2]$---that is, that $\pi$ lives there.}
\end{codepairs}

Interval arithmetic has one famous limitation, the \emph{dependency problem}: because
each occurrence of a variable is treated independently, an expression like $x-x$ is not
recognised as zero, and the enclosure it returns is wider than the true range. This is
not a bug but a property of the method, and it means naive interval evaluation can
over-estimate:

\begin{codepairs}
\pair{arithmetic/interval-advanced/1}{$x-x$ over $[-1,1]$ gives $[-2,2]$, not $\{0\}$:
the two copies are treated as independent.}
\pair{arithmetic/interval-advanced/2}{Iterating the logistic map $4x(1-x)$ on a point
interval, the enclosure widens at each step---a certified picture of how rounding error
grows in a chaotic system, always containing the true orbit.}
\end{codepairs}

\section{Algebra}
\label{sec:algebra}

Arithmetic computes with numbers; algebra computes with the \emph{polynomials} and
\emph{rational functions}\index{rational function} you build when you let symbols stand
in for numbers. These are the central objects of this section, and everything the
previous section established about exactness carries straight over: a polynomial with
rational coefficients keeps those coefficients exact, a factorization is a genuine
identity and not an approximation, and a solution is returned as the algebraic number it
is rather than a decimal that happens to be close.

A polynomial in Mathilda is not a special kind of value. It is an ordinary
expression---a tree of \B{Plus}, \B{Times}, and \B{Power}---and the algebra functions
read its structure off that tree.\calloutnote{Under the hood}{There is no dedicated
polynomial data type: \texttt{x\textasciicircum 2 + 3 x + 2} is the same
\texttt{Plus[\ldots]} expression the evaluator manipulates everywhere else. When an
operation is heavy---expanding a large power, factoring, a polynomial GCD---Mathilda
copies the relevant part into a packed FLINT~\cite{flint} \texttt{fmpq\_mpoly}
(a multivariate polynomial over $\mathbb{Q}$ stored as sorted exponent-coefficient
pairs), does the work there at C speed, and converts the result back to an expression.
The packed form is transient; what you see and store is always an ordinary tree.} That
uniformity is why the same \B{Part}, \B{Map}, and \B{ReplaceAll} you use on any
expression work on a polynomial too.

The section climbs from structure to solution. We begin with the anatomy of a
polynomial and the routines that read it; then expanding and collecting; then factoring,
first over the rationals and then over algebraic number fields. From there we take up
the greatest common divisor, resultants, and elimination---the machinery of eliminating
a variable---and the rational-function operations built on them. Gr\"obner bases
generalise all of this to systems. We close with equation solving: a single polynomial,
then a system, and finally \B{Reduce}, which describes a solution set completely,
including over the real numbers.

\subsection{The anatomy of a polynomial}
\label{ssec:alg-polynomials}

Before manipulating a polynomial it helps to interrogate it, and a handful of functions
read off its structure. \B{PolynomialQ} tests whether an expression \emph{is} a
polynomial in a given variable; \B{Variables} lists the indeterminates; \B{Exponent}
gives the highest power of a variable; \B{Coefficient} and \B{CoefficientList} extract
coefficients; \B{Collect} regroups terms by the powers of a chosen variable; and
\B{HornerForm} rewrites a polynomial in the nested form that evaluates with the fewest
multiplications.

\begin{codepairs}
\pair{algebra/polynomials/1}{\B{PolynomialQ} confirms that $x^3-2x+1$ is a polynomial in
$x$.}
\pair{algebra/polynomials/2}{It is a genuine structural test: $\sin x + x$ is not a
polynomial, because $\sin x$ is not a power of $x$.}
\pair{algebra/polynomials/3}{\B{Variables} returns the independent variables, in
canonical order.}
\pair{algebra/polynomials/4}{\B{Exponent} is the degree in $x$---the largest power that
appears.}
\pair{algebra/polynomials/5}{\B{Coefficient}\texttt{[expr, x, 2]} pulls out the
coefficient of $x^2$; here $\binom{5}{2}=10$.}
\pair{algebra/polynomials/6}{\B{CoefficientList} returns every coefficient at once, from
the constant term upward---the fourth row of Pascal's triangle.}
\pair{algebra/polynomials/7}{\B{Collect} gathers the expression by powers of $x$,
pulling the common factor out of the linear term.}
\pair{algebra/polynomials/8}{\B{HornerForm} nests the polynomial so that evaluating it
costs three multiplications instead of six.}
\end{codepairs}

These are the lenses through which the heavier operations see a polynomial, and each of
those operations---expansion, factoring, division---returns an ordinary expression you
can hand straight back to another.

\subsection{Expanding and collecting}
\label{ssec:alg-expand}

The most basic transformation is to \emph{expand}: to multiply out every product and
positive-integer power until the expression is a flat sum of monomials. \B{Expand} does
this at the top level; \B{ExpandAll} does it everywhere, including inside denominators;
and \B{ExpandNumerator} and \B{ExpandDenominator} expand just one part of a fraction.

\begin{codepairs}
\pair{algebra/expand/1}{\B{Expand} multiplies out a power, giving the binomial
coefficients.}
\pair{algebra/expand/2}{It multiplies out products of factors just as readily.}
\pair{algebra/expand/3}{And it works in any number of variables---here the trinomial
cube, with its multinomial coefficients.}
\pair{algebra/expand/4}{\B{ExpandAll} reaches into the denominator too, distributing the
square over the numerator so the fraction becomes a sum of simpler ones.}
\pair{algebra/expand/5}{\B{ExpandNumerator} expands only the top of a fraction, leaving
the factored denominator intact.}
\pair{algebra/expand/6}{\B{ExpandDenominator} does the reverse, multiplying the
denominator into a single polynomial.}
\end{codepairs}

Expansion is the operation most likely to blow up in size---a product of $k$ factors
each with $t$ terms has up to $t^k$ terms when multiplied out---so it is also where the
packed representation earns its keep.\calloutnote{Performance}{Over the rationals
\B{Expand} routes through FLINT's \texttt{fmpq\_mpoly}~\cite{flint}, whose exponent
vectors are bit-packed and whose term merging is done with a heap; a dense expansion
that would crawl as a tree of \texttt{Plus} nodes finishes in a fraction of a second.
Because the result size is bounded in advance by the Newton polytope of the input,
Mathilda estimates it first and returns an \texttt{Overflow[]} marker---rather than
exhausting memory---when an expansion such as $(a+b+\cdots+g)^{60}$ would be too large to
hold.} The inverse move---writing a sum as a product---is factoring, and it is a great
deal deeper.

\subsection{Factoring over the rationals}
\label{ssec:alg-factor}

To \B{Factor} a polynomial is to write it as a product of \emph{irreducible} factors:
polynomials that cannot themselves be factored further. Over the rational numbers this
factorization is unique up to order and rational multiples, and Mathilda computes it
exactly. \B{FactorList} returns the same information as a list of
$\{\text{factor}, \text{multiplicity}\}$ pairs; \B{FactorTerms} pulls out only the
numeric content; and \B{FactorSquareFree} does a partial job that is worth understanding
because it is the first step of the full algorithm.

\begin{codepairs}
\pair{algebra/factor/1}{\B{Factor} splits $x^6-1$ into four irreducible pieces over
$\mathbb{Q}$: two linear and two quadratic (the cyclotomic factors).}
\pair{algebra/factor/2}{\B{FactorList} gives the factors with their multiplicities, the
leading numeric factor $2$ among them.}
\pair{algebra/factor/3}{\B{FactorSquareFree} groups the repeated factors: it finds that
$x^4+2x^3+x^2$ is a perfect square, but leaves the base $x+x^2$ \emph{unfactored}.}
\pair{algebra/factor/4}{\B{Factor} finishes the job, splitting that base into $x$ and
$1+x$: the full irreducible factorization is $x^2(1+x)^2$.}
\pair{algebra/factor/5}{\B{FactorTerms} extracts only the numeric content $3$, leaving
the primitive part untouched.}
\pair{algebra/factor/6}{Factoring $x^{15}-1$ recovers the cyclotomic polynomials
$\Phi_1$, $\Phi_3$, $\Phi_5$, and $\Phi_{15}$---one for each divisor of $15$.}
\end{codepairs}

The contrast between pairs~3 and~4 is the key to how factoring actually
works.\index{square-free decomposition} A \emph{square-free} decomposition separates a
polynomial's factors by \emph{multiplicity}---which pieces occur once, which twice, and
so on---without yet splitting them into irreducibles.\calloutnote{Theory}{The tool is
the derivative. If an irreducible factor $p$ divides $f$ with multiplicity $m$, then it
divides $f'$ with multiplicity exactly $m-1$, so $\gcd(f, f')$ collects every repeated
factor with its multiplicity dropped by one. \emph{Yun's algorithm}\index{Yun's
algorithm} turns this one observation into a linear number of GCDs that peel off the
multiplicity-$1$ part, the multiplicity-$2$ part, and so on in a single sweep. It is
cheap---only GCDs and divisions---and it guarantees that the harder work to come never
has to cope with a repeated factor. Mathilda's implementation lives in
\texttt{src/poly/facpoly\_squarefree.inc}; the classical reference is
Geddes--Czapor--Labahn~\cite{gcl}.} Once the polynomial is square-free, each factor must
be split into irreducibles, and that is the genuinely hard step.\calloutnote{Under the
hood}{The classical method, \emph{Berlekamp--Zassenhaus}\index{Berlekamp--Zassenhaus
algorithm}, factors modulo a well-chosen prime $p$---a problem that is tractable over a
finite field, via the Cantor--Zassenhaus algorithm---then \emph{Hensel
lifts}\index{Hensel lifting} that factorization from $p$ to $p^2$, $p^4$, and on up to a
modulus large enough to pin down the true integer coefficients. A final recombination
step searches for subsets of the lifted factors whose product has integer coefficients
that divide the original. Mathilda carries a self-contained implementation of this
pipeline (\texttt{facpoly\_bz\_uni.inc}), and uses it, together with a Wang-style
leading-coefficient correction for the multivariate case, whenever FLINT is unavailable.
When FLINT is present---the usual case---the univariate and multivariate factorizations
over $\mathbb{Q}$ are delegated to its state-of-the-art \texttt{fmpz\_poly\_factor} and
\texttt{fmpq\_mpoly\_factor}~\cite{flint}, which turns factorizations that once hung for
tens of seconds into a few milliseconds.}

\subsection{Polynomials over algebraic number fields}
\label{ssec:alg-extensions}

Whether a polynomial is irreducible depends on the numbers you are allowed to use.
$x^2+1$ is irreducible over $\mathbb{Q}$, but the moment you admit the imaginary unit it
splits into $(x-i)(x+i)$. Mathilda lets you extend the coefficient field with the
\texttt{Extension} option to \B{Factor}, adjoining one algebraic number---or a whole
tower of them---and factoring over the resulting \emph{algebraic number
field}\index{algebraic number field} $\mathbb{Q}(\alpha)$.

\begin{codepairs}
\pair{algebra/extensions/1}{Adjoining $i$ with \texttt{Extension -> I} splits $x^2+1$
into its two Gaussian-integer factors.}
\pair{algebra/extensions/2}{A tower: over $\mathbb{Q}(\sqrt2,\sqrt3)$ the quartic
$x^4-10x^2+1$---the minimal polynomial of $\sqrt2+\sqrt3$---splits into all four linear
factors.}
\pair{algebra/extensions/3}{An extension need not split a polynomial completely: over
$\mathbb{Q}(\sqrt2)$, only the $x^2-2$ part of $x^4-5x^2+6$ factors; the $x^2-3$ part
stays irreducible, since $\sqrt3\notin\mathbb{Q}(\sqrt2)$.}
\pair{algebra/extensions/4}{\B{IrreduciblePolynomialQ} confirms $x^2+1$ is irreducible
over the rationals\ldots}
\pair{algebra/extensions/5}{\ldots and that it stops being irreducible once $i$ is
adjoined.}
\pair{algebra/extensions/6}{\B{MinimalPolynomial} runs the other way, recovering the
lowest-degree rational polynomial with $\sqrt2+\sqrt3$ as a root---the very quartic of
pair~2.}
\pair{algebra/extensions/7}{\B{RootReduce} canonicalises an algebraic number; here it
recognises $\sqrt2+\sqrt3$ as the fourth root of that quartic, returned as an exact
\B{Root} object.}
\end{codepairs}

Factoring over a number field looks like it should require a whole new algorithm, but a
classical trick reduces it to factoring over the rationals, which we already
know how to do.\calloutnote{Theory}{\emph{Trager's algorithm}\index{Trager's algorithm}
rests on the \emph{norm}. For $f\in\mathbb{Q}(\alpha)[x]$, the norm
$N(f)=\operatorname{Res}_y\!\big(m_\alpha(y),\,f_y(x)\big)$---the resultant against the
minimal polynomial $m_\alpha$ of $\alpha$---is an ordinary polynomial over $\mathbb{Q}$.
Trager shows that after a suitable integer shift $x\mapsto x-s\alpha$ the norm becomes
square-free, and then each irreducible factor of $f$ over $\mathbb{Q}(\alpha)$ is
recovered as $\gcd\big(f,\,g\big)$ in $\mathbb{Q}(\alpha)[x]$ for one rational-norm
factor $g$. So the recipe is: form the norm, factor it over $\mathbb{Q}$ with the
machinery of the previous subsection, and lift each factor back by a GCD. A tower
$\mathbb{Q}(\alpha_1,\ldots,\alpha_n)$ is first collapsed to a single
\emph{primitive element}\index{primitive element theorem} $\gamma$ by the same
resultant construction, so the multi-generator case needs no new ideas. The code is in
\texttt{src/poly/qafactor.c}.} Mathilda represents an algebraic number not by a special
\texttt{AlgebraicNumber} head but by a \B{Root} object---an exact reference to ``the
$k$-th root of this polynomial''---and \B{ToRadicals} will re-express such a root in
radicals when a radical form exists.

\subsection{GCD, resultants, and elimination}
\label{ssec:alg-gcd}

The greatest common divisor of two polynomials is the highest-degree polynomial dividing
both---exactly the notion that lets you reduce a fraction to lowest terms. \B{PolynomialGCD}
computes it, and \B{PolynomialExtendedGCD} additionally returns the B\'ezout cofactors,
while \B{PolynomialLCM} gives the least common multiple. Division with remainder is
performed by \B{PolynomialQuotient} and \B{PolynomialRemainder} singly, or by
\B{PolynomialQuotientRemainder} for both at once. A second, subtler tool is the
\emph{resultant}\index{resultant}, computed by
\B{Resultant}, which detects a common root of two polynomials \emph{without} finding it,
and its relative the \B{Discriminant}. Together these are the elements of
\emph{elimination theory}, and \B{Eliminate} puts them to work removing a variable from
a system.

\begin{codepairs}
\pair{algebra/gcd-resultant/1}{$x^3-1$ and $x^2-1$ share exactly the factor $x-1$, found
directly by \B{PolynomialGCD}---without factoring either polynomial.}
\pair{algebra/gcd-resultant/2}{\B{PolynomialExtendedGCD} also returns the B\'ezout
cofactors $\{a,b\}$ with $a f + b g = \gcd$.}
\pair{algebra/gcd-resultant/3}{\B{PolynomialQuotientRemainder} divides $x^3-1$ by $x-1$
exactly: quotient $x^2+x+1$, remainder $0$.}
\pair{algebra/gcd-resultant/4}{\B{PolynomialLCM} is the product divided by the GCD.}
\pair{algebra/gcd-resultant/5}{\B{Resultant} vanishes exactly when the two polynomials
share a root. $x^2-2$ and $x^2-3$ share none, so their resultant is the nonzero
constant $1$.}
\pair{algebra/gcd-resultant/6}{With symbolic coefficients the resultant is a polynomial
in them---the condition on $a,b,c,d$ for the two quadratics to share a root.}
\pair{algebra/gcd-resultant/7}{The \B{Discriminant} of the general quadratic is the
familiar $b^2-4ac$.}
\pair{algebra/gcd-resultant/8}{And of the depressed cubic $x^3+px+q$, the classical
$-4p^3-27q^2$.}
\pair{algebra/gcd-resultant/9}{\B{Eliminate} removes $y$ from the two equations, leaving
a single relation in $x$ alone: the projection of their common solutions onto the
$x$-axis.}
\end{codepairs}

The resultant is the determinant of the \emph{Sylvester matrix} of the two polynomials,
and equivalently the product of the differences of their roots; the discriminant of $P$
is (up to a sign and leading coefficient) the resultant of $P$ with its own derivative,
so it vanishes precisely when $P$ has a repeated root. Computing it by that determinant,
though, would be needlessly slow.\calloutnote{Under the hood}{Naive Euclidean division to
find a polynomial GCD suffers from \emph{coefficient explosion}: the intermediate
remainders acquire enormous rational coefficients even when the answer is small. The
\emph{subresultant polynomial remainder sequence}\index{subresultant PRS} tames this by
dividing out, at each step, a predicted content given by Bronstein's $\gamma/\beta/\delta$
recurrence, keeping every coefficient bounded; the resultant and the degree of the GCD
can then be read directly off the sequence. This avoids both the $(m+n)\times(m+n)$
Sylvester determinant and the blow-up of ordinary Euclid, and it is what Mathilda uses
for \B{Resultant}, \B{Subresultants}, and the GCDs beneath \B{Cancel}
(\texttt{src/poly/subresultants.c}; Geddes--Czapor--Labahn~\cite{gcl}). For plain
rational inputs the work is handed to FLINT instead.}

\subsection{Rational functions}
\label{ssec:alg-rational}

A \emph{rational function} is a ratio of polynomials, and three operations put one into a
normal form. \B{Together} combines a sum of fractions over a common denominator;
\B{Cancel} reduces a single fraction to lowest terms by dividing out the GCD of numerator
and denominator; and \B{Apart} performs the partial-fraction decomposition---the reverse
of \B{Together}, and the form you want before integrating. \B{Numerator} and
\B{Denominator} read off the two halves.

\begin{codepairs}
\pair{algebra/rational/1}{\B{Together} puts a sum of fractions over the common
denominator $x(x+1)$.}
\pair{algebra/rational/2}{\B{Numerator} reads off the top of a fraction. Note that it is
\emph{structural}: the shared factor $x+1$ is not cancelled first, so the numerator is
$1+x$\ldots}
\pair{algebra/rational/3}{\ldots and \B{Denominator} likewise returns $x^2-1$, still
carrying that common factor.}
\pair{algebra/rational/4}{\B{Cancel} is what divides the shared $x+1$ out, reducing the
same fraction to $1/(x-1)$.}
\pair{algebra/rational/5}{It works over several variables, cancelling the $x-y$ from
$(x-y)/(x^2-y^2)$ to leave $1/(x+y)$.}
\pair{algebra/rational/6}{\B{Apart} splits a fraction into a sum of pieces with the
simplest possible denominators---the partial-fraction decomposition.}
\end{codepairs}

Two behaviours are worth noting. \B{Together} deliberately does \emph{not}
expand\index{numeric contagion}: it puts the sum over one denominator and cancels, but
leaves the result factored, because expanding is a separate decision you make with
\B{Expand}. And \B{Cancel} works by polynomial GCD rather than by factoring, so it
reduces even fractions whose numerator and denominator are hard to factor.\calloutnote{Under
the hood}{\B{Apart} over $\mathbb{Q}$ uses a fast FLINT path
(\texttt{flint\_apart\_over\_q}): it splits the denominator into coprime factors, solves
the associated B\'ezout relations with an extended-GCD, and reads off each partial
fraction by a $p$-adic expansion---replacing the $O(S^2)$ symbolic linear solve that a
naive method of undetermined coefficients would require. Multivariate and
algebraic-extension cases fall back to the classical elimination approach.} These
partial fractions are exactly what \S\ref{sec:calculus} reaches for when integrating a
rational function.

\subsection{Gr\"obner bases}
\label{ssec:alg-groebner}

Everything so far concerned one or two polynomials. A \emph{system} of polynomials
generates an \emph{ideal}\index{ideal}---the set of all combinations
$\sum h_i f_i$---and the right way to understand that ideal is a
\B{GroebnerBasis}\index{Gröbner basis}: a canonical set of generators that answers, by a
single division, whether a given polynomial lies in the ideal, and from which the
geometry of the common solutions can be read off. The catch, and the craft, is that the
basis depends on a \emph{monomial order}\index{monomial order}---a rule for deciding
which term of a polynomial is ``leading.''

\begin{codepairs}
\pair{algebra/groebner/1}{The circle $x^2+y^2=1$ met with the line $x=y$. In the default
lexicographic order the basis is \emph{triangular}: one generator, $2y^2-1$, involves
$y$ alone, and $x-y$ then recovers $x$.}
\pair{algebra/groebner/2}{For $\{x^2-y,\;y^2-x\}$ the lexicographic basis eliminates $x$
into the single-variable $y^4-y$, whose four roots are the $y$-coordinates of the
solutions.}
\pair{algebra/groebner/3}{The same system in \texttt{DegreeReverseLexicographic} order
gives a different, more balanced basis---cheaper to compute, but no longer triangular.}
\pair{algebra/groebner/4}{Over the finite field $\mathrm{GF}(7)$, with
\texttt{Modulus -> 7}, the whole computation runs in modular arithmetic.}
\end{codepairs}

The choice of order is the whole story. The lexicographic order (pair~1) forces the
basis into a triangular, ``solved'' shape---ideal for eliminating variables and solving,
which is why it underlies system solving in \S\ref{ssec:alg-solve-systems}---but it is
the most expensive to compute, and its coefficients can grow enormously. The degree
orders (pair~3) are far cheaper but do not triangulate.\calloutnote{Under the
hood}{\B{GroebnerBasis} implements \emph{Buchberger's algorithm}\index{Buchberger's
algorithm}: for every pair of basis polynomials it forms their $S$-polynomial---the
combination engineered to cancel the two leading terms---reduces it modulo the current
basis, and adjoins any nonzero remainder, repeating until every $S$-polynomial reduces to
zero. The Gebauer--M\"oller criteria discard, in advance, the many pairs that must reduce
to zero, which is what makes the algorithm practical. Because the lexicographic order is
so costly directly, \texttt{Method -> "GroebnerWalk"} instead computes a cheap
degree-order basis and \emph{walks} it across the Gr\"obner fan to the target order, one
cone wall at a time. The core is \texttt{src/poly/groebner.c}, with the walk in
\texttt{groebnerwalk.c} and a dedicated finite-field engine in \texttt{gbmod.c}; the
standard text is Cox--Little--O'Shea~\cite{clo}.}

\subsection{Solving a polynomial equation}
\label{ssec:alg-solve-poly}

To \B{Solve} a polynomial equation is to describe every value of the unknown that
satisfies it. Mathilda factors the polynomial and dispatches each factor by degree: the
linear, quadratic, cubic, and quartic cases have closed-form solutions in radicals, and
everything else is returned as \B{Root} objects---exact algebraic numbers, indexed
canonically.

\begin{codepairs}
\pair{algebra/solve-poly/1}{The general quadratic returns the quadratic formula, its two
roots exact in $a$, $b$, $c$.}
\pair{algebra/solve-poly/2}{The irreducible cubic $x^3-x-1$ has no rational roots, so by
default its three solutions are returned as \B{Root} objects.}
\pair{algebra/solve-poly/3}{The quintic $x^5-x-1$ returns five \B{Root} objects---and
here that is not a matter of taste but of mathematics: no radical formula exists.}
\pair{algebra/solve-poly/4}{A \B{Root} object is a first-class value: ``the first root of
$\#^5-\#-1$.''}
\pair{algebra/solve-poly/5}{\B{N} evaluates it to any precision, isolating the one real
root near $1.1673$.}
\pair{algebra/solve-poly/6}{\B{ToRadicals} re-expresses a root in radicals where one
exists---here the real cube root of $2$.}
\end{codepairs}

Why are the cubic and quartic returned as \B{Root} objects by default when a radical
formula \emph{does} exist for them? Because the formula is almost always monstrous. Ask
for it explicitly with the \texttt{Cubics} option and you see why:

\mtranscript{algebra/solve-cardano}

That is the same cubic $x^3-x-1$, now via \emph{Cardano's
formula}\index{Cardano's formula}, and the \B{Root} objects of pair~2 are simply a
readable name for those three numbers.\calloutnote{Theory}{Radical formulas exist for the
general polynomial up to degree four---Cardano's for the cubic (1545) and Ferrari's for
the quartic---but no further. \emph{The Abel--Ruffini theorem}\index{Abel--Ruffini
theorem} states that the general polynomial of degree five or more has no solution in
radicals: there is simply no formula in $+,-,\times,\div$ and $\sqrt[n]{\ }$ for its
roots. This is why a \B{Root} object is not a placeholder for an answer Mathilda has not
bothered to find, but frequently \emph{the} answer---the only exact, finite name the
number has. Each \B{Root}$[f, k]$ denotes a specific root of $f$ under a canonical
indexing (real roots in increasing order, then complex ones), so it is a well-defined
value that \B{N} can evaluate and arithmetic can combine.} The quartic behaves the same
way: \texttt{Quartics -> True} produces Ferrari's radical solution, which is larger
still and best left named.

\subsection{Solving a system}
\label{ssec:alg-solve-systems}

For a \emph{system} of polynomial equations, \B{Solve} calls on the Gr\"obner machinery
of \S\ref{ssec:alg-groebner}. It computes a lexicographic basis, which for a system with
finitely many solutions is triangular, then solves the resulting univariate polynomial
and back-substitutes---variable by variable---to recover every solution.

\mtranscript{algebra/solve-systems}

The first system has the two obvious solutions. The second is the beautiful one: the
three symmetric functions of $x,y,z$ are fixed at $6$, $11$, and $6$, which are the
coefficients of $t^3-6t^2+11t-6=(t-1)(t-2)(t-3)$, so the solutions are exactly the six
orderings of $\{1,2,3\}$. The third line exposes the mechanism: the lexicographic
Gr\"obner basis of that system has a generator $z^3-6z^2+11z-6$ in $z$ alone, then one
coupling $y$ to $z$, then one giving $x$ in terms of the rest---the triangular form that
\B{Solve} walks down.\calloutnote{Under the hood}{The pipeline (\texttt{src/solve/solvenlsys.c})
clears each equation to a numerator, computes the lexicographic basis via the Gr\"obner
walk, and---when the ideal is zero-dimensional---peels off the univariate generator in
the last variable, solves it with the single-polynomial solver of the previous
subsection, substitutes each root back, and recurses. It is careful to be
\emph{sound}: a positive-dimensional system (a curve or surface of solutions rather than
finitely many points) is reported as such rather than answered with a misleadingly
finite list, and a branch it cannot triangulate declines rather than silently dropping
solutions.} Parametric systems, with symbols that are not among the solving variables,
are handled by computing the basis over the field of rational functions in those
parameters.

\subsection{Complete solutions and the reals}
\label{ssec:alg-reduce}

\B{Solve} returns the \emph{generic} solutions, as a list of rules. But the complete
truth about an equation is often a case analysis, and that is what \B{Reduce} delivers:
not a list of solutions but a logical combination of equations and inequalities
describing the entire solution set, including the degenerate branches \B{Solve} passes
over.

\begin{codepairs}
\pair{algebra/reduce/1}{The equation $ax=b$ has the solution $x=b/a$---but only when
$a\ne0$. \B{Reduce} states the complete truth, including the degenerate branch where
$a=0$ and the equation holds for all $x$ iff $b=0$.}
\pair{algebra/reduce/2}{For a genuine polynomial it gives the roots as a disjunction, the
two branches of the quadratic formula.}
\pair{algebra/reduce/3}{Over the reals it solves \emph{inequalities}: $x^2>2$ holds
exactly outside $[-\sqrt2,\sqrt2]$.}
\pair{algebra/reduce/4}{And it describes a two-dimensional region---the right half of the
unit circle---as a cell: $x$ ranges over $(0,1]$ and, for each $x$, $y$ takes the two
values $\pm\sqrt{1-x^2}$.}
\pair{algebra/reduce/5}{\B{SolveAlways} answers the dual question---for which parameter
values does $ax+b=0$ hold for \emph{every} $x$?---namely $a=b=0$.}
\end{codepairs}

The distinction is worth stating plainly: \B{Solve} is for the generic answer you want
most of the time, while \B{Reduce} is for the complete, guaranteed-exhaustive
description, degenerate cases and all. Over the reals \B{Reduce} does something genuinely
deep.\calloutnote{Theory}{Tarski proved that the theory of \emph{real-closed fields}
admits \emph{quantifier elimination}\index{quantifier elimination}: any statement built
from polynomial equations and inequalities with $\forall$ and $\exists$ is equivalent to
a quantifier-free one. Collins made it effective with \emph{cylindrical algebraic
decomposition}\index{cylindrical algebraic decomposition} (CAD): project the polynomials
down to one variable---using discriminants, resultants, and leading coefficients to track
where roots appear, merge, or vanish---build a sign diagram there, then lift it back up
one variable at a time, splitting each fibre into cells on which every polynomial has
constant sign. Mathilda uses McCallum's projection (\texttt{src/solve/reduce\_cad.c}),
and it holds itself to a hard soundness rule: if any step is undecidable---an algebraic
sign it cannot resolve, a fibre it cannot isolate---it declines rather than emit a
formula that might be wrong.} The third positional argument chooses the \texttt{Domain}:
\texttt{Reals}, \texttt{Complexes} (the default), or \texttt{Rationals}. Solving over the
\emph{integers}---Diophantine equations---is a different world of methods, and we take it
up in \S\ref{sec:numbertheory}.

That is the arc of algebra in Mathilda: from reading a polynomial's structure, through
expanding and factoring it---over the rationals and over the number fields you choose to
work in---through the GCDs, resultants, and Gr\"obner bases that eliminate variables, to
solving single equations, systems, and finally the complete real solution sets that
\B{Reduce} describes. The rational-function normal forms established here feed directly
into the integration of \S\ref{sec:calculus}, and the more general art of coaxing an
expression into its simplest form---\texttt{Simplify} and its relatives---builds on all
of it, and is taken up later. As always, the exhaustive options for any function named
here are one \texttt{?Name} away at the prompt.

\section{Calculus}
\label{sec:calculus}

Algebra computes with expressions as static objects; calculus computes with the
\emph{change} in them. This is the longest section of the book, because it is
where a computer algebra system does its most characteristic work---turning the
limit processes of the calculus into finite, exact algebra. We climb in the
order the subject is usually taught, but with the machinery on show throughout:
differentiation, limits, power series, residues, then the two faces of
integration---first the indefinite integral and the deep decision problem
beneath it, then the definite integral and the complex-analytic and
transform methods that evaluate it---and finally differential equations, where
all of the preceding tools are put to work at once.

The theme that runs through everything is a single distinction. \emph{Some}
operations of calculus are algorithmic in the strict sense: there is a procedure
that always terminates with the right answer or a proof that none exists.
Differentiation is the obvious one---the chain rule is a finite recursion on the
structure of an expression---but so, remarkably, is integration in elementary
terms, by the Risch algorithm. \emph{Other} operations are not decidable at all,
and there the system is a toolbox of specialised engines, each complete on the
family it recognises and honest enough to decline the rest. You will see both
kinds, and the book will be explicit about which is which.

\subsection{Differentiation}
\label{ssec:calc-deriv}

Differentiation\index{differentiation} is the one operation in this section that
is purely mechanical, and Mathilda treats it accordingly. \B{D} computes partial
derivatives, holding every symbol other than the differentiation variable
constant.

\begin{codepairs}
\pair{calculus/derivatives/1}{\B{D}\texttt{[f, x]} differentiates \(f\) with
respect to \(x\): the power rule, \(3x^2\).}
\pair{calculus/derivatives/2}{The chain rule for a composite, \(\frac{d}{dx}\sin
x^2 = 2x\cos x^2\).}
\pair{calculus/derivatives/3}{\texttt{D[f, \{x, n\}]} takes the \(n\)-th
derivative---here the third, \(-a^3\cos ax\).}
\pair{calculus/derivatives/4}{The base-\(b\) logarithm: \(\frac{d}{dx}\log_b x =
1/(x\log b)\).}
\pair{calculus/derivatives/5}{For an \emph{unknown} function the chain rule
produces a \B{Derivative} object: \(f'(g(x))\,g'(x)\), written
\texttt{Derivative[1][\ldots]}.}
\pair{calculus/derivatives/6}{\B{Derivative}\texttt{[2][f]} is the second-derivative
operator itself, carried symbolically until \(f\) is known.}
\end{codepairs}

The last two lines expose the internal representation. An \(n\)-th derivative of
an unknown function is not left as \(D[\ldots]\); it is stored as the operator
\texttt{Derivative[n][f]} applied to its argument, so that
\texttt{Derivative[1][f][g[x]]} means ``\(f'\) evaluated at \(g(x)\).'' This is
what lets Mathilda differentiate the same expression again and again and always
recognise what it is looking at.

Handed a list of variables, \B{D} produces the objects of vector calculus
directly. A single bracketed variable list is a gradient; a repeated one is a
higher array derivative; a vector argument is differentiated component-wise.

\begin{codepairs}
\pair{calculus/derivatives/7}{\texttt{D[f, \{\{x, y\}\}]} is the
gradient\index{gradient} \(\nabla f = (f_x, f_y)\).}
\pair{calculus/derivatives/8}{Repeating the variable list gives the second array
derivative---the Hessian\index{Hessian matrix} \(\nabla^2 f\).}
\pair{calculus/derivatives/9}{For a \emph{vector} field the same call yields the
Jacobian\index{Jacobian matrix}, one gradient row per component.}
\pair{calculus/derivatives/10}{\B{Dt} is the \emph{total} derivative: every
symbol is an implicit function of everything, so each contributes a
\texttt{Dt} differential.}
\pair{calculus/derivatives/11}{So \texttt{Dt} of \(x^n\) carries a \(\log\) term
from the exponent's own dependence.}
\end{codepairs}

The difference between \B{D} and \B{Dt}\index{total derivative} is exactly the
difference between a partial and a total derivative in an ordinary calculus
course: \texttt{D[y, x]} is \(0\) because \(y\) is held fixed, while
\texttt{Dt[y, x]} is the differential \texttt{Dt[y, x]} because \(y\) may depend
on \(x\). Give a function a definition, and its derivative computes to an
ordinary expression:

\begin{codepairs}
\pair{calculus/derivatives/13}{Once \(f(x)=x^5+6x^3\) is defined,
\texttt{f'[x]} evaluates the operator: \(18x^2+5x^4\).}
\pair{calculus/derivatives/14}{And \texttt{f'[5]} is a plain number.}
\end{codepairs}

\calloutnote{Under the hood}{Differentiation is the textbook example of the
system's two-tier design---C for speed, rewrite rules for mathematics---and of
that design evolving. The derivative rules once lived as pattern-matching
\texttt{DownValues} in \texttt{src/internal/deriv.m}; every call walked a linear
list of some sixty rules through the full rule engine. They were promoted to a
native C dispatcher (\texttt{src/calculus/deriv.c}) that branches on the head
symbol in a single comparison and builds the derivative tree directly, for a
2--8\(\times\) speedup, and \texttt{deriv.m} is now an empty hook you can still
drop custom \texttt{Dt} identities into. The higher tier survives where it earns
its keep: the integral tables (\texttt{CRCMathTablesIntegrals.m}) and the mixed-tower
integrator (\texttt{src/internal/mixed/ParallelMixed.m}) below are still Mathilda
programs, loaded at startup.}

The vector operators \B{Grad}, \B{Div}, \B{Curl}, and \B{Laplacian} are assembled
entirely from \B{D}, so they inherit every rule it knows.\index{vector calculus}

\begin{codepairs}
\pair{calculus/vector-calculus/1}{\B{Grad} of a scalar is its gradient vector.}
\pair{calculus/vector-calculus/2}{\B{Grad} of a vector field is its Jacobian.}
\pair{calculus/vector-calculus/3}{\B{Div} contracts a field against the variables:
\(\nabla\cdot\mathbf{F}\).}
\pair{calculus/vector-calculus/4}{\B{Curl} of \((-y, x, 0)\) is the constant
vorticity \((0,0,2)\).}
\pair{calculus/vector-calculus/5}{\B{Laplacian} of \(x^2+y^2+z^2\) is \(6\).}
\pair{calculus/vector-calculus/6}{With a chart name the operators use the
orthonormal basis---here the Coulomb field \(-\nabla(kq/r)= (kq/r^2)\,\hat r\) in
spherical coordinates.}
\pair{calculus/vector-calculus/7}{And the spherical divergence of \(r^2\hat r\)
is \(4r\), from \(\frac{1}{r^2}\partial_r(r^2\cdot r^2)\).}
\end{codepairs}

\calloutnote{Theory}{The three-argument chart form does the real differential
geometry. In a curvilinear orthogonal chart the operators are not just
partial derivatives: each carries the chart's scale factors (Lam\'e
coefficients) \(h_i\), so the spherical divergence is
\(\frac1{J}\sum_i \partial_i\!\big(\tfrac{J}{h_i}F_i\big)\) with \(J=\prod_i
h_i\). Mathilda knows the Cartesian, polar, cylindrical, and spherical charts;
field ranks that would need Christoffel symbols (a vector Laplacian, the gradient
of a vector field in a curved chart) are left unevaluated rather than guessed.}

\subsection{Limits}
\label{ssec:calc-limit}

\B{Limit} evaluates the limit of an expression as a variable approaches a point,
finite or infinite. The easy cases are the ones a first course drills; the hard
ones are where a CAS earns its name.

\begin{codepairs}
\pair{calculus/limits/1}{The archetypal \(0/0\) form: \(\sin x / x \to 1\).}
\pair{calculus/limits/2}{A removable singularity---the difference quotient of
\(x^2\) at \(1\).}
\pair{calculus/limits/3}{Another \(0/0\), giving \(\tfrac12\).}
\pair{calculus/limits/4}{The definition of \(e\), as a \(1^\infty\) form.}
\pair{calculus/limits/5}{With no direction, \(1/x\) at \(0\) is
\mcode{ComplexInfinity}---the two sides disagree in sign.}
\pair{calculus/limits/6}{Fixing \texttt{Direction} takes the one-sided limit,
here \(+\infty\) from above.}
\end{codepairs}

The default is the two-sided limit; \texttt{Direction -> "FromAbove"} and
\texttt{"FromBelow"} select a one-sided approach, and the answer
\mcode{ComplexInfinity} in pair~5 is Mathilda telling you the unsigned limit is
infinite because the signed one does not exist. So far this is bookkeeping.
The following are not:

\begin{codepairs}
\pair{calculus/limits/7}{A \(\infty^0\) form whose value is the largest base,
\(5\)---no series expansion reaches this.}
\pair{calculus/limits/8}{A tower of nested exponentials whose leading terms
cancel to leave \(-1\).}
\pair{calculus/limits/9}{An oscillation that genuinely has no limit:
\mcode{Indeterminate}, not a guessed value.}
\end{codepairs}

Pairs~7 and~8 are the province of a genuine algorithm.\calloutnote{Theory}{The
last resort for a limit of an exp--log function is \emph{Gruntz's most-rapidly-varying
algorithm}\index{Gruntz's algorithm} (his 1996 ETH thesis, implemented in
\texttt{src/calculus/gruntz.c}). Rather than expand from the bottom up---where the
intermediate expressions swell catastrophically on cancellation-heavy towers like
pair~8---it identifies the single most rapidly varying subexpression \(w\), rewrites
the whole function as a series in \(w\to 0^+\), and reads the limit off the leading
term. Because the rewriting is driven by the comparability classes of the
subexpressions, the swell never happens: \((3^x+5^x)^{1/x}\) collapses at once to
\(5\), the base that dominates. Gruntz also runs as the silent fallback inside the
default cascade, after the substitution, rational, series, and L'Hospital layers
have had their turn.} Pair~9 is the honest counterpart: \(x\sin x\) oscillates
with growing amplitude, so no limit exists, and Mathilda proves this (via an
oscillatory normal-form analysis) rather than returning a plausible-looking
number. That refusal to guess is the through-line of the whole section.

Computing a limit is one thing; \emph{proving} one from its definition is another,
and Mathilda does both. The statement \(\lim_{x\to a} f(x) = L\) unfolds into the
\(\varepsilon\)--\(\delta\) sentence\index{limit!epsilon-delta proof of@\(\varepsilon\)--\(\delta\) proof of}\index{epsilon-delta@\(\varepsilon\)--\(\delta\) definition}
\[
  \forall\,\varepsilon>0 \quad \exists\,\delta>0 \quad \forall x \quad
  \bigl(0<\lvert x-a\rvert<\delta \;\Rightarrow\; \lvert f(x)-L\rvert<\varepsilon\bigr),
\]
a quantified formula over the reals---and \B{Resolve} decides quantified formulas
over the reals. Transcribed verbatim, with \B{ForAll} and \B{Exists} carrying each
bound variable and its restriction just as the quantifiers do on the page, the
definition of \(\lim_{x\to2}(3x-1)=5\) is proved outright---and the same sentence
with the limit misstated is \emph{disproved}, not merely left unconfirmed:

\begin{codepairs}
\pair{calculus/limits/10}{The \(\varepsilon\)--\(\delta\) definition of
\(\lim_{x\to2}(3x-1)=5\), transcribed literally and returned \mcode{True}: a proof.}
\pair{calculus/limits/11}{The identical sentence with \(5\) replaced by the wrong
value \(6\)---\mcode{False}. The decision is two-sided.}
\end{codepairs}

Behind the verdict is quantifier elimination, not sampling.\calloutnote{Under the
hood}{An \(\varepsilon\)--\(\delta\) statement is three alternating quantifiers over
the reals, and settling it is \emph{quantifier
elimination}\index{quantifier elimination}: \B{Resolve} eliminates the innermost
block to a condition on the variables that remain, then the next, until only
\mcode{True} or \mcode{False} is left. The engine underneath
(\texttt{src/solve/reduce\_cad.c}) is a \emph{cylindrical algebraic
decomposition}\index{cylindrical algebraic decomposition}---a complete decision
procedure for the first-order theory of real-closed
fields\index{real-closed field} (Tarski's theorem)---so for any
\emph{semialgebraic} \(f\) (polynomial, rational, or algebraic) the answer is
guaranteed, and each \(\lvert\cdot\rvert\) is absorbed by splitting the absolute
value into its two sign branches. That completeness has a sharp boundary: a
transcendental target such as \(\lim_{x\to0}\sin(x)/x = 1\) has no polynomial
decomposition, so Mathilda \emph{declines}---returns the sentence
unevaluated---rather than guess. \B{Limit} and \B{Resolve} are complementary:
\B{Limit} reaches the transcendental limits a decision procedure cannot, while
\B{Resolve} certifies the algebraic ones with a proof.}

\usagebox{Limit}

\subsection{Series}
\label{ssec:calc-series}

A power series is the local linearisation carried to all orders. \B{Series}
computes the Taylor, Laurent, or Puiseux expansion of an expression about a
point, to a requested order.

\begin{codepairs}
\pair{calculus/series/1}{The Taylor series\index{Taylor series} of \(e^x\) to
order five, with its \(O[x]^6\) remainder term.}
\pair{calculus/series/2}{\B{Normal} strips the \(O\) term, leaving an ordinary
polynomial you can compute with.}
\pair{calculus/series/3}{\B{SeriesCoefficient} extracts a single coefficient
without forming the whole series---here \(1/7!\).}
\pair{calculus/series/4}{A \emph{Laurent}\index{Laurent series} series: \(1/\sin
x\) has a pole at \(0\), so the expansion starts at \(x^{-1}\).}
\pair{calculus/series/5}{A \emph{Puiseux}\index{Puiseux series} series in
fractional powers, for \(\sqrt{\sin x}\).}
\pair{calculus/series/6}{Symbolic exponents survive: the binomial series of
\((1+x)^n\) keeps \(n\) unevaluated.}
\end{codepairs}

The three expansion types are one mechanism seen at three scales: a Taylor series
when the function is analytic, a Laurent series with negative powers at a pole, a
Puiseux series with fractional powers at a branch point. The expansion can also be
taken at infinity, and an unknown function expands formally through its
derivatives:

\begin{codepairs}
\pair{calculus/series/7}{An asymptotic\index{asymptotic expansion} expansion at
\(\infty\): the arctangent tends to \(\pi/2\) with a \(1/x\) tail.}
\pair{calculus/series/8}{For an unknown \(f\) the coefficients appear as
\B{Derivative} values---the Taylor series in symbolic form.}
\pair{calculus/series/9}{The internal object is \B{SeriesData}; constructing one
by hand prints as the same series.}
\end{codepairs}

\calloutnote{Under the hood}{Behind the printed sum is a \B{SeriesData} object
carrying the expansion variable, the centre, the coefficient list, the leading and
trailing exponents, and a denominator that makes fractional powers exact integers
internally. It is a first-class value: \B{D}, \B{Integrate}, \B{Plus}, \B{Times},
and \B{Power} all operate on it term by term and return a new truncated series, so
you can do arithmetic with series as freely as with polynomials. Coefficients
promote automatically to arbitrary-precision rationals when machine integers would
overflow, so a deep Laurent expansion such as \(\operatorname{Series}[1/\sin^{10}x]\)
stays exact.}

\subsection{Residues}
\label{ssec:calc-residue}

The residue of a function at an isolated singularity---the coefficient of
\((z-z_0)^{-1}\) in its Laurent series---is the single number that complex
analysis needs to integrate around a pole. \B{Residue} reads it straight out of
the series.

\begin{codepairs}
\pair{calculus/residues/1}{The residue\index{residue} of \(1/z\) at \(0\) is
\(1\).}
\pair{calculus/residues/2}{A double pole with no simple part: \(1/z^2\) has
residue \(0\).}
\pair{calculus/residues/3}{\(\cot z\) has a simple pole at \(0\) with residue
\(1\).}
\pair{calculus/residues/4}{A higher-order pole: residue \(-3/4\) at the double
pole \(z=0\).}
\pair{calculus/residues/5}{A complex pole---\(1/(z^2+1)\) at \(z=i\) has residue
\(-i/2\).}
\pair{calculus/residues/6}{A fifth-order pole: \(1/\sin^5 z\) at \(0\)
gives \(3/8\).}
\pair{calculus/residues/7}{At a \emph{branch} point the residue is undefined, so
the call is left unevaluated.}
\end{codepairs}

\calloutnote{Under the hood}{For a general integrand the residue is the \(-1\)
coefficient of \B{Series}; the engine raises the expansion order automatically
when the pole's order is hidden behind an unknown factor. But a rational function
with an algebraic (radical) pole location would trip the series inverter, so a
simple pole short-circuits to the exact formula \(P(z_0)/Q'(z_0)\) computed
directly, and higher-order rational poles expand the numerator and denominator
\emph{separately} to keep each algebraic coefficient in its minimal radical
spelling. A fractional-power (Puiseux) expansion signals a branch point---pair~7---where
the notion of residue does not apply, so the call declines rather than return a
Laurent coefficient that is not there.} The residue theorem this feeds---summing
residues to evaluate a real integral---reappears in \S\ref{ssec:calc-definite},
where Mathilda uses it as a definite-integration engine in its own right.

\subsection{Indefinite integration: the Risch story}
\label{ssec:calc-integrate}

Integration is the centre of this section and the deepest single algorithm in the
system. Where differentiation is a finite recursion, integration in closed form
is a genuine decision problem, and the theory that settles it is one of the
triumphs of twentieth-century computer algebra. We start with the cases that look
familiar and work toward the machinery that makes them possible.

\begin{codepairs}
\pair{calculus/integrate-elementary/1}{Term-by-term integration of a polynomial.}
\pair{calculus/integrate-elementary/2}{A logarithmic derivative recognised
directly: \(\int 2x/(x^2+1) = \log(1+x^2)\).}
\pair{calculus/integrate-elementary/3}{The arctangent, from an irreducible
quadratic denominator.}
\pair{calculus/integrate-elementary/4}{A partial-fraction case; the answer combines
two of the three logarithms into a single quadratic \(\log\).}
\pair{calculus/integrate-elementary/5}{Two irreducible quadratic factors give a
mixed \(\log\)/\(\arctan\) answer.}
\end{codepairs}

These are all \emph{rational} integrands, and for them integration is completely
solved.\calloutnote{Under the hood}{The rational integrator
(\texttt{src/calculus/intrat.c}) is a faithful implementation of the modern
pipeline: \emph{Hermite reduction}\index{Hermite reduction} splits off the part
of the answer that is itself rational (the repeated-pole part), leaving a
residual with a square-free denominator; the \emph{Lazard--Rioboo--Trager}
algorithm\index{Rothstein--Trager} then produces the logarithmic part as a sum
\(\sum_i c_i\log(g_i)\) whose constants \(c_i\) are the roots of a
resultant---never by an undetermined-coefficient solve. Rioboo's conversion turns
each complex-conjugate pair of logarithms into a real \(\arctan\), which is why
pair~5 comes out in real form. For plain rational inputs the heavy polynomial
arithmetic is handed to FLINT~\cite{flint}; the classical account is
Geddes--Czapor--Labahn~\cite{gcl}. Every branch ends at the same universal check:
\(\operatorname{Cancel}[\operatorname{Together}[D[\text{result}]-f]]\) must be
\(0\).} Rational functions of radicals are next, rationalised by the right
substitution:

\begin{codepairs}
\pair{calculus/integrate-elementary/6}{A rational function of \(\sqrt x\),
rationalised by \(u=\sqrt x\).}
\pair{calculus/integrate-elementary/7}{A square root of a quadratic, handled by a
Euler substitution.}
\end{codepairs}

Now the mathematics deepens. Many perfectly innocent-looking integrands have
\emph{no} antiderivative among the elementary functions---and this is a theorem,
not a gap. \emph{Liouville's theorem}\index{Liouville's theorem} says that if an
elementary function has an elementary antiderivative, that antiderivative has a
very restricted shape: the original function plus a sum of constant multiples of
logarithms of elementary functions. When no such shape exists, the integral is
provably non-elementary, and the honest answer is a \emph{new} function named for
exactly that integral.

\begin{codepairs}
\pair{calculus/integrate-special/1}{Integration by parts, done for you:
\(\int x e^x = (x-1)e^x\).}
\pair{calculus/integrate-special/2}{\(e^{-x^2}\) has no elementary integral, so
the answer is the error function \B{Erf}.}
\pair{calculus/integrate-special/3}{\(e^x/x\) gives the exponential integral
\B{ExpIntegralEi}.}
\pair{calculus/integrate-special/4}{\(1/\log x\) is the logarithmic integral,
returned as \(\operatorname{Ei}(\log x)\).}
\pair{calculus/integrate-special/5}{\(\sin x/x\) gives the sine integral
\B{SinIntegral}.}
\pair{calculus/integrate-special/6}{A dilogarithm\index{dilogarithm} appears:
\(\int \log x/(1+x) = \log x\log(1+x) + \operatorname{Li}_2(-x)\).}
\pair{calculus/integrate-special/7}{And a genuine polylogarithm ladder for
\(x/(e^x-1)\), the Bose--Einstein integrand.}
\end{codepairs}

Each of these is correct \emph{by construction}: the algorithm certifies that no
elementary form exists before it reaches for the special function, so \B{Erf},
\B{ExpIntegralEi}, \B{SinIntegral}, and \B{PolyLog} are not approximations or
placeholders---they are the exact antiderivatives, and their appearance is a
statement about the integrand.\calloutnote{Theory}{The engine behind pairs~2--7
is the \emph{Risch algorithm}\index{Risch algorithm}
(\texttt{src/calculus/integrate\_risch\_transcendental.c}), the decision
procedure for integration in a differential field built by adjoining
exponentials and logarithms one at a time. Its heart is the
\emph{Risch differential equation}: closing the exponential part of an integral
amounts to solving \(q' + f'q = p\) for a polynomial \(q\), which Mathilda does by
Bronstein's polynomial-time SPDE reduction rather than an undetermined-coefficient
blow-up. Layered above the base algorithm are engines for the special-function
answers: a Gaussian recogniser for \B{Erf}, Cherry's algorithms for the exponential
and logarithmic integrals and the \emph{polylogarithm ladder}\index{polylogarithm
ladder} \(\int x^n/(e^{cx}-\rho)\), and a dilogarithm recogniser. When the field
decision proves an integral non-elementary but the special-function engines cannot
name it, the integral is returned unevaluated with an \texttt{Integrate::nonelem}
message explaining why.}

The book promised to be explicit about what is decidable, and integration lets it
keep the promise directly. The predicate that decides elementary integrability is
exposed:

\begin{codepairs}
\pair{calculus/elementary-q/1}{\(e^x/x\) is \emph{not} elementary---a proof, not a
timeout.}
\pair{calculus/elementary-q/2}{Nor is \(e^{x^2}\).}
\pair{calculus/elementary-q/4}{But \(e^x/(1+e^x)\) is: its integral is
\(\log(1+e^x)\).}
\pair{calculus/elementary-q/5}{And \(\sin x/x\) is not---the sine integral is the
statement that it is not.}
\end{codepairs}

A \texttt{False} here is backed by an exact certificate---a non-constant residue,
or a Risch differential equation with no rational solution---so it means ``no
elementary antiderivative exists,'' never ``I gave up.'' This soundness is
purchased at a real cost, which is worth seeing:

\begin{codepairs}
\pair{calculus/elementary-q/6}{A correct answer in an awkward dress:
\(\int\cos(\log x)\) is real, but printed through complex powers of \(x\).}
\end{codepairs}

\calloutnote{Pitfall}{The output of pair~6 is exactly
\(\tfrac{x}{2}(\cos\log x + \sin\log x)\)---but Mathilda returns it as
\((\tfrac14-\tfrac{i}4)x^{1+i}+(\tfrac14+\tfrac{i}4)x^{1-i}\), a genuinely real
quantity written with conjugate complex exponents. The transcendental integrator
reaches trigonometric integrands by exponentialising with \(\cos =
(e^{i\cdot}+e^{-i\cdot})/2\), and the current simplifier does not always fold the
result back to real form. The value is right and \(\B{N}\) confirms it; the form is
a known cosmetic limitation, not an error. Prefer the algebraic routes when a real
closed form matters.}

Two further capabilities of \B{Integrate} deserve mention, because they reach
integrands the classical Risch algorithm alone cannot. The first is the
\texttt{"ParallelMixedTower"}\index{ParallelMixedTower method} method, a parallel
Risch--Norman integrator over a \emph{mixed} transcendental tower that may contain
a square root anywhere in its structure. It is what closes hard algebraic-function
integrals such as the classic \(\int\sqrt{\tan x}\):

\begin{codepairs}
\pair{calculus/integrate-special/8}{\(\int\sqrt{\tan x}\,dx\): a mixed radical
tower, closed in real \(\arctan\)/\(\log\) form.}
\end{codepairs}

The parallel method treats the whole integrand as living in a differential field
built over rational operations, one square root, and \(\log\)/\(\exp\)/\(\tan\) of
field elements, and solves for the antiderivative by an ansatz whose unknowns are
found by linear algebra---then verifies every result by re-differentiating before
returning it, so a case it cannot realise correctly declines rather than escaping
as a wrong answer. The second is Feynman's trick,
\texttt{"DiffUnderInt"}---differentiation under the integral sign---which belongs
to definite integration and is taken up next.

\usagebox{Integrate}

\subsection{Definite integration}
\label{ssec:calc-definite}

A definite integral is not just an antiderivative evaluated at two points---not
when the integrand has poles, the interval runs to infinity, or the
antiderivative crosses a branch cut. Mathilda's definite integrator is a cascade
of methods, each correct on the family it claims. The workhorse is the
fundamental theorem of calculus, made careful.

\begin{codepairs}
\pair{calculus/definite/1}{The elementary case: \(\int_0^1 x^2 = 1/3\).}
\pair{calculus/definite/2}{An \emph{improper} integral that converges---the
singularity at \(0\) is integrable.}
\pair{calculus/definite/3}{An infinite upper limit, evaluated through the
\B{Limit} engine.}
\pair{calculus/definite/11}{An iterated double integral, reduced innermost-first.}
\end{codepairs}

\calloutnote{Under the hood}{The Newton--Leibniz mechanism
(\texttt{src/calculus/integrate\_newton\_leibniz.c}) antidifferentiates \(f\),
locates the real poles of the \emph{integrand} strictly inside the interval, splits
the interval there, and evaluates each boundary through \B{Limit}---a plain limit at
infinite endpoints, a one-sided limit at interior poles. This is how the improper
integrals above acquire their finite values, and how a genuinely divergent
integral is caught and left unevaluated rather than assigned the difference of two
infinities.} That last safeguard matters:

\begin{codepairs}
\pair{calculus/definite/10}{\(\int_{-1}^1 dx/x\) diverges, and Mathilda says so by
returning the input unevaluated (with an \texttt{Integrate::idiv} message).}
\end{codepairs}

For the classical improper and periodic integrals that complex analysis dispatches
by summing residues, a residue-theorem engine runs \emph{before} Newton--Leibniz.
These answers are correct by construction---once the structural gates hold, the
residue theorem gives the exact value.

\begin{codepairs}
\pair{calculus/definite/4}{A rational function over the whole line, by residues in
the upper half-plane: \(\pi/\sqrt2\).}
\pair{calculus/definite/5}{A Fourier integral, by Jordan's lemma\index{Jordan's
lemma}: \(\int_{-\infty}^\infty \cos x/(1+x^2) = \pi/e\).}
\pair{calculus/definite/9}{A rational function of \(\cos x\) over a full period, via
\(z=e^{ix}\) on the unit circle.}
\end{codepairs}

A third family---half-line integrals of transcendental functions that neither
residues nor the FTC close---is handled by a Mellin-transform engine built on the
Ramanujan Master Theorem\index{Ramanujan Master Theorem}, and a fourth by
symmetry.

\begin{codepairs}
\pair{calculus/definite/6}{The Gaussian \(\int_0^\infty e^{-x^2} = \sqrt\pi/2\).}
\pair{calculus/definite/7}{The sinc integral \(\int_0^\infty \sin x/x = \pi/2\).}
\pair{calculus/definite/8}{The Debye integral \(\int_0^\infty x^3/(e^x-1) =
\pi^4/15\)---the Stefan--Boltzmann law, from the Bose--Einstein kernel.}
\end{codepairs}

\calloutnote{Theory}{Pair~8 is Ramanujan's Master Theorem\index{Mellin transform}
at work. If \(f(x)=\sum_k (-1)^k\varphi(k)x^k/k!\) then its Mellin transform
\(\int_0^\infty x^{s-1}f(x)\,dx\) equals \(\Gamma(s)\varphi(-s)\) on a
fundamental strip, and the Bose--Einstein kernel \(1/(e^{x}-1)\) lands on
\(\Gamma(s)\zeta(s)\)---so \(\int_0^\infty x^3/(e^x-1)=\Gamma(4)\zeta(4)=6\cdot
\pi^4/90=\pi^4/15\). Every application is gated on its convergence strip, checked
by \B{Simplify} or against your \texttt{Assumptions}; when the strip cannot be
proved the value is returned as a \B{ConditionalExpression}, matching Wolfram, and
a provably violated strip declines. The Fermi--Dirac companion
\(\int_0^\infty x^3/(e^x+1)=7\pi^4/120\) comes out the same way.}

When an endpoint is complex, or the spec lists more than two points, \B{Integrate}
switches to a genuine contour integral along the straight segments.

\begin{codepairs}
\pair{calculus/contour/1}{A contour integral of \(1/z\) between two complex
points---a difference of principal logarithms.}
\pair{calculus/contour/2}{\(\int_0^{1+i} z^2\,dz\), path-independent because
\(z^2\) is entire.}
\pair{calculus/contour/3}{A closed square loop around the origin gives \(2\pi
i\)---the residue theorem\index{residue theorem}, numerically.}
\end{codepairs}

The Cauchy principal value\index{Cauchy principal value} is available for the
odd-order interior poles where the symmetric divergences cancel:

\begin{codepairs}
\pair{calculus/contour/4}{The principal value across the pole at \(x=2\):
\(-\log 2\).}
\end{codepairs}

Finally, differentiation under the integral sign---\emph{Feynman's
trick}\index{Feynman's trick}\index{differentiation under the integral sign}---is
a definite-integration method in its own right. For a parameter-dependent integral
\(I(p)\), it differentiates the integrand in \(p\), evaluates the simpler integral,
and integrates back, fixing the constant from an exact base value.

\begin{codepairs}
\pair{calculus/feynman/1}{\(\int_0^1 (x^a-1)/\log x\,dx = \log(1+a)\), the standard
Feynman warm-up.}
\pair{calculus/feynman/2}{And \(\int_0^\infty \log(1+a^2x^2)/(1+x^2)\,dx = \pi\log(1+a)\)
under \(a>0\).}
\end{codepairs}

\subsection{Differential equations}
\label{ssec:calc-dsolve}

The whole apparatus of this section converges on \B{DSolve}, which solves
differential equations symbolically. It returns solutions as rules, with the
constants of integration named \texttt{C[1]}, \texttt{C[2]}, and so on; initial
or boundary conditions, given as equations, are fitted to those constants. Like
\B{Integrate}, it is a cascade polyalgorithm---a long sequence of specialised
methods, each recognising an equation type it can solve completely. We survey it
by that classification, which is also the classification any course on
differential equations teaches.

We begin, as promised, with the warm-up:

\begin{codepairs}
\pair{calculus/dsolve-first-order/1}{The exponential equation \(y'=y\): the
solution is \(C_1 e^x\).}
\pair{calculus/dsolve-first-order/2}{Separable, \(y'=xy\): \(C_1 e^{x^2/2}\).}
\pair{calculus/dsolve-first-order/3}{First-order linear, solved by an integrating
factor\index{integrating factor}.}
\pair{calculus/dsolve-first-order/4}{A Bernoulli\index{Bernoulli equation}
equation \(y'=y+xy^2\), linearised by \(v=y^{-1}\).}
\pair{calculus/dsolve-first-order/5}{An exact\index{exact equation} equation;
the potential inverts to the explicit \(y=C_1/(1+x^2)\).}
\pair{calculus/dsolve-first-order/6}{A Clairaut\index{Clairaut equation} equation,
whose general solution is a family of lines.}
\pair{calculus/dsolve-first-order/7}{A quadratic (Riccati\index{Riccati equation})
right-hand side, whose solution is a shifted tangent.}
\pair{calculus/dsolve-first-order/8}{An initial-value problem---the constant is
fitted to \(y(0)=5\).}
\end{codepairs}

Each line above is a different classical method: separable, linear, Bernoulli,
exact, Clairaut, Riccati. \B{DSolve} recognises the shape of the equation and
dispatches to the matching engine, and every branch it returns is verified by
back-substitution before you see it. The linear constant-coefficient case---the
backbone of a differential-equations course---is next, with the forcing handled by
undetermined coefficients or variation of parameters.

\begin{codepairs}
\pair{calculus/dsolve-linear/1}{A damped oscillator: complex characteristic roots
give \(e^{-2x}\) times sine and cosine.}
\pair{calculus/dsolve-linear/2}{Resonant forcing, \(y''-y=e^x\): the particular
solution carries the tell-tale factor of \(x\).}
\pair{calculus/dsolve-linear/3}{Repeated integration, in \B{Function} form.}
\pair{calculus/dsolve-linear/4}{A boundary-value problem, fitted at two points to
give \(\sin x\).}
\pair{calculus/dsolve-linear/5}{An \emph{inconsistent} BVP: \texttt{\{\}} means no
solution exists, distinct from an unhandled equation.}
\end{codepairs}

The resonance in pair~2 is worth pausing on: when the forcing frequency coincides
with a natural frequency, the naive particular solution fails, and \B{DSolve}
detects the coincidence and multiplies by \(x\) automatically---the same
resonance a physicist watches for in a driven oscillator. Pair~5 shows the
system's care with existence: an over-determined boundary-value problem returns
the empty list as a proof that no solution exists, which is a different statement
from leaving an equation symbolic because no method applied.

Beyond constant coefficients lies the variable-coefficient world, where the
solutions are the named special functions of \S\ref{ssec:calc-series} and the
answer often hinges on a deep classical algorithm.

\begin{codepairs}
\pair{calculus/dsolve-special/1}{The Airy equation \(y''=xy\), whose solutions are
\B{AiryAi} and \B{AiryBi}.}
\pair{calculus/dsolve-special/2}{Bessel's equation of order \(2\), giving
\B{BesselJ} and \B{BesselY}.}
\pair{calculus/dsolve-special/3}{An Euler--Cauchy\index{Euler--Cauchy equation}
equidimensional equation, solved by \(x^r\) with complex exponents.}
\pair{calculus/dsolve-special/4}{A Liouvillian solution found by Kovacic's
algorithm---elementary, not a series.}
\pair{calculus/dsolve-special/5}{No closed form exists, so the fallback is a
Frobenius\index{Frobenius method} power series about the origin.}
\end{codepairs}

\calloutnote{Theory}{Pair~4 is the crown of second-order theory.
\emph{Kovacic's algorithm}\index{Kovacic's algorithm}
(\texttt{src/calculus/dsolve\_kovacic.c}) decides whether a second-order linear
ODE \(y''+Py'+Qy=0\) has a \emph{Liouvillian} solution---one built from exponentials,
integrals, and algebraic functions---and constructs it when it does, by examining
the equation's reduced normal form \(z''=rz\) and the local behaviour of \(r\) at
its poles and at infinity. When it succeeds you get a genuine closed form (here
\((C_1+C_2e^{x^2})e^{-x^2/2}/\sqrt x\)); when it proves none exists, control passes
to the Frobenius\index{Frobenius method} series method, which expands about a regular
singular point with an indicial equation fixing the leading exponents. The two
together---Kovacic for the closed form, Frobenius for its absence---cover the
second-order linear equations a course would set.}

Systems of equations are dispatched to their own cascade, which diagonalises the
coefficient matrix and handles real, complex, and defective spectra.

\begin{codepairs}
\pair{calculus/dsolve-systems/1}{A \(2\times2\) system with complex eigenvalues:
the solution is oscillatory.}
\pair{calculus/dsolve-systems/2}{A system with real eigenvalues \(-1\) and \(3\):
a combination of growing and decaying exponentials.}
\end{codepairs}

Two engines stand out for their generality. The first is the Lie point-symmetry
method, the general backstop for first-order equations.\calloutnote{Theory}{A
\emph{Lie point symmetry}\index{Lie point symmetry} is an infinitesimal
transformation \(V=\xi\,\partial_x+\eta\,\partial_y\) that maps solution curves to
solution curves. Finding one is an underdetermined PDE---so Mathilda's engine
(documented in \texttt{docs/design/dsolve\_lie\_symmetry.md}) runs a table of
ansatz heuristics, each proposing a shape for \((\xi,\eta)\) and testing it against
the linearised symmetry condition. An accepted symmetry hands over an integrating
factor \(\mu=1/(\eta-\xi\omega)\), whose quadrature is the first integral of the
equation. Because a first integral need not admit an explicit \(y(x)\), the answer
is often returned \emph{implicitly}, as a relation \(G(x,y)=C_1\).}

\begin{codepairs}
\pair{calculus/dsolve-lie/1}{A symmetry with a rational infinitesimal yields an
implicit first integral.}
\pair{calculus/dsolve-lie/2}{An irrational right-hand side \(\sqrt{ax+by+c}\),
handled by a similarity-ansatz symmetry.}
\end{codepairs}

The second is the partial-differential-equation solver, which handles first-order
equations by characteristics and the classical second-order equations of
mathematical physics.

\begin{codepairs}
\pair{calculus/dsolve-pde/1}{The transport equation, by the method of
characteristics\index{method of characteristics}: an arbitrary function of
\(x-ct\).}
\pair{calculus/dsolve-pde/2}{A first-order PDE with a lower-order term and forcing.}
\pair{calculus/dsolve-pde/3}{The wave equation as an initial-value problem, by
d'Alembert's formula\index{d'Alembert's formula}.}
\pair{calculus/dsolve-pde/5}{A Sturm--Liouville\index{Sturm--Liouville problem}
eigenvalue problem: the eigenvalues \(n^2\) and eigenfunctions \(\sin(nx)\).}
\end{codepairs}

The arbitrary function \texttt{C[1][\ldots]} in pairs~1 and~2 is the PDE analogue
of an integration constant---a first-order PDE's general solution contains an
arbitrary \emph{function}, not merely a constant, and it is a function of the
characteristic coordinate the equation is constant along. Pair~3 is d'Alembert's
formula for the vibrating string, and pair~5 is the eigenvalue quantisation that
underlies every separated boundary-value problem, returned as a
\B{ConditionalExpression} indexed by an integer. The heat equation is solved by
its Gaussian kernel; because that convolution is non-elementary for a general
initial condition, Mathilda keeps it as an integral rather than inventing a closed
form---the same honesty about non-elementarity that governs \B{Integrate}.

\usagebox{DSolve}

\subsection*{Where this connects}

Calculus is where the rest of Mathilda's mathematics is put to work at once.
Every integral of a rational function leans on the factoring, partial-fraction,
and resultant machinery of \S\ref{sec:algebra}; every \B{Root} object in an
antiderivative is an exact algebraic number the solver returns; the definite
integrals close through the same \B{Limit} engine that opens this section, and the
contour methods rest on \B{Residue}. \B{DSolve} in turn calls on all of it---on
\B{Integrate} for its quadratures, on the eigenvalue routines of \S\ref{sec:linalg}
for its systems, and on the special functions of \S\ref{ssec:calc-series} for its
closed forms. Where a symbolic answer does not exist, or exists but is too costly
to want, the numerical companions \B{NIntegrate}, \B{NDSolve}, and their relatives
in \S\ref{sec:numcalc} take over, and the numeric loops that evaluate a
series or a special function accelerate through the compiler of
Chapter~\ref{ch:compilation}. As always, the exhaustive options for any function
named here are one \texttt{?Name} away at the prompt.

\section{Linear Algebra}
\label{sec:linalg}

Linear algebra is where symbolic computation and fast numerics meet most
directly, and Mathilda meets them with a single idea: a matrix is not a special
kind of value. It is an ordinary expression---a \B{List} of \B{List}s---so the
generic tools you already know (\B{Part}, \B{Map}, \B{Transpose}) apply to it
without ceremony, and a routine like \B{Det} or \B{Inverse} simply reads the
structure off the tree.\index{matrix!as an expression} That uniformity is one of
the two themes of this section. The other is the split that runs underneath every
function here: the \emph{same} call is answered exactly when its entries are exact
and at the speed of C when they are machine numbers. Hand \B{Inverse} an integer
matrix and you get an exact rational inverse; hand it a matrix of measured reals
and the work goes to LAPACK over a packed buffer. The mathematics you type is the
mathematics that runs, whichever world the data lives in.

We climb from structure to spectrum. First the matrix as an expression---its
shape, its constructors, and the packed-array substrate that makes it
fast---then the products and norms that turn vectors into geometry; the
determinant, inverse, and matrix powers; solving linear systems exactly and in
the least-squares sense; rank, row reduction, and the null space, with the
user-facing \texttt{ZeroTest} option that makes exact reduction possible where a
structural pivot would fail; the classical matrix decompositions (LU, QR, SVD,
Schur); and finally the eigenproblem---eigenvalues, eigenvectors, and the Jordan
canonical form---closing with the machine--exact divide that governs all of it.

\subsection{Vectors and matrices as expressions}
\label{ssec:la-matrices}

A vector is a flat list; a matrix is a list of equal-length rows. Because that is
all a matrix is, the everyday expression tools work on it unchanged, and a
handful of predicates and constructors round out the vocabulary.

\begin{codepairs}
\pair{linear-algebra/matrices/1}{A \(3\times3\) matrix is just a list of three
rows.}
\pair{linear-algebra/matrices/2}{\B{Dimensions} reads the shape---here \(3\times3\).}
\pair{linear-algebra/matrices/3}{\B{MatrixQ} confirms it is a rectangular rank-2
array (its siblings \B{VectorQ} and \B{SquareMatrixQ} test the other shapes).}
\pair{linear-algebra/matrices/4}{\B{Part} slices a column: \texttt{m[[All, 2]]} is
the second entry of every row.}
\pair{linear-algebra/matrices/5}{And because rows are ordinary lists, \B{Map}
threads \B{Total} over them to sum each row---no matrix-specific machinery
needed.}
\pair{linear-algebra/matrices/6}{\B{Transpose} exchanges rows and columns.}
\pair{linear-algebra/matrices/7}{\B{ConjugateTranspose} transposes \emph{and}
conjugates---the adjoint \(A^{*}\), the right notion for complex matrices.}
\pair{linear-algebra/matrices/8}{\B{IdentityMatrix} builds the \(n\times n\)
identity.}
\pair{linear-algebra/matrices/9}{\B{DiagonalMatrix} places a list on the
diagonal; the off-diagonal zeros are exact when the diagonal is.}
\pair{linear-algebra/matrices/10}{\B{Diagonal} reads the diagonal back out.}
\end{codepairs}

The lesson of pairs~4 and~5 is worth dwelling on: nothing in \texttt{m[[All, 2]]}
or \texttt{Map[Total, m]} knows it is operating on a ``matrix.'' They are the same
generic operators that act on any expression, and they work here for free because
a matrix \emph{is} an expression. The specialised predicates---\B{SymmetricMatrixQ},
\B{HermitianMatrixQ}, \B{PositiveDefiniteMatrixQ}---layer structural tests on top
of that same representation.

\calloutnote{Under the hood}{Behind the plain list is an invisible fast lane. When
every entry is a machine number, Mathilda may hold the matrix as a
\emph{packed array}\index{packed array}: a flat, row-major buffer of raw
\texttt{double}s carrying its shape on the side, exactly like NumPy's
\texttt{ndarray}, rather than a tree of boxed \texttt{Expr} nodes. \B{Head},
\B{ListQ}, printing, and pattern matching cannot tell a packed list from any
other list---it is the \emph{same value}, only stored densely---so the uniformity
above costs nothing. The visible form of that same storage is \B{NDArray}, which
prints as \texttt{NDArray[\ldots]} and announces itself; the two share every
kernel. Packing and the transparency gate that guards it are covered in
\S\ref{ssec:la-precision}.}

\subsection{Products and norms}
\label{ssec:la-products}

The single operator \B{Dot} (written as an infix period, \texttt{.}) carries the
whole algebra of products: vector--vector, matrix--vector, and matrix--matrix are
one function, contracting the last index of its left argument with the first index
of its right.

\begin{codepairs}
\pair{linear-algebra/products/1}{A matrix times a vector---each output entry is a
row dotted with \(\{x, y\}\).}
\pair{linear-algebra/products/2}{Two matrices multiply to a matrix, row-by-column.}
\pair{linear-algebra/products/3}{Two vectors contract to a scalar, the inner
product \(1\cdot4 + 2\cdot5 + 3\cdot6 = 32\).}
\pair{linear-algebra/products/4}{\B{Cross} is the \(3\)-dimensional cross product:
\(\hat x \times \hat y = \hat z\).}
\pair{linear-algebra/products/5}{In general \B{Cross} takes \(n-1\) vectors of
length \(n\); with one \(2\)-vector it returns the perpendicular in the plane.}
\pair{linear-algebra/products/6}{\B{Tr} sums the diagonal---the trace, invariant
under similarity.}
\end{codepairs}

A \emph{norm} measures length. \B{Norm} defaults to the Euclidean (\(2\)-)norm but
takes any \(p\), and \B{Normalize} scales a vector to unit length.

\begin{codepairs}
\pair{linear-algebra/products/7}{\B{Norm} of a symbolic vector is the Hermitian
\(2\)-norm \(\sqrt{\sum |x_i|^2}\)---note the \B{Abs}, so it is correct for complex
entries.}
\pair{linear-algebra/products/8}{The \(1\)-norm sums absolute values: \(1+2+3=6\).}
\pair{linear-algebra/products/9}{The \(\infty\)-norm takes the largest,
\(\max\{1,2,3\}=3\).}
\pair{linear-algebra/products/10}{\B{Normalize} divides by the norm, sending
\(\{3,4\}\) to the unit vector \(\{3/5, 4/5\}\).}
\end{codepairs}

The \(p\)-norms interpolate a family: \(p=1\) the taxicab length, \(p=2\) the
straight-line length, and \(p=\infty\) the largest coordinate. That the symbolic
answer in pair~7 uses \(\lvert x\rvert^2\) rather than \(x^2\) is not pedantry---for
a complex vector the length-squared is \(\sum x_i\overline{x_i}\), and writing
\(x^2\) would give the wrong answer off the real axis.

\calloutnote{Performance}{On a machine-precision vector or matrix these route
through LAPACK: the \(2\)-norm uses \texttt{dnrm2}'s running-scale recurrence and
the matrix norms use \texttt{dlange}, so a well-scaled vector whose true norm is
representable never overflows to \texttt{inf}. \texttt{Norm[\{1e200, 1e200\}]}
returns \(1.414\times10^{200}\)---the naive \(\sqrt{x_1^2 + x_2^2}\) would square to
infinity first. Both \B{Norm} and \B{Dot} lower inside \B{Compile}, so a norm or
inner product inside a tight numerical loop runs at buffer speed
(\texttt{src/linalg/norm.c}, \texttt{dot.c}).}

\subsection{Determinant, inverse, and matrix powers}
\label{ssec:la-det}

The determinant collapses a square matrix to a single number that vanishes exactly
when the matrix is singular; the inverse undoes the linear map when it does not.

\begin{codepairs}
\pair{linear-algebra/det/1}{\B{Det} of a numeric matrix is exact: \(1\cdot4 -
2\cdot3 = -2\).}
\pair{linear-algebra/det/2}{For a symbolic matrix it returns the Laplace
(cofactor) expansion along the first row.}
\pair{linear-algebra/det/3}{\B{Inverse} of an integer matrix comes back with exact
rational entries.}
\pair{linear-algebra/det/4}{Symbolically, the \(2\times2\) inverse is the classical
adjugate over the determinant \(ad-bc\).}
\pair{linear-algebra/det/5}{Store a matrix to check the defining identity\ldots}
\pair{linear-algebra/det/6}{\ldots{}\(A\,A^{-1}\) is the identity, exactly.}
\pair{linear-algebra/det/7}{\B{MatrixPower} raises a matrix to an integer power by
repeated squaring; the tenth power of \(\left(\begin{smallmatrix}1&1\\1&0\end{smallmatrix}\right)\)
is a block of consecutive Fibonacci numbers.}
\pair{linear-algebra/det/8}{It works symbolically too---here the square of a
general \(2\times2\) matrix.}
\end{codepairs}

Pair~7 is a small gem: the powers of the matrix \(\left(\begin{smallmatrix}1&1\\1&0\end{smallmatrix}\right)\)
generate the Fibonacci sequence, because that matrix encodes the recurrence
\(F_{n+1}=F_n+F_{n-1}\). Its tenth power holds \(F_{11}=89\) and \(F_{10}=55\), and
\B{MatrixPower} reaches them in \(\lceil\log_2 10\rceil=4\) matrix multiplications
rather than nine, by repeated squaring.

\calloutnote{Under the hood}{\B{Det} and \B{Inverse} each choose a path by the kind
of entries they see. An all-integer or rational matrix is handed to FLINT's
\texttt{fmpq\_mat\_det} / \texttt{fmpq\_mat\_inv}, which work in polynomial time and
return the exact answer---so a \(12\times12\) Hilbert determinant is instant, not an
\(O(n!)\) Laplace expansion. A machine-real matrix factorises through LAPACK's
\texttt{dgetrf} LU; a genuinely arbitrary-precision (MPFR) matrix uses an
\(O(n^3)\) MPFR LU. Only a \emph{symbolic} matrix falls to cofactor expansion,
which is why pair~2 shows it. The determinant of a large machine matrix is the
product of the LU pivots and can exceed the IEEE-double range---a
\(200\times200\) random matrix has \(\lvert\det\rvert\approx10^{340}\)---so Mathilda
re-accumulates the pivot product in an MPFR value of unbounded exponent rather than
returning \texttt{inf} (\texttt{src/linalg/det.c}, \texttt{inv.c}, \texttt{matpow.c}).}

\subsection{Solving linear systems}
\label{ssec:la-linsolve}

To solve \(A\mathbf{x}=\mathbf{b}\) is the central task of the subject. \B{LinearSolve}
does it exactly when it can, over- and under-determined systems included; when no
exact solution exists, \B{LeastSquares} returns the closest thing.

\begin{codepairs}
\pair{linear-algebra/linsolve/1}{A square system, solved exactly over the
rationals.}
\pair{linear-algebra/linsolve/2}{With symbolic entries the answer is Cramer's
rule: each component is a ratio of determinants sharing the denominator
\(ru-st=\det A\).}
\pair{linear-algebra/linsolve/3}{An \emph{under}-determined system (more unknowns
than equations) returns one particular solution, with the free variables set to
zero---use \B{Solve} for the full parametrised family.}
\pair{linear-algebra/linsolve/4}{When the system is \emph{over}-determined and
inconsistent, \B{LeastSquares} minimises \(\lVert A\mathbf{x}-\mathbf{b}\rVert\):
the best-fit line \(y = \tfrac{19}{3} + \tfrac12 x\) through \((1,7),(2,7),(3,8)\).}
\pair{linear-algebra/linsolve/5}{And when the system \emph{is} consistent, the
least-squares solution coincides with the exact one.}
\end{codepairs}

The distinction in pairs~3 and~4 is the whole geometry of the problem.
\B{LinearSolve} answers ``is \(\mathbf{b}\) in the column
space\index{Cramer's rule} of \(A\)?''---if so,
which combination of columns produces it. When \(A\) has more columns than rows the
answer is a whole affine subspace, and \B{LinearSolve} hands back one point of it;
when \(A\) has more rows than columns and \(\mathbf{b}\) lies off the column space,
there is no solution at all, and \B{LeastSquares} instead returns the
\(\mathbf{x}\) whose image \(A\mathbf{x}\) is the orthogonal projection of
\(\mathbf{b}\) onto the column space---the least-squares fit\index{least squares} that underlies all of
linear regression (and \B{Fit}, in \S\ref{sec:statistics}).

\calloutnote{Theory}{\B{LeastSquares} minimises \(\lVert A\mathbf{x}-\mathbf{b}\rVert_2\).
Setting the gradient to zero gives the \emph{normal equations} \(A^{*}A\,\mathbf{x} =
A^{*}\mathbf{b}\), whose solution is \(\mathbf{x}=A^{+}\mathbf{b}\) with \(A^{+}\) the
\B{PseudoInverse}\index{Moore-Penrose pseudoinverse}. For a rank-deficient \(A\)
the minimiser is not unique, and \(A^{+}\mathbf{b}\) picks out the one of smallest
norm. A machine-real system takes a thin SVD rather than forming \(A^{+}\)
explicitly---a \(500\times500\) least-squares fit that never finished on the exact
pipeline costs \(77\)~ms this way (\texttt{src/linalg/lstsq.c}).}

\usagebox{LinearSolve}

\subsection{Rank, row reduction, and the null space}
\label{ssec:la-reduction}

Beneath solving lies elimination. \B{RowReduce} brings a matrix to reduced row
echelon form; \B{MatrixRank} counts the independent rows; and \B{NullSpace}
returns a basis for the vectors the matrix annihilates. The three are one
computation seen from three angles.

\begin{codepairs}
\pair{linear-algebra/reduction/1}{\B{RowReduce} of a rank-deficient matrix: the
bottom row of zeros is the dependency made visible.}
\pair{linear-algebra/reduction/2}{\B{MatrixRank} reports \(2\)---the third row is
the sum of a scaling of the first two.}
\pair{linear-algebra/reduction/3}{\B{NullSpace} returns the one vector \(\{1,-2,1\}\)
spanning the kernel.}
\pair{linear-algebra/reduction/4}{Store the matrix\ldots}
\pair{linear-algebra/reduction/5}{\ldots and confirm it: \(A\) times the null
vector is \(\mathbf 0\).}
\pair{linear-algebra/reduction/6}{These are exact even symbolically---row-reducing
a matrix whose third row is \(\text{row}_1+\text{row}_2\) exposes the zero row and
the exact rational entries above.}
\pair{linear-algebra/reduction/7}{\B{MatrixRank} sees the symbolic dependency too:
the second row is \(2\) times the first, so the rank is \(1\).}
\end{codepairs}

The rank-nullity theorem\index{rank-nullity theorem} is on display: a
\(3\times3\) matrix of rank \(2\)
(pair~2) has a null space of dimension \(3-2=1\) (pair~3), and their sum recovers
the column count. That Mathilda finds these exactly, over rationals \emph{and} over
symbols, is the point of pairs~6 and~7---the elimination is fraction-free, so no
spurious denominator ever appears.

\calloutnote{Under the hood}{The default reduction is a Bareiss-like
\emph{fraction-free} Gauss--Jordan elimination\index{Bareiss algorithm}: it clears
each column with cross-multiplications chosen so that the exact divisions leave no
denominator larger than necessary, keeping intermediate coefficients from
exploding the way naive Euclidean elimination would. An all-integer or rational
matrix instead goes to FLINT's \texttt{fmpq\_mat\_rref}, whose reduced echelon form
is unique and therefore matches. \B{MatrixRank} and \B{NullSpace} both call back
into \B{RowReduce} through the evaluator, so every improvement to it propagates to
them for free (\texttt{src/linalg/linsolve.c}, \texttt{matrank.c},
\texttt{nullspace.c}).}

\subsection{Exact reduction and the \texttt{ZeroTest} option}
\label{ssec:la-zerotest}

Elimination lives or dies on one question asked at every pivot: \emph{is this entry
zero?} For numbers it is trivial. For symbols it is the hardest question in the
system, and it is where exact linear algebra quietly goes wrong if you are not
careful. Consider a matrix whose two rows are equal---but equal only after a
trigonometric identity that Mathilda does not apply on its own.

\begin{codepairs}
\pair{linear-algebra/zerotest/1}{The Pythagorean identity is true, but it does not
auto-simplify: Mathilda leaves \texttt{Cos[x]\^{}2 + Sin[x]\^{}2} as written.}
\pair{linear-algebra/zerotest/2}{So these two rows are mathematically identical---the
top-left entry \emph{is} \(1\)---but not \emph{structurally} identical.}
\pair{linear-algebra/zerotest/3}{The default \B{NullSpace} uses a structural zero
test. It cannot see that the rows are dependent, so it wrongly reports full rank
and an empty kernel.}
\pair{linear-algebra/zerotest/4}{Supply \texttt{ZeroTest -> (Simplify[\#] === 0 \&)}
and the reduction consults \B{Simplify} at every pivot---now it finds the
dependency and returns the true null vector.}
\end{codepairs}

This is the reason for the user-facing \texttt{ZeroTest} option on \B{RowReduce}
and \B{NullSpace}.\index{ZeroTest option}\index{zero testing!symbolic} The default
structural test treats each distinct radical or transcendental as an independent
symbol---fast, and correct whenever the answer is structurally apparent, but blind
to any zero that only appears after simplification. Passing a predicate---a body
like \texttt{(Simplify[\#] === 0 \&)} or a head like \B{PossibleZeroQ}---replaces
that test with a semantic one, and the elimination pivots on genuine, not merely
apparent, non-zeros. There is no \emph{perfect} zero test---that this must be so is a
theorem, not a shortfall of effort---so how hard to test at each pivot is a choice left to
you; \S\ref{ssec:alg-possiblezeroq} takes up \mcode{PossibleZeroQ} and the undecidability
theorem behind it.

The same fault line runs through \emph{numeric} reduction, where the question
becomes ``is this entry small enough to treat as zero?'' Here the answer depends on
a tolerance, and the exact and machine paths part company sharply.

\begin{codepairs}
\pair{linear-algebra/zerotest/5}{A matrix whose two off-diagonal pivots are tiny
but exactly non-zero.}
\pair{linear-algebra/zerotest/6}{Exactly, this matrix has three independent
rows---the tiny entries \(10^{-10}\) and \(10^{-20}\) are genuinely non-zero.}
\pair{linear-algebra/zerotest/7}{Convert to machine reals and one pivot falls below
the default tolerance: \B{MatrixRank} now reports \(2\).}
\pair{linear-algebra/zerotest/8}{Force \texttt{Tolerance -> 0} and every non-zero
counts again, recovering the rank of \(3\).}
\end{codepairs}

Neither answer is wrong; they answer different questions. Exact reduction asks
whether an entry is \emph{algebraically} zero and keeps the tiny pivots; machine
reduction asks whether it is zero \emph{to working precision} and discards them,
which is exactly what you want when the small entries are round-off rather than
signal. The lesson is to know which question you are asking---and to reach for
\texttt{Tolerance} (numeric) or \texttt{ZeroTest} (exact) when the default answers
a different one.

\subsection{Matrix decompositions}
\label{ssec:la-decomp}

A decomposition rewrites a matrix as a product of simpler factors---triangular,
orthogonal, diagonal---and almost every serious numerical algorithm is a
decomposition in disguise. Mathilda ships the classical four.

\begin{codepairs}
\pair{linear-algebra/decomp/1}{\B{LUDecomposition} returns the combined
lower/upper factors, the pivot permutation, and a condition estimate---here the
factorisation of a Vandermonde-like matrix.}
\pair{linear-algebra/decomp/2}{\B{QRDecomposition} factors \(A\) into an
orthonormal \texttt{q} and an upper-triangular \texttt{r}. This is the classic
textbook \(3\times3\), and its \(Q\) is exactly rational.}
\pair{linear-algebra/decomp/3}{The orthonormal factor \texttt{q}: its rows are
mutually perpendicular unit vectors.}
\pair{linear-algebra/decomp/4}{The upper-triangular factor \texttt{r}.}
\pair{linear-algebra/decomp/5}{And they reconstruct the original exactly:
\(Q^{\mathsf T} R = A\).}
\pair{linear-algebra/decomp/6}{\B{SingularValueDecomposition} of a rank-\(1\)
matrix: a single non-zero singular value \(\sqrt{10}\) on the diagonal of
\(\Sigma\), the rest zero.}
\pair{linear-algebra/decomp/7}{A symmetric machine matrix for the Schur form\ldots}
\pair{linear-algebra/decomp/8}{\ldots factored by \B{SchurDecomposition}\ldots}
\pair{linear-algebra/decomp/9}{\ldots into an orthonormal \texttt{q} and an
upper-triangular \texttt{t} (diagonal here, since the matrix is symmetric): its
diagonal carries the eigenvalues.}
\pair{linear-algebra/decomp/10}{\B{Chop} confirms \(A = Q\,T\,Q^{*}\) to machine
zero.}
\end{codepairs}

Each factorisation answers a different need. \B{LUDecomposition} is the engine of
\B{LinearSolve}: once \(A=LU\), each right-hand side costs only two triangular
solves. \B{QRDecomposition} orthogonalises the columns and is the stable way to
solve least-squares problems and the workhorse of the eigenvalue algorithm.
\B{SingularValueDecomposition} is the most revealing of all---the singular values
are the semi-axes of the ellipsoid \(A\) maps the unit sphere onto, the number of
non-zero ones is the rank (pair~6 makes a rank-\(1\) matrix show a single
\(\sqrt{10}\)), and truncating the small ones is the mathematics of data
compression and \B{PseudoInverse}. \B{SchurDecomposition}\index{Schur decomposition} is the theoretical keystone:
\emph{every} square matrix is unitarily
similar to a (quasi-)upper-triangular one whose diagonal is the spectrum, so the
eigenvalues can be read off without ever solving a polynomial.

\calloutnote{Under the hood}{The decompositions dispatch on precision. A
machine-real matrix goes to LAPACK---\texttt{dgetrf} (LU), a Householder
\texttt{dgeqrf} (QR)\index{Householder reflection}, divide-and-conquer
\texttt{dgesdd} (SVD), and the Francis QR iteration \texttt{dgees} (Schur)---wired
through a four-tier autodetection ladder (Apple Accelerate, then two flavours of
system LAPACK, then a graceful fallback). An arbitrary-precision matrix runs a
hand-rolled MPFR kernel at the requested precision (Householder QR, one-sided
Jacobi SVD, Hessenberg+Francis Schur). An exact or symbolic matrix stays on a
modified Gram--Schmidt or eigendecomposition path and returns exact
\texttt{Sqrt[\ldots]} forms. \B{SchurDecomposition} is numeric only---a generic
matrix has no closed-form symbolic Schur form---and is left unevaluated on symbolic
input rather than answered wrongly (\texttt{src/linalg/ludecomp.c},
\texttt{qrdecomp.c}, \texttt{svdecomp.c}, \texttt{schurdecomp.c},
\texttt{lapack.c}).}

\subsection{Eigenvalues, eigenvectors, and the Jordan form}
\label{ssec:la-eigen}

An eigenvector is a direction the matrix merely scales; its eigenvalue is the
scale. Finding them is the deepest structural question you can ask of a linear map,
and it is where the exact/numeric split is most vivid.

\begin{codepairs}
\pair{linear-algebra/eigen/1}{\B{CharacteristicPolynomial} gives \(\det(A - xI)\),
whose roots are the eigenvalues: \(x^2 - 5x - 2\) here.}
\pair{linear-algebra/eigen/2}{\B{Eigenvalues} of an exact matrix stays exact,
returning genuine surds---and a \(0\), certifying the matrix is singular.}
\pair{linear-algebra/eigen/3}{Symbolically, the \(2\times2\) eigenvalues are the
trace-and-determinant quadratic formula.}
\pair{linear-algebra/eigen/4}{A \(90^\circ\) rotation has no real eigenvectors, so
its eigenvalues are the complex pair \(\pm i\).}
\pair{linear-algebra/eigen/5}{\B{Eigenvectors} of a defective matrix: the repeated
eigenvalue \(2\) supplies only one eigenvector, so the deficient slot is padded
with a zero vector to keep the list aligned with the eigenvalues.}
\end{codepairs}

Pair~5 is the crack through which the Jordan form enters. The matrix there has
eigenvalue \(2\) with algebraic multiplicity \(2\) but only \emph{one} independent
eigenvector---it is \emph{defective}, not diagonalizable. The zero-vector padding
is Mathilda telling you honestly that an eigenvector is missing. The remedy is the
\emph{Jordan canonical form}\index{Jordan canonical form}: the closest a defective
matrix can come to diagonal, with the missing eigenvectors replaced by
\emph{generalized} ones\index{generalized eigenvector} and a tell-tale \(1\) on the
superdiagonal.

\begin{codepairs}
\pair{linear-algebra/eigen/6}{A matrix with a defective eigenvalue \(12\).}
\pair{linear-algebra/eigen/7}{\B{JordanDecomposition} returns the similarity
\(S\) and the Jordan form \(J\)\ldots}
\pair{linear-algebra/eigen/8}{\ldots whose \(J\) has a \(2\times2\) Jordan block for
\(12\): the eigenvalue twice on the diagonal, a \(1\) above it marking the deficient
eigenspace.}
\pair{linear-algebra/eigen/9}{And the decomposition reconstructs the matrix:
\(A = S\,J\,S^{-1}\).}
\end{codepairs}

The superdiagonal \(1\) in pair~8 is the entire content of the Jordan form: it says
that eigenvalue \(12\) has a two-step chain of generalized eigenvectors rather than
two independent ones. Read \(J\) and you read the matrix's deepest invariant---the
sizes of its Jordan blocks---which no similarity transformation can change.

\calloutnote{Under the hood}{\B{Eigenvalues} does \emph{not} form \(\det(A-xI)\) and
expand a symbolic determinant; it uses the Faddeev--LeVerrier--Souriau
recurrence\index{Faddeev-LeVerrier algorithm}, which produces the characteristic
polynomial in \(O(n^4)\) matrix multiplications, then hands the polynomial to
\B{Solve}---so eigenvalues of a large exact matrix return quickly, and the exact
roots come back as radicals or \B{Root} objects (\S\ref{sec:algebra}). The exact
\B{JordanDecomposition} reads each block's size from the nullity sequence
\(\dim\ker(A-\lambda I)^k\); a machine-real matrix with a distinct spectrum is
diagonalizable and takes a direct LAPACK \texttt{dgeev} path (a \(100\times100\)
in \(\approx\!9.6\)~ms), while a defective or ambiguous numeric matrix is
rationalised, decomposed exactly, and numericalised back to recover the true block
structure (\texttt{src/linalg/eigen.c}, \texttt{charpoly.c}, \texttt{jordandecomp.c}).}

\usagebox{Eigenvalues}

\usagebox{JordanDecomposition}

\subsection{The machine--exact divide}
\label{ssec:la-precision}

Every routine in this section has, in fact, been choosing between worlds all along.
It is worth watching the choice happen on one matrix. The Hilbert matrix---entries
\(1/(i+j-1)\), symmetric and famously ill-conditioned---is the perfect specimen.

\begin{codepairs}
\pair{linear-algebra/precision/1}{Exact: \B{Det} of the \(5\times5\)
\B{HilbertMatrix} is the precise rational \(1/266716800000\).}
\pair{linear-algebra/precision/2}{Machine: the same determinant over doubles is a
single \texttt{Real}---accurate to the digits shown, but a rounded number, not an
exact one.}
\pair{linear-algebra/precision/3}{Arbitrary precision: at \(40\) digits (MPFR) the
determinant of the \(7\times7\) Hilbert matrix comes back to full requested
precision.}
\pair{linear-algebra/precision/4}{Exact inversion of an ill-conditioned matrix is
no trouble at all when the arithmetic is exact---integer entries, computed by
FLINT.}
\pair{linear-algebra/precision/5}{The visible packed form: an \B{NDArray} is a
dense machine buffer that announces itself.}
\pair{linear-algebra/precision/6}{\B{Dot} over two \texttt{NDArray}s runs on the raw
buffers and returns a closed-system \texttt{NDArray}\ldots}
\pair{linear-algebra/precision/7}{\ldots and \B{Det} of one returns a bare machine
scalar, straight off the buffer.}
\end{codepairs}

The three determinants---an exact rational, a machine \texttt{Real}, and a
\(40\)-digit MPFR number---are the same mathematical quantity computed in three
number systems, and Mathilda selects among them purely from the \emph{kind} of the
entries. Exact in, exact out (via FLINT or fraction-free elimination); machine in,
LAPACK out; high-precision MPFR in, an MPFR kernel out. Nothing about the syntax
changes; the representation of the data decides the path.\index{arbitrary precision}\index{LAPACK}

\calloutnote{Performance}{The packed array is what makes the machine path fast, and
the rule that keeps it safe is the \emph{transparency gate}\index{packed array!transparency gate}: a
machine-number list may be held as a raw buffer, and a numeric routine on the
``aware'' list (\B{Dot}, \B{Det}, \B{Inverse}, \B{LinearSolve}, \B{Norm},
\B{MatrixRank}, and their kin) reads that buffer directly---no per-element boxing,
no re-parsing---so a \(300\times300\) \texttt{MatrixPower[A, 4]} is the \(0.69\)~ms of
its two \texttt{dgemm} calls. A routine \emph{not} on that list transparently
un-packs first, trading speed for correctness but never the answer. The same
machine heads lower inside \B{Compile}, so a linear-algebra kernel written at the
prompt runs at full speed inside a compiled numerical loop
(\texttt{src/pack.c}, \texttt{src/linalg/ndlinalg.c}; see also
\S\ref{sec:statistics} on packed reductions).}

\subsection*{Where this connects}

Linear algebra is the substrate the rest of the system computes on. The
least-squares solve of \S\ref{ssec:la-linsolve} is the machinery behind \B{Fit}
and the correlation matrix of \S\ref{sec:statistics}, both of them a \B{Dot} of a
matrix with its \B{Transpose}. The exact eigenvalues of \S\ref{ssec:la-eigen} flow
straight out of the polynomial solving and \B{Root} objects of \S\ref{sec:algebra},
and the fraction-free elimination beneath \B{RowReduce} is the same idea as the
subresultant sequences that tame coefficient growth there. The machine kernels and
the packed substrate reappear wherever numbers are crunched in bulk---the numerical
calculus, the optimizers, the machine-learning primitives---and every machine head
named here lowers inside \B{Compile} (Chapter~\ref{ch:compilation}), so the
linear algebra you write interactively is the linear algebra that runs in a tight
loop. For the exhaustive options of any function above, its reference page is one
\texttt{?Name} away at the prompt.

\section{Numerical Calculus}
\label{sec:numcalc}

Every operation in the previous sections was \emph{exact}: \B{Integrate} returns an
antiderivative, \B{Solve} returns the roots as algebraic numbers, \B{Sum} returns a closed
form. That is the ideal, and when it is attainable it is unbeatable. But most of the
integrals, equations, sums, and differential equations that arise in practice have no closed
form at all---not because Mathilda is not clever enough, but because, as theorems of Liouville,
Abel, and others tell us, none exists. For those problems the honest answer is a
\emph{number}, computed to whatever precision you ask for, and computing it well is a
subject in its own right.

This section is about Mathilda's \emph{numerical} calculus: the family of functions whose
names begin with \texttt{N}. We start with the question of when to reach for them at all, and
with the one contract they all share---how they decide they are ``done.'' Then we take them in
turn: numerical integration (\B{NIntegrate}), summation and products (\B{NSum}, \B{NProduct}),
differentiation (\B{ND}), series and limits (\B{NSeries}, \B{NLimit}), residues (\B{NResidue}),
differential equations (\B{NDSolve}), the roots of polynomials and the solutions of general
equations (\B{NRoots}, \B{NSolve}, \B{FindRoot}), and finally optimization, both local
(\B{FindMinimum}, \B{FindMaximum}) and global (\B{NMinimize}, \B{NMaximize}). For each we show
not just the syntax but the algorithm underneath---because choosing the right method for the
shape of your problem is most of the craft of numerical work.

\subsection{Why and when to go numeric}
\label{ssec:whynumeric}

The clearest way to see the need is to watch a symbolic function stop short and a numerical one
carry on. The Gaussian \(\int_0^1 e^{-x^2}\,dx\) has no elementary antiderivative; \B{Integrate}
does the best exact thing it can, expressing the answer through the error function, while
\B{NIntegrate} simply returns the number:

\begin{codepairs}
\pair{numerical-calculus/motivation/1}{No elementary antiderivative exists, so the exact answer
is written through a special function, \(\tfrac12\sqrt{\pi}\,\operatorname{erf}(1)\).}
\pair{numerical-calculus/motivation/2}{\B{NIntegrate} gives the value you usually wanted in the
first place.}
\pair{numerical-calculus/motivation/3}{The sine integral is worse: it has no antiderivative
elementary in \(x\) at all, so \B{Integrate} names it \(\operatorname{Si}(1)\).}
\pair{numerical-calculus/motivation/4}{Again the number is one step away.}
\end{codepairs}

Roots tell the same story. Abel proved that the general quintic has no solution in radicals, so
\B{Solve} can only name the five roots of \(x^5-x-1\) as abstract \texttt{Root} objects---correct,
but not a number you can plot---whereas \B{NRoots} and \B{FindRoot} deliver digits:

\begin{codepairs}
\pair{numerical-calculus/motivation/5}{There is no radical formula, so each root is returned as
an indexed \texttt{Root} object standing for ``the \(k\)th root of this polynomial.''}
\pair{numerical-calculus/motivation/6}{\B{NRoots} approximates all five at once: one real, two
conjugate pairs.}
\pair{numerical-calculus/motivation/7}{\B{FindRoot} chases a single root of a transcendental
equation---\(\cos x = x\), whose solution (the ``Dottie number'') is provably not expressible in
closed form.}
\end{codepairs}

The rule of thumb: reach for a numerical function when the exact machinery returns something
unevaluated, something written through a special function or \texttt{Root} object you would only
numericalize anyway, or something correct but too large to be useful---and whenever the input is
itself inexact, since there is then no exact answer to find.

\subsection{Precision, and the three goals}
\label{ssec:numgoals}

Every numerical function in this section shares one control panel, so it is worth learning once.
Three options govern how hard the function works and how it judges success:
\texttt{WorkingPrecision}, \B{AccuracyGoal}, and \B{PrecisionGoal}\index{numerical calculus!accuracy
and precision goals}.

\texttt{WorkingPrecision} sets how many digits the internal arithmetic carries. Its default is
\texttt{MachinePrecision}---the roughly \(15.95\) decimal digits of an IEEE double
(\S\ref{ssec:machine-precision})---and a larger value transparently switches the computation onto
Mathilda's MPFR path, so you can ask for forty digits or five hundred:

\begin{codepairs}
\pair{numerical-calculus/goals/1}{At machine precision the answer is a plain double.}
\pair{numerical-calculus/goals/2}{Raising \texttt{WorkingPrecision} to \(40\) recomputes the same
integral in extended precision.}
\pair{numerical-calculus/goals/3}{The same lever works for a sum: forty digits of \(\pi^2/6\).}
\end{codepairs}

The other two options set the \emph{target}, not the arithmetic. \B{AccuracyGoal}~\(a\) asks for a
result whose \emph{absolute} error is below \(10^{-a}\); \B{PrecisionGoal}~\(p\) asks for a
\emph{relative} error below that fraction of the answer's size. A function stops when it has met the
combined tolerance
\[
  \text{error} \;\le\; 10^{-a} \;+\; |x|\,10^{-p},
\]
where \(x\) is the current estimate.\calloutnote{Under the hood}{One small module,
\texttt{src/nc\_accuracy.c}, implements this contract for every routine here, so the semantics never
drift between them. \texttt{Automatic} resolves to two digits below \texttt{WorkingPrecision}; the symbol
\texttt{MachinePrecision} tracks the machine-double budget; and \texttt{Infinity} switches a
criterion off, so \texttt{PrecisionGoal -> Infinity} judges by absolute error alone. The near-universal
defaults are \texttt{AccuracyGoal -> MachinePrecision} and \texttt{PrecisionGoal -> Automatic}.}
The absolute term keeps the relative term from demanding the impossible near a zero answer, and the
relative term keeps the absolute term honest for very large or very small ones.

One behavior is worth internalizing now, because you will meet it throughout the section:
\textbf{these functions never fail with an error.} If a routine cannot reach its goal within its
resource budget, it refines as far as it can, keeps the best estimate it found, prints a
\texttt{Head::accgl} message to the terminal, and returns that estimate
anyway.\calloutnote{Pitfall}{The \texttt{::accgl} message is advisory, not fatal, and it is printed
to standard error---so it appears at your prompt but not inside a captured result. When a numerical
answer looks a digit or two short, it is often because the goal was quietly not met; raise
\texttt{WorkingPrecision}, loosen the goal, or (as we will see) pick a method better suited to the
problem.} We rely on this deliberately below, where forcing the wrong summation method leaves a
visibly wrong last digit.

\subsection{Numerical integration}
\label{ssec:nintegrate}

\B{NIntegrate}\index{NIntegrate@\mcode{NIntegrate}!methods} is the workhorse, and its default
\texttt{Automatic} method is really a small polyalgorithm that inspects the integrand and the
interval and dispatches to a specialist. On a smooth, finite integrand it uses globally-adaptive
Gauss--Kronrod quadrature; on a singular or infinite one it switches to a double-exponential rule;
on an oscillatory one it uses a Levin collocation scheme; and in high dimensions it falls back to
Monte-Carlo sampling. The everyday cases just work:

\begin{codepairs}
\pair{numerical-calculus/nintegrate-basic/1}{An infinite interval: \(\int_0^\infty e^{-x^2}\,dx =
\tfrac12\sqrt{\pi}\).}
\pair{numerical-calculus/nintegrate-basic/2}{An oscillatory tail: the Dirichlet integral
\(\int_0^\infty \tfrac{\sin x}{x}\,dx = \tfrac{\pi}{2}\).}
\pair{numerical-calculus/nintegrate-basic/3}{An integrable endpoint singularity at \(x=0\); the
adaptive rule clusters its samples there.}
\pair{numerical-calculus/nintegrate-basic/4}{Two dimensions: adaptive cubature over the plane gives
\(\int_{\mathbb{R}^2} e^{-x^2-y^2} = \pi\).}
\end{codepairs}

\usagebox{NIntegrate}

When you know the shape of your integrand, naming the \texttt{Method} is both faster and more
accurate. The most useful members of the family are \texttt{"GlobalAdaptive"} (the Gauss--Kronrod
default)\index{Gauss--Kronrod quadrature}, \texttt{"DoubleExponential"} (the tanh-sinh / sinh-sinh /
exp-sinh rules, superb for endpoint singularities and infinite ranges)\index{double-exponential
quadrature}, \texttt{"LevinRule"} (for products with an oscillatory kernel like \(\cos g\), \(\sin
g\), or \(e^{ig}\))\index{Levin collocation}, \texttt{"MonteCarlo"} and its quasi-random variant (for
many dimensions and for region indicators), and \texttt{"PrincipalValue"} (for a Cauchy principal
value around a pole named in \texttt{Exclusions}):

\begin{codepairs}
\pair{numerical-calculus/nintegrate-methods/1}{The double-exponential rule reproduces the Gaussian on
the half-line to machine precision.}
\pair{numerical-calculus/nintegrate-methods/2}{A rapidly oscillating kernel: Levin collocation returns
\(\int_0^\infty \cos(100x)\,e^{-x}\,dx = 1/10001\) without resolving every wiggle.}
\pair{numerical-calculus/nintegrate-methods/3}{A Cauchy principal value straddling the pole at the
origin: the divergent halves cancel to \(\ln 2\).}
\pair{numerical-calculus/nintegrate-methods/4}{Monte-Carlo sampling of a two-dimensional Gaussian; the
seed is fixed, so the estimate is reproducible.}
\end{codepairs}

\calloutnote{Theory}{Gauss--Kronrod quadrature\index{Gauss--Kronrod quadrature} is the engine behind
the default. A Gauss rule of \(n\) points integrates polynomials up to degree \(2n-1\) exactly; the
Kronrod extension adds \(n+1\) points so that the \(15\)-point rule reuses every node of the embedded
\(7\)-point rule. The difference between the two estimates is a cheap, reliable error bar, and the
adaptive driver bisects whichever subinterval carries the largest one until the total error meets the
goal---the classic QUADPACK strategy~\cite{quadpack}.}

\calloutnote{Theory}{The double-exponential (tanh-sinh) rule~\cite{takahasimori} is a change of variable
\(x = \tanh\!\big(\tfrac{\pi}{2}\sinh t\big)\) that makes the integrand decay doubly exponentially at
the ends. A plain trapezoidal rule in \(t\) then converges with astonishing speed, and---crucially---it
places no sample exactly at the endpoints, so an integrable singularity there costs nothing. This is why
it is Mathilda's method of choice for singular and infinite integrals and for high precision.}

Two honesty notes. Several rule names the reference documentation lists---\texttt{"ClenshawCurtisRule"},
\texttt{"LobattoKronrodRule"}, and a few others---are recognized but not yet implemented, and asking for
one yields an \texttt{NIntegrate::method} message rather than a silent substitution. And
\texttt{"AdaptiveMonteCarlo"} currently runs the same sampler as \texttt{"MonteCarlo"} without the
adaptive stratification its name promises. The book prefers to say so plainly.

\subsection{Numerical summation and products}
\label{ssec:nsum}

An infinite sum is an integral's discrete cousin, and \B{NSum}\index{NSum@\mcode{NSum}!methods} carries
the same idea: add the first several terms explicitly, then \emph{extrapolate} the tail rather than
plod through millions of terms. Its \texttt{Automatic} method reads the series and chooses---the
Euler--Maclaurin formula for a smooth monotone tail\index{Euler--Maclaurin summation}, the
Cohen--Villegas--Zagier algorithm for an alternating one\index{alternating series!acceleration}, and
Wynn's epsilon algorithm otherwise\index{Wynn epsilon algorithm}:

\begin{codepairs}
\pair{numerical-calculus/nsum/1}{The Basel sum \(\sum 1/n^2 = \pi^2/6\); Euler--Maclaurin nails it from
a handful of terms.}
\pair{numerical-calculus/nsum/2}{An alternating series for \(\pi/4\); the \texttt{"AlternatingSigns"}
method converges geometrically where naive addition crawls.}
\pair{numerical-calculus/nsum/3}{A slowly converging \(p\)-series at forty digits---compare
\(\zeta(11/10)\). Adding terms directly, you would need more than \(10^{400}\) of them.}
\end{codepairs}

\usagebox{NSum}

Extrapolation is powerful but not free, and the wrong accelerator can hurt. Forcing Wynn's epsilon
onto the Basel sum---whose tail is monotone, not oscillatory---leaves the last digits wrong:

\begin{codepairs}
\pair{numerical-calculus/nsum/4}{Compare with pair~1 above: \(1.62533\) instead of \(1.64493\). Wynn's
epsilon is the wrong tool for a monotone tail, and (because it cannot reach the goal) an
\texttt{NSum::accgl} message accompanies this at the terminal.}
\end{codepairs}

\calloutnote{Theory}{Wynn's epsilon algorithm\index{Wynn epsilon algorithm} is a recursive
implementation of the Shanks transformation, which models a sequence's tail as a sum of decaying
geometric terms and cancels them. It is spectacular on alternating and geometric-like sequences and
weak on the slow algebraic decay of a monotone \(p\)-series---exactly the failure the pair above shows.
The Euler--Maclaurin formula, by contrast, replaces the monotone tail with an integral plus a
correction series in the Bernoulli numbers, which is why \texttt{Automatic} prefers it there.}

\B{NProduct}\index{NProduct@\mcode{NProduct}} handles infinite products by the obvious device---it sums
the logarithms and exponentiates, so \(\prod f = \exp\!\sum \log f\)---and therefore inherits \B{NSum}'s
methods and convergence test wholesale:

\begin{codepairs}
\pair{numerical-calculus/nproduct/1}{\(\prod_{n\ge 2}\!\big(1-\tfrac1{n^2}\big) = \tfrac12\), a
telescoping classic.}
\pair{numerical-calculus/nproduct/2}{\(\prod_{n\ge 2}\tfrac{n^2}{n^2-1} = 2\), its reciprocal.}
\end{codepairs}

\usagebox{NProduct}

\subsection{Numerical differentiation}
\label{ssec:nd}

Differentiation is exact and easy when you have a formula, but a numerical derivative is what you need
when the function is only known through a black box, or is not analytic, or when a symbolic derivative
would be enormous. \B{ND}\index{ND@\mcode{ND}!methods} estimates the derivative of \texttt{expr} at a
point. Its default \texttt{EulerSum} method samples finite differences on a shrinking step and
Richardson-extrapolates them to the limit\index{Richardson extrapolation}:

\begin{codepairs}
\pair{numerical-calculus/nd/1}{\(\tfrac{d}{dx}e^x\) at \(x=1\) is \(e\).}
\pair{numerical-calculus/nd/2}{A high-order, high-precision derivative: the first derivative of
\(\sin(x^2)\) at \(x=3\) to forty digits.}
\pair{numerical-calculus/nd/3}{The integrand \(\operatorname{Re}\,\cos(iy) = \cosh y\) is not
analytic in the complex sense, so only the finite-difference method applies; its derivative at
\(1\) is \(\sinh 1 = 1.1752\).}
\end{codepairs}

\usagebox{ND}

\calloutnote{Under the hood}{\B{ND}'s second method, \texttt{Method -> NIntegrate}, evaluates the
derivative through Cauchy's integral formula, \(f^{(n)}(x_0) = \tfrac{n!}{2\pi i}\oint
\tfrac{f(z)}{(z-x_0)^{n+1}}\,dz\), by calling \B{NResidue} on a small circle. That route needs the
function to be analytic near \(x_0\)---but in return it delivers derivatives of \emph{fractional} and
even complex order, which finite differences cannot. The finite-difference \texttt{EulerSum} method is
the default precisely because it makes no analyticity demand, as pair~3 requires.}

\subsection{Numerical series}
\label{ssec:nseries}

\B{NSeries}\index{NSeries@\mcode{NSeries}} produces a whole Taylor or Laurent expansion numerically:
it samples the function on a circle in the complex plane and recovers the coefficients by a
discrete Fourier transform---Cauchy's integral formula, one contour, all coefficients at once. It is
the natural extension of the contour derivative behind \B{ND}, from a single coefficient to the whole
expansion. Its single algorithm has no \texttt{Method} to choose; you control only the sampling
\texttt{Radius}. Because the DFT carries a little round-off into the coefficients it does not use, it is
idiomatic to \B{Chop} the result:

\begin{codepairs}
\pair{numerical-calculus/nseries/1}{The Taylor series of \(e^x\) about \(0\); the coefficients
\(1,1,\tfrac12,\tfrac16,\tfrac1{24}\) are read straight off.}
\pair{numerical-calculus/nseries/2}{A genuine Laurent series: \(\cot x = \tfrac1x - \tfrac{x}{3} -
\tfrac{x^3}{45} - \cdots\), pole and all.}
\end{codepairs}

\usagebox{NSeries}

\subsection{Numerical limits}
\label{ssec:nlimit}

\B{NLimit}\index{NLimit@\mcode{NLimit}!methods} evaluates a limit by sampling the function along a
sequence of points approaching the target and accelerating that sequence to its limit---the same
extrapolation family as \B{NSum}, since a limit is what a sequence of partial results is converging to.
It is the numerical companion to the symbolic \B{Limit} of \S\ref{sec:calculus}, and it shines when a
limit is delicate or defined only through a hard special function:

\begin{codepairs}
\pair{numerical-calculus/nlimit/1}{A \(0/0\) form: \(\lim_{x\to 0}(1-\cos x)/x^2 = \tfrac12\). Here
\texttt{AccuracyGoal -> 10} simply asks for ten clean digits.}
\pair{numerical-calculus/nlimit/2}{Two diverging pieces whose difference is finite:
\(\lim_{s\to1}\big(\zeta(s)-\tfrac1{s-1}\big) = \gamma\), the Euler--Mascheroni constant.}
\pair{numerical-calculus/nlimit/3}{An \(\infty\)-point limit computed with Wynn's epsilon
(\texttt{"SequenceLimit"}): \(\lim_{n\to\infty} n\,(2^{1/n}-1) = \ln 2\).}
\end{codepairs}

\usagebox{NLimit}

\subsection{Numerical residues}
\label{ssec:nresidue}

The contour trick that powers \B{NSeries} is a tool in its own right. \B{NResidue}\index{NResidue@%
\mcode{NResidue}} computes the residue of a function at a point---the coefficient of \((z-z_0)^{-1}\) in
its Laurent expansion---by integrating around a small circle with the periodic trapezoidal rule, which
is exponentially accurate on a smooth closed contour. Its reach is exactly the symbolic \B{Residue}'s
blind spot: essential singularities, where no Laurent series terminates.

\begin{codepairs}
\pair{numerical-calculus/nresidue/1}{The simplest pole: \(\operatorname{Res}_{z=0}\tfrac1z = 1\).}
\pair{numerical-calculus/nresidue/2}{A pole of \(\tan z\) at \(\pi/2\) has residue \(-1\).}
\pair{numerical-calculus/nresidue/3}{An \emph{essential} singularity: \(\operatorname{Res}_{z=0}
e^{1/z} = 1\). Here the contour \texttt{Radius} is widened to \(1\)---see the pitfall below.}
\pair{numerical-calculus/nresidue/4}{Cauchy coefficient extraction: the residue of \(e^x/x^4\) at
\(0\) is the \(x^3\) Taylor coefficient of \(e^x\), namely \(1/3! = 0.166667\).}
\end{codepairs}

\usagebox{NResidue}

\calloutnote{Pitfall}{The default contour radius is \(1/100\), which is ideal for an ordinary pole but
dangerous for an essential singularity: \(e^{1/z}\) reaches \(e^{100}\) on a circle of radius \(1/100\),
overwhelming the arithmetic. Widening \texttt{Radius} to \(1\), as pair~3 does, keeps the integrand in range.
The lesson generalizes---when a residue near an essential singularity comes back as noise, enlarge the
contour.}

\subsection{Differential equations}
\label{ssec:ndsolve}

Most differential equations have no closed-form solution, which makes \B{NDSolve}\index{NDSolve@%
\mcode{NDSolve}!methods} one of the most important functions in the system. It integrates an initial-value
problem---a single ODE, a coupled system, a higher-order equation reduced automatically to first order,
or a partial differential equation over a rectangle---and returns the solution not as a formula but as an
\B{InterpolatingFunction}\index{InterpolatingFunction@\mcode{InterpolatingFunction}}: a callable object
you can evaluate, differentiate, or plot anywhere in the solved interval.

\begin{codepairs}
\pair{numerical-calculus/ndsolve-numeric/1}{A \emph{stiff} problem, solved with the implicit
\texttt{"BackwardEuler"} method; the answer is an \texttt{InterpolatingFunction} over \([0,3]\).}
\pair{numerical-calculus/ndsolve-numeric/2}{The harmonic oscillator \(y''+y=0\), \(y(0)=1\),
\(y'(0)=0\); assign the solution to reuse it.}
\pair{numerical-calculus/ndsolve-numeric/3}{Evaluating the solution at \(x=2\)\dots}
\pair{numerical-calculus/ndsolve-numeric/4}{\dots reproduces \(\cos 2\) to full machine precision.}
\pair{numerical-calculus/ndsolve-numeric/5}{A coupled first-order system returns one
\texttt{InterpolatingFunction} per dependent variable.}
\end{codepairs}

\usagebox{NDSolve}

The \texttt{Method} option chooses the integrator. The default, \texttt{"DOPRI5"}, is the adaptive
Dormand--Prince \(5(4)\) Runge--Kutta pair\index{Runge--Kutta method!Dormand--Prince}; for a
\emph{stiff} problem---one where the explicit step size is throttled not by accuracy but by stability---you
want an implicit method such as \texttt{"BackwardEuler"} or the variable-order \texttt{"BDF"}. A nonlinear
oscillator makes the interpolating solution vivid. The van der Pol equation \(y'' - \mu(1-y^2)y' + y = 0\)
relaxes onto a \emph{limit cycle}, a closed orbit no linear equation can produce; plotting \(y'\) against
\(y\) draws it directly:

\plotcell{numerical-calculus/vanderpol}

The trajectory starts near \((1,0)\), spirals outward, and locks onto the sharp-cornered limit cycle
characteristic of a relaxation oscillator.

\calloutnote{Theory}{An explicit Runge--Kutta method advances the solution by sampling the slope at
several trial points inside the step and combining them; the Dormand--Prince pair computes two estimates
of different orders from the \emph{same} seven slope evaluations, and their difference drives the adaptive
step size---large where the solution is placid, small where it turns sharply. A \emph{stiff} equation
defeats this: stability, not accuracy, forces a tiny step, and an \emph{implicit} method---one that solves
an algebraic equation for the next value---removes that ceiling. This explicit/implicit split is the
central fact of practical ODE solving~\cite{hairer}.}

For a partial differential equation, \B{NDSolve} uses the \emph{method of lines}: it discretizes the
spatial derivatives on a grid, turning the PDE into a large system of ODEs in time, and integrates that.
The heat equation \(u_t = u_{xx}\) started from a half-sine profile shows the fundamental mode decaying
exponentially:

\plotcell{numerical-calculus/heat-profile}

\calloutnote{Under the hood}{The spatial grid, its resolution (\texttt{"MinPoints"}), and the finite-%
difference accuracy order (\texttt{"DifferenceOrder"}) are sub-options of \texttt{Method -> \{"MethodOfLines",
...\}}. Mathilda builds the spatial derivative stencils by Fornberg's algorithm, compiles the resulting
right-hand side for speed, and---because diffusion is stiff---integrates the semi-discrete system with the
BDF method by default. The returned object is then a two-dimensional \texttt{InterpolatingFunction} in
\(t\) and \(x\).}

\subsection{Roots of polynomials}
\label{ssec:nroots}

For the special case of a single polynomial equation, \B{NRoots}\index{NRoots@\mcode{NRoots}!methods}
finds \emph{all} the roots at once and returns them as a disjunction. This is where Abel's theorem bites
hardest and numerics earn their keep:

\begin{codepairs}
\pair{numerical-calculus/nroots/1}{A cubic: one real root and a conjugate pair.}
\pair{numerical-calculus/nroots/2}{The quintic \(x^5-x-1\)---no radical form---resolved to five roots.}
\pair{numerical-calculus/nroots/3}{Complex coefficients are fine: the two square roots of \(3+4i\) are
\(\pm(2+i)\).}
\pair{numerical-calculus/nroots/4}{A triple root is returned with its multiplicity, three copies of
\(1\).}
\pair{numerical-calculus/nroots/5}{Wilkinson's notoriously ill-conditioned polynomial \(\prod_{k=1}^{15}
(x-k)\): Mathilda recovers all fifteen integer roots cleanly.}
\pair{numerical-calculus/nroots/6}{Precision is set by \B{PrecisionGoal}: thirty digits of \(\sqrt2\).}
\end{codepairs}

\usagebox{NRoots}

\calloutnote{Under the hood}{At machine precision \B{NRoots}'s \texttt{Automatic} method forms the
companion matrix and finds its eigenvalues with LAPACK---exactly what \texttt{numpy.roots} does---then
polishes each root with one Newton step, which is what keeps Wilkinson's polynomial (pair~5) honest. At
higher precision it switches to the Aberth--Ehrlich method\index{Aberth--Ehrlich method}, which refines
\emph{all} the roots simultaneously, each Newton correction repelled from the others so they cannot
collide~\cite{aberth}. A note on options: unlike its siblings, \B{NRoots} takes no \texttt{WorkingPrecision}---%
you steer its precision with \B{PrecisionGoal} alone.}

\subsection{Equations and systems}
\label{ssec:nsolve}

\B{NSolve}\index{NSolve@\mcode{NSolve}!methods} is the general numerical equation solver. For a single
polynomial it defers to \B{NRoots}; for a polynomial \emph{system} it runs a Gröbner-basis elimination and
reads the solutions off an eigenvalue problem; and for a transcendental equation it seeds \B{FindRoot} from
a grid of starting points. A domain argument restricts the answers:

\begin{codepairs}
\pair{numerical-calculus/nsolve/1}{Over the default complexes, \(x^4=1\) has four roots.}
\pair{numerical-calculus/nsolve/2}{Restricting to \texttt{Reals} keeps only \(\pm1\).}
\pair{numerical-calculus/nsolve/3}{A genuine system: \(x^2+y^2=1\), \(x^3-y^3=2\) has six complex
solutions, found by the Gröbner/eigenvalue engine.}
\pair{numerical-calculus/nsolve/4}{A transcendental equation, \(e^x-x=7\), solved over the reals by
grid-seeded \B{FindRoot}: two roots.}
\end{codepairs}

\usagebox{NSolve}

\calloutnote{Theory}{For a zero-dimensional polynomial system Mathilda computes a Gröbner basis
(\S\ref{sec:algebra}), which makes the quotient ring \(\mathbb{Q}[x]/I\) a finite-dimensional vector space.
Multiplication by each variable is then a linear map on that space, and the solutions are read from the
eigenvalues and eigenvectors of those multiplication matrices---the Möller--Stetter method~\cite{clo}. It
is a beautiful bridge: the algebra of \S\ref{sec:algebra} reduces the geometry of a solution set to the
linear algebra of \S\ref{sec:linalg}.}

When you want \emph{one} root and can supply a starting point, \B{FindRoot}\index{FindRoot@%
\mcode{FindRoot}!methods} is the direct tool. Its behavior follows the starting specification: a single
start uses Newton's method, two starts use the secant method, and a bracket \(\{x, x_0, x_{\min},
x_{\max}\}\) uses Brent's method. It also solves systems, with a symbolic Jacobian:

\begin{codepairs}
\pair{numerical-calculus/findroot/1}{Newton's method from \(x=1\) on \(\cos x = x\).}
\pair{numerical-calculus/findroot/2}{The same, to fifty digits---Newton's quadratic convergence makes
extra precision cheap.}
\pair{numerical-calculus/findroot/3}{A bracket \([0,1]\) selects Brent's method, which cannot leave the
bracket.}
\pair{numerical-calculus/findroot/4}{Two starting values select the derivative-free secant method.}
\pair{numerical-calculus/findroot/5}{A \(2\times2\) system solved by Newton with an exact Jacobian.}
\end{codepairs}

\usagebox{FindRoot}

\subsection{Optimization}
\label{ssec:optimization}

The last family finds minima and maxima. The distinction that organizes it is \emph{local} versus
\emph{global}. \B{FindMinimum}\index{FindMinimum@\mcode{FindMinimum}!methods} and \B{FindMaximum} start
from a point and descend to the nearest optimum; they are fast and precise but see only their own valley.
\B{NMinimize}\index{NMinimize@\mcode{NMinimize}!methods} and \B{NMaximize} search the whole region for the
\emph{global} optimum; they are slower but not fooled by local traps.

Start local. \B{FindMinimum}'s default is a quasi-Newton (BFGS) method in more than one variable, Brent's
method in one; it handles box and general constraints, and honors \texttt{WorkingPrecision}:

\begin{codepairs}
\pair{numerical-calculus/findminimum/1}{The Rosenbrock ``banana'' function, a standard hard case: the
minimum sits at \((1,1)\) in a curved valley, and BFGS threads it.}
\pair{numerical-calculus/findminimum/2}{A constrained problem: minimize \(x\cos x\) on \([1,15]\).}
\pair{numerical-calculus/findminimum/3}{Fifty digits of the minimizer of \((x-\pi)^2\)---it is \(\pi\).}
\pair{numerical-calculus/findminimum/4}{\B{FindMaximum} is the mirror image: the first peak of
\(\sin(x)\,e^{-x/10}\).}
\end{codepairs}

\usagebox{FindMinimum}

Beyond the default, \B{FindMinimum} offers a large menu of local methods---\texttt{"Newton"},
\texttt{"ConjugateGradient"}, \texttt{"LBFGSB"}, \texttt{"NelderMead"}, and the constrained solvers
\texttt{"COBYLA"}, \texttt{"COBYQA"}, and \texttt{"SLSQP"}, among others. A few names it recognizes
(\texttt{"LevenbergMarquardt"}, \texttt{"InteriorPoint"}) are not yet implemented and say so.

Now global. A function like Rastrigin's, whose graph is an egg-carton of local minima over a single global
one, would trap any purely local method; \B{NMinimize}'s default differential-evolution search is built to
escape:

\begin{codepairs}
\pair{numerical-calculus/nminimize/1}{The Rastrigin function in two variables has its global minimum
\(0\) at the origin, buried among hundreds of local minima; differential evolution finds it.}
\pair{numerical-calculus/nminimize/2}{A smooth constrained problem: minimize \(x+y\) on the unit disk,
answer \(-\sqrt2\) at \((-\tfrac1{\sqrt2},-\tfrac1{\sqrt2})\).}
\pair{numerical-calculus/nminimize/3}{Integer programming: the same objective restricted to the integer
lattice by \texttt{Element[\{x,y\}, Integers]}.}
\pair{numerical-calculus/nminimize/4}{The deterministic \texttt{"DIRECT"} method solves Rastrigin on a
box exactly, with no randomness at all.}
\end{codepairs}

\usagebox{NMinimize}

\B{NMinimize}'s global methods---\texttt{"DifferentialEvolution"}, \texttt{"SimulatedAnnealing"},
\texttt{"DualAnnealing"}, \texttt{"BasinHopping"}, \texttt{"SHGO"}, \texttt{"DIRECT"}, and others---each
take a bag of sub-options through \texttt{Method -> \{"Name", "SearchPoints" -> \dots\}}; the usage card
above and the reference page enumerate them.

\calloutnote{Theory}{Differential evolution\index{differential evolution} keeps a population of candidate
points and breeds new ones by adding the scaled difference of two random members to a third---a
gradient-free move that adapts its own step size to the population's spread~\cite{stornprice}. Simulated
annealing~\cite{kirkpatrick} instead walks one point randomly, accepting uphill moves with a probability
that cools over time, so it can climb out of a local basin early and settle late. And DIRECT~\cite{direct}
is deterministic: it repeatedly bisects the most promising rectangles of the domain, balancing global reach
against local refinement---which is why pair~4 lands exactly on the origin every time.}

\calloutnote{Performance}{Global optimizers evaluate the objective thousands of times, so Mathilda compiles
it to bytecode (Chapter~\ref{ch:compilation}) before the search begins. Supplying an \texttt{EvaluationMonitor}
forces every evaluation back through the interpreter and disables that fast path---useful for debugging a
search, costly in a production run.}

\subsection*{Where this connects}

Numerical calculus is the pragmatic mirror of the exact mathematics in the rest of this chapter. Each
function here has an exact sibling in \S\ref{sec:calculus}---\B{ND} to \B{D}, \B{NIntegrate} to
\B{Integrate}, \B{NLimit} to \B{Limit}, \B{NResidue} to \B{Residue}---and reaches for a number precisely
where the exact route runs out. Underneath, the numerical solvers lean on the machinery of the neighboring
sections: \B{NSolve} on the Gröbner bases of \S\ref{sec:algebra}, and \B{NRoots} and \B{NSolve} on the
eigenvalue kernels of \S\ref{sec:linalg}. For the exhaustive options and corner cases of any function named
here, its reference page is one \texttt{?Name} away at the prompt.

\section{Number Theory}
\label{sec:numbertheory}

The previous section computed \emph{numbers}: quantities carried to whatever precision
you asked for, where the last digit is always an approximation and the question is only
how good an approximation. This section returns to the opposite pole---the exact
arithmetic of the whole numbers, where nothing rounds and nothing overflows. An integer
in Mathilda is not a 64-bit machine word straining against its limits but, the moment it
outgrows one, a GMP\index{GMP} big integer that grows without bound, so
\(\B{Fibonacci}[200]\) and \(2^{1000}\) are computed to the last digit as readily as
\(2+2\). Number theory is the mathematics of exactly these objects, and it is where a
computer algebra system shows a face very unlike a calculator's: not \emph{how much} but
\emph{how}---how an integer factors, which integers are prime, how they behave under a
modulus, how a rational unfolds into a continued fraction.

We take the subject in the classical order. First the greatest common divisor and the
Euclidean algorithm that computes it (\B{GCD}, \B{LCM}, \B{ExtendedGCD}, \B{CoprimeQ},
\B{Divisible}); then primes and how Mathilda tests, counts, and finds them (\B{PrimeQ},
\B{Prime}, \B{PrimePi}, \B{NextPrime}); then the factorization pipeline that breaks a
composite into primes (\B{FactorInteger}, \B{IntegerExponent}, \B{SquareFreeQ}); the
classical arithmetic functions built on that factorization (\B{Divisors},
\B{DivisorSigma}, \B{EulerPhi}, \B{MoebiusMu}, \B{LiouvilleLambda}, \B{PrimeOmega},
\B{PrimeNu}); modular arithmetic and the structure of the residues (\B{PowerMod},
\B{MultiplicativeOrder}, \B{PrimitiveRoot}, \B{JacobiSymbol}); continued fractions
(\B{ContinuedFraction}, \B{FromContinuedFraction}); the digit and combinatorial functions
(\B{IntegerDigits}, \B{Binomial}, \B{PartitionsP}); and we close by assembling almost all
of it into a working RSA cryptosystem.

\subsection{The Euclidean algorithm and B\'ezout}
\label{ssec:nt-euclid}

The greatest common divisor is the foundation stone. \B{GCD} returns the largest integer
dividing all of its arguments and \B{LCM}\index{least common multiple} the smallest they all divide; both take any
number of arguments and both were met in passing in \S\ref{sec:arithmetic}. What makes
them a starting point rather than a footnote is the algorithm underneath---Euclid's, the
oldest nontrivial algorithm still in daily use\index{Euclidean algorithm}---and the
identity that ties the two functions together:

\begin{codepairs}
\pair{number-theory/euclid/1}{The largest integer dividing both \(24\) and \(36\).}
\pair{number-theory/euclid/2}{And the smallest that \(4\) and \(6\) both divide.}
\pair{number-theory/euclid/3}{Their product, \(\gcd\cdot\operatorname{lcm}\), \dots}
\pair{number-theory/euclid/4}{\dots equals the product of the numbers themselves---a
classic identity, true for any pair.}
\end{codepairs}

\B{ExtendedGCD} does more than report the gcd: it returns a \emph{B\'ezout
certificate}\index{Bezout's identity@B\'ezout's identity}, integers expressing the gcd as
a linear combination of the inputs. For \(\B{ExtendedGCD}[a,b]=\{g,\{s,t\}\}\) you always
have \(sa+tb=g\):

\begin{codepairs}
\pair{number-theory/euclid/5}{Here \(\gcd(240,46)=2\), with the certificate coefficients
\(s=-9\), \(t=47\).}
\pair{number-theory/euclid/6}{Read in the right order they check out: \((-9)\cdot240 +
47\cdot46 = 2\). These coefficients are exactly what inverts one number modulo another---
the engine behind the modular division of \S\ref{ssec:nt-modular}.}
\pair{number-theory/euclid/7}{\B{CoprimeQ}\index{coprime} asks the yes/no question directly: \(14\) and
\(15\) share no factor.}
\pair{number-theory/euclid/8}{\(14\) and \(21\) share a \(7\), so they are not coprime.}
\pair{number-theory/euclid/9}{\B{Divisible} tests one-way divisibility: \(4\mid100\).}
\pair{number-theory/euclid/10}{\(7\) does not divide \(100\).}
\end{codepairs}

\usagebox{ExtendedGCD}

\calloutnote{Under the hood}{Euclid's algorithm replaces the pair \((a,b)\) with
\((b,\,a\bmod b)\) until the second entry is zero; the extended form carries the two
B\'ezout coefficients along the same recurrence. Mathilda does not run this in
interpreted code---it hands the integers to GMP's \texttt{mpz\_gcd} and
\texttt{mpz\_gcdext}, whose subquadratic ``half-GCD'' handles thousand-digit inputs in
microseconds. The partial quotients Euclid discards on the way are, as
\S\ref{ssec:nt-contfrac} will show, precisely a continued fraction~\cite{knuth-taocp2}.}

\subsection{Primes and primality}
\label{ssec:nt-primes}

A prime is an integer greater than \(1\) with no divisors but \(1\) and itself, and the
primes are the multiplicative atoms from which every integer is built. \B{PrimeQ} is the
test:

\begin{codepairs}
\pair{number-theory/primes/1}{\(97\) is prime.}
\pair{number-theory/primes/2}{\(91=7\cdot13\) is not---a number that looks prime until you
try to factor it.}
\pair{number-theory/primes/3}{\B{Prime} indexes the primes directly: the hundredth prime
is \(541\).}
\pair{number-theory/primes/4}{\B{PrimePi} counts them: there are \(25\) primes up to
\(100\).}
\end{codepairs}

The counting function \(\pi(x)\) is where number theory meets analysis. The Prime Number
Theorem\index{Prime Number Theorem} says \(\pi(x)\sim x/\ln x\), and Mathilda lets you
watch the approximation at \(x=10^{6}\):

\begin{codepairs}
\pair{number-theory/primes/5}{The exact count below a million.}
\pair{number-theory/primes/6}{The theorem's estimate \(x/\ln x\)---low by about eight
percent here, an error the theorem promises shrinks (slowly) to zero.}
\pair{number-theory/primes/7}{\B{NextPrime} finds the next prime above its argument,
instantly even past a million.}
\pair{number-theory/primes/8}{A negative second argument walks the other way, so this is
the largest prime \emph{below} \(100\)---there is no separate \texttt{PreviousPrime}.}
\end{codepairs}

\usagebox{PrimeQ}
\usagebox{PrimePi}

\B{PrimePi} is not one algorithm but a family, selectable through \texttt{Method}. Below a
threshold Mathilda simply sieves; above it, counting every prime is hopeless and it uses
a combinatorial method that never enumerates them---all agreeing, of course, on the
answer:

\begin{codepairs}
\pair{number-theory/primes/9}{Legendre's sieve-free recurrence.}
\pair{number-theory/primes/10}{The modern Del\'eglise--Rivat method---same count,
reachable to far larger \(x\).}
\end{codepairs}

A last twist: the notion of ``prime'' depends on the ring. In the Gaussian
integers\index{Gaussian integers} \(\mathbb{Z}[i]\) a rational prime may split, and
\B{PrimeQ} tests this too under \texttt{GaussianIntegers -> True}:

\begin{codepairs}
\pair{number-theory/primes/11}{\(5=(2+i)(2-i)\) is \emph{not} a Gaussian prime---it splits.}
\pair{number-theory/primes/12}{\(3\) stays prime in \(\mathbb{Z}[i]\); it is inert.}
\pair{number-theory/primes/13}{And \(2+3i\) is a Gaussian prime, its norm \(13\) being a
rational prime.}
\end{codepairs}

\calloutnote{Under the hood}{\B{PrimeQ} is not trial division. For machine-sized integers
it runs a deterministic Miller--Rabin test over a fixed witness set; for larger integers
it uses the Baillie--PSW test (a strong Fermat test composed with a Lucas test), for which
no counterexample is known below the ranges Mathilda will ever see. Both are
probabilistic in theory and, at these sizes, certain in practice~\cite{crandallpomerance}.}

\subsection{Integer factorization}
\label{ssec:nt-factor}

Testing primality is easy; \emph{factoring} a composite is the hard problem on which much
of modern cryptography rests. \B{FactorInteger} returns the prime factorization as a list
of \(\{\text{prime},\text{exponent}\}\) pairs:

\begin{codepairs}
\pair{number-theory/factorization/1}{\(360=2^{3}\cdot3^{2}\cdot5\).}
\pair{number-theory/factorization/2}{The fifth Fermat number \(2^{32}+1\). Fermat
conjectured all such numbers prime; Euler found this factor \(641\) in
\(1732\)\index{Fermat number}, and Mathilda recovers his result at once.}
\pair{number-theory/factorization/3}{A pandigital composite peeled apart into eight primes.}
\pair{number-theory/factorization/4}{A \(19\)-digit semiprime---two ten-digit primes---
factored in milliseconds. Its size is trivial for Mathilda but hints at the wall that RSA
hides behind: multiply the primes to hundreds of digits and no known method finishes.}
\end{codepairs}

\usagebox{FactorInteger}

The single call above is a whole cascade of algorithms, and \texttt{Method} lets you
name one explicitly:

\begin{codepairs}
\pair{number-theory/factorization/5}{The same semiprime, factored by Pollard's rho method
on demand\index{Pollard rho method}.}
\pair{number-theory/factorization/6}{\B{IntegerExponent} reads off the multiplicity of a
base directly: \(10^{6}\parallel10^{6}\).}
\pair{number-theory/factorization/7}{The power of \(2\) dividing \(360\).}
\pair{number-theory/factorization/8}{\B{SquareFreeQ}\index{squarefree}: \(30=2\cdot3\cdot5\) has no repeated
prime.}
\pair{number-theory/factorization/9}{\(12=2^{2}\cdot3\) does, so it is not square-free.}
\end{codepairs}

\calloutnote{Under the hood}{\B{FactorInteger} runs a pipeline, each stage suited to a
different size of factor. It strips small primes by trial division, then hunts medium
factors with Pollard's rho (Brent's variant) and the \(p\pm1\) methods, which succeed when
\(p-1\) or \(p+1\) is smooth. For the genuinely hard factors it turns to Lenstra's
elliptic-curve method (ECM)\index{elliptic curve method}, whose running time depends on
the size of the factor found rather than the number
factored~\cite{lenstraecm,crandallpomerance}, with continued-fraction and quadratic-sieve
routines behind that. The default chooses among them by size; \texttt{Method} overrides
the choice.}

\calloutnote{Performance}{The asymmetry between \S\ref{ssec:nt-primes} and this section is
the whole point. Deciding that a \(300\)-digit number is composite takes milliseconds;
producing its factors is, as far as anyone knows, infeasible. RSA
(\S\ref{ssec:nt-rsa}) is exactly this gap, weaponised.}

\subsection{Arithmetic functions}
\label{ssec:nt-arith}

Once an integer is factored, a whole family of \emph{arithmetic
functions}\index{arithmetic function}---functions of a single integer, defined through its
factorization---becomes computable~\cite{hardywright}. Most are \emph{multiplicative}\index{multiplicative
function}: their value on \(mn\) is the product of their values on \(m\) and \(n\) when
\(m,n\) are coprime. The most basic list or count the divisors:

\begin{codepairs}
\pair{number-theory/arithmetic-functions/1}{\B{Divisors} lists every positive divisor of
\(36\), in order.}
\pair{number-theory/arithmetic-functions/2}{\B{DivisorSigma} with first argument \(0\)
counts them: \(36\) has nine divisors.}
\pair{number-theory/arithmetic-functions/3}{With first argument \(1\) it sums them;
\(\sigma_1(28)=56\), a fact we return to below.}
\end{codepairs}

Euler's totient \(\varphi(n)\) counts the integers in \(1,\dots,n\) coprime to \(n\)---
equivalently, the size of the group of units modulo \(n\):

\begin{codepairs}
\pair{number-theory/arithmetic-functions/4}{\(\varphi(36)=12\).}
\pair{number-theory/arithmetic-functions/5}{The same number from the product formula
\(n\prod_{p\mid n}(1-1/p)\)---this is what ``multiplicative'' buys you.}
\end{codepairs}

The remaining functions read off the shape of the factorization. The M\"obius function
\(\mu(n)\)\index{Mobius function@M\"obius function} is \(0\) when a square divides \(n\)
and \(\pm1\) otherwise; Liouville's \(\lambda(n)\), \(\Omega(n)\), and \(\nu(n)\) count
prime factors with and without multiplicity:

\begin{codepairs}
\pair{number-theory/arithmetic-functions/6}{\(\mu(30)=(-1)^{3}\): three distinct primes,
none repeated.}
\pair{number-theory/arithmetic-functions/7}{\(\mu(12)=0\): the square \(4\) divides it.}
\pair{number-theory/arithmetic-functions/8}{\(\lambda(12)=(-1)^{\Omega(12)}=(-1)^{3}=-1\).}
\pair{number-theory/arithmetic-functions/9}{\B{PrimeOmega} counts primes \emph{with}
multiplicity: \(360=2^3\cdot3^2\cdot5\) has \(\Omega=6\).}
\pair{number-theory/arithmetic-functions/10}{\B{PrimeNu} counts \emph{distinct} primes:
\(\nu(360)=3\).}
\pair{number-theory/arithmetic-functions/11}{Like the divisor functions, these extend to
\(\mathbb{Z}[i]\): the Gaussian divisors of \(5\) exhibit its splitting \(5=(2+i)(2-i)\).}
\end{codepairs}

\usagebox{DivisorSigma}
\usagebox{EulerPhi}
\usagebox{MoebiusMu}

These functions are not curiosities; the divisor sum is the gateway to one of the oldest
questions in mathematics. A \emph{perfect number}\index{perfect number} equals the sum of
its proper divisors---equivalently, \(\sigma_1(n)=2n\)---and Mathilda finds the first four
by brute force in an instant:

\begin{codepairs}
\pair{number-theory/perfect/1}{Every perfect number below \(10^{4}\): \(6,28,496,8128\).}
\pair{number-theory/perfect/2}{Euclid's formula \(2^{p-1}(2^{p}-1)\) for \(p=2,3,5,7\)
reproduces exactly that list\index{Euclid--Euler theorem}.}
\end{codepairs}

The match is the Euclid--Euler theorem: \(2^{p-1}(2^{p}-1)\) is perfect precisely when
\(2^{p}-1\) is a \emph{Mersenne prime}\index{Mersenne prime}. So the even perfect numbers
are in exact correspondence with the Mersenne primes, and testing one is testing the
other:

\begin{codepairs}
\pair{number-theory/perfect/3}{\(2^{13}-1=8191\) is a Mersenne prime, giving the fifth
perfect number.}
\pair{number-theory/perfect/4}{But \(2^{11}-1\) is \emph{not} prime---the exponent being
prime is necessary, not sufficient.}
\pair{number-theory/perfect/5}{Indeed \(2^{11}-1=2047=23\cdot89\).}
\end{codepairs}

\subsection{Modular arithmetic and congruences}
\label{ssec:nt-modular}

To work modulo \(n\) is to care only about remainders. \B{Mod} (met in
\S\ref{sec:arithmetic}, and always returning a residue with the sign of the modulus) is
the entry point; the number-theoretic weight sits in what you can do inside a modulus. The
first essential is fast modular exponentiation, because forming a huge power and
\emph{then} reducing it is wasteful when you can reduce at every step:

\begin{codepairs}
\pair{number-theory/modular/1}{Forming \(7^{100}\) in full and reducing---correct, but it
builds an \(85\)-digit integer first.}
\pair{number-theory/modular/2}{\B{PowerMod} gives the same residue without ever forming
the giant power.}
\pair{number-theory/modular/3}{A far larger exponent, still immediate.}
\end{codepairs}

\B{PowerMod} carries two overloads that pack a surprising amount of number theory. A
negative exponent asks for the modular \emph{inverse}, and a fractional one for a modular
\emph{root}:

\begin{codepairs}
\pair{number-theory/modular/4}{The inverse of \(3\) modulo \(7\) is \(5\), since
\(3\cdot5=15\equiv1\).}
\pair{number-theory/modular/5}{A modular square root: \(3^{2}=9\equiv2\pmod 7\), computed
by the Tonelli--Shanks algorithm\index{Tonelli--Shanks algorithm}.}
\end{codepairs}

The \emph{multiplicative order} of \(a\) modulo \(n\) is the least power returning \(a\) to
\(1\); when it reaches the full \(\varphi(n)\), \(a\) is a \emph{primitive root}\index{primitive
root}, a generator of the units:

\begin{codepairs}
\pair{number-theory/modular/6}{\(2\) has order \(3\) modulo \(7\) (\(2^{3}=8\equiv1\)).}
\pair{number-theory/modular/7}{\(3\) has order \(6\)---the full group---so it is a
primitive root.}
\pair{number-theory/modular/8}{\B{PrimitiveRoot} returns the least one, \(3\).}
\pair{number-theory/modular/9}{\B{PrimitiveRootList} gives all of them; there are
\(\varphi(\varphi(7))=\varphi(6)=2\).}
\pair{number-theory/modular/10}{Euler's theorem in one line: \(3^{\varphi(10)}\equiv1\),
because \(3\) is coprime to \(10\).}
\end{codepairs}

The order also frames the hardest question in this circle of ideas. Reversing it---given a
base \(a\) and a target \(b\), find the exponent \(x\) with \(a^{x}\equiv b\pmod n\)---is the
\emph{discrete logarithm problem}\index{discrete logarithm problem}, and the three-argument
form of \B{MultiplicativeOrder} solves exactly that, returning the least exponent whose power
lands in a given residue set:

\begin{codepairs}
\pair{number-theory/modular/13}{The discrete logarithm of \(5\) to base \(3\) modulo \(7\):
the least \(x\) with \(3^{x}\equiv5\), namely \(x=5\) (indeed \(3^{5}=243\equiv5\)).}
\end{codepairs}

\calloutnote{Theory}{Computing the order of \(a\) is easy---it divides \(\varphi(n)\) and is
found by stripping prime factors from it. The discrete logarithm is the \emph{reverse}
question, and it is believed \emph{hard}: no polynomial-time algorithm is known over a
general prime field, and that presumed hardness is the foundation of Diffie--Hellman key
exchange and the ElGamal and DSA signature schemes\index{discrete logarithm problem!in
cryptography}. The best general methods---baby-step giant-step, Pollard's rho for logarithms,
Pohlig--Hellman (which reduces the problem to the prime-power factors of the group order),
and index calculus---are surveyed in~\cite{hac}. Mathilda's three-argument form simply walks
the orbit \(a^{1},a^{2},\dots\), so it is a teaching tool for small moduli, not a break of
the problem at cryptographic sizes.}

Finally, the \B{JacobiSymbol}\index{Jacobi symbol} \(\left(\tfrac{a}{n}\right)\) encodes quadratic
residues\index{quadratic reciprocity}---whether \(a\) is a square modulo \(n\):

\begin{codepairs}
\pair{number-theory/modular/11}{\(\left(\tfrac{2}{7}\right)=+1\): \(2\) is a square mod
\(7\) (indeed \(3^{2}\equiv2\), as above).}
\pair{number-theory/modular/12}{\(\left(\tfrac{3}{7}\right)=-1\): \(3\) is a non-residue.}
\end{codepairs}

\usagebox{PowerMod}
\usagebox{MultiplicativeOrder}
\usagebox{PrimitiveRoot}
\usagebox{JacobiSymbol}

\calloutnote{Theory}{Two theorems govern this section. Fermat's little
theorem\index{Fermat's little theorem}---\(a^{p-1}\equiv1\pmod p\) for a prime \(p\nmid
a\)---is the special case \(n=p\) of \emph{Euler's theorem}\index{Euler's theorem},
\(a^{\varphi(n)}\equiv1\pmod n\) whenever \(\gcd(a,n)=1\). Together with the fact that the
units modulo a prime form a \emph{cyclic} group---generated by a primitive root---they are
the whole basis of what follows, and of RSA in \S\ref{ssec:nt-rsa}~\cite{irelandrosen}.}

Mathilda has no single \texttt{ChineseRemainder} builtin, and it is instructive to see
that none is needed: the Chinese Remainder Theorem\index{Chinese Remainder Theorem} is a
short composition of the tools above. Given residues modulo pairwise-coprime moduli, we
reconstruct the unique solution modulo their product with cofactors and \B{PowerMod}
inverses---the classical recipe of Sun Zi:

\begin{codepairs}
\pair{number-theory/crt/3}{The combined modulus \(3\cdot5\cdot7=105\).}
\pair{number-theory/crt/4}{Each cofactor \(N/n_i\).}
\pair{number-theory/crt/5}{Each cofactor's inverse modulo its own \(n_i\).}
\pair{number-theory/crt/6}{The weighted sum, reduced---the unique residue \(23\).}
\pair{number-theory/crt/7}{And it checks: \(23\) leaves remainders \(2,3,2\) modulo
\(3,5,7\), as required.}
\end{codepairs}

\subsection{Continued fractions}
\label{ssec:nt-contfrac}

A continued fraction\index{continued fraction} writes a number as
\(a_0+\cfrac{1}{a_1+\cfrac{1}{a_2+\cdots}}\), a form that captures rational approximation
better than any decimal. \B{ContinuedFraction} produces the list of partial quotients and
\B{FromContinuedFraction} folds it back:

\begin{codepairs}
\pair{number-theory/continued-fractions/1}{\(415/93=[4;2,6,7]\). These quotients are
exactly Euclid's from \S\ref{ssec:nt-euclid}, run on \(415\) and \(93\).}
\pair{number-theory/continued-fractions/2}{Folding the list back recovers the rational
exactly---the two operations are inverse.}
\end{codepairs}

Irrationals have \emph{infinite} expansions, so you ask for a fixed number of terms---and
the classical patterns appear at once:

\begin{codepairs}
\pair{number-theory/continued-fractions/3}{\(\sqrt2=[1;\overline{2}]\), purely periodic.}
\pair{number-theory/continued-fractions/4}{\(\sqrt7\) has period \(1,1,1,4\).}
\pair{number-theory/continued-fractions/5}{Asked without a length, Mathilda returns the
period in a nested list---every quadratic irrational is \emph{eventually periodic}, which
is Lagrange's theorem\index{quadratic irrational}.}
\pair{number-theory/continued-fractions/6}{\(\pi\) shows no pattern, but its early terms
are the key to good approximations.}
\end{codepairs}

Truncating \(\pi\)'s expansion just after the large term \(15\) and folding it back gives
the approximation known since antiquity:

\begin{codepairs}
\pair{number-theory/continued-fractions/7}{\(355/113\), from only four terms.}
\pair{number-theory/continued-fractions/8}{It agrees with \(\pi\) to better than
\(3\times10^{-7}\)---the best rational approximation with any denominator under a thousand.}
\end{codepairs}

\usagebox{ContinuedFraction}
\usagebox{FromContinuedFraction}

\calloutnote{Theory}{A large partial quotient signals an unusually good rational
approximation just before it: truncating there makes the discarded tail
\(1/(a_{k+1}+\cdots)\) tiny. \(\pi\)'s term \(292\) is why \(355/113\) is so
extraordinarily accurate. The convergents obtained by truncation are, in a precise sense,
the \emph{best} rational approximations to a number\index{continued fraction!convergent}~\cite{khinchin}---the theory that also governs Pell's
equation, which \S\ref{ssec:alg-diophantine} solves through the continued fraction of
\(\sqrt D\).}

\subsection{Digits and bases}
\label{ssec:nt-digits}

The written form of an integer---its digits in a given base---is itself a number-theoretic
object, and Mathilda exposes it directly. \B{IntegerDigits} and \B{FromDigits} convert both
ways, in any base:

\begin{codepairs}
\pair{number-theory/digits/1}{The base-ten digits of \(123456\).}
\pair{number-theory/digits/2}{In base two, \(255\) is eight ones.}
\pair{number-theory/digits/3}{In base sixteen it is two ``fifteens''.}
\pair{number-theory/digits/4}{\B{FromDigits} rebuilds the integer from its base-two digits.}
\pair{number-theory/digits/5}{\B{IntegerString} renders the digits as text, here the
familiar hexadecimal \texttt{ff}.}
\pair{number-theory/digits/6}{\B{IntegerLength} counts the digits without listing them.}
\end{codepairs}

The digit sum is the basis of the schoolroom rule of \emph{casting out nines}---a number
is congruent to its digit sum modulo \(9\), because \(10\equiv1\):

\begin{codepairs}
\pair{number-theory/digits/7}{\B{DigitSum} adds the base-ten digits: \(1+2+3+4+5+6=21\).}
\pair{number-theory/digits/8}{Its residue modulo \(9\) \dots}
\pair{number-theory/digits/9}{\dots equals the number's own residue modulo \(9\). That is
casting out nines.}
\pair{number-theory/digits/10}{\B{DigitCount} tallies how often each digit occurs---one
\(1\), two \(2\)s, three \(3\)s, four \(4\)s (the final slot counts the zeros).}
\end{codepairs}

\usagebox{IntegerDigits}

\subsection{Combinatorial number theory}
\label{ssec:nt-combinatorial}

Number theory shades into combinatorics through the factorials, binomial coefficients, and
partition functions---all computed exactly, all growing explosively. The factorial family
comes first:

\begin{codepairs}
\pair{number-theory/combinatorial/1}{\(20!\), exact to the last of its nineteen digits.}
\pair{number-theory/combinatorial/2}{The double factorial \(9!!=9\cdot7\cdot5\cdot3\cdot1\).}
\pair{number-theory/combinatorial/3}{The falling factorial \(10\cdot9\cdot8\), via
\B{FactorialPower}.}
\pair{number-theory/combinatorial/4}{\B{Binomial}: the number of \(3\)-subsets of a
\(10\)-set.}
\pair{number-theory/combinatorial/5}{Reduced modulo \(7\) it is \(1\)---as Lucas's theorem
predicts from the base-\(7\) digits of \(10\) and \(3\)\index{Lucas's theorem}.}
\end{codepairs}

The partition functions count the ways to write \(n\) as a sum of positive integers.
\B{IntegerPartitions} lists them and \B{PartitionsP} counts them, and the count grows so
fast that its exact evaluation is a small miracle of analysis:

\begin{codepairs}
\pair{number-theory/combinatorial/6}{The seven partitions of \(5\).}
\pair{number-theory/combinatorial/7}{Their count, \(p(10)=42\).}
\pair{number-theory/combinatorial/8}{\(p(100)=190569292\)---already far past listing.}
\pair{number-theory/combinatorial/9}{\(p(1000)\), a \(32\)-digit integer computed exactly.}
\pair{number-theory/combinatorial/10}{\B{PartitionsQ} counts partitions into
\emph{distinct} parts.}
\end{codepairs}

Finally the Fibonacci and Lucas numbers, the archetypal integer sequences:

\begin{codepairs}
\pair{number-theory/combinatorial/11}{\(F_{50}\).}
\pair{number-theory/combinatorial/12}{\(F_{200}\), a \(42\)-digit number, still instant.}
\pair{number-theory/combinatorial/13}{The tenth Lucas number.}
\pair{number-theory/combinatorial/14}{A jewel of an identity: \(\gcd(F_m,F_n)=F_{\gcd(m,n)}\),
here \(\gcd(F_{12},F_{18})=F_{6}\).}
\end{codepairs}

\usagebox{Binomial}
\usagebox{PartitionsP}
\usagebox{Fibonacci}

\calloutnote{Under the hood}{\B{PartitionsP} does not count partitions one by one---that
would be hopeless at \(n=1000\). It evaluates the Hardy--Ramanujan--Rademacher
series\index{Hardy--Ramanujan formula}~\cite{andrews}, a convergent sum of the analytic terms that gives
\(p(n)\) exactly after enough of them, rounded to the nearest integer. \B{Fibonacci} uses
fast doubling---computing \(F_{2k}\) and \(F_{2k+1}\) from \(F_k\) and \(F_{k+1}\)---so
\(F_n\) costs \(O(\log n)\) big-integer multiplications rather than \(n\) additions. The
close relatives \B{BernoulliB}, \B{EulerE}, and \B{Zeta}, whose values encode still deeper
number theory, belong to the special functions of \S\ref{sec:specialfns}.}

\subsection{Putting it together: textbook RSA}
\label{ssec:nt-rsa}

Almost every tool of this section meets in one place: the RSA cryptosystem\index{RSA},
whose security is exactly the factoring gap of \S\ref{ssec:nt-factor}. We use the historic
textbook primes \(p=61\), \(q=53\) so every step is legible; the session is continuous,
each line building on the last. First the public modulus and its secret totient:

\begin{codepairs}
\pair{number-theory/rsa/1}{Pick the first prime, \(p=61\)\dots}
\pair{number-theory/rsa/2}{\dots and the second, \(q=53\).}
\pair{number-theory/rsa/3}{The modulus \(n=pq=3233\) is public; its factors are the secret.}
\pair{number-theory/rsa/4}{The totient \(\varphi(n)=3120\)\dots}
\pair{number-theory/rsa/5}{\dots which for a product of two primes is \((p-1)(q-1)\), as
knowing the factorization lets us confirm.}
\end{codepairs}

Next a public exponent coprime to \(\varphi(n)\), and the private exponent that inverts it:

\begin{codepairs}
\pair{number-theory/rsa/6}{The conventional public exponent \(e=17\)\dots}
\pair{number-theory/rsa/7}{\dots coprime to \(\varphi(n)\), so it is invertible.}
\pair{number-theory/rsa/8}{\B{ExtendedGCD} already exposes the inverse in its B\'ezout
coefficient: \((-367)\cdot17+2\cdot3120=1\).}
\pair{number-theory/rsa/9}{The private key \(d\) is that inverse as a positive residue,
\(2753\) (indeed \(-367+3120\)).}
\pair{number-theory/rsa/10}{Confirming the defining relation \(ed\equiv1\pmod{\varphi(n)}\).}
\end{codepairs}

Now encryption raises the message to \(e\) modulo \(n\), and decryption raises the
ciphertext to \(d\):

\begin{codepairs}
\pair{number-theory/rsa/11}{The plaintext number \(42\).}
\pair{number-theory/rsa/12}{Encrypted, it becomes \(2557\)---bearing no visible relation
to the original.}
\pair{number-theory/rsa/13}{Decryption with the private exponent returns \(42\) exactly.}
\pair{number-theory/rsa/14}{The round trip closes.}
\end{codepairs}

Why it works is Euler's theorem from \S\ref{ssec:nt-modular}: because \(ed\equiv1\pmod{\varphi(n)}\),
raising the message to \(e\) and then to \(d\) raises it to a power \(\equiv1\), which
leaves it unchanged modulo \(n\). Anyone holding \((n,e)\) can encrypt; only the holder of
\(d\)---recoverable only by factoring \(n\)---can decrypt. Every step ran on integers no
larger than a few thousand, but nothing about the method changes when the primes have a
hundred digits each; only the factoring stays out of reach.

\subsection*{Where this connects}

Number theory is the exact spine of Mathilda's arithmetic. Its foundation is the integer
model of \S\ref{sec:arithmetic}---the same floored \B{Mod} and GMP big integers, here put
to work. Its natural sequel is the arithmetic of \emph{solutions}: the classical Diophantine
questions---Pell's equation, Pythagorean triples, sums of squares, and the rest---are
answered by \B{Solve} and \B{Reduce} over the integers, treated in
\S\ref{ssec:alg-diophantine}, where the continued fractions of \S\ref{ssec:nt-contfrac}
reappear as the engine behind Pell. The polynomial \B{Factor} of \S\ref{sec:algebra} is the
algebraic mirror of \B{FactorInteger}, square-free decomposition and Hensel lifting
standing in for trial division and rho. And the integer-valued special functions that shade
number theory into analysis---\B{BernoulliB}, \B{Zeta}, and their kin---await in
\S\ref{sec:specialfns}. Underneath every result here is GMP, and the same exact integers
feed the compiler of Chapter~\ref{ch:compilation} when a number-theoretic loop needs to run
at machine speed.

\section{Special Functions}
\label{sec:specialfns}

The elementary functions---polynomials, the exponential, the trigonometric and
their inverses---run out long before mathematics does. The moment you write down
an integral like \(\int_0^\infty t^{z-1}e^{-t}\,dt\), or a power series whose
coefficients grow by a fixed ratio, or the solution of a second-order linear
differential equation with a variable coefficient, you have left the elementary
world and entered the realm of the \emph{special functions}: the gamma and zeta
functions, the error function, the Bessel and Airy functions, the Legendre and
hypergeometric families.\index{special functions} They are ``special'' only in
the sense that they were named and tabulated before computers existed; today they
are as ordinary as \(\sin\), and Mathilda treats them so.

What makes a special function useful in a computer algebra system is that it wears
three faces at once, and Mathilda shows you all three. There are the
\emph{exact symbolic values}---\(\Gamma(1/2)=\sqrt\pi\), \(\zeta(2)=\pi^2/6\)---which a
numeric library can only approximate. There are the \emph{symbolic identities}: the
recurrences, reflection and derivative rules, series expansions, and orthogonality
relations that let these functions participate in algebra rather than merely
returning numbers. And there is \emph{rigorous arbitrary-precision numerics}, for
when a number really is what you want, computed to any precision you ask for and,
where the kernel is backed by FLINT's ball arithmetic, with a guaranteed error
bound. This section takes the major families in turn, and the same three faces
recur in each: what closes into a clean symbol, what identity the system knows, and
how it puts a number on the value.

\subsection{The gamma family}
\label{ssec:sf-gamma}

Everything begins with the \B{Gamma} function, \(\Gamma(z)=\int_0^\infty
t^{z-1}e^{-t}\,dt\), the continuous interpolation of the
factorial.\index{gamma function} Its defining property, \(\Gamma(z+1)=z\,\Gamma(z)\)
with \(\Gamma(1)=1\), forces \(\Gamma(n)=(n-1)!\) at the integers, and from the
single exact value \(\Gamma(1/2)=\sqrt\pi\) the recurrence generates the whole
ladder of half-integers.

\begin{codepairs}
\pair{special-functions/gamma/1}{At the positive integers \B{Gamma} is the shifted
factorial: \(\Gamma(n)=(n-1)!\), so this is \(0!\) through \(5!\).}
\pair{special-functions/gamma/2}{The one transcendental seed value,
\(\Gamma(1/2)=\sqrt\pi\), returned exactly.}
\pair{special-functions/gamma/3}{The recurrence climbs from it:
\(\Gamma(7/2)=\tfrac{15}{8}\sqrt\pi\), a rational multiple of \(\sqrt\pi\).}
\pair{special-functions/gamma/4}{And descends to the negative half-integers---here
\(\Gamma(-1/2)=-2\sqrt\pi\), finite because the poles fall only on the
\emph{integers}.}
\pair{special-functions/gamma/5}{At the non-positive integers \(\Gamma\) has simple
poles, reported as \mcode{ComplexInfinity}.}
\pair{special-functions/gamma/6}{The two-argument form is the upper incomplete
gamma \(\Gamma(a,z)\); at a positive integer \(a\) it closes into elementary
form.}
\pair{special-functions/gamma/7}{With a numeric second argument that closed form is
an exact number: \(\Gamma(5,2)=168/e^{2}\).}
\pair{special-functions/gamma/8}{Where no symbol exists, \B{N} delivers digits---here
\(\Gamma(i)\) to thirty places, real and imaginary parts together.}
\pair{special-functions/gamma/9}{And \(\Gamma\) knows its own derivative, which
closes through the digamma function \(\psi=\Gamma'/\Gamma\).}
\end{codepairs}

The half-integer ladder (pairs~2--4) is the gamma function's signature, and it is
worth seeing why it terminates in \(\sqrt\pi\) rather than another special value.
\calloutnote{Theory}{The reflection formula
\(\Gamma(z)\,\Gamma(1-z)=\pi/\sin(\pi z)\)\index{gamma function!reflection formula}
evaluated at \(z=1/2\) gives \(\Gamma(1/2)^2=\pi\), hence \(\Gamma(1/2)=\sqrt\pi\).
Combined with the recurrence \(\Gamma(z+1)=z\,\Gamma(z)\), it pins every
half-integer to a rational multiple of \(\sqrt\pi\). The same reflection formula
explains the pole structure: since \(1/\Gamma\) is entire with zeros exactly at
\(0,-1,-2,\dots\), \(\Gamma\) itself has simple poles there and nowhere else, which
is why pair~4 is finite but pair~5 is not.} The numeric face (pairs~7--8) is the
one a differential-equation solver or a statistics routine actually leans
on.\calloutnote{Under the hood}{\B{Gamma} lives in
\texttt{src/special\_functions/gamma.c}. A machine-precision real routes to the C
library's \texttt{tgamma}; an arbitrary-precision real to MPFR's correctly rounded
\texttt{mpfr\_gamma}; and an arbitrary-precision \emph{complex} argument to Spouge's
approximation, whose coefficients are recomputed at runtime to the precision you
request, with the reflection formula handling \(\operatorname{Re}(z)<\tfrac12\). The
incomplete form uses \texttt{mpfr\_gamma\_inc} on the real axis and a lower-incomplete
series or a Lentz continued fraction in the complex plane.}

\usagebox{Gamma}

The \B{Beta} function \(B(a,b)=\Gamma(a)\Gamma(b)/\Gamma(a+b)\) is the gamma
function's closest relative---the normalizing constant of the beta distribution and
the value of the integral \(\int_0^1 t^{a-1}(1-t)^{b-1}\,dt\).\index{beta function}
Because it is built entirely from gammas, its exact reductions follow from theirs.

\begin{codepairs}
\pair{special-functions/beta/1}{A positive integer argument collapses the gamma
ratio to a rational: \(B(2,3)=1/12\).}
\pair{special-functions/beta/2}{Even with a rational partner it stays exact,
\(B(3,1/3)=27/14\).}
\pair{special-functions/beta/3}{Half-integers fold to rational multiples of \(\pi\)
instead---\(B(1/2,1/2)=\pi\), the arcsine normalization.}
\pair{special-functions/beta/4}{And \(B(5/2,7/2)=\tfrac{3}{256}\pi\).}
\pair{special-functions/beta/5}{When no such collapse is available, the residual
gamma form survives rather than being forced to a number.}
\pair{special-functions/beta/6}{A fully numeric call evaluates like any other.}
\end{codepairs}

Pair~5 is a small lesson in the system's temperament: \(B(1/3,1/3)\) is left as
\(\Gamma(1/3)^2/\Gamma(2/3)\) because no further exact simplification is available
without the reflection formula, which Mathilda does not apply automatically to a
product of gammas at complementary rationals.\calloutnote{Pitfall}{Do not expect
\(\Gamma(1/3)\,\Gamma(2/3)\) to collapse to \(2\pi/\sqrt3\) on its own. The gamma
reflection and multiplication identities are known to the theory but are \emph{not}
auto-applied---applying them blindly would fight as often as it helps---so products
of gammas at rational points are carried symbolically until you rewrite them by
hand. This is the honest cost of not over-simplifying.}

Two functions read the shape of the gamma function near its poles and along the
real axis. \B{LogGamma} is the log-gamma function \(\log\Gamma(z)\): the analytic
continuation of \(\log(\Gamma(z))\), not the same thing as \(\log(\Gamma(z))\)
term by term, and the right primitive whenever \(\Gamma\) itself would overflow.
\B{PolyGamma} is the family of logarithmic derivatives: \(\psi(z)=\psi^{(0)}(z)=
\Gamma'(z)/\Gamma(z)\) is the digamma function,\index{digamma function} and
\(\psi^{(n)}\) its \(n\)-th derivative.

\begin{codepairs}
\pair{special-functions/polygamma/1}{\B{LogGamma} keeps integers exact as the log of
a factorial: \(\log\Gamma(10)=\log(362880)=\log(9!)\).}
\pair{special-functions/polygamma/2}{At \(z=171\) the gamma function has already
overflowed a machine double---}
\pair{special-functions/polygamma/3}{---but \B{LogGamma} sails on, returning
\(\approx706.6\). This is why factorial ratios and log-likelihoods are computed
through it, never through \(\Gamma\) directly.}
\pair{special-functions/polygamma/4}{The digamma at a positive integer is a rational
minus Euler's constant: \(\psi(5)=\tfrac{25}{12}-\gamma\).}
\pair{special-functions/polygamma/5}{Gauss's digamma theorem closes rational
arguments into elementary form: \(\psi(1/2)=-\gamma-\log 4\).}
\pair{special-functions/polygamma/6}{The trigamma at \(1\) is \(\zeta(2)=\pi^2/6\).}
\pair{special-functions/polygamma/7}{\B{LogGamma} differentiates to digamma\ldots}
\pair{special-functions/polygamma/8}{\ldots and each further derivative raises the
polygamma order by one.}
\end{codepairs}

The contrast between pairs~2 and~3 is the entire reason \B{LogGamma} exists as a
first-class head rather than a shorthand for \texttt{Log[Gamma[z]]}: it stays finite
across a range where the gamma function runs from \(10^{-300}\) to \(10^{300}\), and
its imaginary part tracks the \emph{continued} argument past \(\pi\), where a naive
\(\log(\Gamma(z))\) would wrap.\calloutnote{Under the hood}{\B{PolyGamma}
(\texttt{src/special\_functions/polygamma.c}) handles order \(n=0\) at rational
arguments through the Gauss digamma theorem---a finite trigonometric sum---and closes
odd orders through \(\zeta(n{+}1)\). Numerically it uses MPFR's digamma for
\(n=0\) and a recurrence-shift plus a Bernoulli asymptotic expansion for higher
orders; complex arbitrary-precision values are handed to FLINT's \texttt{acb\_polygamma}.
The negative order \(\psi^{(-1)}(z)\) is \B{LogGamma} itself.}

The last member of the family is the \B{Pochhammer} symbol \((a)_n=
a(a+1)\cdots(a+n-1)=\Gamma(a+n)/\Gamma(a)\), the \emph{rising factorial} that
supplies the coefficients of every hypergeometric series we meet at the end of this
section.\index{Pochhammer symbol}

\begin{codepairs}
\pair{special-functions/pochhammer/1}{A symbolic base gives an explicit polynomial
product of \(n\) linear factors.}
\pair{special-functions/pochhammer/2}{A numeric base gives an exact value,
\((10)_6=3603600\).}
\pair{special-functions/pochhammer/3}{\((1)_{25}\) is just \(25!\), carried exactly
as a GMP big integer.}
\pair{special-functions/pochhammer/4}{Half-integer arguments reduce through the gamma
ratio: \((1/2)_3=15/8\).}
\pair{special-functions/pochhammer/5}{And \((3/2)_{1/2}=2/\sqrt\pi\), a gamma ratio
that happens to close.}
\end{codepairs}

\subsection{The zeta family}
\label{ssec:sf-zeta}

The Riemann zeta function \(\zeta(s)=\sum_{n\ge1}n^{-s}\) is the deepest object in
this section, and \B{Zeta} carries its exact arithmetic.\index{Riemann zeta
function} At the even positive integers it closes into powers of \(\pi\); its values
at the non-positive integers, defined by analytic continuation, are rational; and
its values at the odd positive integers are believed to be new transcendentals with
no closed form at all.

\begin{codepairs}
\pair{special-functions/zeta/1}{The even values are rational multiples of powers of
\(\pi\): \(\zeta(2)=\pi^2/6\)\ldots}
\pair{special-functions/zeta/2}{\ldots and \(\zeta(4)=\pi^4/90\).}
\pair{special-functions/zeta/3}{This is exactly the Basel sum, which \B{Sum}
recognizes independently and returns identically.}
\pair{special-functions/zeta/4}{Analytic continuation gives \(\zeta(0)=-1/2\)\ldots}
\pair{special-functions/zeta/5}{\ldots and the famous \(\zeta(-1)=-1/12\).}
\pair{special-functions/zeta/6}{The negative even integers are the \emph{trivial
zeros}: \(\zeta(-2)=0\).}
\pair{special-functions/zeta/7}{The odd values have no known closed form, so
\(\zeta(3)\) is returned as itself---a genuine constant, not a failure.}
\pair{special-functions/zeta/8}{But it is a \emph{number}: \B{N} gives Apéry's
constant to thirty digits.}
\pair{special-functions/zeta/9}{On the critical line, at the height \(t=14.1347\ldots\)
of the first nontrivial zero, \(\zeta(1/2+it)\) collapses to the order of
\(10^{-7}\)---the Riemann hypothesis made visible.}
\pair{special-functions/zeta/10}{The Hurwitz zeta \(\zeta(s,a)=\sum_{n\ge0}(n+a)^{-s}\)
shifts the denominators, and still closes: \(\zeta(2,3)=\pi^2/6-5/4\).}
\pair{special-functions/zeta/11}{The Stieltjes constants are inert symbols; the
zeroth is Euler's constant \(\gamma\).}
\end{codepairs}

Pairs~4--6 are not sleight of hand but the content of the \emph{functional
equation}, and pair~9 is a numerical glimpse of the most famous open problem in
mathematics.\calloutnote{Theory}{Riemann's functional
equation\index{Riemann zeta function!functional equation}
\(\zeta(s)=2^s\pi^{s-1}\sin(\pi s/2)\,\Gamma(1-s)\,\zeta(1-s)\) relates the value at
\(s\) to the value at \(1-s\), and with it the whole left half-plane is determined by
the right. The factor \(\sin(\pi s/2)\) vanishes at the negative even integers,
producing the trivial zeros of pair~6; the values \(\zeta(0)=-1/2\) and
\(\zeta(-1)=-1/12\) fall out of the same equation. The \emph{nontrivial} zeros all
lie in the strip \(0<\operatorname{Re}(s)<1\), and the \emph{Riemann hypothesis}
\index{Riemann hypothesis} asserts they lie exactly on the line
\(\operatorname{Re}(s)=1/2\); pair~9 evaluates \(\zeta\) at the first of them, where
it is zero to the precision requested.} The even values, by contrast, are elementary
once you know the Bernoulli numbers.\calloutnote{Theory}{Euler's formula
\(\zeta(2n)=(-1)^{n+1}\dfrac{B_{2n}(2\pi)^{2n}}{2\,(2n)!}\) expresses every even zeta
value through a \B{BernoulliB} number\index{Bernoulli numbers}; for \(n=1\),
\(B_2=1/6\) gives \(\zeta(2)=\pi^2/6\). No analogous formula is known for the odd
values, which is why \(\zeta(3)\) (pair~7) has no closed form while \(\zeta(2)\) and
\(\zeta(4)\) do.} Underneath the exact layer is a rigorous numeric
one.\calloutnote{Under the hood}{\B{Zeta}, \B{HurwitzZeta} and \B{StieltjesGamma}
(\texttt{src/special\_functions/zeta.c}, and \texttt{src/flint\_num\_bridge.c} for
the numerics) delegate arbitrary-precision \emph{complex} evaluation to FLINT's Arb
library---\texttt{acb\_dirichlet\_zeta}, \texttt{acb\_dirichlet\_hurwitz},
\texttt{acb\_dirichlet\_stieltjes}. Arb works in \emph{ball arithmetic}\index{ball
arithmetic}: every value carries a midpoint and a rigorous radius, so pair~9's
near-zero comes with a proof that the true value lies within a known ball of the
printed digits---not a hope that rounding behaved.}

\usagebox{Zeta}

The Stieltjes constants \(\gamma_n\) (pair~11) are the coefficients of the Laurent
expansion of \(\zeta\) about its pole at \(s=1\); the zeroth is Euler's constant
\(\gamma\), surfaced here as \B{EulerGamma}, and the rest are carried as inert
\B{StieltjesGamma} symbols.

\subsection{Error functions and the integral functions}
\label{ssec:sf-erf}

The error function \(\operatorname{erf}(z)=\frac{2}{\sqrt\pi}\int_0^z e^{-t^2}\,dt\)
is the antiderivative of the Gaussian, and no elementary function equals
it---which is precisely why it needs a name.\index{error function} \B{Erf}, its
complement \B{Erfc}, the imaginary variant \B{Erfi}, and the inverse \B{InverseErf}
form a small closed world.

\begin{codepairs}
\pair{special-functions/erf/1}{\(\operatorname{erf}(0)=0\)\ldots}
\pair{special-functions/erf/2}{\ldots and \(\operatorname{erf}(\infty)=1\): the total
mass of the normalized Gaussian.}
\pair{special-functions/erf/3}{It is an odd function, and the odd symmetry is applied
symbolically.}
\pair{special-functions/erf/4}{Arbitrary-precision numerics track the digits you ask
for.}
\pair{special-functions/erf/5}{The complement \(\operatorname{erfc}=1-\operatorname{erf}\)
vanishes at infinity, computed cancellation-free so the tail keeps its
significance.}
\pair{special-functions/erf/6}{The imaginary error function
\(\operatorname{erfi}(z)=-i\,\operatorname{erf}(iz)\) has its own kernel.}
\pair{special-functions/erf/7}{The derivative is the Gaussian itself---so the whole
Taylor series is generated from this one rule.}
\pair{special-functions/erf/8}{The inverse function has exact endpoints:
\(\operatorname{erf}^{-1}(1)=\infty\).}
\pair{special-functions/erf/9}{And numeric values in between, on \([-1,1]\).}
\end{codepairs}

The split between \B{Erf} and \B{Erfc} (pairs~2 and~5) is not redundancy but
numerics: for large positive \(z\), \(\operatorname{erf}(z)\) is so close to \(1\)
that \(1-\operatorname{erf}(z)\) loses every significant digit to cancellation, so
\(\operatorname{erfc}\) is computed by its own cancellation-free
route.\calloutnote{Under the hood}{The kernels live in
\texttt{src/special\_functions/erf.c} (and \texttt{inverf.c} for the inverse). Real
arguments use the C library or MPFR's \texttt{erf}/\texttt{erfc}; complex arguments
sum the cancellation-aware Maclaurin series of DLMF~7.6.2 in MPFR with
\(|z|^2/\ln2\) guard bits, so even a machine-precision complex result carries full
accuracy. Neither the C library nor MPFR ships an inverse error function, so
\B{InverseErf} is Newton's iteration on \(\operatorname{erf}(z)-s\), started from a
Winitzki seed and polished to the working precision.}

The same principle---name the antiderivative when it is not elementary---generates a
whole class of \emph{integral functions}, and Mathilda produces them directly when
you integrate.

\begin{codepairs}
\pair{special-functions/integrals/1}{The sine integral \(\operatorname{Si}\) is the
antiderivative of the sinc function, and \B{Integrate} returns it by name.}
\pair{special-functions/integrals/2}{The cosine integral \(\operatorname{Ci}\)
likewise.}
\pair{special-functions/integrals/3}{The exponential integral \(\operatorname{Ei}\)
names \(\int e^x/x\,dx\).}
\pair{special-functions/integrals/4}{And \(\int dx/\log x\) is the same function in
disguise---the logarithmic integral \(\operatorname{li}(x)=\operatorname{Ei}(\log x)\).}
\end{codepairs}

Each of these is a genuine function with the full three-faced treatment: it
evaluates, it differentiates, it has special values.

\begin{codepairs}
\pair{special-functions/integrals/5}{\B{SinIntegral} to twenty digits\ldots}
\pair{special-functions/integrals/6}{\ldots \B{CosIntegral}\ldots}
\pair{special-functions/integrals/7}{\ldots \B{ExpIntegralEi}\ldots}
\pair{special-functions/integrals/8}{\ldots and \B{LogIntegral}, whose growth rate is
the density of the primes.}
\pair{special-functions/integrals/9}{\(\operatorname{Si}\) saturates to \(\pi/2\) at
infinity---the Dirichlet integral \(\int_0^\infty \sin(x)/x\,dx\).}
\pair{special-functions/integrals/10}{The Fresnel integral \(C\) of optics and
diffraction, \B{FresnelC}\ldots}
\pair{special-functions/integrals/11}{\ldots and its partner \B{FresnelS}.}
\pair{special-functions/integrals/12}{And \(\operatorname{Si}\) differentiates back
to the cardinal sine \B{Sinc}---the loop closes.}
\end{codepairs}

That \(\int dx/\log x\) and \(\int e^x/x\,dx\) come back as the \emph{same} function
(pairs~3--4) is the logarithmic integral's claim to fame: \(\operatorname{li}(x)\)
is the leading term of the prime-counting function, and the substitution
\(t=\log x\) turns one integral into the other.\calloutnote{Under the hood}{On the
real axis \B{ExpIntegralEi} uses MPFR's correctly rounded \texttt{mpfr\_eint} for
\(x>0\) and an on-cut convergent series for \(x<0\); large \(|z|\) switches to the
asymptotic expansion, and the two regimes overlap and agree. \B{LogIntegral} is
literally \(\operatorname{Ei}(\log z)\), reusing the same kernel. The sine, cosine,
and Fresnel integrals sum their Maclaurin series for moderate arguments and their
trigonometric-prefactored asymptotic forms for large ones, with the complex plane
handled by the shared MPFR-complex toolkit.}

\subsection{Bessel and Airy functions}
\label{ssec:sf-bessel}

Bessel functions are what solutions to the wave and heat equations look like in
cylindrical coordinates, and they are the archetype of a function defined by a
differential equation rather than a formula.\index{Bessel functions} \B{BesselJ}
is the function of the first kind, regular at the origin; \B{BesselY} the second
kind, singular there; and \B{BesselI}, \B{BesselK} the modified pair for the
hyperbolic version of the equation.

\begin{codepairs}
\pair{special-functions/bessel/1}{\(J_0(0)=1\); every higher \(J_n\) vanishes at the
origin.}
\pair{special-functions/bessel/2}{At half-integer order the Bessel function is
\emph{elementary}: \(J_{1/2}(x)=\sqrt{2/\pi x}\,\sin x\).}
\pair{special-functions/bessel/3}{And \(J_{-1/2}(x)=\sqrt{2/\pi x}\,\cos x\)---the
spherical Bessel functions in disguise.}
\pair{special-functions/bessel/4}{Numerics to any precision.}
\pair{special-functions/bessel/5}{The Frobenius series about the origin, generated on
demand.}
\pair{special-functions/bessel/6}{The recurrence shows itself in the derivative:
\(J_0'=-J_1\).}
\pair{special-functions/bessel/7}{The second kind \(Y_0\), singular at the
origin\ldots}
\pair{special-functions/bessel/8}{\ldots the modified \(I_0\)\ldots}
\pair{special-functions/bessel/9}{\ldots and \(K_0\), the exponentially decaying
modified solution.}
\end{codepairs}

The half-integer reductions (pairs~2--3) are a real theorem, not a lookup: they are
why a quantum-mechanics problem in a spherical well produces \(\sin\) and \(\cos\)
over \(\sqrt r\).\calloutnote{Theory}{The Bessel functions \(J_\nu\) and \(Y_\nu\)
are the two independent solutions of Bessel's equation
\(x^2y''+xy'+(x^2-\nu^2)y=0\). When \(\nu\) is a half-integer the equation degenerates
just enough for the solution to close in elementary terms: \(J_{1/2}\) and
\(J_{-1/2}\) are \(\sin\) and \(\cos\) scaled by \(\sqrt{2/\pi x}\), and the
three-term recurrence \(J_{\nu-1}(x)+J_{\nu+1}(x)=\tfrac{2\nu}{x}J_\nu(x)\) generates
every other half-integer order from them. Mathilda applies exactly this reduction
(\texttt{src/special\_functions/bessel.c}).}

\usagebox{BesselJ}

The Airy functions are the simplest functions that oscillate on one side of the
origin and decay on the other, the model for a wave near a turning
point.\index{Airy functions} \B{AiryAi} and \B{AiryBi} solve \(y''=xy\), and their
values at the origin tie the whole family back to the gamma function.

\begin{codepairs}
\pair{special-functions/airy/1}{\(\operatorname{Ai}(0)\) closes through the gamma
function: \(1/(3^{2/3}\Gamma(2/3))\).}
\pair{special-functions/airy/2}{\(\operatorname{Bi}(0)\) shares the same
\(\Gamma(2/3)\), scaled differently.}
\pair{special-functions/airy/3}{The derivative at the origin brings in
\(\Gamma(1/3)\) instead.}
\pair{special-functions/airy/4}{The Airy value numerically, to thirty digits\ldots}
\pair{special-functions/airy/5}{\ldots and at \(x=1\), already deep in the decaying
region.}
\pair{special-functions/airy/6}{\B{AiryAi} differentiates to its own derivative head,
\B{AiryAiPrime}, closing the family under \B{D}.}
\end{codepairs}

That \(\operatorname{Ai}(0)\) and \(\operatorname{Ai}'(0)\) involve \(\Gamma(2/3)\)
and \(\Gamma(1/3)\) respectively (pairs~1--3) is the thread tying this subsection
back to the first: the special functions are a single web, and the gamma function
sits near its center.\calloutnote{Theory}{The Airy equation \(y''=xy\) has no
singular points at finite \(x\), so both solutions are entire. The initial values
\(\operatorname{Ai}(0)=3^{-2/3}/\Gamma(2/3)\) and
\(\operatorname{Ai}'(0)=-3^{-1/3}/\Gamma(1/3)\) come from the Mellin-transform
representation of the solution, and every Taylor coefficient follows from the
equation by \(y^{(n+2)}=x\,y^{(n)}+n\,y^{(n-1)}\). The kernels are in
\texttt{src/special\_functions/airyai.c} and \texttt{airybi.c}.}

\subsection{Legendre functions of the first and second kind}
\label{ssec:sf-legendre}

Legendre's equation \((1-x^2)y''-2xy'+n(n+1)y=0\) is the angular part of Laplace's
equation on the sphere, and it too has two independent
solutions.\index{Legendre functions} \B{LegendreP} is the first: at a non-negative
integer degree it is the Legendre \emph{polynomial}, the workhorse of Gaussian
quadrature and multipole expansions. \B{LegendreQ} is the second kind---the other
solution of the same equation---and it is new in this release.\index{Legendre
functions!second kind}

\begin{codepairs}
\pair{special-functions/legendre/1}{The Legendre polynomials \(P_0\) through \(P_3\),
generated on demand.}
\pair{special-functions/legendre/2}{\(P_4\), with its characteristic rational
coefficients.}
\pair{special-functions/legendre/3}{Orthogonality on \([-1,1]\): the squared norm is
\(\int_{-1}^1 P_n^2 = 2/(2n{+}1)\), here \(2/5\).}
\pair{special-functions/legendre/4}{And distinct degrees are orthogonal---the
integral of \(P_2P_3\) is exactly zero.}
\pair{special-functions/legendre/5}{The second-kind solution \(Q_0\) is the
logarithm \(\tfrac12\log\frac{1+x}{1-x}=\operatorname{arctanh}x\).}
\pair{special-functions/legendre/6}{\(Q_1\) carries that same logarithm, times
\(x\), minus one.}
\pair{special-functions/legendre/7}{\(Q_2\) pairs the logarithm with the polynomial
\(P_2\)---the general pattern.}
\pair{special-functions/legendre/8}{And \(Q\) evaluates numerically like any other
function.}
\end{codepairs}

Pairs~3--4 are the property that makes the Legendre polynomials indispensable, and
the system computes them exactly rather than numerically.\calloutnote{Theory}{The
polynomials \(P_n\) are orthogonal on \([-1,1]\):
\(\int_{-1}^1 P_m(x)P_n(x)\,dx = \frac{2}{2n+1}\delta_{mn}\). This single relation is
what lets any function on the interval be expanded in a Legendre series, and it is the
foundation of Gauss--Legendre quadrature. The second-kind solution \(Q_n\) satisfies
the same differential equation but is singular at \(x=\pm1\)---the endpoints where the
logarithm in pairs~5--7 blows up---so a physical solution regular on the whole sphere
keeps only the \(P_n\) part. \B{LegendreQ} follows the pattern
\(Q_n(x)=P_n(x)\,Q_0(x)-W_{n-1}(x)\) with \(Q_0=\operatorname{arctanh}x\); the
implementation is in \texttt{src/special\_functions/legendre.c}.} The visible
structure of \(Q_2\) in pair~7---a Legendre polynomial multiplying
\(\operatorname{arctanh}x\), minus a lower-degree polynomial---is exactly that
recurrence made concrete.

\subsection{The hypergeometric family: one function to rule them all}
\label{ssec:sf-hypergeometric}

Almost every function in this section is a special case of one
master series.\index{hypergeometric function} The generalized hypergeometric
function
\[
  {}_pF_q\!\left(\begin{matrix}a_1,\dots,a_p\\ b_1,\dots,b_q\end{matrix};z\right)
  = \sum_{k=0}^\infty
    \frac{(a_1)_k\cdots(a_p)_k}{(b_1)_k\cdots(b_q)_k}\,\frac{z^k}{k!}
\]
is a power series whose term ratio is a fixed rational function of the summation
index---and that one structural constraint is general enough to contain the
exponential, the logarithm, the trigonometric functions, the Bessel functions, the
error function, and the Legendre functions as special cases. Mathilda exposes it as
\B{HypergeometricPFQ}, with the classical low-order cases \B{Hypergeometric2F1}
(Gauss), \B{Hypergeometric1F1} (Kummer's confluent), and \B{Hypergeometric0F1}
(the confluent limit) as named heads. The system's job is to \emph{recognize} when a
hypergeometric collapses to something elementary, and it does.

\begin{codepairs}
\pair{special-functions/hypergeometric/1}{Gauss's \({}_2F_1(1,1;2;x)\) is a
logarithm in disguise: \(-\log(1-x)/x\).}
\pair{special-functions/hypergeometric/2}{A negative-integer parameter terminates the
series into a polynomial---here \((1-x)^3\).}
\pair{special-functions/hypergeometric/3}{Kummer's \({}_1F_1(1;2;x)\) is
\((e^x-1)/x\)\ldots}
\pair{special-functions/hypergeometric/4}{\ldots and when its two parameters agree it
is the exponential itself.}
\pair{special-functions/hypergeometric/5}{The confluent \({}_0F_1(;1/2;x)\) is a
hyperbolic cosine of a square root.}
\pair{special-functions/hypergeometric/6}{The generalized \({}_pF_q\) carries any
number of parameters, and reduces by the same rules---here matching pair~1.}
\pair{special-functions/hypergeometric/7}{Where no elementary form exists, it is a
number: this one is \(2\log 2\).}
\end{codepairs}

The unification is not just aesthetic; it is how a CAS avoids re-implementing the
same reductions a dozen times.\calloutnote{Theory}{A function is
\emph{hypergeometric} when the ratio of consecutive series coefficients,
\(c_{k+1}/c_k\), is a rational function of \(k\)---and astonishingly many functions
are. The exponential is \({}_0F_0\); \(\log(1-x)=-x\,{}_2F_1(1,1;2;x)\) (pair~1);
\(J_\nu(x)\) is a \({}_0F_1\) up to a power factor (compare pair~5's \(\cosh\), which
is \(J_{-1/2}\) territory); \(\operatorname{erf}\), the Legendre and Gegenbauer
polynomials, and the Airy functions are all \({}_pF_q\) with particular parameters.
The two levers that produce a closed form are visible above: a \emph{negative
integer} upper parameter truncates the infinite series to a polynomial (pair~2), and
special parameter relations trigger classical summation theorems (pairs~1, 3--5).
Mathilda's reductions live in \texttt{src/special\_functions/hypergeopfq.c}.}

Three functions extend the reach of the elementary world one step further and round
out the family. The polylogarithm \B{PolyLog} generalizes the logarithm to a whole
ladder of orders;\index{polylogarithm} the Lerch transcendent \B{LerchPhi} contains
the polylogarithm, the Hurwitz zeta, and much else as special
cases;\index{Lerch transcendent} and the Lambert \(W\) function \B{ProductLog}
inverts \(w\,e^w=z\), the transcendental equation that appears whenever an unknown
sits both inside and outside an exponential.\index{Lambert W function}

\begin{codepairs}
\pair{special-functions/polylog/1}{Order one is the logarithm:
\(\operatorname{Li}_1(x)=-\log(1-x)\).}
\pair{special-functions/polylog/2}{The dilogarithm at \(1\) is \(\zeta(2)=\pi^2/6\).}
\pair{special-functions/polylog/3}{At \(-1\) it is \(-\pi^2/12\).}
\pair{special-functions/polylog/4}{And the celebrated closed form
\(\operatorname{Li}_2(1/2)=\pi^2/12-\tfrac12\log^2 2\).}
\pair{special-functions/polylog/5}{At argument \(1\) the order-\(s\) polylog is
\(\zeta(s)\), so \(\operatorname{Li}_3(1)=\zeta(3)\).}
\pair{special-functions/polylog/6}{The Lerch transcendent reduces to the same
logarithm for these arguments---the master function containing all the rest.}
\pair{special-functions/polylog/7}{The Lambert \(W\): \(W(e)=1\), since
\(1\cdot e^1=e\).}
\pair{special-functions/polylog/8}{And \(W(-1/e)=-1\), the branch point.}
\pair{special-functions/polylog/9}{Its value at \(1\) is the omega constant,
\(0.5671\ldots\).}
\pair{special-functions/polylog/10}{Its derivative is written implicitly in terms of
itself---the signature of a function defined by inversion.}
\end{codepairs}

The dilogarithm identities (pairs~2--4) are among the most beautiful closed forms in
analysis, and that Mathilda knows \(\operatorname{Li}_2(1/2)=\pi^2/12-\tfrac12\log^2
2\) exactly rather than as a decimal is the whole point of a symbolic
system.\calloutnote{Under the hood}{\B{PolyLog}
(\texttt{src/special\_functions/polylog.c}), \B{LerchPhi} (\texttt{lerchphi.c}), and
\B{ProductLog} (\texttt{productlog.c}) each carry their own special-value tables and
their own numeric kernels. The Lerch transcendent
\(\Phi(z,s,a)=\sum_{k\ge0}z^k/(k+a)^s\) is the genuine master function of the three:
\(\operatorname{Li}_s(z)=z\,\Phi(z,s,1)\) and \(\zeta(s,a)=\Phi(1,s,a)\), so pair~6's
reduction and the Hurwitz zeta of \S\ref{ssec:sf-zeta} are two faces of one series.
\B{ProductLog}'s implicit derivative (pair~10) comes from differentiating
\(w\,e^w=z\) directly, since no explicit formula for \(W\) exists.}

\subsection*{Where this connects}

The special functions are the point where the three layers of Mathilda meet on a
single object. Their \emph{exact} values are computed by the same rational and
symbolic arithmetic as \S\ref{sec:arithmetic} and \S\ref{sec:algebra}; their
\emph{identities}---the derivatives that close the gamma, Bessel, Airy, and Legendre
families under \B{D}, the series expansions, the antiderivatives that \B{Integrate}
returns by name---are what let them take part in the calculus of
\S\ref{sec:calculus} rather than merely being evaluated at its end; and their
\emph{numerics}, backed by MPFR on the real line and by FLINT's rigorous Arb ball
arithmetic in the complex plane, are what the numerical methods of
\S\ref{sec:numcalc} call when an integral or a differential equation resolves into
one of these functions. The orthogonality of the Legendre polynomials underwrites
Gaussian quadrature; the gamma and zeta functions thread through analytic number
theory; and the hypergeometric family is the organizing idea that keeps the whole
collection from being a bag of unrelated formulas. For the exhaustive options and
corner cases of any function named here---the branch-cut conventions, the incomplete
and generalized forms, the full list of recognized reductions---its reference page
is one \texttt{?Name} away at the prompt.

\section{Statistics}
\label{sec:statistics}

Every section so far has computed with objects given exactly: an integer, a
polynomial, a function with a known derivative. Statistics begins at the other
end, with \emph{data}---a list of measured numbers, each carrying error, none of
them the answer on its own---and asks what can be said about the whole. The
questions are the classical ones. Where is the data centered? How far does it
spread? What is the shape of its tail? Do two quantities rise and fall together?
And how does a sample, once in hand, speak for the population it was drawn
from?\index{statistics!descriptive}

Mathilda answers these with the same principle that governs the rest of the
system: \emph{stay exact while the data is exact}. Hand \B{Variance} a list of
integers and it returns a rational number, not a rounded decimal; hand it a list
of measured reals and it returns a machine number, because that is all the data
supports. The estimators are the textbook ones---the sample variance carries the
unbiased \(n-1\) in its denominator, kurtosis is Pearson's (not the excess
form)---and every one threads columnwise over a matrix, so a table of several
variables is summarized in a single call. Underneath, when the data is a packed
array of machine reals, these reductions run on the buffer at the speed of C, and
each lowers inside \B{Compile}; the mathematics you type at the prompt is the
mathematics that runs in a tight numerical loop.

We take the subject in the order a data analyst would. First \emph{location}---the
mean and the median, and why the choice between them matters (\B{Mean},
\B{Median}, \B{Commonest}); then \emph{spread} (\B{Variance},
\B{StandardDeviation}, \B{RootMeanSquare}); the \emph{five-number summary} that
sketches a distribution with five order statistics (\B{Min}, \B{Quartiles},
\B{Max}); the \emph{shape} of the distribution through its moments (\B{Moment},
\B{CentralMoment}, \B{Skewness}, \B{Kurtosis}); the relationship between two
variables (\B{Covariance}, \B{Correlation}); the smoothing of a time series
(\B{MovingAverage}, \B{MovingMedian}, \B{ExponentialMovingAverage}); the tallying
of a sample into an empirical distribution (\B{Tally}, \B{Counts}); and we close
by crossing from description to \emph{probability}, drawing samples from a
distribution and measuring them back (\B{RandomVariate}, \B{PDF},
\B{LearnDistribution}).

\subsection{Location: the mean and the median}
\label{ssec:stat-location}

The oldest question in statistics is where a batch of numbers sits. The
\emph{mean} \(\bar x = \tfrac1n\sum x_i\) is the balance point; the \emph{median}
is the middle value once the data is sorted. On clean data they nearly agree, and
the choice between them looks like a matter of taste. It is not. Consider Simon
Newcomb's 1882 measurements of the time light takes to traverse a fixed
baseline\index{Newcomb, Simon}, recorded as deviations from \(24800\)
nanoseconds---and note the two readings, \(-44\) and \(-2\), that any modern eye
flags as blunders:

\begin{codepairs}
\pair{statistics/location/1}{Newcomb's first twenty timings, two of them wild.}
\pair{statistics/location/2}{\B{Mean} of integer data is exact: the rational
\(87/4\).}
\pair{statistics/location/3}{As a decimal, \(21.75\)---dragged down by the two
outliers.}
\pair{statistics/location/4}{\B{Median} sorts and takes the center: \(51/2 = 25.5\),
a full three units higher and untouched by the blunders.}
\pair{statistics/location/5}{\B{Commonest} reports the most frequent value---here a
tie between \(24\) and \(29\).}
\end{codepairs}

The gap between \(21.75\) and \(25.5\) is the whole lesson. The mean feels every
value in proportion to its size, so a single gross outlier moves it without limit;
the median feels only \emph{rank}, so up to half the data can be corrupted before
it budges. When a distribution is symmetric and clean the two coincide and the
mean is the more efficient estimate; when it is skewed or contaminated, the median
is the honest one.\index{median!robustness}

\calloutnote{Pitfall}{The mean is not robust. One fat-fingered data-entry error---a
misplaced decimal, a sign flip like Newcomb's \(-44\)---can shift a mean
arbitrarily far, while leaving the median essentially fixed. Report both when the
data might be dirty, and worry when they disagree.}

\subsection{Spread: variance, standard deviation, and RMS}
\label{ssec:stat-spread}

Location is half the story; the other half is how widely the data scatters about
it. The \emph{sample variance} averages the squared deviations from the mean, and
the \emph{standard deviation} is its square root, back in the original units.
Albert Michelson's 1879 speed-of-light experiment\index{Michelson, Albert} gives a
clean integer batch to work with---measurements in km/s above \(299000\):

\begin{codepairs}
\pair{statistics/spread/1}{Michelson's first twenty runs.}
\pair{statistics/spread/2}{Their mean, exact at \(909\) km/s above the baseline.}
\pair{statistics/spread/3}{\B{Variance} stays exact too: the rational
\(209180/19\).}
\pair{statistics/spread/4}{\B{StandardDeviation} is its exact square root, the surd
\(2\sqrt{52295/19}\).}
\pair{statistics/spread/5}{A decimal is one \B{N} away: about \(105\) km/s.}
\pair{statistics/spread/6}{\B{RootMeanSquare} is the quadratic mean
\(\sqrt{\tfrac1n\sum x_i^2}\)---magnitude about the origin, not about the mean.}
\end{codepairs}

Two details of the exact output repay attention. The variance is
\(209180/\mathbf{19}\), not over \(20\): the sample variance divides by \(n-1\),
not \(n\). And the standard deviation is returned as a genuine surd, carried
symbolically until you ask for digits---the same refusal to round prematurely that
runs through all of Mathilda.

\calloutnote{Theory}{Dividing by \(n-1\) rather than \(n\) is \emph{Bessel's
correction}\index{Bessel's correction}. The deviations are taken about the
\emph{sample} mean, which is itself pulled toward the data, so squared deviations
about it underestimate the true spread; the sample mean costs one degree of
freedom, and dividing by \(n-1\) restores an unbiased estimate of the population
variance.}

\usagebox{Variance}

\subsection{The five-number summary}
\label{ssec:stat-fivenum}

Before any model is fitted, a distribution can be sketched by five order
statistics: the minimum, the three quartiles, and the maximum. Together they
locate the center, bound the range, and measure the spread of the middle
half---the raw material of a box plot. Henry Cavendish's 1798 determination of the
mean density of the Earth\index{Cavendish, Henry} (relative to water), twenty-nine
measured values, is a natural specimen:

\begin{codepairs}
\pair{statistics/fivenum/1}{Cavendish's twenty-nine density measurements.}
\pair{statistics/fivenum/2}{\B{Min}: the lightest reading, \(4.88\).}
\pair{statistics/fivenum/3}{\B{Quartiles} returns \(\{Q_1, Q_2, Q_3\}\)---the
lower quartile, the median, and the upper quartile.}
\pair{statistics/fivenum/4}{\B{Max}: the heaviest, \(5.85\).}
\pair{statistics/fivenum/5}{The median again, \(5.46\), matching the middle
quartile.}
\pair{statistics/fivenum/6}{And the mean, \(5.448\)---almost equal to the median,
the signature of a roughly symmetric batch.}
\end{codepairs}

The interquartile range \(Q_3 - Q_1 = 5.6125 - 5.2975 = 0.315\) captures the
spread of the central half while ignoring the tails entirely, which is why it,
like the median, is a robust measure. That the mean and median nearly coincide
here---\(5.448\) against \(5.46\)---is the numerical fingerprint of symmetry;
where they part company, the data is skewed, and measuring that skew is the next
step.

\subsection{Shape: moments, skewness, and kurtosis}
\label{ssec:stat-shape}

Location and spread are the first two \emph{moments} of a distribution. The higher
moments describe its shape. The \(r\)-th raw moment is \(\tfrac1n\sum x_i^r\); the
\(r\)-th \emph{central} moment takes the deviations about the mean,
\(\tfrac1n\sum (x_i-\bar x)^r\). From these, two standardized quantities read off
the shape directly. Take a small, deliberately lopsided batch---the citation
counts of ten papers, nine ordinary and one runaway hit:

\begin{codepairs}
\pair{statistics/shape/1}{Ten citation counts, with one paper at \(20\).}
\pair{statistics/shape/2}{\B{Moment} gives the raw second moment
\(\tfrac1n\sum x_i^2 = 493/10\).}
\pair{statistics/shape/3}{\B{CentralMoment} of order \(2\)---the variance about the
mean, dividing by \(n\).}
\pair{statistics/shape/4}{The third central moment, \(43452/125\); its sign is the
sign of the skew.}
\end{codepairs}

\emph{Skewness} is the third central moment standardized by the cube of the
standard deviation; \emph{kurtosis} is the fourth, standardized by the square of
the variance. Both are pure numbers, free of units:

\begin{codepairs}
\pair{statistics/shape/5}{\B{Skewness}, exact---a positive surd, so the tail runs
to the right.}
\pair{statistics/shape/6}{About \(2.45\): strongly right-skewed, as the lone hit
demands.}
\pair{statistics/shape/7}{\B{Kurtosis}, the fourth standardized moment.}
\pair{statistics/shape/8}{About \(7.44\)---far above the \(3\) of a normal
distribution, the mark of a heavy tail.}
\end{codepairs}

A symmetric distribution has zero skewness; a positive value means a long right
tail, a negative value a long left one. Kurtosis measures how much of the variance
comes from rare, extreme deviations: Mathilda reports the \emph{Pearson} kurtosis,
for which the normal distribution sits at exactly \(3\), so the value \(7.44\) here
says the single outlier dominates the fourth moment.

\calloutnote{Theory}{Skewness and kurtosis are \emph{standardized} moments,
\(\mu_3/\sigma^3\) and \(\mu_4/\sigma^4\), built to be dimensionless and
invariant under shifting and rescaling the data. The ``excess kurtosis''
reported by some packages is simply Mathilda's kurtosis minus \(3\), centering
the normal distribution at zero.\index{kurtosis}\index{skewness}}

\subsection{Two variables: covariance and correlation}
\label{ssec:stat-correlation}

With two measurements on each subject, the question becomes whether they move
together. \B{Covariance} averages the product of paired deviations; \B{Correlation}
normalizes it by the two standard deviations to a pure number in \([-1, 1]\), so
that \(+1\) is perfect increasing linear agreement, \(-1\) perfect decreasing, and
\(0\) no linear relation at all.

No example makes the point---and its limits---better than \emph{Anscombe's
quartet}\index{Anscombe's quartet}, four data sets that Frank Anscombe constructed
in 1973 to be statistically indistinguishable by every summary yet obviously
different to the eye~\cite{anscombe1973}:

\begin{codepairs}
\pair{statistics/anscombe/1}{The common predictor \(x\) shared by the first three
sets.}
\pair{statistics/anscombe/2}{Set~I's response.}
\pair{statistics/anscombe/5}{Set~IV instead holds \(x\) nearly constant, with one
point far out.}
\pair{statistics/anscombe/6}{Set~IV's response.}
\end{codepairs}

Now compute the summaries of all four at once. They agree to the digits that
matter:

\begin{codepairs}
\pair{statistics/anscombe/7}{All four means of \(y\) sit at \(7.5\).}
\pair{statistics/anscombe/8}{All four variances of \(y\) sit at \(4.13\).}
\pair{statistics/anscombe/9}{And all four correlations equal \(0.816\).}
\pair{statistics/anscombe/10}{The covariance of set~I, \(5.501\)---unnormalized, in
the product of the units.}
\pair{statistics/anscombe/11}{Given the two columns as a matrix, \B{Correlation}
returns the correlation matrix: unit diagonal, \(0.816\) off it.}
\end{codepairs}

Yet the four data sets could hardly be more different: one is a clean line, one a
smooth curve, one a perfect line derailed by a single outlier, and one a vertical
stack whose slope is decided by one lone point. Identical means, identical
variances, identical correlations---and four unrelated stories. The moral, as old
as it is ignored, is that summary statistics never substitute for looking at the
data.

\calloutnote{Pitfall}{A correlation coefficient compresses a whole scatter of
points into one number, and much is lost in the compression. Anscombe's quartet is
the standing proof: equal correlations can hide a line, a curve, and an outlier.
Before trusting a coefficient, plot the points---\S\ref{sec:intro-stats} and the
graphics chapter show how.}

\usagebox{Correlation}

\subsection{Time series: moving statistics}
\label{ssec:stat-moving}

When the data is a sequence in time, its raw values jitter and the trend hides
underneath. A \emph{moving statistic} smooths the jitter by summarizing a sliding
window rather than the whole series at once. Take a short, noisy monthly series:

\begin{codepairs}
\pair{statistics/moving/1}{Twelve monthly readings, trending up through the noise.}
\pair{statistics/moving/2}{\B{MovingAverage} over a window of three---each output is
the mean of three neighbors, so the run is two shorter than the input.}
\pair{statistics/moving/3}{A \emph{weighted} moving average: weights \(\{1,2,1\}\)
lean on the center of each window.}
\pair{statistics/moving/4}{\B{MovingMedian} smooths the same way but resists a lone
spike, exactly as the median did in \S\ref{ssec:stat-location}.}
\pair{statistics/moving/5}{\B{ExponentialMovingAverage} weights every past value,
decaying geometrically; with \(\alpha=1/3\) the result stays exact.}
\pair{statistics/moving/6}{The same recurrence numerically, with \(\alpha=0.2\):
each output is \(0.2\) of the new value plus \(0.8\) of the running average.}
\end{codepairs}

The three smoothers trade off differently. The plain moving average is simplest
but lags the trend and reacts sharply as a window enters or leaves an outlier. The
moving median inherits the median's robustness, flattening isolated spikes
outright. The exponential moving average never forgets---every past point
contributes, weighted by \((1-\alpha)^k\)---so it responds smoothly and needs only
the previous output to advance, which is why it is the workhorse of streaming and
financial data.

\calloutnote{Performance}{These reductions are exactly the kind of numerical
loop Mathilda accelerates. On a packed array of machine reals the moving
statistics run on the raw buffer without boxing a single element, and the same
kernels lower inside \B{Compile}, so a smoothing pass over a million-point series
costs a single sweep of memory rather than a million interpreter steps.}

\subsection{Frequencies and the empirical distribution}
\label{ssec:stat-frequencies}

For discrete or categorical data the natural summary is not a mean but a
\emph{count}: how often each distinct value occurs. This tabulation is the
\emph{empirical distribution}, the sample's own stand-in for a probability
distribution. Fifteen rolls of a die will serve:

\begin{codepairs}
\pair{statistics/frequencies/1}{Fifteen die rolls.}
\pair{statistics/frequencies/2}{\B{Tally} pairs each distinct value with its count,
in order of first appearance.}
\pair{statistics/frequencies/3}{\B{Counts} returns the same tabulation as an
association, keyed by value---ready for \(O(1)\) lookup.}
\pair{statistics/frequencies/4}{\B{Commonest} extracts the mode, the most frequent
value.}
\pair{statistics/frequencies/5}{Asked for two, it returns the top two by
frequency.}
\end{codepairs}

Normalizing the counts by the sample size turns the tally into an empirical
probability mass function, and its running sum into an empirical cumulative
distribution---the discrete companions of the density and distribution functions
of the next subsection. For continuous data the same idea, applied to bins rather
than exact values, is a histogram, which the graphics chapter draws directly with
\B{Histogram}.

\subsection{From data to distributions}
\label{ssec:stat-distributions}

Descriptive statistics run in one direction, from a sample to its summaries.
Probability runs in the other, from a model to the samples it generates---and the
two meet in the estimators of this section, which are exactly the quantities that
converge to a distribution's parameters as the sample grows. Mathilda lets you
watch that convergence. Draw a thousand values from a normal population with mean
\(100\) and standard deviation \(15\)---the scale of an IQ test---and measure the
sample back:

\begin{codepairs}
\pair{statistics/distributions/1}{\B{SeedRandom} fixes the stream, so the draw is
reproducible from run to run.}
\pair{statistics/distributions/2}{A thousand draws from \B{NormalDistribution};
the semicolon holds back the long list.}
\pair{statistics/distributions/3}{The sample mean, \(100.20\)---within a fifth of a
unit of the true \(100\).}
\pair{statistics/distributions/4}{The sample standard deviation, \(15.03\), all but
recovering the true \(15\).}
\pair{statistics/distributions/5}{\B{PDF} evaluates the density; it is symmetric
about the mean, so \(85\) and \(115\) share a value and \(100\) peaks.}
\pair{statistics/distributions/6}{\B{LearnDistribution} runs the inference the
other way, fitting a distribution object back to the data.}
\end{codepairs}

That the sample mean and standard deviation land so near \(100\) and \(15\) is the
\emph{law of large numbers}\index{law of large numbers} in miniature: the sample
estimators converge to the population parameters as \(n\) grows, and a thousand
draws already pin them to a fraction of a unit. \B{RandomVariate} draws from a
distribution, \B{PDF} evaluates its density, and \B{LearnDistribution} inverts the
arrow---inferring a distribution from data---closing the loop between the
descriptive statistics of this section and the probability that underwrites them.

\calloutnote{Under the hood}{\B{RandomVariate} shares its stream with
\B{RandomReal}, so a single \B{SeedRandom} makes an entire session's simulation
reproducible---indispensable when a result must be regenerated exactly, as every
example in this book is. The normal draws come from a standard transform of that
uniform stream.}

\subsection*{Where this connects}

Statistics is where Mathilda's numerical substrate meets its data. Its estimators
lean directly on the reductions of the array engine---a \B{Mean} or \B{Variance}
of a packed column is a single buffer sweep, and each lowers inside \B{Compile} so
it runs at full speed inside a numerical loop (Chapter~\ref{ch:compilation}). The
correlation matrix of \S\ref{ssec:stat-correlation} is a Gram matrix, computed
through the same \B{Dot} and \B{Transpose} that drive the linear algebra of
\S\ref{sec:linalg}, and it is the first step of principal component analysis. The
sampling of \S\ref{ssec:stat-distributions} rests on the random-number generators,
and the fit performed by \B{LearnDistribution} opens into the machine-learning
family---classifiers, clustering, and density estimation---that treats a data set
not as something to summarize but as something to learn from. For the exhaustive
options of any function named here, its reference page is one \texttt{?Name} away
at the prompt.


\include{chapters/05-graphics}
\chapter{Data Structures}
\label{ch:datastructures}

The previous chapters computed with mathematics; this one steps back to ask what,
exactly, a value \emph{is} inside Mathilda. The answer is unusually simple, and it
is the foundation everything else rests on: there is only one kind of value, the
\emph{expression}, and every list, matrix, rule, pattern, and association you will
ever type is one. That uniformity is what lets a single handful of tools---\B{Part},
\B{Map}, \B{ReplaceAll}---work on everything without special cases. But simplicity at
the surface is not the same as naivety underneath. Behind the uniform tree, Mathilda
keeps two specialized structures that make real work fast: the \emph{association},
a hash-backed key--value store with amortized \(O(1)\) lookup, and the \emph{packed
array}, a dense machine-precision buffer that runs numeric code at the speed of C.
The theme of the chapter is precisely this tension---one uniform data model on top,
invisible acceleration beneath---and how Mathilda keeps the two from ever
contradicting each other.

We build up in four steps. First the expression tree itself: head and arguments,
atoms and compounds, and the tools that read a value's shape. Then lists---the
workhorse aggregate---their parts, spans, and the newer array operations. Then
associations, the key--value store. And finally the packed-array substrate, where
most of the chapter's ``under the hood'' weight lies, because that is where a
transparent list quietly becomes a C buffer and back again.

\section{Everything is an expression}
\label{sec:ds-expressions}

Type \(a + b\,c\) and it looks like an infix formula. It is not stored as one.
Every compound value in Mathilda is a \emph{function application}\index{expression!head
and arguments}---a head applied to a sequence of arguments---and \B{FullForm} prints
that underlying tree with the surface syntax stripped away.

\begin{codepairs}
\pair{06-data-structures/expressions/1}{The sum-of-products \(a+bc\) is
\texttt{Plus} applied to \(a\) and \texttt{Times[b, c]}: operators are just
shorthand for named heads.}
\pair{06-data-structures/expressions/2}{A list is no different---\(\{1,2,3\}\) is
\texttt{List[1, 2, 3]}, the head \B{List} applied to three arguments.}
\pair{06-data-structures/expressions/3}{Even a fraction is a compound: \(1/2\) is
\texttt{Rational[1, 2]}. What prints as a number is a two-argument tree.}
\end{codepairs}

Because there is only one representation, the head of a value tells you what kind of
thing it is, and \B{Head} reads it off for anything at all---an operator expression,
an atom, a list, or a call to an undefined function.

\begin{codepairs}
\pair{06-data-structures/expressions/4}{The head of \(a+bc\) is \B{Plus}\ldots}
\pair{06-data-structures/expressions/5}{\ldots and of the sub-product \(bc\) it is
\B{Times}.}
\pair{06-data-structures/expressions/6}{A machine integer has head \texttt{Integer}\ldots}
\pair{06-data-structures/expressions/7}{\ldots a machine real, head \texttt{Real}\ldots}
\pair{06-data-structures/expressions/8}{\ldots a fraction, head \B{Rational}\ldots}
\pair{06-data-structures/expressions/9}{\ldots and a complex number, head \B{Complex}.}
\pair{06-data-structures/expressions/10}{Strings carry head \texttt{String}.}
\pair{06-data-structures/expressions/11}{A list's head is \texttt{List}\ldots}
\pair{06-data-structures/expressions/12}{\ldots and an unknown \(f[x]\) has head
\(f\), the symbol you named. Nothing is privileged.}
\end{codepairs}

The values that have no internal parts to descend into---numbers, symbols,
strings---are the \emph{atoms}\index{atom}; everything else is a compound with a
head and arguments. \B{AtomQ} draws the line, and it is worth watching where it
falls, because ``looks compound'' and ``is compound'' are not the same.

\begin{codepairs}
\pair{06-data-structures/expressions/13}{An integer is an atom\ldots}
\pair{06-data-structures/expressions/14}{\ldots so is a symbol\ldots}
\pair{06-data-structures/expressions/15}{\ldots and so, despite printing as
\(1/2\), is a \B{Rational}: it is indivisible as a value even though its
\texttt{FullForm} shows two integers.}
\pair{06-data-structures/expressions/16}{A function call is not an atom---you can
take it apart.}
\pair{06-data-structures/expressions/17}{Neither is a list.}
\end{codepairs}

\calloutnote{Under the hood}{An \texttt{Expr} is a tagged union
(\texttt{src/expr.h}): the tag is one of \texttt{EXPR\_INTEGER} (a machine
\texttt{int64}), \texttt{EXPR\_REAL} (a \texttt{double}), \texttt{EXPR\_SYMBOL},
\texttt{EXPR\_STRING}, \texttt{EXPR\_BIGINT} (a GMP integer), or
\texttt{EXPR\_FUNCTION} (a head pointer plus an argument array). The first five are
the atoms; only \texttt{EXPR\_FUNCTION} has parts. Expressions are treated as
immutable during evaluation---a transformation returns a new tree rather than
editing one in place---and \texttt{expr\_copy} is a reference-count bump, not a deep
clone, so this costs far less than it sounds. A \texttt{Rational} and a
\texttt{Complex} are \texttt{EXPR\_FUNCTION} nodes at the level of storage, but
\B{AtomQ} reports them atomic because they behave as single numbers; that deliberate
override is the gap between the two answers above.}

The shape of a tree---how deeply it nests, and what lives at each level---is itself
queryable. \bidx{Depth}\mcode{Depth} gives one more than the longest path from the
root to a leaf, and \B{Level} collects the subexpressions at a chosen depth.

\begin{codepairs}
\pair{06-data-structures/expressions/18}{\mcode{Depth} of \(f[g[h[x]]]\) is
\(4\): three nested heads plus the atom \(x\) at the bottom.}
\pair{06-data-structures/expressions/19}{A \(2\times2\) matrix has depth \(3\)---the
outer list, the rows, and the numbers.}
\pair{06-data-structures/expressions/20}{\B{Level} with the specification
\(\{-1\}\) gathers the atoms (level \(-1\) counts up from the leaves): every
number and symbol at the bottom of \(a + f[x, y^n]\).}
\pair{06-data-structures/expressions/21}{Level \(2\) instead gathers everything
within two steps of the root, so \(y^n\) appears whole rather than split into
\(y\) and \(n\).}
\end{codepairs}

A negative level counts distance from the leaves and a positive one distance from
the root, so \(\{-1\}\) is ``the atoms'' and \(2\) is ``down to two levels
deep.''\index{levelspec} These level specifications are shared by \B{Map},
\B{Cases}, \B{Position}, and the rest of the structural operators---learn them once
and they apply everywhere, which is the recurring dividend of the uniform model.

\section{Lists}
\label{sec:ds-lists}

The list is the aggregate you will reach for most, and because it is just an
\texttt{EXPR\_FUNCTION} with head \texttt{List}, every generic tool applies to it
unchanged. What lists add is a rich vocabulary for addressing their elements.

\subsection{Parts and spans}
\label{ssec:ds-parts}

\B{Part}---written \(\mathit{expr}\)\texttt{[[\ldots]]}---extracts an element by
position, counting from \(1\), or from the end with negative indices. A
\emph{span}\index{list!span}---the \B{Span} operator, written \texttt{i ;; j}---extracts
a contiguous run, optionally with a step; and for a nested list a comma-separated
index descends level by level.

\begin{codepairs}
\pair{06-data-structures/lists/2}{The third element of the vector of primes.}
\pair{06-data-structures/lists/3}{A negative index counts from the end: \(-1\) is
the last.}
\pair{06-data-structures/lists/4}{A span \texttt{2 ;; 4} takes elements two through
four inclusive.}
\pair{06-data-structures/lists/5}{With a step, \texttt{1 ;; -1 ;; 2} takes every
second element from first to last.}
\pair{06-data-structures/lists/7}{On a matrix, \texttt{[[2, 3]]} descends: row
two, then column three.}
\pair{06-data-structures/lists/8}{\texttt{All} in a slot keeps that whole axis, so
\texttt{[[All, 2]]} is the second column.}
\end{codepairs}

The same \texttt{[[\ldots]]} notation reads a single element, a slice, or a whole
axis depending on what you put in the slot, and it composes down through any depth
of nesting---a direct consequence of a matrix being nothing more special than a list
of lists.

\subsection{Shape and nesting}
\label{ssec:ds-shape}

\B{Length} counts the top-level arguments of any expression, and \B{Dimensions}
reports the rectangular shape of a nested list---the lengths at each level, read
outward in.

\begin{codepairs}
\pair{06-data-structures/lists/9}{Six primes, length six.}
\pair{06-data-structures/lists/10}{\B{Length} of the matrix is \(2\)---it has two
rows; length is always the count of immediate arguments.}
\pair{06-data-structures/lists/11}{\B{Dimensions} gives the full shape:
\(2\times3\).}
\pair{06-data-structures/lists/12}{And for a rank-\(3\) array, all three axis
lengths.}
\pair{06-data-structures/lists/13}{Because \B{Length} is generic, it counts the
arguments of any head---here the four arguments of \(f\).}
\end{codepairs}

The contrast between the last two pairs is the point: \B{Dimensions} presumes the
rectangular, tensor-like structure a numeric array has, while \B{Length} makes no
such assumption and works on \(f[a,b,c,d]\) exactly as on a list. One is for data;
both are for expressions.

\subsection{Reshaping, padding, and gradients}
\label{ssec:ds-arrayops}

Four operations turn a flat list into a shaped one, resize it, differentiate it, and
locate a value inside it. \B{ArrayReshape} pours the elements of a list, in
row-major order, into an array of whatever shape you request; \B{ArrayPad} grows (or,
with negative amounts, trims) an array at its edges under a choice of fill scheme;
\B{ListGradient} takes the numerical derivative of sampled data by finite
differences; and \B{FirstPosition} finds where a pattern first matches.

\begin{codepairs}
\pair{06-data-structures/arrayops/1}{\B{ArrayReshape} lays \(1\ldots12\) into a
\(3\times4\) grid, filling row by row.}
\pair{06-data-structures/arrayops/2}{When the data runs short, the third argument
supplies the fill: seven elements into a \(3\times3\) grid, padded with \(x\).}
\pair{06-data-structures/arrayops/3}{\B{ArrayPad} adds two elements of the default
fill (\(0\)) at each end.}
\pair{06-data-structures/arrayops/4}{An asymmetric \(\{1,2\}\) pad with the
\texttt{"Fixed"} scheme repeats the edge values rather than inserting zeros.}
\pair{06-data-structures/arrayops/5}{A negative amount trims: \(-2\) drops two from
each end.}
\end{codepairs}

\B{ListGradient} is the most mathematically interesting of the four---a
from-first-principles port of \texttt{numpy.gradient}. On the interior of the array
it takes a second-order central difference; at each endpoint, where a centered
stencil would run off the edge, it falls back to a one-sided difference.

\begin{codepairs}
\pair{06-data-structures/arrayops/6}{Sample \(x^2\) at \(x=1,\dots,5\)
(\(\{1,4,9,16,25\}\)); the interior values \(4,6,8\) are exactly the derivative
\(2x\) at \(x=2,3,4\), and the ends \(3,9\) are the one-sided approximations.}
\pair{06-data-structures/arrayops/7}{With symbolic input the finite-difference
formula itself is returned---halved central differences inside, full one-sided
differences at the ends.}
\pair{06-data-structures/arrayops/8}{Sampling \(x^3\) and asking for a fourth-order
stencil recovers \(3x^2\) exactly at every point, endpoints included.}
\end{codepairs}

\calloutnote{Theory}{The single kernel behind \B{ListGradient} is Fornberg's
finite-difference-weight recurrence\index{finite differences!Fornberg weights},
which produces the exact weights for a derivative of any order on any (even
non-uniform) grid. A second-order interior stencil is
\((f_{i+1}-f_{i-1})/2h\)---which is why sampling a quadratic gives the derivative
back with no error, since the central difference of \(x^2\) is exact---and the
endpoints drop to first order (numpy's \texttt{edge\_order = 1}). Because the weights
are exact rationals, integer input yields exact \B{Rational} output and symbolic
input a symbolic formula, as the middle pair shows; only a machine-real array takes
the packed floating-point fast path of \S\ref{sec:ds-packed}.}

\B{FirstPosition} reports the position of the first match as a list of indices---the
very indices \B{Part} would use to extract it---searching depth-first through the
whole expression.

\begin{codepairs}
\pair{06-data-structures/arrayops/9}{The value \(7\) sits at position \(\{3\}\).}
\pair{06-data-structures/arrayops/10}{In a nested list the position is a full index
path: \(4\) is at row two, column two.}
\pair{06-data-structures/arrayops/11}{No match returns \(\mcode{Missing["NotFound"]}\)\ldots}
\pair{06-data-structures/arrayops/12}{\ldots unless you supply a default, which is
returned in its place.}
\end{codepairs}

\calloutnote{Pitfall}{\B{FirstPosition} takes its arguments with the attribute
\texttt{HoldRest}, so the default is evaluated only if it is actually
returned---a deliberate design that lets you write an expensive fallback without
paying for it on a hit. The search follows \B{Position}'s conventions exactly
(levels \(\{0,\infty\}\), \texttt{Heads -> True}), which is why a position can be
the empty list \(\{\}\), meaning the whole expression matched.}

\section{Associations}
\label{sec:ds-associations}

A list indexes by position. Often what you have instead is a set of named
fields---a record, a lookup table, a histogram---and forcing that onto positional
storage means remembering that ``price is element three.'' The \B{Association}\index{association!keyed
store}, written \texttt{<|key -> value, \ldots|>}, indexes by \emph{key}: any expression can
be a key, keys are unique, and they keep their insertion order.

\begin{codepairs}
\pair{06-data-structures/associations/1}{An inventory keyed by fruit name. It prints
back in the order written.}
\pair{06-data-structures/associations/2}{Applied like a function,
\(\mathit{assoc}[\mathit{key}]\) looks the key up---the idiomatic accessor.}
\pair{06-data-structures/associations/3}{The \texttt{[[key]]} form does the same
through \B{Part}.}
\pair{06-data-structures/associations/4}{\B{Keys} lists the keys, in order\ldots}
\pair{06-data-structures/associations/5}{\ldots and \B{Values} the corresponding
values.}
\pair{06-data-structures/associations/6}{\B{Lookup} is the explicit accessor.}
\pair{06-data-structures/associations/7}{Its three-argument form supplies a default
for an absent key---the safe way to read data that may be incomplete.}
\pair{06-data-structures/associations/8}{\B{KeyExistsQ} tests for a key's
presence\ldots}
\pair{06-data-structures/associations/9}{\ldots reporting \texttt{False} when it is
missing.}
\pair{06-data-structures/associations/10}{A bare lookup of an absent key returns a
\mcode{Missing} marker naming what was not found, rather than an error.}
\end{codepairs}

Two behaviors set associations apart from a bag of rules. Keys are unique with
last-value-wins, and functions that consume a collection act on the
\emph{values}\index{association!value threading} while keeping keys aligned.

\begin{codepairs}
\pair{06-data-structures/associations/11}{A repeated key collapses: the later value
wins, but the key keeps its first position.}
\pair{06-data-structures/associations/12}{\B{Map} threads over values and returns an
association---keys ride along untouched.}
\pair{06-data-structures/associations/13}{\B{Total} reduces over the values, exactly
as it would over \texttt{Values[inv]}.}
\pair{06-data-structures/associations/14}{\B{AssociationThread} zips a list of keys
against a list of values into an association---the constructor to reach for when the
two arrive separately.}
\end{codepairs}

\calloutnote{Under the hood}{Single-key access is amortized \(O(1)\). A persistent
key\(\rightarrow\)position hash index (\texttt{src/assoc\_index.h}) is cached on the
association and built on the first single-key read, so \B{Lookup}, \B{KeyExistsQ},
\(\mathit{a}[\mathit{key}]\), and \(\mathit{a}[[\mathit{key}]]\) never scan the
entries. The index reuses spare space already on the expression node, so it costs no
extra memory, and it is transparent to \B{SameQ}, \B{Sort}, and \B{Union}, which see
the association exactly as before. The value-threading rule---that \B{Map}, \B{Total},
\B{Select}, \B{Cases}, and their kin operate on values and keep keys aligned---runs
through one shared helper, which is why these operations compose predictably into
data pipelines.}

\mcode{Part} addresses an association by key exactly as it addresses a list by position,
and a part that selects \emph{several} entries returns a \emph{sub-association}---the
chosen entries, keys and all, in their original order.

\begin{codepairs}
\pair{06-data-structures/assoc-part/1}{A small keyed record to slice.}
\pair{06-data-structures/assoc-part/2}{A span selects a contiguous run of entries and
keeps them keyed.}
\pair{06-data-structures/assoc-part/3}{A list of keys picks exactly those, in the order
requested.}
\pair{06-data-structures/assoc-part/4}{Integer positions return the same
sub-association---position and key are two names for one slot.}
\end{codepairs}

\subsection{Associations are atoms}
\label{ssec:ds-assoc-atoms}

An association looks like a container, but to the structural machinery it is a single,
indivisible object---an \emph{atom}. That is the rule that lets one travel safely through
code that knows nothing about associations. Traversal and testing see the \emph{values};
the keys are addressing metadata, reached only through an explicit \B{Key}.

\begin{codepairs}
\pair{06-data-structures/assoc-atoms/1}{A three-entry association.}
\pair{06-data-structures/assoc-atoms/2}{\B{AtomQ} reports \texttt{True}: structurally it
is one node, not a compound of parts to be taken apart.}
\pair{06-data-structures/assoc-atoms/3}{\mcode{Cases} scans the \emph{values}, so a
pattern filter returns matching values, not entries.}
\pair{06-data-structures/assoc-atoms/4}{\mcode{Position} likewise finds a value, and
names where it lives by its key, wrapped in \B{Key}.}
\pair{06-data-structures/assoc-atoms/5}{And \B{ReplaceAll} cannot touch a key: a rule
that would rewrite \texttt{"y"} leaves the association unchanged, because the keys are not
part of the expression it walks.}
\end{codepairs}

\calloutnote{Under the hood}{Atomicity is a deliberate---and recent---design choice.
\mcode{AtomQ}, \mcode{Depth}, \mcode{LeafCount}, \mcode{Level}, \mcode{Cases},
\mcode{Count}, \mcode{MemberQ}, \mcode{FreeQ}, \mcode{Position}, and the level
specifications of \mcode{Map} and \mcode{Apply} all agree to see only an association's
values. The pay-off is safety: because \texttt{/.} can never merge or rewrite keys, and a
pattern scan can never mistake a key for a value, an association carried through a generic
pipeline---as a rule set, a \mcode{Cases} filter, or a user-defined transformation---behaves
predictably and cannot be silently corrupted. When you \emph{do} need to name a key
structurally, \mcode{Position} hands back a \B{Key}$[\,k\,]$ wrapper that \mcode{Part}
accepts directly.}

\subsection{Operator forms, slots, and pipelines}
\label{ssec:ds-assoc-pipeline}

Data work is a pipeline: filter the rows, sort them, pull a field, summarise. Associations
are built for it, through two conveniences. A pure function can read a field by \emph{name}
with a named slot, \texttt{\#field}, and most operators carry a \emph{curried} form that
takes their data last, so stages compose left to right.

\begin{codepairs}
\pair{06-data-structures/assoc-pipeline/2}{\B{Select} keeps the rows whose \texttt{qty}
field---read by the named slot \texttt{\#qty}---exceeds two.}
\pair{06-data-structures/assoc-pipeline/3}{\B{SortBy} orders the same rows by that field.}
\pair{06-data-structures/assoc-pipeline/5}{\mcode{Lookup["pears"]} is an operator form:
it waits for an association, then looks the key up.}
\pair{06-data-structures/assoc-pipeline/6}{Listable heads thread over the values, so
adding a scalar adds it to each value and keeps the keys aligned.}
\pair{06-data-structures/assoc-pipeline/7}{And \mcode{Total} of a \emph{list} of
associations sums them key by key---the reduction a table of records calls for.}
\end{codepairs}

\calloutnote{Under the hood}{A named slot \texttt{\#field} parses as
\B{Slot}\texttt{["field"]}, which on an association reads that key---so
\texttt{\#qty\,>\,2\,\&} is an ordinary pure function that happens to index by name. The
curried forms come from one table-driven mechanism (\texttt{src/opform.c}) that gives some
forty-three heads---\mcode{Lookup}, \mcode{Select}, \mcode{Map}, \mcode{GroupBy},
\mcode{SortBy}, and their kin---an operator form \texttt{f[args][data]} equivalent to
\texttt{f[data, args]}. Together they let a data pipeline be written point-free, each stage
a self-contained operator applied in turn.}

\subsection{Aggregating: Counts, GroupBy, and Merge}
\label{ssec:ds-aggregate}

The reason associations matter for data work is that they are what the aggregation
operators \emph{produce}. \B{Counts} tallies a list into an association from value to
frequency; \B{GroupBy} partitions a list by a key function; and \B{Merge} combines
several associations, reconciling clashing keys with a function of your choosing.

\begin{codepairs}
\pair{06-data-structures/aggregate/1}{\B{Counts} turns a list of observations into
an empirical frequency table---the discrete distribution of the sample.}
\pair{06-data-structures/aggregate/2}{\B{GroupBy} partitions \(1\ldots10\) by
\texttt{EvenQ}: two groups, keyed \texttt{False} and \texttt{True}.}
\pair{06-data-structures/aggregate/3}{Its three-argument form applies a reducer to
each group---here \B{Total}, collapsing the split-apply-combine in one call.}
\pair{06-data-structures/aggregate/5}{With a \(\mathit{key}\to\mathit{value}\) rule,
\B{GroupBy} groups the rows by their \B{First} field but collects the \B{Last}: sales
gathered per category.}
\pair{06-data-structures/aggregate/6}{Add a reducer and the same call groups,
extracts, and totals---a one-line group-and-sum.}
\pair{06-data-structures/aggregate/7}{\B{Merge} unions the keys of two associations
and applies \B{Total} to the values collected under each: the shared key \(y\) sums
its two contributions, the unshared keys pass through.}
\end{codepairs}

The middle pairs are the split-apply-combine idiom\index{split-apply-combine} that
data analysis lives on: \emph{split} the data into groups, \emph{apply} a summary to
each, \emph{combine} the summaries into one table. \B{GroupBy} does all three, and
because the result is an association it feeds straight into another aggregation---the
frequency table from \B{Counts} is exactly the empirical distribution the statistics
chapter (\S\ref{sec:statistics}) builds its estimators on.

\B{Merge} is one of a family that treats an association as a keyed \emph{set}. The keys
can be intersected, subtracted, and unioned; entries can be dropped by a predicate on
their values; and two tables can be joined on a shared field the way a database would.

\begin{codepairs}
\pair{06-data-structures/assoc-keyalgebra/1}{\B{KeyIntersection} restricts each
association to the keys they all share---here only \texttt{"b"}, each keeping its own
value.}
\pair{06-data-structures/assoc-keyalgebra/2}{\B{KeyComplement} keeps the first
association's keys that appear in no other.}
\pair{06-data-structures/assoc-keyalgebra/3}{\B{KeyUnion} aligns several associations to
their combined key set, filling a gap with a \mcode{Missing} marker.}
\pair{06-data-structures/assoc-keyalgebra/4}{\B{Discard} is \mcode{Select}'s complement:
it drops the entries whose value satisfies the test.}
\pair{06-data-structures/assoc-keyalgebra/5}{\B{JoinAcross} joins two lists of records on
a shared key, merging each matched pair into one row---a relational join.}
\pair{06-data-structures/assoc-keyalgebra/6}{\B{PositionIndex} inverts a list into an
association from each value to the positions it occupies.}
\pair{06-data-structures/assoc-keyalgebra/7}{And \mcode{Merge} accepts a plain list of
rules, collecting clashing keys before it reduces---so it doubles as a grouping tally.}
\end{codepairs}

\usagebox{Association}

\section{NDArrays and packed arrays}
\label{sec:ds-packed}

Everything so far treated a list as a tree of independent expression nodes. For
symbolic work that is exactly right. For numeric work it is ruinous: a list of a
million reals is a million separate nodes, and the evaluator sweeps every argument of
every call on every pass, so merely \emph{holding} a large list costs time on each
operation, whether or not anything is computed. Mathilda's answer is the
\emph{packed array}\index{packed array}---the same list, stored as one dense
machine-precision buffer instead of a million nodes---and the discipline that keeps
that storage change completely invisible.

\subsection{The invisible buffer}
\label{ssec:ds-invisible}

\B{ToNDArray} packs a list into a buffer and hands it back looking exactly like the
list it was. The only tool that can tell the difference is \B{NDArrayQ}, which asks
about storage rather than about value.

\begin{codepairs}
\pair{06-data-structures/packing/1}{\B{ToNDArray} returns what looks like an ordinary
list of reals---because that is what it still is.}
\pair{06-data-structures/packing/2}{Its head is \texttt{List}, it is a list, and it
is not an atom: every structural predicate answers as it would for the unpacked
list.}
\pair{06-data-structures/packing/3}{Only \B{NDArrayQ} reveals that it is packed.}
\pair{06-data-structures/packing/4}{And it is \B{SameQ} to the plain list---same
value, same identity.}
\pair{06-data-structures/packing/5}{\B{DataType} reports the element type\index{data
type!machine (dtype)} of the buffer: a \texttt{"float64"} of doubles.}
\end{codepairs}

The contract is exact: packing changes representation and \emph{nothing else}---not
the value, not an element's head, not the printed form, not the ordering. So the
representation Mathilda picks on its own can never be observed. That single rule is
what the rest of this section is really about.

\subsection{The exactness contract}
\label{ssec:ds-exact}

The sharpest test of the contract\index{packed array!exactness contract} is exact
integers. A list of integers packs to an \texttt{int64} buffer---and every element
must still come back an \emph{Integer}, not a real, because \(1 \ne 1.0\) is a
difference the buffer is forbidden to introduce.

\begin{codepairs}
\pair{06-data-structures/packing/7}{An all-integer list packs to \texttt{"int64"},
not to floats.}
\pair{06-data-structures/packing/8}{And an element read back out has head
\texttt{Integer}, exactly as from the unpacked list.}
\pair{06-data-structures/packing/9}{A reduction stays exact: the sum of \(1\ldots
1000\) is the integer \(500500\), computed on the buffer but with no loss of
exactness.}
\pair{06-data-structures/packing/10}{Where a buffer cannot hold the exact answer, it
abandons: \(\mathit{Range}[10]/2\) is a list of exact \B{Rational}s, which no
machine buffer represents, so the ordinary list gives them.}
\pair{06-data-structures/packing/11}{And a list holding a \B{Rational} refuses to
pack at all---\B{ToNDArray} returns it unchanged rather than round it.}
\end{codepairs}

The rule is the same one that governs the numeric tower everywhere in Mathilda:
compute the buffer answer only when it is \emph{the} answer the interpreter would
give, and otherwise fall back to the slower, exact path. An overflow past
\texttt{int64}, a rational, a radical, a symbol---each makes the buffer decline and
the ordinary list answer.

\calloutnote{Pitfall}{Only \B{ToNDArray} widens a mixed integer/real list to a float
buffer, because you explicitly asked for a single machine type. \emph{Automatic}
packing never does: it leaves \(\{1, 2, 3.\}\) an ordinary list, since silently
turning the exact \(1\) into \(1.0\) would be an observable change and the system is
forbidden to make one on its own. The asymmetry is the contract in action---a
representation you request may cost exactness; one the system chooses for you never
may.}

\subsection{When packing happens}
\label{ssec:ds-auto}

You rarely call \B{ToNDArray} yourself. The common case is that a \emph{producer}
builds a packed array directly: \B{Range}, \B{ConstantArray}, \B{Table} with a
machine-real body, and a handful of others write the buffer without ever
materializing the nodes---which is where most of the win is, since building a million
nodes and then packing them costs far more than writing a million doubles.

\begin{codepairs}
\pair{06-data-structures/packing/12}{\B{Range} of reals packs automatically once the
result is big enough.}
\pair{06-data-structures/packing/13}{A list written out literally in the source never
packs, at any size---there is no producer to opt in, and a source literal is left
exactly as typed.}
\pair{06-data-structures/packing/14}{The master switch \B{\$AutoArrayPacking} turns
automatic packing off\ldots}
\pair{06-data-structures/packing/15}{\ldots and back on. It never touches an explicit
\B{ToNDArray}---that is the user asking, not the system guessing.}
\end{codepairs}

The distinction in the middle pairs surprises people: \mcode{Range[1., 300.]} packs
but the identical-looking literal \(\{1., 2., 3., 4.\}\) does not. Packing is opt-in
per producer, and a hand-written literal has no producer to opt in. A threshold of
four elements (the size of a \(2\times2\) matrix, so that any matrix packs) keeps the
mechanism from paying its own overhead on tiny lists.

\subsection{The visible NDArray}
\label{ssec:ds-visible}

The packed list is one of two faces of the same storage. The other is the
\emph{visible} \B{NDArray}---the identical dense buffer, but wearing its own head, so
it announces itself as a distinct, purely numeric object modelled on numpy's
\texttt{ndarray}. The two share every byte of representation and differ only in what
they claim to be.

\begin{codepairs}
\pair{06-data-structures/ndarray/2}{\B{NDArray} prints wrapped in its own head\ldots}
\pair{06-data-structures/ndarray/3}{\ldots whose \B{Head} is \texttt{NDArray}, not
\texttt{List}.}
\pair{06-data-structures/ndarray/4}{So it is \emph{not} \B{SameQ} to the plain
list---the head is part of a value's identity.}
\pair{06-data-structures/ndarray/5}{\B{NDArrayQ} is \texttt{True} for both faces: it
asks about storage.}
\pair{06-data-structures/ndarray/6}{\B{PackedArrayQ} is the narrower test---\texttt{True}
only for the packed \emph{list}, \texttt{False} for the visible \texttt{NDArray},
which is an atom rather than a list that happens to be packed.}
\pair{06-data-structures/ndarray/7}{\B{Normal} converts a visible array back to an
ordinary nested list.}
\end{codepairs}

The difference is one of intent. A packed list is Mathilda saying ``this is an
ordinary list; I have merely chosen to store it efficiently.'' A visible
\B{NDArray} is you saying ``this is a numeric array; treat it as one''---so it
combines with scalars by numpy-style broadcasting, refuses symbolic operands, and
stays a closed system under the linear-algebra routines. \B{ToPackedArray} and
\B{FromPackedArray} are Mathematica's names for \B{ToNDArray} and \B{FromNDArray},
the same builtins under a second registration.

\begin{codepairs}
\pair{06-data-structures/ndarray/8}{\B{ToPackedArray} widens a mixed integer/real
list to a single \texttt{"float64"} buffer\ldots}
\pair{06-data-structures/ndarray/9}{\ldots the integers becoming doubles, since one
buffer holds one machine type.}
\pair{06-data-structures/ndarray/11}{A derived array stays packed: adding a scalar to
a packed list gives a packed list.}
\pair{06-data-structures/ndarray/12}{So does mapping an elementary function over
it---\B{Sin} of a packed buffer is a packed buffer.}
\pair{06-data-structures/ndarray/13}{And \B{N} widens a packed \texttt{int64} buffer
to a packed \texttt{float64} one in a single pass, rather than boxing a million
reals.}
\pair{06-data-structures/ndarray/15}{The visible \B{NDArray} carries its rank and
shape directly, so \B{Dimensions} reads them in \(O(1)\).}
\end{codepairs}

\usagebox{NDArray}

\subsection{The transparency gate}
\label{ssec:ds-gate}

If a packed list can flow into any operation, what stops an operation that knows
nothing about buffers---a pattern match, a \B{Cases} scan, a user-defined rule---from
reading garbage out of one? The answer is the \emph{transparency
gate}\index{packed array!transparency gate}, and it is the mechanism that makes the
whole abstraction safe.

\begin{codepairs}
\pair{06-data-structures/gate/2}{Start with a packed range.}
\pair{06-data-structures/gate/3}{\B{Cases}---which has no idea packing exists---filters
it with a pattern and gets exactly the right answer.}
\pair{06-data-structures/gate/4}{\B{Reverse} understands buffers, so it reverses in
place\ldots}
\pair{06-data-structures/gate/5}{\ldots and keeps the result packed.}
\pair{06-data-structures/gate/6}{And the buffer path and the ordinary path agree on
the answer, by construction.}
\end{codepairs}

\calloutnote{Under the hood}{The gate lives in the evaluator (\texttt{src/eval.c}),
just before a call reaches its builtin. Its rule inverts the usual default: a packed
buffer is \emph{materialized} back into ordinary nodes for any head that has not
explicitly declared itself packed-aware on the \texttt{AWARE} list in
\texttt{src/pack.c}. So \B{Cases}, \B{Position}, \B{ReplaceAll}, pattern matching,
and every user rule are correct with no code of their own---they simply never see a
buffer. A head like \B{Reverse}, \B{Total}, \B{Sin}, or \B{Dot} that \emph{has} opted
in and been checked receives the buffer intact and runs its C kernel. Forgetting to
opt a head in costs a materialization---a slowdown---never a wrong answer. The
storage overlay is arranged so that a walker which slipped past the gate and iterated
a buffer as if it were a node array would segfault immediately rather than return
something plausible, turning the dangerous silent-wrong-answer class into a loud
crash.}

\calloutnote{Performance}{The payoff is large and the design record measures it on a
\(10^6\)-element list of reals: \B{Length} falls from about \(30.6\)
milliseconds on an ordinary list to \(2\) microseconds packed, \B{Total} from
\(95.8\) milliseconds to \(4.0\), and \B{Sin} from \(320\) milliseconds to
\(16.8\)---the \(30\)-nanosecond-per-element-per-pass evaluator sweep replaced by a
tight loop over contiguous memory (\texttt{docs/design/packed\_arrays.md}). Because a
packed buffer is exactly the representation \B{Compile} wants, these same arrays cross
the compiled-code boundary without being unpacked and repacked at each end---the
subject of Chapter~\ref{ch:compilation}.}

\section*{Where this connects}

Data structures are the substrate the rest of Mathilda stands on. The uniform
expression tree of \S\ref{sec:ds-expressions} is what makes the pattern matcher and
rule engine of Chapter~\ref{ch:programming} possible: a pattern is an expression, a
rule is an expression, and matching one tree against another is the single mechanism
behind \B{Map}, \B{Cases}, \B{ReplaceAll}, and every definition you write. The
levelspecs and structural operators introduced here reappear there as the working
vocabulary of functional programming. Associations are the natural output of the
grouping and counting that descriptive statistics (\S\ref{sec:statistics}) runs on,
and their \(O(1)\) keyed store underlies any code that indexes data by name. The
packed-array substrate of \S\ref{sec:ds-packed} is where numeric Mathilda gets its
speed: it feeds the machine-precision linear algebra of \S\ref{sec:linalg}, and it is
half of the story of \B{Compile} (Chapter~\ref{ch:compilation}), the other half being
the bytecode VM that consumes those buffers. For the exhaustive options of any
function named here, its reference page is one \texttt{?Name} away at the prompt---and
the full machinery of the evaluator, the matcher, and the transparency gate is laid
out in Chapter~\ref{ch:internals}.

\chapter{Programming in Mathilda}
\label{ch:programming}

Everything in the previous chapters was, in a sense, one long computation: you
typed an expression, Mathilda rewrote it until it stopped changing, and printed
the result. This chapter is about writing \emph{your own} rewrites---about
turning Mathilda from a calculator you drive into a language you program. It
supports three styles, and the point of the chapter is that they are not three
separate languages bolted together but three faces of one mechanism. A
\emph{pattern-based} program describes what an expression looks like and how to
transform it. A \emph{procedural} program says do this, then that, changing
variables as it goes. A \emph{functional} program builds a result by applying
and composing functions rather than by mutating anything. You can mix all three
in a single line, because underneath they are the same thing: rules, applied by
the evaluator to expressions, under the control of a handful of
attributes.\index{programming language}

That last clause is the load-bearing one, and it is worth stating before any
syntax. Two facts from the system's design govern everything here. First,
\emph{everything is an expression}\index{expression!uniformity}---a list, a rule,
a pattern, a function definition, a loop are all trees with a head and
arguments, so the same tools operate on all of them. Second, \emph{attributes
drive evaluation}\index{attribute}: before Mathilda touches a call \texttt{f[\ldots]}
it reads the attributes of \texttt{f}, and those bits decide whether the
arguments are evaluated or held, threaded over lists, flattened, or sorted. Once
you see the paradigms as different ways of installing and firing rules, and the
attributes as the switches that shape when rules fire, the language stops being a
grab-bag of commands and becomes a single coherent machine. We build up to that
machine deliberately, closing the chapter with the attributes that were quietly
running the whole time.

We take pattern matching first, because it is the most characteristic style and
the one the other two lean on. Then procedural programming---scoping, branching,
and loops---for the moments when a plain imperative recipe is the clearest thing
to write. Then functional programming, the shortest and often the fastest code.
And finally evaluation control: how to stop the evaluator, restart it, and hold a
head inert, which is where the attribute story comes home.

\section{Pattern matching and rules}
\label{sec:prog-patterns}

A \emph{pattern}\index{pattern matching} is a template that describes a class of
expressions. Matching a pattern against a subject is the act of unifying the two
structurally---checking that the subject has the described shape and, along the
way, binding names to the pieces that filled the blanks. This one operation is
the engine beneath \B{Cases}, \B{Replace}, every rule you write with
\texttt{->}, and every function you define with \texttt{:=}. Learn it once and
you have learned the spine of the language.

\subsection{The blank and its kin}
\label{ssec:prog-blank}

The atom of every pattern is the \emph{blank}, written \texttt{\_}
(\B{Blank}\index{pattern matching!Blank (underscore)}): it matches any single
expression. A blank can be \emph{typed} by writing a head after it---\texttt{\_Integer}
matches only an expression whose head is \texttt{Integer}---and it can be
\emph{named}, as in \texttt{x\_}, which matches anything and binds the name
\texttt{x} to whatever it matched. Two relatives match \emph{runs} of arguments
rather than a single one: \B{BlankSequence} (\texttt{\_\_}) matches one or more
expressions, and \B{BlankNullSequence} (\texttt{\_\_\_}) matches zero or more.
\B{MatchQ}\texttt{[expr, patt]} asks the yes/no question directly, and is the
quickest way to build an intuition for what fits what.

\begin{codepairs}
\pair{07-programming/patterns/1}{The bare typed blank: \(42\) has head
\texttt{Integer}.}
\pair{07-programming/patterns/2}{\(3.14\) does not---its head is \texttt{Real}, so
\texttt{\_Integer} fails.}
\pair{07-programming/patterns/3}{\ldots but \texttt{\_Real} matches it. A typed
blank is a structural test on the head.}
\pair{07-programming/patterns/4}{Head-typed blanks reach the atomic numeric heads:
\(1/2\) is \texttt{Rational[1, 2]}.}
\pair{07-programming/patterns/5}{And \(3 + 4\,i\) is \texttt{Complex[3, 4]}.}
\pair{07-programming/patterns/6}{Structure, not just heads: \texttt{f[\_, \_]}
matches a two-argument \texttt{f}.}
\pair{07-programming/patterns/7}{\texttt{\_\_} (\B{BlankSequence}) matches a run of
one or more arguments.}
\pair{07-programming/patterns/8}{An empty \texttt{f[]} has no arguments, so
\texttt{\_\_} (one or more) fails\ldots}
\pair{07-programming/patterns/9}{\ldots while \texttt{\_\_\_}
(\B{BlankNullSequence}, zero or more) matches it.}
\pair{07-programming/patterns/10}{Sequence blanks match \emph{positions}: some
prefix, a \(3\), some suffix.}
\pair{07-programming/patterns/11}{A named typed blank binds: the rule fires only on
the integer-argument \texttt{f} and rewrites it, leaving the others untouched.}
\end{codepairs}

The last line is the whole idea in miniature. The named blank \texttt{n\_Integer}
does two jobs at once---it \emph{restricts} the match to an integer argument and
it \emph{captures} that argument as \texttt{n}, ready to be used on the
right-hand side. The two calls that do not fit, \texttt{f[x]} and \texttt{f[1, 2]},
are simply left alone; a pattern that does not match is not an error, it is a
non-event.\index{pattern matching!named pattern}

\calloutnote{Under the hood}{Matching is structural tree unification with
backtracking, in \texttt{src/match.c}. A named pattern binds its variable in a
\texttt{MatchEnv}; a sequence blank (\texttt{\_\_}, \texttt{\_\_\_}) forces a
\emph{search} over how to partition the remaining arguments, and the matcher
backtracks through the partitions until the rest of the pattern also fits. A
non-linear pattern---the same name used twice, as in \texttt{f[x\_, x\_]}---simply
checks that the two bindings are structurally equal, which is why it can express
``two equal arguments'' with no extra syntax.}

\subsection{Selecting by pattern}
\label{ssec:prog-select}

Once a pattern can say \emph{what to look for}, a family of functions puts it to
work searching a larger expression. \B{Cases} keeps the elements that match;
\B{Count} tallies them; \B{Position} reports where they are; \B{DeleteCases}
removes them; and \B{Select}, the workhorse's cousin, keeps elements for which a
\emph{predicate function} returns \texttt{True}. A rule attached to \B{Cases}
transforms each match as it is collected.

\begin{codepairs}
\pair{07-programming/cases/1}{A deliberately mixed bag---integers, a string, a
real, symbols, a power.}
\pair{07-programming/cases/2}{\B{Cases} with \texttt{\_Integer} keeps the three
integers.}
\pair{07-programming/cases/3}{With the predicate pattern \texttt{\_?NumberQ} it
keeps every number, real included.}
\pair{07-programming/cases/4}{\B{Count} returns how many matched rather than which.}
\pair{07-programming/cases/5}{\B{Position} gives the paths of the matches---here the
two places \texttt{b} occurs.}
\pair{07-programming/cases/6}{\B{Select} filters by a predicate: the primes up to
\(20\).}
\pair{07-programming/cases/7}{\B{DeleteCases} is the complement of \B{Cases}: drop
the \texttt{Missing[]} markers.}
\pair{07-programming/cases/8}{A rule on \B{Cases} transforms each match---here each
pair to its sum.}
\end{codepairs}

The division of labor is worth noticing. \B{Cases} and its relatives take a
\emph{pattern} and match structure; \B{Select} takes a \emph{function} and tests
each element. They overlap---\texttt{Cases[list, \_?EvenQ]} and
\texttt{Select[list, EvenQ]} agree---but the pattern form can do things a
predicate cannot, such as destructure (\texttt{\{a\_, b\_\}}) and rewrite in the
same breath, as pair~8 does.

\subsection{Constraining a match}
\label{ssec:prog-constrain}

A bare blank is often too generous. Three constructs narrow it. A
\emph{pattern test}, \texttt{patt\,?\,f}, keeps the match only if the function
\texttt{f} applied to it returns \texttt{True}. A \emph{condition},
\texttt{patt\,/;\,cond}, keeps it only if the boolean expression \texttt{cond}---which
may mention the pattern's bound names---is \texttt{True}. And
\emph{alternatives}, \texttt{p1\,|\,p2\,|\,\ldots}, match any one of several
patterns. The complementary \texttt{Except[c]} matches anything that does
\emph{not} match \texttt{c}.

\begin{codepairs}
\pair{07-programming/constrain/1}{\texttt{\_?EvenQ}: the blank, then the predicate
\B{EvenQ} applied to what it matched.}
\pair{07-programming/constrain/2}{\texttt{n\_ /; n > 7}: a condition on the bound
name \texttt{n}---an arbitrary boolean, not just a head test.}
\pair{07-programming/constrain/3}{\texttt{2 | 4 | 6}: alternatives match any one of
the listed forms.}
\pair{07-programming/constrain/4}{A condition drives a rewrite: negate exactly the
negative elements. This is \(\lvert\cdot\rvert\) written as a rule.}
\pair{07-programming/constrain/5}{A condition may relate two bound names:
\texttt{a < b} keeps only the increasing pairs.}
\pair{07-programming/constrain/6}{\texttt{Except[\_String]} keeps everything that is
not a string.}
\end{codepairs}

The difference between \texttt{?} and \texttt{/;} is the difference between a
test on the \emph{whole match} and a test on its \emph{named parts}.
\texttt{\_?f} hands the matched expression to \texttt{f} and asks for a verdict;
\texttt{patt /; cond} evaluates an arbitrary predicate that can see every name the
pattern bound, so \texttt{\{a\_, b\_\} /; a < b} can compare two pieces at once, a
thing no single-argument predicate test can do. Reach for the test when the
condition is ``this one thing passes \texttt{f}'', and for the condition when the
constraint couples several bindings.

\calloutnote{Theory}{A condition can make matching \emph{undecidable} in general,
because \texttt{cond} is an arbitrary Mathilda computation---the matcher cannot
know whether it halts. This is the price of the expressiveness: the same
generality that lets \texttt{/;} express any side condition means a pattern is a
tiny program, not a finite automaton. In practice conditions are cheap boolean
checks, but it is worth remembering that ``does this match?'' is, with
\texttt{/;} in play, as hard as ``does this program return \texttt{True}?''}

\subsection{Rules and replacement}
\label{ssec:prog-rules}

A \emph{rule} pairs a pattern with a replacement. \B{Rule} (\texttt{lhs -> rhs})
is \emph{immediate}: its right-hand side is evaluated once, when the rule is
built. \B{RuleDelayed} (\texttt{lhs :> rhs}) is \emph{delayed}: the right-hand
side is held and re-evaluated afresh on every match, after the pattern's names
are bound. Rules are applied by three operators that differ only in \emph{where}
and \emph{how often} they look: \B{Replace} tries the rule at one level (the
whole expression by default), \B{ReplaceAll} (\texttt{/.}) sweeps the entire tree
top-down, and \B{ReplaceRepeated} (\texttt{//.}) applies \texttt{/.} over and over
until the expression stops changing.

\begin{codepairs}
\pair{07-programming/rules/1}{\B{Replace} rewrites the whole expression when it
matches the rule's left-hand side.}
\pair{07-programming/rules/2}{\B{ReplaceAll} (\texttt{/.}) finds the \(x^2\)
\emph{inside} the sum and rewrites it in place.}
\pair{07-programming/rules/3}{It rewrites every occurrence in one downward sweep:
both \texttt{x} and the \texttt{x} inside \(x^2\).}
\pair{07-programming/rules/4}{A list of rules is applied together---a simultaneous
substitution.}
\pair{07-programming/rules/5}{Given a list of \emph{rule lists}, \B{Replace} maps
over them, giving one answer per alternative.}
\pair{07-programming/rules/6}{\texttt{/.} applies a rule once per part and does not
re-descend into what it just built---so one \texttt{f} is peeled off.}
\pair{07-programming/rules/7}{\texttt{//.} (\B{ReplaceRepeated}) repeats to a fixed
point, peeling every \texttt{f} until nothing changes.}
\pair{07-programming/rules/8}{\B{Rule} (\texttt{->}) evaluates its RHS \emph{once};
the counter is bumped at rule-build time, so all three parts get the same \(1\).}
\pair{07-programming/rules/9}{\B{RuleDelayed} (\texttt{:>}) re-evaluates its RHS per
match, so the counter climbs \(1, 2, 3\).}
\end{codepairs}

Pairs~6 and~7 draw the line between \texttt{/.} and \texttt{//.}: the single
sweep peels one layer, the repeated sweep runs to a fixed point. Pairs~8 and~9
draw the sharper and more consequential line between \texttt{->} and
\texttt{:>}.\index{rule!immediate vs delayed} With a passive right-hand side the
two are indistinguishable, but the moment the RHS \emph{does} something---calls a
random generator, reads a counter, evaluates a costly function---the choice
matters. Use \texttt{:>} whenever the right-hand side should be recomputed for
each match, which is nearly always the case for rules with bound pattern
variables.

\calloutnote{Pitfall}{\texttt{/.}\ descends into \emph{heads} as well as
arguments, because a head is just another part of the tree. A blank rule like
\texttt{\_ -> new} applied to \texttt{\{a, b\}} will therefore also match the head
\texttt{List} and replace it---rarely what you want. When you mean ``the elements,
not the structure,'' constrain the pattern (\texttt{\_Symbol}, say) or apply the
rule at a definite level with \texttt{Replace[expr, rule, \{1\}]}, as pairs~8 and~9
do to keep the counter demonstration clean.}

\usagebox{ReplaceAll}

\subsection{Functions as rules: DownValues}
\label{ssec:prog-downvalues}

Here the whole apparatus pays off. Defining a function with \texttt{:=}
(\B{SetDelayed}) does not create a new kind of object---it \emph{installs a
delayed rule} on a symbol. \texttt{square[x\_] := x\^{}2} attaches to the symbol
\texttt{square} the rule ``a call matching \texttt{square[x\_]} rewrites to
\texttt{x\^{}2}''. These rules are the symbol's
\emph{DownValues}\index{DownValues}, and every time you write \texttt{square[7]}
the evaluator matches the call against them, exactly as \texttt{/.}\ would. A
function \emph{is} a bundle of pattern rules.

\begin{codepairs}
\pair{07-programming/downvalues/1}{Install a rule on \texttt{square} with
\texttt{:=}.}
\pair{07-programming/downvalues/2}{It fires by matching: on a number, a symbolic
sum, and a constant alike.}
\pair{07-programming/downvalues/3}{A second symbol, \texttt{area}, gets its own
rule\ldots}
\pair{07-programming/downvalues/4}{\ldots and a \emph{second} rule with a different
argument count.}
\pair{07-programming/downvalues/5}{Dispatch is by pattern: one argument gives the
disk area, two the rectangle.}
\pair{07-programming/downvalues/6}{The base cases of \texttt{fib} are installed with
\texttt{=} (immediate).}
\pair{07-programming/downvalues/7}{The second base case.}
\pair{07-programming/downvalues/8}{The recursive case, with \texttt{:=}, refers to
\texttt{fib} of smaller arguments.}
\pair{07-programming/downvalues/9}{\texttt{fib[10]} unwinds through the rules to
\(55\).}
\pair{07-programming/downvalues/10}{\texttt{collatz} on even integers---the guard
\texttt{\_?EvenQ} selects the rule\ldots}
\pair{07-programming/downvalues/11}{\ldots and a separate rule handles the odd case.}
\pair{07-programming/downvalues/12}{\B{NestList} then iterates the two-rule function:
the Collatz orbit of \(27\).}
\end{codepairs}

The \texttt{fib} definition shows why the ordering of rules matters and how
recursion falls out for free. The two base cases \texttt{fib[0]} and
\texttt{fib[1]} are \emph{more specific} than the general \texttt{fib[n\_]}, and
Mathilda tries specific rules before general ones, so the recursion bottoms out
correctly. The \texttt{collatz} pair shows dispatch by \emph{guard}: two rules
with the same shape but disjoint conditions act as a single case-split function.

\calloutnote{Under the hood}{DownValues live on the head symbol in the symbol
table (\texttt{src/symtab.c}) as a linked list, kept in order of
\emph{specificity}---genuinely more-specific patterns ahead of more-general ones,
with insertion order breaking ties---so that \texttt{fib[0]} is tried before
\texttt{fib[n\_]} however the two were entered (\S\ref{sec:int-symtab}). Evaluating
\texttt{square[7]} is then just \texttt{apply\_down\_values}: walk the list, and
the first pattern that matches wins---the same matcher, the same first-match
rule, as \texttt{/.}. There is no separate ``function call'' machinery in the
kernel; a user function and a rewrite rule are one mechanism.}

This unity buys a classic trick almost for free. A rule whose right-hand side
\emph{installs another rule} turns a function into its own cache:

\mtranscript{07-programming/memoize}

The definition \texttt{fib[n\_] := fib[n] = fib[n-1] + fib[n-2]} reads oddly at
first. Each time it fires for a new \texttt{n}, the inner \texttt{fib[n] = \ldots}
is an immediate assignment that installs a \emph{new}, specific DownValue
\texttt{fib[n] -> value}. The next request for that same \texttt{n} matches the
new specific rule instead of re-entering the recursion. Naive \texttt{fib} makes
exponentially many calls; this one makes \(O(n)\), which is why \texttt{fib[100]}
returns a twenty-one-digit integer instantly rather than after the heat death of
the machine.\index{memoization}\index{dynamic programming}

\calloutnote{Performance}{Memoization by self-installing rules is the CAS
idiom for dynamic programming. The cost is memory---each computed value becomes a
permanent DownValue on the symbol---so it is a space-for-time trade, not a free
lunch; a long-running kernel that memoizes over a wide domain will accumulate
rules until the symbol is cleared. The win is asymptotic and often enormous, but
the stored table is real, and \B{Clear} is how you reclaim it.}

\section{Procedural programming}
\label{sec:prog-procedural}

Sometimes the clearest description of a computation is a recipe: set this up, do
that repeatedly, choose between these branches. Mathilda has the full
imperative apparatus, and because assignments and loops are themselves
expressions with attributes, they compose with the other two paradigms rather
than sitting apart from them.

\subsection{Assignment and scoping}
\label{ssec:prog-scoping}

Two assignment operators anchor procedural code: \B{Set} (\texttt{=}) evaluates
its right-hand side \emph{now} and binds the result, while \B{SetDelayed}
(\texttt{:=}) holds the right-hand side and re-evaluates it on each use---the same
immediate/delayed distinction as \texttt{->} versus \texttt{:>}, and for the same
reason. A sequence of steps is joined with \B{CompoundExpression}
(\texttt{;}), which evaluates each in turn and returns the last.

The subtler question is \emph{scoping}: how to introduce a variable that is local
to a block of code, so it neither sees nor disturbs the outside world. Mathilda
offers three constructs, and the differences between them are exactly the
differences a working programmer needs to understand. \B{Module} creates genuinely
\emph{fresh} local variables (lexical scoping). \B{With} performs lexical
\emph{substitution}, replacing a name by a value textually before the body runs.
\B{Block} does \emph{dynamic} scoping, temporarily overriding a symbol's global
value for the duration of the body and restoring it afterward.

\begin{codepairs}
\pair{07-programming/scoping/1}{A global \texttt{x} exists in the session.}
\pair{07-programming/scoping/2}{\B{Module} makes a fresh, private \texttt{x}; inside,
it is \(2\).}
\pair{07-programming/scoping/3}{The global \texttt{x} is untouched---still \(100\).}
\pair{07-programming/scoping/4}{\B{With} substitutes \texttt{x -> 5} into the body
before evaluating it.}
\pair{07-programming/scoping/5}{Because \B{With} substitutes \emph{textually}, it
reaches inside \B{Hold}: \texttt{a} becomes \(2\) even though the body is held.}
\pair{07-programming/scoping/6}{A \B{Module} local can be read and assigned from
inside a nested pure function---here a running counter.}
\pair{07-programming/scoping/7}{A global \texttt{n}\ldots}
\pair{07-programming/scoping/8}{\ldots and a function \texttt{g} that reads it.}
\pair{07-programming/scoping/9}{\B{Block} overrides \texttt{n} \emph{dynamically}:
\texttt{g[]} sees the temporary \(3\), so \(3^2 = 9\).}
\pair{07-programming/scoping/10}{\B{Module} does \emph{not}: its private \texttt{n}
is invisible to \texttt{g}, which still reads the global \(100\).}
\pair{07-programming/scoping/11}{Afterward the global \texttt{n} is restored, and
\texttt{g[]} is \(100\) again.}
\end{codepairs}

Pairs~9 and~10 are the crux, and the single most useful thing to internalize
about scoping. \B{Block} localizes a \emph{value}: for the duration of its body,
the symbol \texttt{n} has a new value, and \emph{any} code that reads
\texttt{n}---including a function \texttt{g} defined elsewhere---sees it. \B{Module}
localizes a \emph{name}: it invents a brand-new symbol so its \texttt{n} cannot
possibly collide with anyone else's, which is exactly why \texttt{g}, whose body
refers to the ordinary global \texttt{n}, does not see the Module's version.
Module is what you want almost always, because it is hygienic; Block is the right
tool precisely when you \emph{do} want to reach into other code's use of a symbol---temporarily
raising \texttt{\$RecursionLimit}, say, or overriding a constant for one
computation.\index{scoping!lexical}\index{scoping!dynamic}

\calloutnote{Under the hood}{\B{Module} implements freshness by renaming: its
local \texttt{x} becomes a uniquely numbered symbol \texttt{x\$nnn} (drawn from
\texttt{\$ModuleNumber}) carrying the \texttt{Temporary} attribute, which is why
it cannot clash with a global \texttt{x} or with another Module's. \B{With} needs
no rename because it substitutes and then vanishes---and because the substitution
happens \emph{before} the body evaluates, it is the one scoping construct that
penetrates a \B{Hold} (pair~5). \B{Block} saves the symbol's current value, installs
the new one, and restores the old value on exit even if the body throws. All
three are \texttt{HoldAll}, so the body is not evaluated until the locals are in
place.}

\usagebox{Module}

\subsection{Branching}
\label{ssec:prog-branching}

Three constructs choose among alternatives. \B{If}\texttt{[cond, t, f]} is the
two-way branch. \B{Which}\texttt{[t1, v1, t2, v2, \ldots]} walks a list of
test/value pairs and returns the value of the first test that is \texttt{True}---a
multi-way branch, with a trailing \texttt{True} acting as the catch-all.
\B{Switch}\texttt{[e, form1, v1, \ldots]} instead \emph{pattern-matches} the
value \texttt{e} against a sequence of forms, so it dispatches on shape rather
than on a chain of booleans, with a trailing \texttt{\_} as the default.

\begin{codepairs}
\pair{07-programming/branching/1}{\B{If} chooses the true branch.}
\pair{07-programming/branching/2}{\B{Which} as a multi-way classifier\ldots}
\pair{07-programming/branching/3}{\ldots dispatched over a list with
\texttt{/@}.}
\pair{07-programming/branching/4}{The sign function, three cases with a \texttt{True}
default.}
\pair{07-programming/branching/5}{Its values on a negative, zero, and positive
input.}
\pair{07-programming/branching/6}{\B{Switch} dispatches on \emph{pattern}, not on a
boolean chain.}
\pair{07-programming/branching/7}{So it separates an integer, a real, and a symbol
by their heads.}
\pair{07-programming/branching/8}{\B{Switch} inside a pure function, mapped over a
list.}
\end{codepairs}

Which and Switch are the same idea approached from the two sides of the chapter:
Which tests \emph{predicates} (procedural), Switch matches \emph{patterns}
(the rewriting style). Choose Switch when the distinction you care about is
structural---a head, a shape---and Which when it is a numeric or logical
condition.

\calloutnote{Under the hood}{\B{If} is \texttt{HoldRest}, \B{Which} is
\texttt{HoldAll}, and \B{Switch} is \texttt{HoldRest}: the branches are held and
only the chosen one is ever evaluated. Without that, \texttt{If[x > 0, expensive,
cheap]} would compute both arms before choosing---the attribute is not an
optimization but a correctness requirement. This is the first place the attribute
system does visible work; \S\ref{sec:prog-attributes} makes the pattern explicit.
The three lower to machine branches inside \B{Compile} as well, so a classifier
used in a numeric loop costs nothing (Chapter~\ref{ch:compilation}).}

\subsection{Loops}
\label{ssec:prog-loops}

For genuine iteration Mathilda has \B{Do} (repeat over a range), \B{While}
(repeat while a test holds), and \B{For} (the C-style initialize/test/increment
loop). \B{Break} exits the enclosing loop early and \B{Continue} skips to its next
iteration; both take effect the instant they are evaluated.

\begin{codepairs}
\pair{07-programming/loops/1}{\B{Do} accumulates \(1 + 2 + \cdots + 100 = 5050\).}
\pair{07-programming/loops/2}{The same shape builds \(6! = 720\).}
\pair{07-programming/loops/3}{\B{While} doubles until it passes \(100\).}
\pair{07-programming/loops/4}{\B{For} sums the first ten squares, \(385\).}
\pair{07-programming/loops/5}{\B{Break} leaves the loop as soon as \texttt{i}
exceeds \(3\), so the last value recorded is \(3\).}
\pair{07-programming/loops/6}{\B{Continue} skips the even \texttt{i}, summing only
the odds: \(1+3+5+7+9 = 25\).}
\end{codepairs}

Each loop here returns a value through an accumulator variable, which is the
honest way to use them: a loop's own return value is \texttt{Null}. That is the
tell that loops are about \emph{effects}---changing \texttt{s}, \texttt{p},
\texttt{n}---not about producing a value, and it is the seam where procedural code
gives way to the functional style of the next section, which produces values
without mutating anything.

\calloutnote{Performance}{A loop whose body is entirely machine arithmetic takes
an automatic fast path (\texttt{src/numloop.c}): it is compiled to a small
double-stack program and run with no \texttt{Expr} allocation per iteration, so
\texttt{Do[s = s + i, \{i, 1, 10\^{}7\}]} is a tight machine loop, not ten million
evaluator steps. The moment the body does something non-numeric---a \B{Print}, a
symbolic head, a call to a user function---the fast path silently declines and the
interpreter runs the body exactly as written, so the answer is always the
interpreter's. An exact-integer counter accumulates in overflow-checked
\texttt{int64} and promotes to a bignum rather than wrapping, which is why
\texttt{Do[p = p\,i, \{i, 1, 6\}]} is exactly \(720\) and would be exactly
\(25!\) at the top end.}

\section{Functional programming}
\label{sec:prog-functional}

Functional code builds a result by applying and composing functions instead of
mutating variables. It is usually the shortest way to say what you mean, and---because
it avoids explicit loops---often the fastest, since the same numeric
fast path that accelerates \B{Do} also drives \B{Map}, \B{Nest}, and \B{Fold}.

\subsection{Pure functions}
\label{ssec:prog-purefunc}

A \emph{pure} (anonymous) function is one written without a name.
\B{Function}\texttt{[x, body]} names its parameter explicitly; the terser
\texttt{body\,\&} leaves the parameters implicit, referring to them by
\B{Slot}: \texttt{\#} or \texttt{\#1} is the first argument, \texttt{\#2} the
second, and \B{SlotSequence} \texttt{\#\#} is all of them at once.\index{pure function}\index{anonymous function}

\begin{codepairs}
\pair{07-programming/purefunc/1}{\texttt{\#\^{}2 \&} is the squaring function;
\texttt{\#} is its argument.}
\pair{07-programming/purefunc/2}{\texttt{\#1} and \texttt{\#2} are the first and
second arguments.}
\pair{07-programming/purefunc/3}{The named form \B{Function}\texttt{[u, u\^{}3]} is
the same idea, spelled out.}
\pair{07-programming/purefunc/4}{A named function of two parameters.}
\pair{07-programming/purefunc/5}{Pure functions are values: pass one straight to
\B{Map}.}
\pair{07-programming/purefunc/6}{And to \B{Select} as a predicate.}
\pair{07-programming/purefunc/7}{\texttt{\#\#} (\B{SlotSequence}) is the whole
argument run, spliced---here twice.}
\pair{07-programming/purefunc/8}{\texttt{\#2} before \texttt{\#1} reorders the
arguments.}
\end{codepairs}

The \texttt{\&} notation is the connective tissue of functional Mathilda: it lets
you write a one-off function exactly where it is needed---as the thing to map, the
predicate to select by, the step to iterate---without the ceremony of naming it.
Named \B{Function} earns its keep when the body is long enough that \texttt{\#1},
\texttt{\#2}, \texttt{\#3} become hard to read, or when you need its third-argument
form to give the function evaluation attributes of its own.

\calloutnote{Under the hood}{A pure function binds its parameters by
\emph{lexical substitution} (\texttt{src/purefunc.c}), not through the global
symbol table: when \texttt{(\#\^{}2 \&)[a + b]} fires, \texttt{\#} is replaced by
\texttt{a + b} in the body before it evaluates. Nested pure functions scope
correctly---an inner \texttt{\&} establishes a fresh set of slots that shadow the
outer ones---but named-parameter \B{Function}s are deliberately \emph{not}
closures: an inner \texttt{Function[e, \ldots]} does not capture an outer
function's parameter. When you need that capture, inject the value with \B{With}.}

\usagebox{Function}

\subsection{Mapping and applying}
\label{ssec:prog-map}

Two operators cover most higher-order work. \B{Map} (\texttt{/@}) applies a
function to each element of an expression. \B{Apply} (\texttt{@@}) does something
subtler and often surprising to newcomers: it replaces the \emph{head} of an
expression. \texttt{Plus @@ \{1, 2, 3, 4\}} turns the list's head \texttt{List}
into \texttt{Plus}, giving \texttt{Plus[1, 2, 3, 4]}, which evaluates to the sum.
\B{MapThread} and \B{Thread} generalize mapping to several lists at once.

\begin{codepairs}
\pair{07-programming/mapapply/1}{\B{Map} wraps each element in \texttt{box}.}
\pair{07-programming/mapapply/2}{\texttt{/@} is the same operator in infix form.}
\pair{07-programming/mapapply/3}{\B{Apply} replaces the head \texttt{List} with
\texttt{Plus}\ldots}
\pair{07-programming/mapapply/4}{\ldots which \texttt{@@} writes infix; the list
becomes a sum.}
\pair{07-programming/mapapply/5}{With a symbolic head the mechanism is bare:
\texttt{List} becomes \texttt{f}.}
\pair{07-programming/mapapply/6}{\texttt{@@@} applies at level one---the head of each
sublist.}
\pair{07-programming/mapapply/7}{\B{MapThread} zips several lists through a function
of several arguments.}
\pair{07-programming/mapapply/8}{\B{Thread} spreads a function over the list
argument, copying the scalar \texttt{u} across.}
\pair{07-programming/mapapply/9}{\B{Thread} over \texttt{==} turns one list equation
into a list of equations.}
\end{codepairs}

\B{Map} changes the \emph{elements}, \B{Apply} changes the \emph{head}: holding
that distinction straight is most of what it takes to read functional Mathilda
fluently. The two combine constantly---\texttt{Total} is essentially
\texttt{Plus @@}, and a great many one-liners are a \texttt{Map} to transform
elements followed by an \texttt{Apply} to combine them.

\subsection{Building by iteration}
\label{ssec:prog-iterate}

Where a procedural program loops, a functional one iterates a function. \B{Nest}
applies a function \(n\) times; \B{NestList} records every intermediate result.
\B{Fold} threads a binary function through a list, carrying an accumulator;
\B{FoldList} keeps the running history---which is exactly a cumulative reduction.
\B{Composition} (\texttt{@*}) builds the function that applies several functions
in succession.

\begin{codepairs}
\pair{07-programming/iterate/1}{\B{Nest} applies \texttt{f} three times,
symbolically.}
\pair{07-programming/iterate/2}{\B{NestList} keeps every stage of the same
iteration.}
\pair{07-programming/iterate/3}{With a numeric function it collapses to a number:
\((1+\cdot)^2\) thrice from \(1\).}
\pair{07-programming/iterate/4}{A trivial counter, to read the shape of the
output.}
\pair{07-programming/iterate/5}{Newton's iteration for \(\sqrt 2\)---each step
\((x + 2/x)/2\)---converging in a handful of steps.}
\pair{07-programming/iterate/6}{\B{Fold} sums a list through an accumulator, giving
\(10\).}
\pair{07-programming/iterate/7}{\B{FoldList} keeps the partial sums---a cumulative
total.}
\pair{07-programming/iterate/8}{\B{Fold} by Horner's rule assembles the digits
\(1, 9, 2, 6\) into the base-ten number \(1926\).}
\pair{07-programming/iterate/9}{\B{Composition} applies its functions
innermost-first.}
\pair{07-programming/iterate/10}{\texttt{@*} is the same, in infix.}
\end{codepairs}

The Newton example (pair~5) is worth dwelling on, because it shows functional
iteration doing real numerical work: each application of the pure function
\texttt{(\# + 2/\#)/2 \&} is one Newton step for \(\sqrt 2\), and \B{NestList}
exhibits the quadratic convergence---the digits roughly double each step until the
machine-precision limit is reached and the value stops moving. The same
computation as a \B{While} loop would be several lines and a mutable variable;
here it is one expression, and it takes the numeric fast path automatically.

\calloutnote{Performance}{\B{Nest}, \B{NestList}, \B{Fold}, \B{FoldList}, \B{Map}
and their relatives share the loop fast path (\texttt{src/numloop.c}): when the
iterated function is machine arithmetic throughout, the whole iteration runs as
compiled \texttt{double} bytecode with no per-step allocation. And crucially, the
same functional heads \emph{also} lower inside \B{Compile}, so a \B{Fold} or
\B{Nest} written at the prompt is the same \B{Fold} or \B{Nest} that runs inside a
compiled kernel (Chapter~\ref{ch:compilation}). Exactness is preserved position
by position: a value the iteration merely passes through keeps its exact type
even when computed neighbors have gone inexact.}

\section{Controlling evaluation}
\label{sec:prog-evaluation}

All three paradigms rest on the evaluator, and to program well you sometimes need
to \emph{stop} it: to keep an expression in the form you typed rather than the
form it would reduce to. Mathilda gives you fine control over evaluation, and the
tools reveal the machinery that has been running underneath the whole chapter.

\subsection{Holding and releasing}
\label{ssec:prog-hold}

\B{Hold}\texttt{[expr]} keeps \texttt{expr} unevaluated. \B{HoldForm} does the
same but prints without the wrapper showing, so held expressions display
naturally. \B{Evaluate}, inside a held argument, forces \emph{that} argument to
evaluate anyway. And \B{ReleaseHold} strips the wrappers, letting the held
expression finally run.

\begin{codepairs}
\pair{07-programming/holds/1}{\B{Hold} keeps \(1+1\) from evaluating.}
\pair{07-programming/holds/2}{\B{HoldForm} holds too, but prints without showing the
wrapper.}
\pair{07-programming/holds/3}{\B{Hold} shields all of its arguments.}
\pair{07-programming/holds/4}{\B{Evaluate} punches through the hold for one
argument---\(1+1\) becomes \(2\), the other stays.}
\pair{07-programming/holds/5}{\B{ReleaseHold} removes the wrapper and lets
evaluation resume.}
\pair{07-programming/holds/6}{\B{Hold} holds because it \emph{carries the attribute}
\texttt{HoldAll}.}
\pair{07-programming/holds/7}{So does \B{SetDelayed}---which is why \texttt{:=} can
capture an unevaluated right-hand side.}
\end{codepairs}

Pairs~6 and~7 are the reveal. \B{Hold} is not special-cased in the evaluator; it
holds its arguments because it has the attribute \texttt{HoldAll}, and
\emph{anything} with that attribute holds the same way. \B{SetDelayed} has it too,
which is the whole reason \texttt{f[x\_] := body} can store \texttt{body}
unevaluated instead of computing it once at definition time. Holding is not a
feature of a few privileged heads---it is an attribute, and the next subsection
shows how to make it, and the others, do work.

\subsection{Inactive heads}
\label{ssec:prog-inactive}

Holding an \emph{argument} is one thing; sometimes you want to hold a
\emph{head}---to build \texttt{Plus[2, 3]} or \texttt{Integrate[f, x]} as an inert
object that will not fire until you say so. \B{Inactive}\texttt{[f]} is that inert
form of the head \texttt{f}: \texttt{Inactive[f][args]} evaluates its arguments
but does not run \texttt{f}'s own rules, so the call stays symbolic.
\B{Activate} reverses it, turning every \texttt{Inactive[h]} back into \texttt{h}
and letting evaluation proceed.

\begin{codepairs}
\pair{07-programming/inactive/1}{\B{Inactive}\texttt{[Plus]} is an inert head:
\texttt{Plus} does not fire, so the sum stays symbolic.}
\pair{07-programming/inactive/2}{The \emph{arguments} still evaluate---\(1+1\) and
\(2+2\) reduce---but the outer \texttt{Plus} does not.}
\pair{07-programming/inactive/3}{\B{Activate} releases the head and the addition
finally happens.}
\pair{07-programming/inactive/4}{The chief use: holding an integral inert, so no
integration cascade runs.}
\pair{07-programming/inactive/5}{\B{D} applies the fundamental theorem of calculus
to the inert integral \emph{without} evaluating it.}
\pair{07-programming/inactive/6}{And \B{Activate} runs the held integral when you
want the answer.}
\end{codepairs}

The distinction in pairs~1 and~2 is precise and worth stating: \B{Inactive} freezes
the \emph{head}, not the arguments, so \texttt{Inactive[Plus][1+1, 2+2]} evaluates
its arguments to \texttt{Inactive[Plus][2, 4]} and stops there. This is more than
a curiosity. An inert \texttt{Inactive[Integrate][g, x]} lets you write down an
integral you do not (yet) want to compute---perhaps because the integrand is
non-elementary and the attempt would be expensive---and still manipulate it
symbolically. Pair~5 shows the payoff: \B{D} knows that differentiating an integral
gives back the integrand, and applies that rule to the inert form directly, so an
inert first integral can be verified without ever paying for the
integration.\index{Inactive!fundamental theorem of calculus}

\calloutnote{Under the hood}{Inertness is automatic, not a special rule.
\texttt{Inactive[f][args]} is an expression whose head is the compound
\texttt{Inactive[f]}, and that compound head carries no DownValues, so the call is
already a fixed point of the evaluator---there is nothing to fire. It is the same
structural trick as \texttt{Derivative[n][f][x]}, a call through a compound head.
\B{Activate} is then just a targeted replacement of \texttt{Inactive[h]} by
\texttt{h} followed by re-evaluation.}

\subsection{Attributes drive it all}
\label{sec:prog-attributes}

We can now name the mechanism that has been at work in every section. Before
Mathilda processes a call \texttt{f[a1, \ldots, aN]}, it reads the
\emph{attributes} of \texttt{f} and lets them shape the step: which arguments to
evaluate (the \texttt{Hold} family), whether to thread over lists
(\texttt{Listable}), whether to flatten nested same-head calls (\texttt{Flat}),
whether to sort the arguments (\texttt{Orderless}). \B{Attributes} reports them,
\B{SetAttributes} adds them, and \B{ClearAttributes} removes them---and you can set
them on your own symbols, watching the evaluator's behavior change.

\begin{codepairs}
\pair{07-programming/attributes/1}{\B{Plus} is \texttt{Flat} (associative),
\texttt{Listable}, \texttt{Orderless} (commutative), and more---its algebra is its
attributes.}
\pair{07-programming/attributes/2}{Give a fresh symbol \texttt{circ} the attribute
\texttt{Orderless}\ldots}
\pair{07-programming/attributes/3}{\ldots and its arguments are sorted into canonical
order automatically.}
\pair{07-programming/attributes/4}{Make \texttt{gp} \texttt{Flat}\ldots}
\pair{07-programming/attributes/5}{\ldots and nested \texttt{gp} calls flatten into
one.}
\pair{07-programming/attributes/6}{Make \texttt{hl} \texttt{Listable}\ldots}
\pair{07-programming/attributes/7}{\ldots and it threads element-wise over a list
argument, with no loop written.}
\pair{07-programming/attributes/8}{Over several list arguments it threads them in
parallel.}
\pair{07-programming/attributes/9}{Remove the attribute\ldots}
\pair{07-programming/attributes/10}{\ldots and the threading stops---\texttt{hl}
holds the list whole again.}
\end{codepairs}

This is the design at the heart of Mathilda, and it is why the language feels so
uniform. The evaluator itself is small and generic; it does not know that
\B{Plus} is commutative or that \B{Sin} threads over lists. It knows only how to
read attribute bits and act on them. Commutativity is \texttt{Orderless};
associativity is \texttt{Flat}; element-wise threading is \texttt{Listable};
holding an argument is one of the \texttt{Hold} bits. A new builtin opts into any
of these behaviors by setting a flag, and the same seven or eight attributes
compose to describe everything from arithmetic to your own functions. The three
paradigms of this chapter are three ways of installing rules; the attributes are
the dial that sets \emph{when and how} those rules fire.\index{attribute!Flat}\index{attribute!Orderless}\index{attribute!Listable}\index{attribute!HoldAll}

\calloutnote{Theory}{The evaluator's per-call sequence is fixed and worth
memorizing: evaluate the head; read its attributes; evaluate the arguments except
those held; thread if \texttt{Listable}; flatten if \texttt{Flat}; sort if
\texttt{Orderless}; then try the builtin, then the DownValues. Each attribute is a
step in that pipeline, which is why \texttt{Orderless} on \texttt{circ} and
\texttt{Flat} on \texttt{gp} take effect the instant they are set, with no change
to the symbol's rules---they change how the evaluator \emph{prepares} the call
before any rule is consulted. The mechanism lives in \texttt{src/eval.c} and
\texttt{src/attr.c}.}

\section*{Where this connects}

Programming in Mathilda is not a layer on top of the algebra system; it is the
same system, used deliberately. A function you define with \texttt{:=} is a rule
in the symbol table, matched by the same engine that powers \B{Solve}'s
case analysis and the derivative rules of Chapter~\ref{ch:mathematics}; the
pattern language of \S\ref{sec:prog-patterns} is the language those internal rule
sets are written in. The functional iteration of \S\ref{ssec:prog-iterate} and the
loops of \S\ref{ssec:prog-loops} share one numeric fast path, and both lower into
the bytecode compiler---so the \B{Map}, \B{Fold}, and \B{Nest} you write here are
exactly what \B{Compile} accelerates in Chapter~\ref{ch:compilation}, where the
compilable subset and its cliff are the subject. The attributes of
\S\ref{sec:prog-attributes} are the same bits Chapter~\ref{ch:internals} traces
through the evaluator, and the ``everything is an expression'' model that lets
\B{Map} and \B{ReplaceAll} work on rules, patterns, and functions alike is the
data model of Chapter~\ref{ch:datastructures}. As always, the exhaustive options
of any function named here are one \texttt{?Name} away at the prompt.

\chapter{Compilation and the Compiler}
\label{ch:compilation}

Everything you have done so far has gone through the evaluator: you type an
expression, Mathilda rewrites it until it stops changing, and prints the result.
That machine is wonderfully general---it will happily add a symbol to a bignum,
thread a function over a ragged list, or apply a rule you wrote a moment ago---but
generality has a price. To evaluate \(x^2 + 1\) at a single machine number, the
evaluator still walks a tree, reads attributes, checks for rules, and dispatches
on the head at every node, allocating and freeing expression objects as it goes.
For one evaluation that cost is invisible. Inside a loop that runs a million
times over plain \texttt{double}s, it is almost all of the running time.

This chapter is about the other gear. \B{Compile} takes a numeric expression and
lowers it, once, to typed \emph{bytecode}\index{compilation!bytecode}---a short
program of monomorphic instructions run by a small register machine that never
allocates an expression, never consults the symbol table, and never asks what
type a value has, because the types were all settled at compile time. The answer
is exactly the evaluator's; the speed is eight to several hundred times
greater\index{Compile!performance}. The same machine runs quietly behind
\B{Table}, \B{Plot}, \B{NIntegrate} and \B{Nest} without your asking, a feature
called \emph{auto-compilation}\index{auto-compilation} that this chapter also
takes apart and shows you how to switch off.

A word to prevent one confusion. This is the \mcode{Compile} \emph{builtin}---a
tool inside the running system---not the business of compiling Mathilda itself
from C source, which is Chapter~\ref{ch:building}. Here nothing touches a C
compiler; the ``compiler'' is a component of the evaluator, and its target is a
bytecode of its own.

\section{From an expression tree to bytecode}
\label{sec:compile-basics}

\B{Compile} takes two arguments: a list of parameters and a body to evaluate over
them. The result is a \mcode{CompiledFunction}, an opaque object you call exactly
like any function.

\begin{codepairs}
\pair{08-compilation/basics/1}{\B{Compile} holds its arguments and returns a
\texttt{CompiledFunction}---the body is not evaluated symbolically, it is
\emph{compiled}. It echoes showing its parameters and the source body.}
\pair{08-compilation/basics/2}{Bind it to a name and it behaves like a function.
Here the parameter is declared \mcode{\{x, \_Real\}}: a machine real.}
\pair{08-compilation/basics/3}{Called on a real, it runs the bytecode and returns
a machine real.}
\pair{08-compilation/basics/4}{An exact integer argument is accepted and coerced
to the declared type---\(3\) enters as \(3.0\), so the result is \(10.0\).}
\pair{08-compilation/basics/5}{Parameters can be typed \mcode{\_Integer} \dots}
\pair{08-compilation/basics/6}{\dots or \mcode{\_Complex}; the compiler emits
integer or complex opcodes accordingly.}
\pair{08-compilation/basics/7}{And when an argument is \emph{not} a machine
number, the compiled function quietly hands the work back to the evaluator: a
symbolic \(y\) gives the ordinary symbolic answer. Compilation is an
optimisation, never a restriction on what you may pass.}
\end{codepairs}

The parameter list is the compiler's type declaration. A bare symbol means a
machine real; \mcode{\{x, \_Real\}}, \mcode{\{n, \_Integer\}}, and
\mcode{\{z, \_Complex\}} name the three scalar types explicitly; and a third
entry gives an \emph{array} rank, so \mcode{\{v, \_Real, 1\}} is a rank-1 vector
of reals and \mcode{\{m, \_Real, 2\}} a matrix (\S\ref{sec:compile-arrays}). Every
downstream instruction is chosen from these declared types, which is why the
compiler must know them up front.

\usagebox{Compile}

That last row of the example is the single most important thing to understand
about \mcode{Compile}: its \emph{fallback contract}\index{Compile!interpreter
fallback}. A \texttt{CompiledFunction} always agrees with the evaluator. If you
hand it an argument it cannot represent as a machine number, or if the body uses
something the compiler cannot lower (\S\ref{sec:compile-cliff}), it does not fail
and it does not lie---it re-runs the expression through the interpreter and
returns that. You never trade correctness for speed; at worst you fail to gain
the speed.

\section{What the compiler produces}
\label{sec:compile-bytecode}

The best way to see what compilation \emph{is} is to look at what it produces.
\B{CompilePrint} disassembles a compiled function and writes its bytecode to the
terminal:

\compileprint{08-compilation/scalar}

Read it from the top. The \emph{signature} records the parameter types and the
source body. Then three \emph{registers}: \texttt{R0} holds the argument, the
result comes back in \texttt{R1}, and one scratch register was enough. Then eight
instructions. Each is \emph{monomorphic}\index{compilation!monomorphic opcodes}:
\texttt{ADD\_RK} is ``add a real constant \(K\) to a real register,''
\texttt{MUL\_R} is ``multiply two real registers.'' There is no run-time test of
what type the operands are, because the answer was fixed when the program was
built---that erasure of per-node type dispatch is the whole source of the
speed-up. Notice too that Horner's nested form fell straight out of the source:
the compiler folded each numeric coefficient into an \emph{immediate} operand
(the \(K\) in \texttt{ADD\_RK}) rather than loading it from a register. The
program ends, as every compiled program does, with \texttt{RET}.

Not every body is straight-line. Control flow compiles to
jumps\index{compilation!control flow}: a conditional becomes a compare into a
Boolean register and a \emph{jump-if-zero} (\texttt{JZ}) around the branch not
taken, with the jump targets marked \texttt{>}.

\compileprint{08-compilation/branch}

And a counted loop compiles to a single back-edge---here \texttt{LOOP}, over
\emph{integer} opcodes (\texttt{POWI\_I}, \texttt{ADD\_I}), because every value in
this body is an integer and the ``all-Real fast path'' note is absent:

\compileprint{08-compilation/loop}

\calloutnote{Under the hood}{The register machine is a threaded-code interpreter:
on GCC and Clang each opcode's handler ends by computing the address of the next
handler and jumping straight to it (a computed \texttt{goto}), so there is no
per-instruction \texttt{switch}. Registers are a flat array of 16-byte slots on
the C stack, banked into scalar, array and ``tile'' regions, so a scalar program
allocates nothing at all per call. Elementary and special functions
(\texttt{Sin}, \texttt{Gamma}, \texttt{Erf}, \dots) do not each get their own
opcode; they share a \texttt{KERN} instruction that carries a pointer to the same
machine kernel the packed-array code uses (\S\ref{sec:ds-packed}), so a new
special function becomes compilable for free.}

You do not usually need the disassembly. What you do need, often, is a quick
answer to ``did this body actually compile, and to what?''---and for that there
is \B{CompileDiagnostics}, which type-checks a body without building the object
and reports what it found:

\begin{codepairs}
\pair{08-compilation/diagnostics/1}{The result type is \texttt{Real}; the
optimised program is three instructions, down from five before the optimiser
(constant folding, common-subexpression elimination and dead-code removal run on
the bytecode).}
\pair{08-compilation/diagnostics/2}{A trigonometric identity the compiler does
\emph{not} know: it dutifully computes both \(\sin^2\) and \(\cos^2\) in six
instructions rather than returning \(1\). The compiler is a lowering, not a
simplifier---simplify first, then compile.}
\pair{08-compilation/diagnostics/3}{Array reductions compile too:
\(\mathtt{Total}[v] + \mathtt{Max}[v]\) over a real vector is twelve
instructions, all machine code.}
\end{codepairs}

\usagebox{CompileDiagnostics}
\usagebox{CompilePrint}

\section{How fast? Compiled against the tree}
\label{sec:compile-speed}

Numbers make the case. Here is one polynomial body run over a million points two
ways: once as a \texttt{CompiledFunction}, and once as an ordinary definition
\mcode{treePoly[x\_] := \dots} that the evaluator must rewrite node by node. We
time each with \B{AbsoluteTiming}, which reports elapsed wall-clock seconds, and
take the ratio:

\mtranscript{08-compilation/speed-scalar}

\noindent The compiled loop is more than ten times faster here, for identical
output. (The exact ratio is re-measured every time this book is built and depends
on the machine; it is the \emph{order} of the thing that is stable.)

\calloutnote{Pitfall}{Measure with \B{AbsoluteTiming} (wall-clock), not \B{Timing}
(CPU seconds \emph{summed over threads}): a compiled array kernel that uses
several cores can report a \mcode{Timing} larger than the wall-clock time it
actually took, which reads backwards. And time distinct expressions, or discard a
warm-up run: the first evaluation of anything pays one-off costs (compilation
itself, cold caches) that a fair comparison should not count.}

How does that hold up against a mature numerical stack? The table below pits
Mathilda's compiled kernels against Python---vectorised NumPy where a vectorised
form exists, a plain CPython loop where one does not\calloutnote{Performance}{From
\texttt{docs/design/performance.md}, measured 2026-07-30 on an Intel Core
i9-9880H with NumPy 2.4.4 / SciPy 1.17.1 on CPython 3.11 (both linked against
Apple Accelerate). Lennard-Jones, Mandelbrot and Monte~Carlo are vectorised NumPy
and are genuine library comparisons; the logistic map has no vectorised form
(numba was not installed), so its row is ``what a Python user without a JIT
measures,'' not a library result.}:

\begin{center}
\begin{tabular}{lrr}
\hline
Benchmark & Mathilda & Python / NumPy \\
\hline
Lennard-Jones energy, 1452 bodies (all pairs) & 133.19 ms & 91.25 ms \\
Mandelbrot, \(800\times800\), 100 iterations & 769.23 ms & 980.20 ms \\
Monte Carlo \(\pi\), \(10^{7}\) samples (vectorised) & 162.36 ms & 204.53 ms \\
Logistic map, \(10^{7}\) iterations & 183.76 ms & 460.57 ms \\
\hline
\end{tabular}
\end{center}

\noindent Mathilda's compiled code wins three of the four rows and is close on the
fourth. It is worth being precise about that fourth: Lennard-Jones is where
NumPy's vectorised all-pairs form, running a hand-tuned C loop under one Python
call, comes out ahead---91~ms against 133. That is the honest shape of it. Where
the work is a genuine scalar recurrence with no array form, such as the logistic
map, the gap the other way is largest, because that is exactly the case a
vectorised library cannot help with and a bytecode VM can.

\section{The compilable subset is a cliff}
\label{sec:compile-cliff}

The compiler does not handle every function---and the way it fails is the single
most important operational fact in this chapter. The compilable
subset\index{compilation!compilable subset} is a \emph{cliff, not a slope}: if
\emph{one} head anywhere in the body cannot be lowered, the \emph{whole} body
falls back to the interpreter. Not the offending sub-expression---all of it. And
because the fallback is silent and gives the right answer, nothing tells you it
happened except the clock.

\B{CompileDiagnostics} is how you see the edge of the cliff. Give it the argument
types and a body, and it tells you whether the body compiles and, if not, the
exact innermost sub-expression that stopped it:

\begin{codepairs}
\pair{08-compilation/cliff/1}{\(\mathtt{Sin}[x] + \mathtt{Gamma}[x]\) compiles:
both heads have machine kernels.}
\pair{08-compilation/cliff/2}{Replace \texttt{Gamma} with \B{BarnesG}, which has
no machine kernel, and the whole body is rejected---the report names
\texttt{BarnesG[x]} as the culprit even though \(\sin\) beside it was perfectly
compilable.}
\pair{08-compilation/cliff/3}{The same over arrays: \B{Total} lowers, but
\B{Riffle} does not, so \(\mathtt{Total}[\mathtt{Riffle}[v, 0.]]\) does not
compile at all.}
\pair{08-compilation/cliff/4}{And an integer body: \B{MoebiusMu} has no machine
lowering, so it sinks the sum with \B{EulerPhi} regardless of the latter.}
\end{codepairs}

The lesson is a workflow. When a numeric loop is slower than you expected, do not
guess---ask \mcode{CompileDiagnostics} whether the body even compiled, fix or
remove the one head it names, and ask again. The reward for staying on the subset
is total; the penalty for stepping off it, by one head, is also total.

\calloutnote{Under the hood}{Which heads are on the subset is not a fixed list but
a moving frontier, and it is guarded. A build-time audit,
\texttt{make check-compile-coverage}, cross-checks every head that has a fast
numeric path in the interpreter against whether it also compiles, and fails the
build if a newly-added numeric builtin forgets its compiled lowering---so the
cliff can only ever recede. When an auto-compiled builtin
(\S\ref{sec:compile-auto}) falls back, setting the environment variable
\texttt{MATHILDA\_COMPILE\_DIAG=1} makes it name the head that forced it, on
standard error.}

\section{Auto-compilation}
\label{sec:compile-auto}

Most of the time the compiler works for you without being asked. When you write a
\B{Table} over a real range, or a \B{Plot}, or an \B{NIntegrate}, or a \B{Nest}
with a machine-number seed, Mathilda tries to compile the body once and run the
compiled form at every sample point or iteration. This is \emph{auto-compilation},
and it is why numeric work in Mathilda is fast without a single call to
\mcode{Compile} in sight.

The one switch that governs it is \B{\$AutoCompilation}\index{auto-compilation}.
It is \texttt{True} by default; set it to \texttt{False} and every such body goes
through the interpreter instead. Because a compiled body is contracted to give the
interpreter's answer, the switch changes speed and \emph{nothing else}---which is
exactly what makes it useful for measuring the effect. The session below turns it
off and on around the same \B{Table}, confirms the result is byte-for-byte the
same either way, and then times a million-step logistic-map iteration each way:

\mtranscript{08-compilation/autocompile}

\noindent The two tables (\texttt{r1} and \texttt{r2}) are equal: the switch is a
pure performance control, never a semantic one. The final ratio---dozens to
one---is the interpreter tax that auto-compilation removes from the
\(3.5\,x(1-x)\) recurrence. (As before, the exact figure is re-measured on each
build; its order is the point.)

Two mechanisms hide behind the one switch. An \emph{adapter} compiles a held body
once and reuses it across many sample points---this is what \B{Table}, \B{Plot},
\B{Plot3D}, \B{NIntegrate}, \B{NSum}, \B{Sum}, \B{Product} and \B{FindRoot} use.
A separate \emph{loop compiler} handles the bodies of \B{Do}, \B{While},
\B{Nest}, \B{Fold}, \B{Map} and their kin. Both are gated by the same
\mcode{\$AutoCompilation} flag, and both are conservative in the same way: an
exact-integer \mcode{Table} iterator such as \mcode{\{i, 1, n\}} is left to the
interpreter, because exactness must be preserved, while an inexact one like
\mcode{\{i, 1., n\}} is fair game.

\calloutnote{Pitfall}{Auto-compilation never wraps a machine integer
(\S\ref{sec:compile-precision}): you did not ask it to compile anything, so it may
not turn a transparent speed-up into a silently different answer. If an
auto-compiled integer computation would overflow, it abandons the compiled run and
lets the interpreter promote to a bignum---the safe default is not negotiable
behind your back.}

\section{Arrays, packing, and fusion}
\label{sec:compile-arrays}

A parameter with a rank compiles to code that operates on a whole array in one
call. Declare \mcode{\{u, \_Real, 1\}} and the body is compiled to run over the
backing buffer of a packed array (\S\ref{sec:ds-packed}), start to finish, with no
per-element expression ever built:

\mtranscript{08-compilation/arrays}

\noindent The result of a million-element call is itself a packed array
(\B{NDArrayQ} confirms it), ready to feed the next operation without unpacking.

What makes the array path fast is worth stating precisely, because it is not what
one first assumes. Look again at the disassembly of exactly this body:

\compileprint{08-compilation/array}

\noindent The element-wise chain \(u^2 + 1\) was \emph{fused}\index{compilation!fusion}
into a single strip-mined pass: one parallel loop (\texttt{APAR}) walks the buffer
in 64-wide tiles, and within each tile the square and the add happen on tile
registers (\texttt{VPOWI\_R}, \texttt{VADD\_R}) before the result is written back
(\texttt{VSTORE\_R}). There is no intermediate array holding \(u^2\). Array speed
in Mathilda does not come from ``removing interpretation''---the packed
interpreter is already fast, because it too calls machine kernels
(\S\ref{sec:ds-packed}). It comes from \emph{removing per-element dispatch and
intermediate buffers}, which is why compiling a length-16 vector barely helps
while compiling a long one, or a longer chain of operations, helps a great deal.

\calloutnote{Performance}{This is the same lesson the packed-array chapter
teaches from the other side. Packing decides whether the data is a dense buffer at
all; compilation, and the fusion it enables, decides how few passes over that
buffer the arithmetic takes. The two compose: a compiled body over packed
arguments is the fast path, and it is the one \B{Table}, \B{Plot} and the numeric
solvers of Chapter~\ref{ch:datastructures}'s sequel put you on automatically.}

\section{Machine integers and arbitrary precision}
\label{sec:compile-precision}

By default the compiler's integers are machine \texttt{int64}s and its reals are
machine \texttt{double}s---that is the point, and it is what makes the code fast.
But two questions follow immediately: what happens when a machine integer
overflows, and what if machine precision is not enough? Mathilda has a careful
answer to each.

\begin{codepairs}
\pair{08-compilation/precision/1}{\(3000000^3\) is far past the \texttt{int64}
range. The compiled code detects the overflow, abandons the call, and the
interpreter re-runs it and promotes to a bignum \dots}
\pair{08-compilation/precision/2}{\dots giving exactly the interpreter's answer.
Overflow is caught, never wrapped: the compiled result is always the true one.}
\pair{08-compilation/precision/3}{For more digits than a \texttt{double} holds,
\mcode{WorkingPrecision -> n} compiles the arithmetic in extended precision (MPFR)
at \(n\) digits---here \(\sqrt{2}\) to forty \dots}
\pair{08-compilation/precision/4}{\dots identical to what the interpreter computes
at the same precision.}
\pair{08-compilation/precision/5}{The diagnostics confirm the result type is now
\texttt{MPFRReal}, not \texttt{Real}.}
\pair{08-compilation/precision/6}{Alternatively \mcode{"BigIntegers" -> True}
keeps integer arithmetic exact (GMP) \emph{inside} the compiled code, so the cube
stays compiled and correct rather than bailing out on overflow \dots}
\pair{08-compilation/precision/7}{\dots with result type \texttt{BigInteger}.}
\end{codepairs}

There is a third mode, for the rare case where you know your integers stay within
range and want to skip even the overflow check: \mcode{RuntimeOptions -> "Speed"}.
It removes the per-operation check, so an overflowing result is left wrapped
modulo \(2^{64}\) instead of being promoted:

\begin{codepairs}
\pair{08-compilation/precision/8}{The same cube, now in \texttt{"Speed"} mode: no
promotion, so the answer is the low 64 bits of the true result---a different, and
here plainly wrong, number.}
\pair{08-compilation/precision/9}{Which is exactly \(3000000^3 \bmod 2^{64}\): the
machine word has simply wrapped.}
\end{codepairs}

\noindent This is the one place in the chapter where the compiled answer is
\emph{not} contracted to the true one, and it is off by default for that reason.
Reach for it only when you have measured that the check matters and proven that
your values cannot overflow.

\calloutnote{Under the hood}{The extended-precision modes are deliberately narrow:
\mcode{WorkingPrecision} and \mcode{"BigIntegers"} cover straight-line arithmetic
and the elementary functions on scalars---enough for the sample-point bodies the
numeric solvers of \S\ref{sec:numcalc} compile at high precision---but not arrays
or control flow, which stay on the machine path. The machine path itself is
untouched by these options: switch them off and the bytecode is identical to the
last byte.}

\section{When to compile}
\label{sec:compile-when}

The rule of thumb is short. If a numeric expression is evaluated many times---in a
loop, over an array, at every abscissa of an integration---compilation is a large,
free win, and most of the time (\S\ref{sec:compile-auto}) you get it without
asking. Reach for \mcode{Compile} explicitly when you want to build the fast
function once and reuse it, when you want to hand it to \B{Map} or a solver, or
when you simply want to see with \mcode{CompileDiagnostics} and \mcode{CompilePrint}
what the machine is doing. Do \emph{not} compile a symbolic computation, a
one-shot evaluation, or a body full of heads off the subset---there the fallback
gives you nothing but the overhead of trying.

The threads pulled together here run through the rest of the book. The packed
arrays the compiler operates on are Chapter~\ref{ch:datastructures}; the numeric
solvers that auto-compile their integrands and right-hand sides are
\S\ref{sec:numcalc}; and the evaluator whose generality the compiler trades away
for speed is the subject of the internals to come.

\chapter{Data I/O}
\label{ch:dataio}

Every computation so far has begun and ended inside the session: you typed an
expression, Mathilda evaluated it, and the result scrolled past. This chapter is
about the doorway to the world outside that session---reading numbers a
colleague measured and left in a file, saving a result so the next run can pick
it up, turning a table on disk into a list you can compute with. It is the least
mathematical chapter in the book and, for day-to-day work, one of the most
useful.\index{file I/O}

A word of honesty before we begin. Input and output is a \emph{young} part of
Mathilda, and this chapter is deliberately narrow: it documents what the system
actually does today and is explicit about what it does not. What exists is a
genuine, Mathematica-shaped \emph{stream layer}---you open a file, get a stream
object, and read or write through it---together with \B{ReadList} for reading a
data file field by field, \B{Get} and \B{Put} for saving and reloading whole
Mathilda expressions, and a handful of pure path utilities. What does \emph{not}
yet exist is just as important to know: there is no general \B{Import}/\B{Export}
of data formats (those two are confined to raster images and graphics,
\S\ref{sec:io-importexport}), no binary read/write, and no file-management verbs
such as copying or deleting. Where a gap matters, this chapter names it rather
than papering over it.

Because every example in this book is re-run from scratch at build time (this
chapter is no exception), each one here is self-contained: it writes a
file under \texttt{/tmp} and then reads that same file back, so nothing depends
on data that happens to be lying around. That discipline is worth adopting in
your own scripts, and it makes the round trip---write, then read---the natural
place to start.

\section{The round trip: writing and reading a file}
\label{sec:io-roundtrip}

The smallest complete story in file I/O is to write something and read it back.
Three functions open the door. \B{OpenWrite} opens a file for writing and hands
back a stream object; \B{WriteString} sends raw text to that stream; and \B{Close}
finishes, flushing and releasing the file. To read the data back as numbers,
\mcode{ReadList} does the whole job in one call.

\begin{codepairs}
\pair{09-data-io/roundtrip/1}{\mcode{OpenWrite} truncates the file and returns an
\B{OutputStream} object. The integer is an internal handle---the first stream
opened this session, so \(1\).}
\pair{09-data-io/roundtrip/2}{\mcode{WriteString} sends two lines of text to the
stream; the trailing semicolon suppresses its (empty) result.}
\pair{09-data-io/roundtrip/3}{\mcode{Close} flushes and releases the file,
returning its name to confirm the stream is shut.}
\pair{09-data-io/roundtrip/4}{\B{FileExistsQ} confirms the file is really on
disk.}
\pair{09-data-io/roundtrip/5}{\mcode{ReadList} with the read type \texttt{Number}
scans the whole file and returns its six numbers as a flat list.}
\end{codepairs}

Two things are worth pausing on. First, the newline in the string
\texttt{"88 91 79\textbackslash n95 82 100\textbackslash n"} is a real newline:
Mathilda's string lexer honors the C escapes \mcode{\textbackslash n} and
\mcode{\textbackslash t}, so you can lay out a file's lines directly in a string
literal.\index{string!escape sequences} Second, \mcode{ReadList} did not care that
the six numbers were split across two lines---read as \texttt{Number}, whitespace
and line breaks alike simply separate one token from the next. Where the numbers
sit in the file is a formatting detail; what you get back is the data.

\calloutnote{Under the hood}{A stream reaches the language as an inert object,
\mcode{OutputStream}\texttt{["name", }\emph{id}\texttt{]} or
\B{InputStream}\texttt{["name", }\emph{id}\texttt{]}, whose integer is a handle
into a process-global registry (\texttt{src/io/streams.c}). An output stream
holds an open \texttt{FILE*}; an input stream slurps the whole file into a buffer
with a moving read point. The registry is freed by \mcode{Close} or, for anything
left open, by an \texttt{atexit} hook at program exit, so a session that forgets
to close a stream still leaks nothing under valgrind.}

\usagebox{OpenWrite}

\section{Type-directed reading}
\label{sec:io-readtypes}

The power of \mcode{ReadList} is that you tell it \emph{what kind} of object to read,
and it interprets the bytes accordingly. The same file can be read as numbers, as
words, as whole lines, or as raw bytes---you choose by naming a \emph{read type}.
Take a small table with integers, decimals, and scientific-notation tokens:

\begin{codepairs}
\pair{09-data-io/readtypes/4}{\texttt{Number} reads each token as a number:
integers stay exact, and \texttt{2e5}/\texttt{1.5e-3} become reals.}
\pair{09-data-io/readtypes/5}{\texttt{Real} forces every token to an approximate
number, so even \(1\) comes back as \(1.0\).}
\pair{09-data-io/readtypes/6}{\texttt{Word} keeps each token as a string, doing no
numeric parsing at all.}
\pair{09-data-io/readtypes/7}{\texttt{String} reads a whole line at a time---three
lines, three strings.}
\pair{09-data-io/readtypes/8}{\texttt{Byte} reads raw byte codes; the first five
bytes of \texttt{"1 2 3"} are the ASCII codes \(49, 32, 50, 32, 51\).}
\pair{09-data-io/readtypes/9}{\texttt{Character} reads one-character strings,
byte by byte.}
\pair{09-data-io/readtypes/10}{A third argument caps the read: only the first
three numbers.}
\end{codepairs}

So the same bytes on disk become, depending on the type you ask for, seven
numbers, seven words, three whole lines, or a stream of raw byte codes. The read
type is the lens.\index{ReadList@\mcode{ReadList}!read types}

A \emph{list} of types reads a heterogeneous record: one object of each type per
pass, grouped into a sublist, repeated to end of file. This is exactly how you
read a columnar file whose fields have different types---a name and an age, say:

\begin{codepairs}
\pair{09-data-io/readtypes/14}{\texttt{\{Word, Number\}} reads a string then a
number on each pass, giving one \(\{\text{name}, \text{age}\}\) pair per line.}
\pair{09-data-io/readtypes/18}{When end of file arrives partway through a
pass---here after the lone \(3\)---the unread slot is filled with \B{EndOfFile}
rather than dropped, so the record structure stays intact.}
\end{codepairs}

That \mcode{EndOfFile} in the last pair is not an error; it is how the reader tells
you a record was cut short, so a ragged file never silently loses its shape.

\calloutnote{Theory}{The read types divide cleanly by what counts as a boundary.
\texttt{Byte} and \texttt{Character} consume exactly one byte and never look for a
separator. \texttt{Word}, \texttt{Number}, and \texttt{Real} scan up to the next
\emph{word separator} (whitespace by default). \texttt{Record} and \texttt{String}
scan up to the next \emph{record separator} (a line break). \texttt{Expression}
reads one complete Mathilda expression---the type that makes \mcode{Get} and
\mcode{ReadList} interchangeable for saved results (\S\ref{sec:io-putget}).}

\usagebox{ReadList}

\section{Separators and delimited data}
\label{sec:io-separators}

The word and record separators are not fixed; you set them with options, and that
is what lets \mcode{ReadList} read comma-separated values, tab-separated files, and
other delimited formats. \texttt{WordSeparators} names the strings that break one
token from the next:

\begin{codepairs}
\pair{09-data-io/separators/4}{With \texttt{WordSeparators -> \{","\}}, the comma
becomes the delimiter and a CSV of numbers reads as a flat list.}
\pair{09-data-io/separators/5}{Pair the comma delimiter with a type list to read
the CSV row by row into sublists---a small numeric matrix.}
\end{codepairs}

The complementary option, \texttt{TokenWords}, forces a given string to be split
off as its own word even when nothing surrounds it---useful when an operator or
punctuation mark is glued to its neighbors:

\begin{codepairs}
\pair{09-data-io/separators/9}{Read as plain \texttt{Word}s, \texttt{"a+b"} is a
single token---there is no separator inside it.}
\pair{09-data-io/separators/10}{\texttt{TokenWords -> \{"+"\}} makes the plus sign
a word in its own right, splitting the run into three.}
\end{codepairs}

Alongside these, \texttt{RecordSeparators} sets what ends a line (the default is
the usual trio of \texttt{"\textbackslash r\textbackslash n"},
\texttt{"\textbackslash n"}, and \texttt{"\textbackslash r"}), and the boolean
options \texttt{NullWords} and \texttt{NullRecords} decide whether empty fields
between adjacent separators are kept or skipped---the difference between reading
\texttt{a,,b} as two fields or three.\index{separator!word}\index{separator!record}
Between them, these options cover the common shapes of plain-text tabular
data.

\calloutnote{Pitfall}{Mathilda has no dedicated CSV importer---a comma-delimited
file is just a file you hand the right \texttt{WordSeparators}. That is flexible,
but it also means the quoting and escaping conventions of the CSV standard
(quoted fields containing commas, doubled quotes) are \emph{not} understood: a
field is whatever lies between separators, quotes and all. For clean numeric
tables this is exactly right; for messy real-world CSV with embedded commas, it
is a sharp edge to know about.}

\section{The stream as a cursor}
\label{sec:io-streams}

\mcode{ReadList} reads a whole file in one gulp. When you want to read a file
\emph{incrementally}---a header, then a body sized by what the header said, or a
record at a time in a loop---you open the file yourself and read through the
stream. \B{OpenRead} returns an \mcode{InputStream} that carries a persistent
\emph{current point}, and each \B{Read} advances it, so successive reads return
successive objects.

\begin{codepairs}
\pair{09-data-io/streams/4}{\mcode{OpenRead} returns an \mcode{InputStream}. Its handle
is \(2\): the write stream that created the file was handle \(1\).}
\pair{09-data-io/streams/5}{\mcode{Read} takes one object and moves the point past it.}
\pair{09-data-io/streams/6}{The next \mcode{Read} continues from where the last one
stopped---the cursor persists.}
\pair{09-data-io/streams/7}{A read type may itself be a list: this reads the next
two numbers as a pair.}
\pair{09-data-io/streams/8}{\B{StreamPosition} reports the current point as a byte
offset into the file.}
\pair{09-data-io/streams/9}{\B{Streams} lists every stream still open.}
\pair{09-data-io/streams/10}{\B{SetStreamPosition} moves the point---here back to
the start, returning the new offset.}
\pair{09-data-io/streams/11}{So the very next \mcode{Read} yields the first number
again.}
\end{codepairs}

Positioning turns a stream into random-access storage: \mcode{StreamPosition} tells
you where you are, \mcode{SetStreamPosition} takes you anywhere in the file (and
\texttt{SetStreamPosition[}stream\texttt{, Infinity]} jumps to the end). Because a
type structure can nest, a single \mcode{Read} can pull in a whole shaped object---a
matrix, for instance, read as a list of rows:

\begin{codepairs}
\pair{09-data-io/streams/17}{A nested type structure fills depth-first, so this
reads the file's four numbers into a \(2\times 2\) matrix in one call.}
\pair{09-data-io/streams/21}{Once the point is past the last object, \mcode{Read}
returns \mcode{EndOfFile}---the sentinel a read loop watches for to stop.}
\end{codepairs}

\calloutnote{Under the hood}{\mcode{ReadList} is not a separate engine---it is literally
a loop of the same \texttt{read\_structure()} routine that \mcode{Read} calls once
(\texttt{src/io/read.c}), run to end of file. That is why the two share their read
types and every separator option to the letter: there is one reader, exposed two
ways. The \texttt{Expression} leaf is read \emph{unevaluated} and handed to the
ordinary evaluator, which is why \texttt{Read[s, Expression]} evaluates what it
read while \texttt{Read[s, Hold[Expression]]} keeps it held.}

\section{Persisting expressions: Put, Get, and Write}
\label{sec:io-putget}

The reading so far has treated a file as \emph{data}---numbers and words. But a
file can also hold Mathilda \emph{expressions}, written in the same input syntax
you type at the prompt, and this is how you save a result and reload it in a later
session. \B{Write} sends an expression to a stream in input form; \mcode{Put} (and its
operator \mcode{>>}) writes to a named file; \mcode{Get} reads a file back, evaluating
as it goes.

First, the distinction between the two ways of writing. \mcode{Write} renders each
argument in input form and---crucially---\emph{evaluates} it first, whereas
\mcode{WriteString} writes its text raw:

\begin{codepairs}
\pair{09-data-io/putget/6}{Reading the file back as lines shows what was written:
\mcode{Write}\texttt{[s, 2 + 2]} evaluated to \texttt{4} before writing, \(x^2+1\) was
written in input form, and \mcode{WriteString}'s text went out verbatim.}
\end{codepairs}

For saving results, the \mcode{>>} operator and \mcode{Get} are the everyday tools.
\mcode{>>} is shorthand for \mcode{Put}: it writes one expression to a file, replacing
whatever was there. \mcode{>>>} is \B{PutAppend}: it adds to the file instead.

\begin{codepairs}
\pair{09-data-io/putget/7}{\texttt{FactorInteger[40320] >> file} writes the
factorization of \(8!\) to a file (the result is suppressed).}
\pair{09-data-io/putget/8}{\mcode{Get} reads the file, evaluating each expression, and
returns the value of the last---here the one factorization we wrote.}
\pair{09-data-io/putget/11}{After two \mcode{>>>} appends, \mcode{Get} returns only
the \emph{last} expression in the file: the factorization of \(100\).}
\pair{09-data-io/putget/12}{\mcode{ReadList} reads the same file as
\texttt{Expression}s and \emph{collects} all three, rather than keeping only the
last.}
\end{codepairs}

That contrast between pairs is the whole design. \mcode{Get} is for \emph{loading}: it
runs a file as a script and gives you the final value, exactly as the REPL
bootstrap loads Mathilda's own startup definitions. \mcode{ReadList} is for
\emph{collecting}: it gathers every expression into a list. \mcode{Put} writes several
expressions at once, and they round-trip through \mcode{ReadList} unchanged:

\begin{codepairs}
\pair{09-data-io/putget/13}{\mcode{Put} writes two expressions, one per line.}
\pair{09-data-io/putget/14}{\mcode{ReadList} reads them straight back---a symbolic
result saved to disk and recovered intact.}
\end{codepairs}

\calloutnote{Under the hood}{\mcode{Get} and \mcode{Put} predate the stream layer and live
in \texttt{src/readwrite.c}; \mcode{Get} shares its file-reading core with
\texttt{LoadModule}, the same machinery that loads \texttt{src/internal/init.m} at
startup. Because \mcode{Put} renders with the standard printer, a saved file is
ordinary Mathilda source---you can open it in an editor, and \mcode{Get} it back like
any script.}

\usagebox{Get}

\section{File names and existence}
\label{sec:io-filenames}

The last group of functions is the plumbing around the files themselves: testing
whether a path exists and taking a file name apart. \mcode{FileExistsQ} is the one
that touches the disk; the rest are pure string operations on the path, and never
open anything.

\begin{codepairs}
\pair{09-data-io/filenames/4}{\mcode{FileExistsQ} is \texttt{True} for a file we just
wrote\ldots}
\pair{09-data-io/filenames/5}{\ldots and \texttt{False} for one that is not there.}
\pair{09-data-io/filenames/6}{\B{FileExtension} returns the trailing extension,
without the dot.}
\pair{09-data-io/filenames/7}{\B{FileBaseName} returns the leaf with that one
extension stripped---note it splits off only the last suffix.}
\pair{09-data-io/filenames/10}{\B{FileNameSplit} breaks a path into its
components; the leading \texttt{""} marks an absolute path.}
\pair{09-data-io/filenames/11}{\B{FileNameJoin} is the inverse, assembling
components with the OS separator.}
\pair{09-data-io/filenames/12}{Split then join round-trips to a canonical form of
the original path.}
\end{codepairs}

These are string surgery, portable and side-effect-free---\mcode{FileNameJoin} and
\mcode{FileNameSplit} even take an \texttt{OperatingSystem} option so you can
manipulate Windows paths from Unix or vice versa.\index{file path!components} One
more console convenience rounds out the set: \B{FilePrint} streams the raw bytes
of a file straight to your terminal, the equivalent of the shell's
\texttt{cat}, with optional line ranges (\texttt{FilePrint["f", 1;;10]}). It
writes as a side effect and returns nothing, which is precisely why it does not
appear in the verified transcripts of this chapter---there is no value to
capture. Throughout, we have shown a file's contents by reading them back as data
with \mcode{ReadList}; \mcode{FilePrint} is the quick look when you are at the prompt.

\section{Import and Export: images and graphics}
\label{sec:io-importexport}

A reader coming from Mathematica will expect \mcode{Import} and \mcode{Export} to be the
universal front door for data---CSV, JSON, spreadsheets, and a hundred other
formats. In Mathilda today they are far narrower, and it is important to be plain
about it: \mcode{Import} and \mcode{Export} handle \emph{raster images and graphics
only}. An image round-trips cleanly through a file---

\begin{codepairs}
\pair{09-data-io/importexport/2}{\mcode{Export} writes an image to a PNG and returns
the file name.}
\pair{09-data-io/importexport/3}{\mcode{Import} reads it back; the dimensions survive
the round trip.}
\end{codepairs}

---but a data file is not something they claim, and they say so by declining
rather than guessing:

\begin{codepairs}
\pair{09-data-io/importexport/7}{Handed a plain text file, \mcode{Import} stays
unevaluated---its head comes back as \texttt{Import}, the sign of a call left
alone, not a \mcode{\$Failed}.}
\pair{09-data-io/importexport/8}{Likewise \mcode{Export} declines to write a bare list
to a \texttt{.csv}: there is no data-format writer behind it yet.}
\end{codepairs}

The image and graphics side of \mcode{Import}/\mcode{Export} is real and useful---PNG,
JPEG, BMP, and a resolution-independent vector PDF for plots---and it belongs to
the graphics story, so the graphics chapter (Chapter~\ref{ch:graphics}) develops
it in full. For structured \emph{data}, the honest advice is to reach for the
tools of this chapter instead: \mcode{ReadList} with the right separators for tabular
text, and \mcode{Put}/\mcode{Get} for anything expressible as Mathilda expressions. A
general format layer is future work, not a present feature.

\calloutnote{Pitfall}{The difference between an unevaluated result and
\mcode{\$Failed} is a real signal, not a quirk. \mcode{\$Failed} means ``I tried
and could not''---a missing file, a malformed number. A call that comes back with
its own head, like \texttt{Import[\ldots]} above, means ``I do not handle this at
all,'' leaving the expression untouched so a rule of your own could still act on
it. Reading that distinction saves you from debugging a failure that never
happened.}

\subsection*{Where this connects}

Data I/O is the seam between Mathilda and everything outside it, and its two
halves lead in different directions. The \emph{data} side---\mcode{ReadList} and the
stream layer---feeds the rest of the book: a file of measurements read as a packed
list of numbers is exactly the input the statistics of \S\ref{sec:statistics}
summarize and the array engine of Chapter~\ref{ch:datastructures} accelerates,
and reading a matrix with a nested \mcode{Read} structure produces the same lists that
the linear algebra of \S\ref{sec:linalg} operates on. The \emph{expression}
side---\mcode{Put} and \mcode{Get}---is how a symbolic result outlives its session and how
Mathilda loads its own definitions at startup, a thread the internals chapter
(Chapter~\ref{ch:internals}) picks up. The image and graphics forms of \mcode{Import}
and \mcode{Export} are taken up in Chapter~\ref{ch:graphics}. And for the exhaustive
options of any function named here---every read type, every separator, every
edge case of a path utility---its reference page is one \texttt{?Name} away at the
prompt.

\chapter{The Internals of Mathilda}
\label{ch:internals}

Everything in this book so far has looked at Mathilda from the outside: you type an
expression, and an answer comes back. This chapter turns the machine over and looks
at the works. It is the one place in the book where we get into the C-level detail
of how the system is actually built---how an expression is laid out in memory, what
the evaluator does on every step and in exactly what order, how the pattern matcher
searches, how memory is owned and reclaimed, and how the parser, printer and the
interactive loop bracket the whole thing.

Appendix~\ref{app:architecture} is the aerial photograph: the pipeline in one
figure, each subsystem named in a paragraph. This chapter is the walk on the
ground. If you have not read that appendix, Figure~\ref{fig:pipeline} there is worth
a glance first---it shows how the five components fit together and it is the frame we
now fill in. Where the appendix says ``the evaluator reads the head's attributes,''
this chapter tells you which field it reads, when, and what it costs.

Two conventions for what follows. The \emph{transcripts} are, as everywhere in this
book, real output from the running binary---the internals are observable from the
REPL more often than you might expect, and we lean on that. The \emph{C excerpts},
tinted blue and titled with their source file, are the exception: they are quoted by
hand from the source tree to show a data structure or a control path, abridged for
the page and \emph{not} run by the build the way the transcripts are. Treat them as a
faithful sketch and open the named file for the whole truth.

\section{Expression trees}
\label{sec:int-expr}

The foundational decision, the one from which everything else follows, is that
\emph{every value is the same kind of object}: an \texttt{Expr}. A number, a symbol,
a string, a list, a rule, a pattern, a matrix, a held piece of code---all of them are
\texttt{Expr} nodes, and the whole system is machinery for building, comparing,
rewriting and freeing them. You saw this from the user's side in
Chapter~\ref{ch:datastructures}; here is the node itself.

An \texttt{Expr} is a \emph{tagged union}. A single enum says which kind of value the
node holds, and a union holds the payload for that kind:

\begin{ccode}[src/expr.h (abridged)]
typedef enum {
    EXPR_INTEGER, EXPR_REAL, EXPR_SYMBOL, EXPR_STRING,
    EXPR_FUNCTION, EXPR_BIGINT, EXPR_NDARRAY, EXPR_COMPILED
#ifdef USE_MPFR
    , EXPR_MPFR
#endif
} ExprType;

typedef struct Expr {
    ExprType type;
    unsigned refcount;            /* shared ownership / copy-on-write         */
    uint64_t last_evaluated_at;   /* eval-clock stamp; top bit = GROUND flag  */
    uint64_t hash_cache;          /* memoized structural hash; 0 = unset      */
    union {
        int64_t integer;
        double  real;
        struct { char* name; struct SymbolDef* def; } symbol;
        char*   string;
        struct {
            struct Expr*  head;
            struct Expr** args;
            size_t        arg_count;
            struct AssocIndex*  index;         /* Association key->position   */
            struct SymSetCache* symset_cache;  /* symbol set of the subtree  */
        } function;
        mpz_t bigint;                          /* GMP arbitrary-precision int */
        NDArrayData ndarray;                   /* dense machine-number buffer */
        struct CompiledFunction* compiled;
#ifdef USE_MPFR
        mpfr_t mpfr;                           /* arbitrary-precision real    */
#endif
    } data;
} Expr;
\end{ccode}

The classic six kinds are the ones you meet first---integer, real, symbol, string,
function, and the arbitrary-precision \texttt{EXPR\_BIGINT}. Two more,
\texttt{EXPR\_NDARRAY} and \texttt{EXPR\_COMPILED}, are the dense numeric buffer of
\S\ref{sec:int-packed} and the compiled-code object of
Chapter~\ref{ch:compilation}; a ninth, \texttt{EXPR\_MPFR}, appears only in builds
with extended-precision reals. All of the last three are \emph{atomic} to the
evaluator---it never looks inside them.

A \emph{function} node is the compound case, and it is the interesting one: a head
(itself an \texttt{Expr*}) and a heap array of argument pointers with an explicit
count. There is no separate capacity field---the array is allocated to exactly
\texttt{arg\_count} entries---which matters later, when we pool these arrays by
length (\S\ref{sec:int-memory}). The compound form \(\mathtt{f[x, y]}\) is a function
node whose head is the symbol \texttt{f}; a list is a function node with head
\texttt{List}. You can X-ray any surface syntax into this tree with \B{FullForm}, and
confirm each node's kind with \B{Head} and \B{AtomQ}:

\mtranscript{10-internals/expr-fullform}

Read those as the tree behind the syntax. A sum is \texttt{Plus[a, Times[b, c]]}; a
fraction is a single \texttt{Rational} node; a division is a \texttt{Times} against a
\texttt{Power} with exponent \(-1\); a complex literal is \texttt{Complex[2, 3]}; and
even a pattern, \(\mathtt{x\_}\), is an ordinary expression,
\texttt{Pattern[x, Blank[]]}. \B{Head} of an integer is \texttt{Integer}, of a list
is \texttt{List}: atoms carry a head too, which is why the same tools---\B{Part},
\B{Map}, \B{Replace}---work on everything. \B{AtomQ} draws the one structural line
that matters internally: an atom has no arguments to descend into, a function node
does.

\subsection*{The three metadata words}

Before the union sit three machine words that are not part of the value at all. They
are pure acceleration, invisible to equality and ordering, and understanding them
explains most of what makes Mathilda fast.

\texttt{refcount} is the reference count: shared ownership, discussed in full in
\S\ref{sec:int-memory}. \texttt{last\_evaluated\_at} is a timestamp the evaluator
stamps on a node when it reaches a fixed point, so a later evaluation of the same
node can return instantly (\S\ref{sec:int-eval}); its top bit is a flag for
loop-invariant values. \texttt{hash\_cache} memoizes the node's structural hash the
first time anything asks for it, turning the repeated subtree walks in \texttt{Plus}
grouping, \texttt{Association} indexing and set operations into
\(O(1)\) reads.\calloutnote{Under the hood}{Because these three fields are metadata,
not value, \texttt{expr\_eq}, \texttt{expr\_hash} and \texttt{expr\_compare} all
ignore them---two nodes are equal when their \emph{structure} matches, whatever their
refcounts or stamps. The flip side is a strict rule the whole tree obeys: a cached
hash is only safe because a shared node is never mutated in place (\S\ref{sec:int-memory}),
and any code that does rewrite a node's fields must reset \texttt{hash\_cache} to 0
first. The \texttt{MATHILDA\_HASH\_VERIFY} build recomputes and asserts every cache
hit to catch a missed reset.}

Integers show a fourth theme---representations that adapt. A value begins life as a
machine \texttt{int64}, and is promoted to a GMP bignum only when an operation would
overflow, demoting back when a result fits again. There is only ever ``an integer,''
never a user-visible big-integer mode:

\mtranscript{10-internals/expr-bignum}

\(2^{62}\) is a machine word; \(2^{63}\) has already overflowed into a
\texttt{EXPR\_BIGINT}; \(100!\) is a 158-digit bignum---and \B{Head} says
\texttt{Integer} for all of them. The promotion is a change of internal
representation, not of type, and it is completely invisible from the language.

\section{The evaluator}
\label{sec:int-eval}

The evaluator is the heart of the system, and its governing idea is
\emph{fixed-point rewriting}: \texttt{evaluate(e)} calls \texttt{evaluate\_step(e)}
over and over, each call rewriting the expression by one level, and stops when a step
changes nothing. There is no separate ``simplifier'' pass and no explicit chaining of
rules; composition is a consequence of running the step to a fixed point. You can
watch the process directly with \B{Trace}, which records every intermediate expression:

\mtranscript{10-internals/eval-trace}

A definition and an arithmetic fold each show the shape of it---\texttt{f[3]} becomes
\texttt{3\^{}2} becomes \texttt{9}; the three-term sum collapses to \texttt{9} in one
step. And because a step just rewrites and the loop merely repeats, composition needs
no plumbing at all: define one function in terms of another and the answer falls out
over successive passes.

\mtranscript{10-internals/eval-fixedpoint}

Each element of the \B{Trace} list is one intermediate expression on the way to the
answer: \texttt{g[10]} rewrites to \texttt{h[10]}, then to \texttt{10 + 1}, then to
\texttt{11}, at which point a step produces nothing new and the loop halts. No code in
\texttt{g} or \texttt{h} knows about the other; the loop is the only glue.

\subsection*{One step, in order}

The real content is what \texttt{evaluate\_step} does for a function call
\(\mathtt{f[a_1, \dots, a_n]}\). The appendix gives the eight-line summary; the true
sequence, in \texttt{src/eval.c}, is longer and every stage earns its place.
Figure~\ref{fig:evalstep} lists it in order.

\begin{figure}[h]
\begin{framed}
\small
\begin{verbatim}
   evaluate_step( f[a1, ..., an] ):

    1.  evaluate the head f
    2.  resolve f's SymbolDef once, cache it on the head node;
        read attributes as (base_attributes | attributes)  -- O(1)
    3.  evaluate each argument UNLESS a Hold attribute holds that slot
          (Evaluate[] punches through a hold; atoms skip the recursion)
    4.  reassemble f[...] from the evaluated arguments
    5.  flatten Sequence[...] arguments        (unless SequenceHold)
    6.  strip Nothing from a List
    7.  strip Unevaluated[...] wrappers in non-held slots
   7.5  transparency gate: materialize any packed buffer whose head
          is not packed-aware                  (see 10.9)
    8.  Flat      -- flatten nested f[f[...]]   if ATTR_FLAT
    9.  Listable  -- thread over list arguments if ATTR_LISTABLE
   10.  Orderless -- sort arguments canonically if ATTR_ORDERLESS
   11.  try f's DOWNVALUES  (user/library rules)  -- before the builtin
   12.  try f's NDArray element-wise machine kernel, if registered
   13.  call f's builtin C function, if any
   14.  Set / SetDelayed and the composite-head cases (pure Function, ...)
\end{verbatim}
\end{framed}
\caption{One rewrite level. Steps 8--10 apply the associative, threading and
commutative attributes in that fixed order; the surprise for a newcomer is step 11,
where user \texttt{DownValues} are tried \emph{before} the builtin C code, so a rule
you write can override a primitive.}
\label{fig:evalstep}
\end{figure}

Three things about this sequence are worth dwelling on.

\emph{Attributes are read once, and cheaply.} Step~2 resolves the head symbol to its
definition record and caches the pointer on the head node, so subsequent touches skip
the symbol-table lookup entirely. The attributes themselves are a single field read
and a bitwise OR---\texttt{base\_attributes} (the immutable floor a builtin is born
with) combined with \texttt{attributes} (whatever the user has since set). This
replaced an older design that walked a 140-entry table of name comparisons on every
call. The attribute bits are the small vocabulary that steers the generic loop; the
full set is in Table~\ref{tab:attrs}.

\emph{Holding happens in step~3, bit by bit.} Whether an argument is evaluated is
decided per slot, by three lines:

\begin{ccode}[src/eval.c (the argument loop)]
bool hold = hold_all_complete;                        /* HoldAllComplete: every slot */
if (i == 0 && (attrs & ATTR_HOLDFIRST)) hold = true;  /* HoldFirst: slot 0           */
if (i >  0 && (attrs & ATTR_HOLDREST))  hold = true;  /* HoldRest:  slots 1..n       */
\end{ccode}

\texttt{HoldAll} needs no case of its own: it is simply \texttt{HoldFirst \& HoldRest}
as a combined mask. Everything about non-standard evaluation---\texttt{Table} not
evaluating its iterator, \texttt{If} not evaluating the branch it will not take,
\texttt{Set} not evaluating its left side---reduces to which of these bits the head
carries. That is the whole subject of \S\ref{sec:int-nonstd}.

\emph{DownValues come before the builtin.} Step~11 is ahead of step~13 on purpose: a
rule the user or a library has installed on \texttt{f} is tried before \texttt{f}'s
own C code, so mathematics expressed as a rewrite rule (the integral tables, the
derivative rules of \S\ref{sec:int-boot}) can extend or override a primitive without
touching C.

The attribute bits do all the steering, and you can read a head's bits from the
language with \B{Attributes} and change them with \B{SetAttributes}:

\mtranscript{10-internals/eval-attributes}

\texttt{Plus} carries five attributes at once; \texttt{Sin} only two; \texttt{Table}
is \texttt{HoldAll}; \texttt{Set} is \texttt{HoldFirst}. Watch the effect of each: an
\texttt{Orderless} head sorts \texttt{b + a + c} into canonical order; adding
\texttt{Flat} to \texttt{f} makes \texttt{f[f[a, b], c]} flatten to
\texttt{f[a, b, c]}; adding \texttt{Listable} to \texttt{h} makes it thread over a
list. Nothing about the main loop changed between those heads---only their bits did.

\begin{table}[h]
\centering
\small
\begin{tabular}{lll}
\hline
Attribute & Bit & Effect \\
\hline
\texttt{HoldFirst} & \texttt{1<<0} & do not evaluate argument 0 \\
\texttt{HoldRest} & \texttt{1<<1} & do not evaluate arguments 1..n \\
\texttt{HoldAll} & (both) & do not evaluate any argument \\
\texttt{HoldAllComplete} & \texttt{1<<2} & hold all, and suppress Sequence/Flat/Unevaluated \\
\texttt{Flat} & \texttt{1<<3} & associative: flatten nested same-head calls \\
\texttt{Orderless} & \texttt{1<<4} & commutative: sort arguments canonically \\
\texttt{Listable} & \texttt{1<<5} & thread element-wise over list arguments \\
\texttt{NumericFunction} & \texttt{1<<6} & numeric under \texttt{N}; threads over intervals \\
\texttt{Protected} & \texttt{1<<7} & values and attributes cannot be changed \\
\texttt{OneIdentity} & \texttt{1<<8} & \(\mathtt{f[x]}\equiv\mathtt{x}\) for matching (\S\ref{sec:int-match}) \\
\texttt{SequenceHold} & \texttt{1<<13} & do not flatten \texttt{Sequence} in arguments \\
\texttt{Constant} & \texttt{1<<14} & derivative is zero \\
\hline
\end{tabular}
\caption{The core attribute bits, from \texttt{src/attr.h}. \texttt{Locked},
\texttt{Temporary} and the \texttt{NHold*} family occupy the remaining bits; bit 11,
once \texttt{ReadProtected}, was removed---Mathilda never hides a definition.}
\label{tab:attrs}
\end{table}

\subsection*{Knowing when to stop}

A fixed-point loop needs to know when to quit, and doing that naively---comparing the
whole tree before and after every step---would be quadratic on deep expressions.
Mathilda avoids it with a cheap boolean: \texttt{evaluate\_step} reports whether it
actually rewrote anything, and only when it claims a change does the loop fall back to
a structural comparison. When a step fires no rule, the loop stops without walking the
tree at all.

Stability across \emph{calls} is handled by the \texttt{last\_evaluated\_at} stamp. A
global evaluation clock ticks forward on every definition change---every \texttt{Set},
\texttt{SetDelayed}, \texttt{Clear} or \texttt{SetAttributes}. When
\texttt{evaluate} reaches a fixed point it stamps the result with the current clock;
a later \texttt{evaluate} of a node already carrying the current stamp returns an
inc-ref'd view immediately, skipping the loop and all of \texttt{evaluate\_step}. A
single definition change, by bumping the clock, invalidates every cached evaluation in
the system at once.\calloutnote{Under the hood}{The stamp's top bit is a second,
finer cache. A node built entirely from literals under one of six pure structural
constructors (\texttt{List}, \texttt{Association}, \texttt{Rule},
\texttt{RuleDelayed}, \texttt{Complex}, \texttt{Rational}) is marked \emph{ground}:
its value cannot depend on any symbol's definition, so it survives clock churn and can
be re-validated as a fixed point in \(O(1)\) even after an ordinary assignment. This
is what keeps a loop that rebuilds a constant list on every iteration from
re-evaluating it every time.}

Two limits bound the process, and they are deliberately asymmetric.
\(\mathtt{\$RecursionLimit}\) caps how deep evaluation may nest; it is a real,
user-settable symbol (default 1024, with a floor of 20), and on overflow the offending
expression is returned wrapped in \texttt{Hold} so it cannot immediately re-enter.
\(\mathtt{\$IterationLimit}\), by contrast, is \emph{not} a settable symbol at all: the
number of fixed-point iterations for a single expression is capped by a compile-time
constant (4096). It is an honest simplification---the constant is generous enough that
only a pathological non-terminating rewrite reaches it---but it is worth knowing that,
unlike in some systems, you cannot raise it from the language.

\section{Non-standard evaluation}
\label{sec:int-nonstd}

``Non-standard evaluation'' is the umbrella term for everything that departs from the
default ``evaluate all the arguments, then call the head.'' In many languages this
would need special forms baked into the interpreter. In Mathilda it is one mechanism:
the \texttt{Hold} attributes of \S\ref{sec:int-eval}, step~3. A head that wants to
control its arguments simply carries \texttt{HoldFirst}, \texttt{HoldRest} or
\texttt{HoldAll}, and the generic loop hands it the arguments unevaluated. Everything
in this section is built from that one bit-test.

\usagebox{Hold}

The primitives that manage held code are small and composable:

\mtranscript{10-internals/nonstandard-hold}

\B{Hold} keeps its contents literally---it is nothing but a head carrying
\texttt{HoldAll}, so \texttt{Hold[1 + 1]} never adds. \B{Evaluate} punches a hole in a
hold, forcing just its argument while its siblings stay held. \B{Unevaluated} is the
mirror image: it passes an argument through unevaluated and then dissolves, so
\texttt{Length[Unevaluated[1 + 2 + 3]]} sees the three-element \texttt{Plus} and
answers \texttt{3}. \B{ReleaseHold} strips a wrapper so its contents can finally run.
The distinction between \B{Set} and \B{SetDelayed} is the same mechanism: \texttt{=}
is \texttt{HoldFirst} (it holds the left side but evaluates the right side now), while
\texttt{:=} is \texttt{HoldAll} (it holds both and re-evaluates the right side on every
use). And \B{If}, carrying \texttt{HoldRest}, never touches the branch it does not
take---which is why the untaken \texttt{boom = 1/0} above caused no error and left
\texttt{boom} unassigned.

\subsection*{Scoping without a stack}

The scoping constructs are the richest users of held arguments, and they achieve their
effects in three distinct ways---worth contrasting, because the differences are
observable:

\mtranscript{10-internals/nonstandard-scoping}

\B{Module} implements lexical locals by \emph{alpha-renaming}: each local \texttt{x} is
renamed to a globally fresh symbol \texttt{x\$}\textit{n}, and you can see the counter
advance---the two identical \texttt{Module} calls return \texttt{Hold[x\$1]} and
\texttt{Hold[x\$2]}. \B{Block}, by contrast, is \emph{dynamic} scoping: it does not
rename anything but temporarily saves the global value of \texttt{n}, runs the body,
and restores it---so inside the block \texttt{n} is 2 and after it \texttt{n} is 10
again. \B{With} does neither: it substitutes its constants directly into the body
before evaluation, a purely lexical rewrite. Pure functions
(\(\mathtt{\#^2\, \&}\)) bind their \texttt{Slot} arguments by the same
substitution.\calloutnote{Under the hood}{None of these three touches a call stack of
activation records---Mathilda has no such thing. \texttt{Module} works by renaming and
installing a temporary \texttt{OwnValue}; \texttt{Block} by saving and restoring a
symbol's value cell; \texttt{With} and pure functions by lexical substitution into a
held body. \texttt{Table} and the other iterators are dynamic scoping too, saving the
iterator symbol's value around the loop. The uniform ``everything is an expression''
model means a binding is just a value on a symbol, and scoping is just deciding when to
put it there and when to take it away.}

\section{The pattern matcher}
\label{sec:int-match}

Pattern matching is how Mathilda decides whether a rule applies, and it is the most
combinatorially delicate subsystem in the tree. Chapter~\ref{ch:programming} taught
the pattern \emph{language}---blanks, sequences, tests, conditions; this section is the
engine beneath it, \texttt{src/match.c}. From the language side you drive it with
\B{MatchQ} and its relatives:

\mtranscript{10-internals/matcher-basics}

Everything in that session is the matcher at work: a typed blank checks a head; a
repeated name enforces equality; sequence blanks match runs of arguments; an
\texttt{Orderless} head matches against reordered arguments; a test or condition
gates a match; an \texttt{Optional} supplies a default. We now take the mechanism
apart.

\subsection*{Bindings and backtracking}

Pattern variables are held in a \texttt{MatchEnv}: two parallel arrays, one of names
and one of bound values, with a linear scan for lookup. A handful of variables is the
common case, so a scan beats the overhead of a hash.

\begin{ccode}[src/match.h and match.c]
typedef struct MatchEnv {
    char** symbols;     /* borrowed interned names -- NOT strdup'd     */
    Expr** values;      /* the env owns a copy of each bound value     */
    size_t count;
    size_t capacity;
    bool (*callback)(struct MatchEnv*, void*);  /* ReplaceList: see every match */
    void* callback_data;
} MatchEnv;

/* Backtracking is a single integer. Bindings are only ever appended, so
 * truncating the count back to a checkpoint is a complete, cheap undo. */
static void env_rollback(MatchEnv* env, size_t saved_count) {
    while (env->count > saved_count)
        expr_free(env->values[--env->count]);
}
\end{ccode}

That \texttt{env\_rollback} is the whole backtracking discipline. At every branch
point the matcher records \texttt{env->count}, tries an alternative, and if it fails
truncates the bindings back to the saved count and tries the next. Because a name is
never overwritten---only appended after a check---rolling the count back is an exact
undo. \emph{Non-linear} patterns fall straight out of this: when the matcher is about
to bind a name that already has a value, it does not re-bind; it compares the new
candidate to the existing binding with \texttt{expr\_eq} and fails if they differ.
That single rule is why \texttt{f[x\_, x\_]} matches \texttt{f[a, a]} but not
\texttt{f[a, b]}.

The recursion itself threads a continuation---a \texttt{ParentMatch} chain
representing ``the sibling patterns still to match''---so that a deep sub-match, on
success, can carry on matching the arguments that follow it. There is no explicit work
stack; the C call stack, guarded by the same \(\mathtt{\$RecursionLimit}\) the
evaluator uses, is the stack.

\subsection*{Sequences: the combinatorial heart}

Ordinary structural matching is easy: same head, same arity, recurse pairwise. The
hard part is sequence patterns, \(\mathtt{x\_\_}\) and \(\mathtt{x\_\_\_}\), which
match \emph{runs} of arguments of unknown length. The matcher searches the partitions:
for the first sequence pattern it tries letting it consume \(k\) arguments for each
legal \(k\), and for each choice recurses on the rest. Figure~\ref{fig:backtrack}
shows the search for \texttt{f[x\_\_, y\_\_]} against \texttt{f[1, 2, 3]}.

\begin{figure}[h]
\begin{framed}
\small
\begin{verbatim}
   match  f[x__, y__]   against   f[1, 2, 3]        (each __ needs >= 1 arg)

     x__ = {1}          x__ = {1,2}         x__ = {1,2,3}
       |                   |                    |
     y__ = {2,3}  OK     y__ = {3}   OK       y__ = {}   FAIL
       ^                                        (y__ requires >= 1)
       first success returned
\end{verbatim}
\end{framed}
\caption{Sequence matching is a partition search with backtracking. Each split is
tried in turn; the first that lets every remaining pattern match wins.}
\label{fig:backtrack}
\end{figure}

Two forces make this affordable in practice. The first is a hard short-circuit: when a
sequence pattern is the \emph{last} one in the argument list, it must consume all the
remaining arguments, so there is nothing to search---\(k\) is pinned, not enumerated.
This one rule is what keeps \texttt{Plus[x\_\_, y\_\_]} from hanging on a wide sum. The
second is that for an ordinary (ordered) head the consumed run and the remainder are
handed to the recursion as \emph{slices} of the caller's array---no copying.

\texttt{Orderless} heads are where the real cost lives, because a commutative match may
pair a pattern with arguments in any position, not just a contiguous run. There the
matcher enumerates \(k\)-element \emph{subsets} of the arguments, which is genuinely
combinatorial; it prunes hard (moving fixed elements ahead of variable-length blanks
so the selective patterns filter first), but the worst case is why matching a
complicated pattern against a very large sum can be slow. \texttt{Flat} heads add a
third wrinkle: a single pattern element may consume several arguments, which the
matcher assembles back into an \texttt{f[\dots]} sub-expression to match against.

\subsection*{When the matcher calls the evaluator}

A \mcode{PatternTest} (\(\mathtt{p\,?\,test}\)) and a \mcode{Condition}
(\(\mathtt{p\,/;\,c}\)) are not structural: they are predicates. When the matcher reaches one, it substitutes
the current bindings into the test or condition and \emph{calls
\texttt{evaluate}}---the matcher and the evaluator are mutually recursive. A test that
does not return \texttt{True} fails the match and triggers the same
\texttt{env\_rollback} as a structural mismatch, so the partition search moves on to
the next candidate. This is why \texttt{\_?EvenQ} and \texttt{x\_ /; x > 5} work: the
condition is arbitrary Mathilda code, run mid-match with the tentative bindings in
place.

One attribute is a matcher-only concern: \texttt{OneIdentity}. It is \emph{not}
applied during evaluation (\(\mathtt{f[x]}\) is never silently collapsed to
\texttt{x}); it exists solely so that, when a structural match of \texttt{f[\dots]}
fails, the matcher may retry treating the whole subject as a single argument---the
behaviour that lets a pattern like \(\mathtt{a\_.\, x\_ + b\_.}\) match a bare
\texttt{x}. Once a match succeeds, \texttt{replace\_bindings} substitutes the captured
values into the rule's right-hand side, splicing any \texttt{Sequence} bindings into
place, and returns the rewritten expression.

\section{Rules and the symbol table}
\label{sec:int-symtab}

A rule is stored, retrieved and applied through the symbol table, and the design of
that table is what makes rule dispatch fast even for a symbol with hundreds of
definitions. Every symbol resolves to a \texttt{SymbolDef}: its \texttt{OwnValues}
(bare assignments like \texttt{x = 5}), its \texttt{DownValues} (rules on a head, like
\(\mathtt{f[x\_] := x^2}\)), an optional builtin C function pointer, its attribute
bits, its docstring, and its NDArray kernels. A rule is a \texttt{Rule} node:

\begin{ccode}[src/symtab.h (abridged)]
typedef struct Rule {
    Expr* pattern;
    Expr* replacement;
    int32_t     dispatch_arity;        /* arg count, or -1 for variable arity   */
    const char* first_arg_head_canon;  /* interned head of arg 1, or NULL        */
    int32_t     specificity;           /* larger = more specific; sorts the list */
    struct Rule* next;
} Rule;
\end{ccode}

\subsection*{Specificity, not recency}

\texttt{DownValues} are kept in a linked list, and the order of that list is the order
the matcher tries them---so it decides which of several applicable rules wins. The
ordering is by \emph{specificity}: a more specific pattern (a literal
\texttt{fib[0]}) sorts ahead of a more general one (\texttt{fib[n\_]}), with insertion
order breaking ties. It is a common misconception that new definitions simply go to
the front of the list. They do not. A general rule entered \emph{after} a specific one
still yields to it:

\mtranscript{10-internals/matcher-downvalues}

The Fibonacci base cases \texttt{fib[0]} and \texttt{fib[1]}, though entered after the
recursive rule, are tried first, and \texttt{DownValues[fib]} shows them at the front
of the list. The decisive case is \texttt{r}: the specific \texttt{r[0]} is entered
first, the general \texttt{r[x\_]} second, and yet \texttt{r[0]} still returns
\texttt{special}---specificity beat recency.

\usagebox{DownValues}

Before the matcher is ever invoked, a cheap filter discards rules that cannot possibly
apply. The \texttt{dispatch\_arity} and \texttt{first\_arg\_head\_canon} fields on each
rule are a precomputed dispatch key: a rule whose pattern takes three arguments is
skipped outright for a two-argument call, and a rule whose first argument must have
head \texttt{Plus} is skipped when the input's first argument is an
\texttt{Integer}---an interned-pointer comparison, no matching attempted. Only the
survivors reach the backtracking matcher of \S\ref{sec:int-match}.
\texttt{OwnValues} get an even shorter path: the overwhelmingly common
\texttt{x = value} binding is recognised and returned without building a match
environment at all.

\subsection*{Interning and contexts}

The symbol table is also the string interner. Every symbol name exists exactly once as
a canonical \texttt{const char*}, so two references to \texttt{Sin} share one pointer
and identity tests reduce to pointer equality---which is why the evaluator's hot paths
can write \texttt{head == SYM\_List} instead of a string comparison. Those
\texttt{SYM\_*} constants are the interned pointers for the kernel's own symbols,
filled in once at startup.

Layered over interning is the \B{Context} system: symbols live in namespaces, and
\(\mathtt{\$Context}\) is where new ones are created.

\mtranscript{10-internals/symtab-context}

A fresh symbol is born in \texttt{Global`}; a builtin lives in \texttt{System`};
\B{Begin} switches the current context so that symbols mentioned inside it are
namespaced (\texttt{here} becomes \texttt{MyPkg`here}), and \B{End} restores the
previous context. This is the machinery packages are built on.

The rule \emph{engine}---\B{Replace}, \B{ReplaceAll} (\texttt{/.}) and
\B{ReplaceRepeated} (\texttt{//.})---applies rules without installing them on a symbol:

\mtranscript{10-internals/rules-replace}

Note the semantics the internals enforce. \texttt{/.} rewrites top-down and does
\emph{not} re-descend into a subexpression it has just rewritten---so
\texttt{f[f[x]] /. f[a\_] -> a} yields \texttt{f[x]}, not \texttt{x}. \texttt{//.}
repeats to a fixed point and does reach \texttt{x}. \B{Replace} acts only at the top
level by default, which is why it leaves the list of \texttt{f}'s untouched. And
\B{ReplaceList} returns \emph{every} way the pattern can match---the direct use of the
matcher's success callback, enumerating all six orderings of the commutative sum.

\section{Memory management}
\label{sec:int-memory}

Mathilda manages memory by hand---there is no garbage collector---and the entire
scheme rests on one fact you have already met: \texttt{expr\_copy} is not a deep copy.
It is a reference-count bump.

\begin{ccode}[src/expr.c (abridged)]
Expr* expr_copy(Expr* e) {          /* a logical copy: no tree is duplicated */
    if (!e) return NULL;
    e->refcount++;
    return e;
}

void expr_free(Expr* e) {
    if (!e) return;
    if (e->refcount > 1) { e->refcount--; return; }  /* still shared: keep it */
    /* last reference: release the payload by type -- mpz_clear a bignum, free
     * an NDArray's buffer, recurse into a function node's head and args -- then
     * recycle the node struct. Interned symbol names are never freed. */
}
\end{ccode}

Sharing a subtree costs one increment; freeing decrements and reclaims only on the
last reference. The invariant that makes this safe is that a node with
\texttt{refcount > 1} is \emph{immutable}---any code wanting to rewrite a shared node
first calls \texttt{expr\_unshare}, which returns a private one-level copy (its
children still shared) and resets the metadata caches. This copy-on-write discipline
is what lets the cached hash and evaluation stamp of \S\ref{sec:int-expr} be trusted.

\subsection*{The builtin ownership contract}

Every builtin obeys one contract, and getting it right is the single most important
rule for anyone extending the system. A builtin \emph{takes ownership} of its argument
expression. It returns either a freshly built \texttt{Expr*}---in which case the
evaluator frees the original---or \texttt{NULL} to mean ``not my case,'' in which case
the evaluator keeps the expression and leaves it unevaluated for a rule to handle. That
\texttt{NULL} convention is exactly what lets a symbolic argument flow untouched
through a head that only knows how to handle numbers. A builtin that wants to reuse a
piece of its input must null the slot out before returning it, so the evaluator's
subsequent free does not double-free the reused child.\calloutnote{Under the hood}{Two
free-lists keep allocation cache-hot. Fixed-size \texttt{Expr} node structs are pooled
in a bounded list (the dead node's link is stored in a union field it no longer needs),
and the small argument arrays are pooled in per-length buckets for lengths 1 through 8
---which is why the argument array is sized exactly, never over-allocated. Both pools
are drained by an \texttt{atexit} hook, so a clean run leaves no reachable-but-unfreed
memory and \texttt{valgrind} stays quiet.}

You can watch the footprint from the language. \B{MemoryInUse} and \B{MaxMemoryUsed}
report the process heap in bytes:

\mtranscript{10-internals/memory-inuse}

Building a million-element \B{Range} raises the footprint by about eight
megabytes---a dense \texttt{int64} buffer of \(10^6\) elements
(\S\ref{sec:int-packed}), not a million boxed nodes. \B{Clear} releases the binding,
but the resident figure does not fall: freed memory returns to the process's own pools
and allocator, not immediately to the operating system, so \B{MaxMemoryUsed} records
the high-water mark. The numbers here are re-measured on every build and depend on the
machine; it is their shape that is the lesson.

\section{The parser and the printer}
\label{sec:int-parse}

The parser and printer are the two ends of the pipeline---the translation between the
string a human types and the tree the evaluator rewrites.

The parser (\texttt{src/parse.c}) is a \emph{Pratt}, or precedence-climbing, parser
with no separate lexer: tokens are read inline as parsing proceeds. Every operator
sits on a numeric precedence ladder that runs from 1 to 10000 in round steps, with
generous gaps so a new operator can be slotted between two existing ones without
renumbering. The best way to see the parser's decisions is to X-ray them with
\B{FullForm}, which shows the tree the grouping produced:

\mtranscript{10-internals/parser-fullform}

Every line is a precedence decision made visible. \texttt{a + b*c} groups the
multiplication first; \texttt{a\^{}b\^{}c} associates to the right; subtraction and
division are not primitive at all---\texttt{a - b} is \texttt{Plus[a, Times[-1, b]]}
and \texttt{a/b} is \texttt{Times[a, Power[b, -1]]}, so the evaluator needs rules only
for \texttt{Plus}, \texttt{Times} and \texttt{Power}. A chain of equalities folds into
a single \texttt{Inequality} node; \texttt{\&\&} binds tighter than \texttt{||}; the
prefix \texttt{\#\^{}2 \&} is \texttt{Function[Power[Slot[1], 2]]}; and
\texttt{a[[1]]} is \texttt{Part[a, 1]}. Implicit multiplication---the space in
\texttt{b c}---is injected by the parser as an ordinary \texttt{Times} at the
multiplication rung.

The printer (\texttt{src/print.c}) runs the same ladder in reverse, adding parentheses
exactly where a subexpression's precedence is lower than its context so that the output
re-parses to the same tree:

\mtranscript{10-internals/printer-forms}

\B{InputForm} prints a form suitable for re-reading, \B{FullForm} the raw tree, and
the default \texttt{OutputForm} the familiar two-dimensional-ish notation. The printer
restores \texttt{Times[a, Power[b, -1]]} to \texttt{a/b}, \texttt{Rational[3, 4]} to
\texttt{3/4}, and parenthesises \texttt{(a + b) c} because \texttt{Plus} binds looser
than \texttt{Times}. It is the same precedence table the parser used, which is what
guarantees the round trip.

\section{The read--eval--print loop and session state}
\label{sec:int-repl}

The evaluator is wrapped by the interactive loop (\texttt{src/repl.c}), and it is
worth being precise that the loop is a distinct layer---it does several things
\emph{around} each call to \texttt{evaluate} that the evaluator itself knows nothing
about. For one input line, in order, the loop:

\begin{enumerate}
  \item sets \(\mathtt{\$Line}\) to the current line number;
  \item applies the \(\mathtt{\$PreRead}\) hook to the raw input \emph{text};
  \item parses the text to an \texttt{Expr};
  \item stores it as \texttt{In[\textit{n}]} (a \texttt{DownValue} on \texttt{In}), so
        the literal input can be recalled;
  \item applies the \(\mathtt{\$Pre}\) hook to the parsed expression;
  \item calls \texttt{evaluate};
  \item applies the \(\mathtt{\$Post}\) hook to the result;
  \item stores the result as \texttt{Out[\textit{n}]} (a \texttt{DownValue} on
        \texttt{Out});
  \item applies the \(\mathtt{\$PrePrint}\) hook to a copy for display;
  \item prints \texttt{Out[\textit{n}]= } and the expression.
\end{enumerate}

The history reach-backs are built on step 8: \texttt{\%} parses to
\texttt{Out[-1]}, \texttt{\%\%} to \texttt{Out[-2]}, and \texttt{Out[\textit{n}]}
resolves through the same \texttt{DownValue} store, a negative index counting back from
\(\mathtt{\$Line}\). Because all of this lives in the loop and not in \texttt{evaluate},
it is present only in an interactive session. The proof is to run the same lines
\emph{non-interactively}---which is exactly how every transcript in this book is
produced, through a batch pipe:

\mtranscript{10-internals/repl-history}

Outside the interactive loop there is no \texttt{In}/\texttt{Out} store and no running
\(\mathtt{\$Line}\): \texttt{\%} stays the symbolic \texttt{Out[-1]}, \texttt{Out[1]}
does not resolve, and \(\mathtt{\$Line}\) is just an unbound symbol. Nothing is broken
---this is the same evaluator---but the session bookkeeping simply is not there,
because the loop that maintains it is not running. That is the clearest possible
demonstration that history is a service of the REPL, layered on top of a stateless
evaluator.\calloutnote{Under the hood}{The loop chooses one of three modes at startup.
A script given with \texttt{-file} is evaluated statement by statement with nothing
echoed. A non-terminal standard input (a pipe) selects a line-oriented NDJSON protocol
---the one that drives this book's build. Only a real terminal runs the interactive
loop, with GNU Readline and the completeness-driven multi-line editing described in
\texttt{SPEC.md}~\S11, where pressing Return submits only when brackets, strings and
comments are all balanced.}

\section{Numbers and the packed-array substrate}
\label{sec:int-packed}

Numbers deserve a closer look, because ``number'' is really several internal
representations that the system chooses between automatically. The
\texttt{Head}, though, is uniform:

\mtranscript{10-internals/numbers-tower}

An exact integer is a machine \texttt{int64} promoted to a GMP bignum as needed
(\S\ref{sec:int-expr}); a rational is a \texttt{Rational} of two integers; a machine
real is a \texttt{double}; and at higher precision the arithmetic moves onto MPFR,
which is why \texttt{Precision[N[Pi, 40]]} reports about forty digits rather than the
\texttt{MachinePrecision} of a bare \texttt{N}. Exactness is contagious in the usual
way: \texttt{1/2 + 1/3} stays the exact \texttt{5/6}, but a single machine real, as in
\texttt{1/2 + 0.5}, makes the whole result inexact. This numeric tower is the subject
of \S\ref{sec:numcalc} from the mathematics side; here it is enough to see that one
\texttt{Head} hides several machine types.

The last representation is the one that makes numeric work fast: the packed array. A
list all of whose elements are machine numbers can be stored not as a tree of boxed
nodes but as a single \texttt{EXPR\_NDARRAY}---a dense, flat buffer. Crucially this is
\emph{one node type with two surfaces}. A visible \texttt{NDArray[\dots]} that a user
typed is atomic: its head is \texttt{NDArray}, and \texttt{AtomQ} is true. A
\emph{packed list} is the same buffer wearing the head \texttt{List}: it is
indistinguishable from the nested list it stores---\texttt{ListQ} is true, patterns
match it, it prints as a list---and only \B{NDArrayQ} can tell the two apart.

\mtranscript{10-internals/packed-gate}

That indistinguishability is a promise the evaluator has to keep, and it keeps it with
the \emph{transparency gate} of \S\ref{sec:int-eval}, step~7.5. There are thousands of
places in the source that read a function node's argument array directly; a packed
buffer has no such array. So before the evaluator hands a packed argument to any head
that has not explicitly declared itself packed-aware, it \emph{materializes} the buffer
back into an ordinary list. A head opts in only after it has been checked to read the
buffer correctly, and the gate additionally materializes an integer buffer for any head
not proven exact on \texttt{int64}---which is why \texttt{Total[Range[5]]} stays the
exact integer \texttt{15} and never becomes \texttt{15.}. The full story of packing,
the exactness contract, and the audit that guards the list of aware heads is in
Chapter~\ref{ch:datastructures} (\S\ref{sec:ds-packed}); what matters here is that
packing is a performance representation the evaluator makes invisible, at the cost of
one gate check per call.

\section{Bootstrapping: startup and the two-tier design}
\label{sec:int-boot}

When Mathilda starts, \texttt{main()} in \texttt{src/repl.c} builds the world in a
fixed order: \texttt{symtab\_init()} creates the symbol table and interns the kernel's
\texttt{SYM\_*} names, then \texttt{core\_init()} in \texttt{src/core.c} registers the
core builtins and calls each subsystem's \texttt{*\_init()} function in turn---several
dozen of them, from arithmetic through the special functions and the packed-array
substrate to the higher-level algebra and graphics services. That ordered list of calls
in \texttt{core.c} is the authoritative source of truth for initialization order.

Then comes the seam that gives the system its character. After the C side is built,
\texttt{main()} loads \texttt{src/internal/init.m}, which in turn reads the remaining
bootstrap \texttt{.m} files. This is Mathilda's \emph{two-tier design}:
performance-sensitive primitives are written in C, while a great deal of the actual
mathematics is expressed as ordinary rewrite rules in Mathilda's own language. The
derivative rules live in \texttt{deriv.m}; the integral tables are \texttt{.m} data;
and the same fixed-point evaluator that runs your rules runs these:

\mtranscript{10-internals/twotier-rules}

\texttt{D[Sin[x], x]}, the chain rule for \texttt{D[f[g[x]], x]}, and the
antiderivative of \(1/(1+x^2)\) are not special-cased in C---they are
\texttt{DownValues} loaded at startup, tried at step~11 of every relevant evaluation
(\S\ref{sec:int-eval}) exactly as a rule you write would be. This is the deepest
reason the architecture holds together: because everything is an expression and the
evaluator is a generic fixed-point loop, extending the system's mathematics most often
means adding a rule, not writing C. The recipes for doing so---adding a builtin, a
module, a pattern, an operator---are Chapter~\ref{ch:contributing}, and the build,
test and audit machinery that keeps all of this honest is
Chapter~\ref{ch:development}.

\include{chapters/11-development}
\include{chapters/12-contributing}
\include{chapters/13-about-the-author}

\appendix
\chapter{System Architecture at a Glance}
\label{app:architecture}

This appendix is a map. It gives a high-level tour of how Mathilda is put
together---the handful of core pieces every computation flows through---and points
at where in the source tree each piece lives, so that a curious reader can open the
right file. It is deliberately compact; the authoritative engineering reference is
\texttt{SPEC.md} at the root of the source tree, and Chapter~\ref{ch:internals}
develops these ideas in depth.

\section{The pipeline}
\label{app:pipeline}

At the largest scale Mathilda is a loop that turns a string into an expression,
rewrites that expression until it stops changing, and turns the result back into a
string. Five components do the work, and the evaluator consults the middle three on
every step.

\begin{figure}[h]
\begin{framed}
\small
\begin{verbatim}
    "Integrate[1/(1+x^2), x]"          input string
             |
             v
      [ Parser ]      parse.c      Pratt / operator-precedence
             |
             v   Expr* tree  ->  Integrate[Times[...], x]
             |
      [ Evaluator ]   eval.c       fixed-point loop
        |  ^  ^  ^
        |  |  |  +--- Built-in C library    (plus.c, calculus/, ...)
        |  |  +------ Pattern matcher (match.c) + Rule engine (replace.c)
        |  +--------- Attribute system (attr.c)
        |  +--------- Symbol table (symtab.c)  OwnValues / DownValues
             |
             v
      [ Printer ]     print.c      precedence-aware formatting
             |
             v
        "ArcTan[x]"                  output string
\end{verbatim}
\end{framed}
\caption{The evaluation pipeline. The parser and printer bracket a fixed-point
evaluator that, on every step, consults the symbol table, the attribute system,
the pattern matcher, and the built-in library.}
\label{fig:pipeline}
\end{figure}

The remaining sections walk these pieces in turn.

\section{Everything is an expression}
\label{app:expr}

The single most important design decision in Mathilda is that \emph{everything is an
expression}. A number, a symbol, a list, a matrix, a rule, even a pattern---all are
the same kind of tree, an \texttt{Expr} (defined in \texttt{src/expr.\{c,h\}}). An
\texttt{Expr} is a tagged union: an integer (machine \texttt{int64} or a GMP
bignum), a real, a symbol, a string, or a \emph{function} node carrying a head and a
list of argument expressions. The compound form \texttt{f[x, y]} is a function node
with head \texttt{f}; a list \texttt{\{a, b\}} is really \texttt{List[a, b]}
internally.

You can see the tree behind any surface syntax with \B{FullForm}:

\mtranscript{appendix-architecture/expr-model}

Read those lines as X-rays. The sum \texttt{a + b c} is
\texttt{Plus[a, Times[b, c]]}; a fraction is a \texttt{Rational}; a division is a
\texttt{Times} against a \texttt{Power} with exponent $-1$; and even a pattern like
\texttt{x\_} is an ordinary expression, \texttt{Pattern[x, Blank[]]}. Because atoms
carry a head too---\B{Head} of an integer is \texttt{Integer}, of a list is
\texttt{List}---one uniform set of tools (\B{Part}, \B{Map}, \B{Replace}, and the
rest) operates on \emph{everything}. This uniformity is why the evaluator and the
pattern matcher can be small: they only ever manipulate one data structure.

Integers illustrate a second theme: representations that adapt. A value begins as a
fast machine word and is promoted to a GMP bignum only when an operation would
overflow, demoting back when a result fits again---so there is only ever ``an
integer,'' never a separate big-integer mode.

\section{How a call is evaluated}
\label{app:eval}

Evaluation is a fixed-point process (\texttt{src/eval.c}): \texttt{evaluate} applies
one \texttt{evaluate\_step} after another until the expression stops changing. For a
function call \texttt{f[a1, \dots, aN]} a single step proceeds, in order:

\begin{enumerate}
  \item evaluate the head \texttt{f};
  \item read \texttt{f}'s \textbf{attributes};
  \item evaluate the arguments---except those held back by a \texttt{Hold} attribute;
  \item if \texttt{f} is \textbf{Listable}, thread over any list arguments;
  \item if \texttt{f} is \textbf{Flat}, flatten nested \texttt{f[\dots]} calls;
  \item if \texttt{f} is \textbf{Orderless}, sort the arguments canonically;
  \item call \texttt{f}'s built-in C function, if it has one---this returns a new
        \texttt{Expr*} on success, or \texttt{NULL} to mean ``I can't evaluate this,
        leave it alone'';
  \item otherwise try \texttt{f}'s \textbf{DownValues}---the user- or library-defined
        rewrite rules---and take the first that matches.
\end{enumerate}

The attribute bits (\texttt{src/attr.c}) are the small vocabulary that steers all of
this. They are what let the evaluator stay generic while each head opts into exactly
the behaviour it needs\calloutnote{Under the hood}{\B{Plus} is \texttt{Flat}
(associative), \texttt{Orderless} (commutative), \texttt{Listable} (threads over lists),
\texttt{OneIdentity}, and a \texttt{NumericFunction}; \B{Sin} is merely \texttt{Listable}
and numeric; \B{Table} is \texttt{HoldAll}, so its iterator specification is \emph{not}
evaluated before \texttt{Table} runs. Nothing about the evaluator's main loop changes
between these heads---only their attribute bits do.}:

\mtranscript{appendix-architecture/attributes}

Because a step just rewrites and the loop runs to a fixed point, rule chains compose
on their own, with no explicit plumbing. Define \texttt{f} in terms of \texttt{g} and
\texttt{g} in terms of squaring, and the answer falls out over successive passes:

\mtranscript{appendix-architecture/fixed-point}

\section{Matching, rules, and the built-in library}
\label{app:match}

Two subsystems implement the last two steps above. The \textbf{pattern matcher}
(\texttt{src/match.c}) unifies an expression against a pattern with backtracking,
handling blanks (\texttt{\_}), typed blanks (\texttt{\_h}), sequence patterns
(\texttt{\_\_}, \texttt{\_\_\_}), tests, and conditions, binding pattern variables in
a match environment. The \textbf{rule engine} (\texttt{src/replace.c}) applies
rules---\texttt{Rule} (\texttt{->}, immediate) and \texttt{RuleDelayed}
(\texttt{:>}, delayed)---through \B{Replace}, \B{ReplaceAll} (\texttt{/.}), and
\B{ReplaceRepeated} (\texttt{//.}).

The \textbf{built-in C library} supplies the primitives. Each builtin obeys a single
ownership contract: it takes ownership of its argument expression and returns either
a freshly built \texttt{Expr*} (success) or \texttt{NULL} (``not my case,'' leaving
the expression for a rule to handle). That \texttt{NULL} convention is what lets
symbolic arguments flow untouched through heads that only know how to handle numbers.

\section{The source tree}
\label{app:tree}

The C sources live under \texttt{src/}. The core is a small number of files; the bulk
is the mathematical library, organised so that---as a rule of thumb---each builtin
gets its own file, grouped into a subsystem directory.

\begin{figure}[h]
\begin{framed}
\small
\begin{verbatim}
src/
  expr.{c,h}     expression representation (the tagged union)
  parse.{c,h}    Pratt parser            print.{c,h}   printer
  eval.{c,h}     fixed-point evaluator    symtab.{c,h}  symbol table
  attr.{c,h}     attribute bitflags       match.{c,h}   pattern matcher
  replace.{c,h}  rule engine              core.{c,h}    builtin + init hub
  plus.c times.c power.c arithmetic.c     arithmetic heads
  ndarray.* ndkernels.* pack.*            packed / dense-array substrate
  compile/       Compile[] bytecode compiler and auto-compilation
  calculus/      D, Series, Limit, Integrate, Risch-Norman
  poly/  rat/  simp/  linalg/  solve/     algebra & solving
  numerical_calculus/  numerical_roots/   N-prefixed numerics
  sum/  product/  special_functions/      series, products, transcendentals
  numbertheory/  stats/  graph/  ml/      further subsystems
  strings/  list/  graphics/              text, lists, 2D/3D graphics
  internal/      init.m and .m rule files (derivatives, integral tables)
tests/           CMake-built unit suite
makefile         primary build system
\end{verbatim}
\end{framed}
\caption{A condensed view of the source tree. The full layout, with a line for every
directory, is in \texttt{SPEC.md}~\S2.}
\label{fig:tree}
\end{figure}

\section{Startup and module initialization}
\label{app:init}

When Mathilda starts, \texttt{main()} in \texttt{src/repl.c} calls
\texttt{symtab\_init()} to create the symbol table, then \texttt{core\_init()} in
\texttt{src/core.c}. \texttt{core\_init()} registers the core builtins and then calls
each subsystem's \texttt{*\_init()} function in a fixed order---several dozen of them,
spanning arithmetic, special functions, the packed-array substrate, the numeric
subsystems, and the higher-level algebra and graphics services. \texttt{core.c} is
the source of truth for that order.

After the C side is initialised, \texttt{main()} loads
\texttt{src/internal/init.m}, which in turn \texttt{Get[]}s the remaining bootstrap
\texttt{.m} files. This is the seam of Mathilda's \textbf{two-tier design}:
performance-sensitive primitives are written in C, while higher-level mathematical
identities---the derivative rules in \texttt{deriv.m}, the integral tables---are
expressed as ordinary rewrite rules in Mathilda's own language. The same fixed-point
evaluator that runs a user's rules runs these, so extending the system's mathematics
often means adding a rule, not writing C.

\bigskip
\noindent
For the authoritative detail behind every claim here, see \texttt{SPEC.md}; for
recipes on adding builtins, modules, patterns, and operators, see
\texttt{docs/extending.md} and Chapter~\ref{ch:contributing}.


\backmatter
\chapter{Colophon}

This book was written in \LaTeX\ using the \textsf{memoir} document class, with
\textsf{listings} for the verified transcripts, \textsf{tcolorbox} for the
callout boxes, and \textsf{biblatex} for the references. Body text is set in
Latin Modern; transcripts are set in Inconsolata.

Every \texttt{In[\,]}/\texttt{Out[\,]} transcript is generated at build time by
running its input through the Mathilda binary via the system's NDJSON pipe
protocol, using the same driver (\texttt{site/verify\_tutorial.py}) that verifies
the online tutorials. Every built-in name printed as a hyperlink is resolved
from Mathilda's own builtin index (\texttt{site/docs/assets/builtins.json}) to
its reference page. The tools that do this---\texttt{book/tools/gen\_links.py},
\texttt{book/tools/build\_examples.py}, and \texttt{book/tools/check\_links.py}
---are small and are meant to be read.

The book's guiding principles, and instructions for extending it, are recorded in
\texttt{book/CONTEXT.md}; its section-by-section roadmap is in
\texttt{book/ROADMAP.md}.

\printbibliography[heading=bibintoc,title={References}]
\printindex

\end{document}
